\documentclass[fs85,seriesWit]{korfarm-base}
\usepackage{korfarm-witt}

% ============================================================
% passage-note.tex (v11)
% 지문 박스 + 첫 페이지 우측 메모란 (동적 높이)
% 정책: PAGE_BREAK_POLICY.md §3, §4 참조
%
% 변경 (v11):
%  · 지문 박스 폭 확대: enlarge right by=-3.7cm  (메모 3.4cm + gap 3mm)
%  · 메모란 동적 높이: frame.north east ~ frame.south east 사이 자동 채움
%  · 메모 줄선: clip 영역으로 박스 내부만 채움 (박스 높이 모르고도 작동)
%  · 문단 들여쓰기 활성화: tcolorbox 내부 \parindent=1em, \noindent 제거
%
% 사용:
%   \kfPassageNote{제목}{지문 본문}      → 메모란 포함
%   \kfPassageNoteNT{지문 본문}          → 메모란 포함, 제목 없음
%   \kfPassageFull{제목}{본문}           → 메모란 없음 (활동 내 보기)
% ============================================================

\makeatletter

% 메모란 규격 (cm 고정)
\newlength{\kf@memoW}
\setlength{\kf@memoW}{3.4cm}
\newlength{\kf@memoGap}
\setlength{\kf@memoGap}{3mm}

% --- 메인: 제목 있는 버전 ---
\newcommand{\kf@passagenote@with}[2]{%
  \par\ifkfPassageNewPage\ifdim\pagetotal>0.4\textheight\clearpage\fi\fi\smallskip\needspace{6\baselineskip}%
  \begin{tcolorbox}[
    enhanced, breakable,
    width=\dimexpr\linewidth-4cm\relax,
    colback=kfPassageBg, colframe=kfPrimary!70,
    boxrule=0.6pt, arc=3pt,
    left=\kfBoxPad, right=\kfBoxPad, top=\dimexpr\kfBoxPad+8pt\relax, bottom=\kfBoxPad,
    title={\small\bfseries\color{kfPrimary} #1},
    coltitle=kfPrimary,
    colbacktitle=kfPrimaryLight!60,
    attach boxed title to top left={yshift=-2.4mm, xshift=8pt},
    boxed title style={colframe=kfPrimary!50, boxrule=0.5pt, arc=2pt},
    parbox=false,
    overlay unbroken and first={%
      \begin{scope}
        \draw[kfRule, line width=0.4pt, fill=kfNoteBg, rounded corners=2pt]
          ([xshift=\kf@memoGap]frame.north east) rectangle
          ([xshift=\kf@memoGap+\kf@memoW]frame.south east);
        \node[anchor=north east, font=\footnotesize\bfseries\color{kfMute}, inner sep=3pt]
          at ([xshift=\kf@memoGap+\kf@memoW, yshift=-2pt]frame.north east) {메모};
        \node[anchor=north east, fill=kfPrimary!15, text=kfPrimary,
              font=\scriptsize\bfseries, rounded corners=2pt,
              inner xsep=4pt, inner ysep=2pt]
          at ([xshift=\kf@memoGap+\kf@memoW-3pt, yshift=-19pt]frame.north east) {\kf@memocat};
        \begin{scope}
          \path[clip]
            ([xshift=\kf@memoGap]frame.north east) rectangle
            ([xshift=\kf@memoGap+\kf@memoW]frame.south east);
          \foreach \i in {1,...,40}{%
            \draw[kfRule, line width=0.25pt]
              ([xshift=\kf@memoGap+2mm, yshift=-7mm - \i*5mm]frame.north east) --
              ([xshift=\kf@memoGap+\kf@memoW-2mm, yshift=-7mm - \i*5mm]frame.north east);%
          }
        \end{scope}
      \end{scope}
    }
  ]
    \setlength{\parindent}{1em}%
    \hspace*{1em}\ignorespaces#2%
  \end{tcolorbox}\par\smallskip
}

% --- 제목 없는 버전 ---
\newcommand{\kf@passagenote@notitle}[1]{%
  \par\ifkfPassageNewPage\ifdim\pagetotal>0.4\textheight\clearpage\fi\fi\smallskip\needspace{6\baselineskip}%
  \begin{tcolorbox}[
    enhanced, breakable,
    width=\dimexpr\linewidth-4cm\relax,
    colback=kfPassageBg, colframe=kfPrimary!70,
    boxrule=0.6pt, arc=3pt,
    left=\kfBoxPad, right=\kfBoxPad, top=\dimexpr\kfBoxPad+8pt\relax, bottom=\kfBoxPad,
    parbox=false,
    overlay unbroken and first={%
      \begin{scope}
        \draw[kfRule, line width=0.4pt, fill=kfNoteBg, rounded corners=2pt]
          ([xshift=\kf@memoGap]frame.north east) rectangle
          ([xshift=\kf@memoGap+\kf@memoW]frame.south east);
        \node[anchor=north east, font=\footnotesize\bfseries\color{kfMute}, inner sep=3pt]
          at ([xshift=\kf@memoGap+\kf@memoW, yshift=-2pt]frame.north east) {메모};
        \node[anchor=north east, fill=kfPrimary!15, text=kfPrimary,
              font=\scriptsize\bfseries, rounded corners=2pt,
              inner xsep=4pt, inner ysep=2pt]
          at ([xshift=\kf@memoGap+\kf@memoW-3pt, yshift=-19pt]frame.north east) {\kf@memocat};
        \begin{scope}
          \path[clip]
            ([xshift=\kf@memoGap]frame.north east) rectangle
            ([xshift=\kf@memoGap+\kf@memoW]frame.south east);
          \foreach \i in {1,...,40}{%
            \draw[kfRule, line width=0.25pt]
              ([xshift=\kf@memoGap+2mm, yshift=-7mm - \i*5mm]frame.north east) --
              ([xshift=\kf@memoGap+\kf@memoW-2mm, yshift=-7mm - \i*5mm]frame.north east);%
          }
        \end{scope}
      \end{scope}
    }
  ]
    \setlength{\parindent}{1em}%
    \hspace*{1em}\ignorespaces#1%
  \end{tcolorbox}\par\smallskip
}

\newcommand{\kfPassageNote}[2]{%
  \def\kf@tmptitle{#1}%
  \ifx\kf@tmptitle\@empty
    \kf@passagenote@notitle{#2}%
  \else
    \kf@passagenote@with{#1}{#2}%
  \fi
}
\newcommand{\kfPassageNoteNT}[1]{\kf@passagenote@notitle{#1}}
\makeatother

% ============================================================
% 풀폭 지문 박스 (메모란 없음) — 활동 내 보기 박스
% PAGE_BREAK_POLICY.md §3
% ============================================================
\makeatletter
\newcommand{\kf@passagefull@with}[2]{%
  \par\ifkfPassageNewPage\ifdim\pagetotal>0.4\textheight\clearpage\fi\fi\smallskip\needspace{6\baselineskip}%
  \begin{tcolorbox}[
    enhanced, breakable,
    colback=kfPassageBg, colframe=kfPrimary!70,
    boxrule=0.6pt, arc=3pt,
    left=\kfBoxPad, right=\kfBoxPad, top=\dimexpr\kfBoxPad+8pt\relax, bottom=\kfBoxPad,
    title={\small\bfseries\color{kfPrimary} #1},
    coltitle=kfPrimary,
    colbacktitle=kfPrimaryLight!60,
    attach boxed title to top left={yshift=-2.4mm, xshift=8pt},
    boxed title style={colframe=kfPrimary!50, boxrule=0.5pt, arc=2pt},
    parbox=false
  ]
    \setlength{\parindent}{1em}%
    \hspace*{1em}\ignorespaces#2%
  \end{tcolorbox}\par\smallskip
}
\newcommand{\kf@passagefull@notitle}[1]{%
  \par\ifkfPassageNewPage\ifdim\pagetotal>0.4\textheight\clearpage\fi\fi\smallskip\needspace{6\baselineskip}%
  \begin{tcolorbox}[
    enhanced, breakable,
    colback=kfPassageBg, colframe=kfPrimary!70,
    boxrule=0.6pt, arc=3pt,
    left=\kfBoxPad, right=\kfBoxPad, top=\dimexpr\kfBoxPad+8pt\relax, bottom=\kfBoxPad,
    parbox=false
  ]
    \setlength{\parindent}{1em}%
    \hspace*{1em}\ignorespaces#1%
  \end{tcolorbox}\par\smallskip
}
\newcommand{\kfPassageFull}[2]{%
  \def\kf@tmptitle{#1}%
  \ifx\kf@tmptitle\@empty
    \kf@passagefull@notitle{#2}%
  \else
    \kf@passagefull@with{#1}{#2}%
  \fi
}
\newcommand{\kfPassageFullNT}[1]{\kf@passagefull@notitle{#1}}
\makeatother

% ============================================================
% 작은 문장 박스 (문장 독해용) — PAGE_BREAK_POLICY.md §2.5
% ============================================================
\newcommand{\kfSentenceBox}[2]{%
  \par\smallskip\needspace{3\baselineskip}%
  \noindent\begin{tcolorbox}[
    enhanced, breakable=false,
    colback=kfPassageBg, colframe=kfMute,
    boxrule=0.4pt, arc=2pt,
    left=8pt, right=6pt, top=4pt, bottom=4pt,
    title={\footnotesize\bfseries\color{kfPrimary} #1},
    coltitle=kfPrimary, colbacktitle=kfPaper,
    attach boxed title to top left={yshift=-1.6mm, xshift=4pt},
    boxed title style={colframe=kfMute, boxrule=0.3pt, arc=1pt}
  ]%
    \setstretch{1.3}#2%
  \end{tcolorbox}\par\smallskip%
}

% ============================================================
% 지문 본문 아래 안내/정리 박스 (v16 신규)
% 사용: \kfPassageHint{한 줄 정리}{본문 안내 텍스트}
% ============================================================
\newcommand{\kfPassageHint}[2]{%
  \par\smallskip\needspace{4\baselineskip}%
  \begin{tcolorbox}[
    enhanced, breakable=false,
    colback=kfPrimary!5, colframe=kfPrimary!40,
    boxrule=0.4pt, arc=2pt,
    left=8pt, right=6pt, top=4pt, bottom=4pt,
    title={\footnotesize\bfseries\color{kfPrimary} #1},
    coltitle=kfPrimary, colbacktitle=kfPaper,
    attach boxed title to top left={yshift=-1.5mm, xshift=4pt},
    boxed title style={colframe=kfPrimary!40, boxrule=0.3pt, arc=1pt}
  ]%
    \small\setstretch{1.3}#2%
  \end{tcolorbox}\par\smallskip%
}

% ============================================================
% 환경 버전 (호환성) — 메모란 포함
% ============================================================
\makeatletter
\newenvironment{kfPassageNoteEnv}[1]%
  {\par\needspace{6\baselineskip}%
   \def\kf@psrc{#1}%
   \begin{tcolorbox}[
     enhanced, breakable,
     width=\dimexpr\linewidth-4cm\relax,
     colback=kfPassageBg, colframe=kfPrimary!70,
     boxrule=0.6pt, arc=3pt,
     left=\kfBoxPad, right=\kfBoxPad, top=\dimexpr\kfBoxPad+8pt\relax, bottom=\kfBoxPad,
     title={\small\bfseries\color{kfPrimary} \kf@psrc},
     coltitle=kfPrimary, colbacktitle=kfPrimaryLight!60,
     attach boxed title to top left={yshift=-2.4mm, xshift=8pt},
     boxed title style={colframe=kfPrimary!50, boxrule=0.5pt, arc=2pt},
     parbox=false,
     overlay unbroken and first={%
       \begin{scope}
         \draw[kfRule, line width=0.4pt, fill=kfNoteBg, rounded corners=2pt]
           ([xshift=\kf@memoGap]frame.north east) rectangle
           ([xshift=\kf@memoGap+\kf@memoW]frame.south east);
         \node[anchor=north east, font=\footnotesize\bfseries\color{kfMute}, inner sep=3pt]
           at ([xshift=\kf@memoGap+\kf@memoW, yshift=-2pt]frame.north east) {메모};
         \begin{scope}
           \path[clip]
             ([xshift=\kf@memoGap]frame.north east) rectangle
             ([xshift=\kf@memoGap+\kf@memoW]frame.south east);
           \foreach \i in {1,...,40}{%
             \draw[kfRule, line width=0.25pt]
               ([xshift=\kf@memoGap+2mm, yshift=-7mm - \i*5mm]frame.north east) --
               ([xshift=\kf@memoGap+\kf@memoW-2mm, yshift=-7mm - \i*5mm]frame.north east);%
           }
         \end{scope}
       \end{scope}
     }
   ]%
   \setlength{\parindent}{1em}%
  }%
  {\end{tcolorbox}\par\smallskip}
\makeatother

% ============================================================
% condition-box.tex
% <조건> 박스 컴포넌트 (tcolorbox)
% 사용:
%   \begin{kfCondition}
%     \item 첫째 조건
%     \item 둘째 조건
%   \end{kfCondition}
% ============================================================

\newenvironment{kfCondition}%
  {\par\smallskip\needspace{4\baselineskip}%
   \begin{tcolorbox}[
     enhanced, breakable=false,
     colback=kfPrimaryLight!40, colframe=kfPrimary,
     boxrule=0.5pt, arc=2pt,
     left=\kfBoxPad, right=\kfBoxPad, top=\kfBoxPad, bottom=\kfBoxPad,
     title={\small\bfseries\color{kfPrimary} \,조건\,},
     coltitle=kfPrimary, colbacktitle=kfPaper,
     attach boxed title to top left={yshift=-2mm, xshift=6pt},
     boxed title style={colframe=kfPrimary, boxrule=0.4pt, arc=1pt}
   ]%
   \begin{enumerate}[leftmargin=16pt, itemsep=1pt, topsep=1pt,
                    label=\textbf{\arabic*.}]}%
  {\end{enumerate}\end{tcolorbox}\par\smallskip}

% ============================================================
% evidence-box.tex
% <보기> 박스 컴포넌트
% 사용:
%   \begin{kfEvidence}
%     보기 본문 ...
%   \end{kfEvidence}
% ============================================================

\newenvironment{kfEvidence}%
  {\par\smallskip\needspace{4\baselineskip}%
   \begin{tcolorbox}[
     enhanced, breakable=false,
     colback=kfPaper, colframe=kfMute,
     boxrule=0.5pt, arc=2pt,
     left=\kfBoxPad, right=\kfBoxPad, top=\kfBoxPad, bottom=\kfBoxPad,
     title={\small\bfseries\color{kfInk} \,보기\,},
     coltitle=kfInk, colbacktitle=kfPrimaryLight!50,
     attach boxed title to top left={yshift=-2mm, xshift=6pt},
     boxed title style={colframe=kfMute, boxrule=0.4pt, arc=1pt}
   ]%
   \leavevmode\mbox{}\ignorespaces%
  }%
  {\end{tcolorbox}\par\smallskip}

% ============================================================
% choice-list.tex
% 5선지 (① ② ③ ④ ⑤) 자동 들여쓰기 컴포넌트
% 사용:
%   \begin{kfChoices}
%     \kfChoice 첫째 선택지
%     \kfChoice 둘째 선택지
%     \kfChoice 셋째 선택지
%     \kfChoice 넷째 선택지
%     \kfChoice 다섯째 선택지
%   \end{kfChoices}
% ============================================================

\newcounter{kfChoiceCnt}
\newcommand{\kfChoiceMark}[1]{%
  \ifcase#1\relax%
  \or ①\or ②\or ③\or ④\or ⑤\or ⑥\or ⑦\or ⑧\or ⑨%
  \fi
}

% 환경: 자동 번호
\newenvironment{kfChoices}%
  {\setcounter{kfChoiceCnt}{0}%
   \par\smallskip%
   \begin{list}{}{%
     \setlength{\leftmargin}{20pt}%
     \setlength{\itemindent}{0pt}%
     \setlength{\labelwidth}{18pt}%
     \setlength{\labelsep}{6pt}%
     \setlength{\itemsep}{2pt}%
     \setlength{\topsep}{1pt}%
     \setlength{\parsep}{0pt}%
     \renewcommand{\makelabel}[1]{##1\hss}%
   }}%
  {\end{list}\par\smallskip}

\newcommand{\kfChoice}{%
  \stepcounter{kfChoiceCnt}%
  \item[\textbf{\kfChoiceMark{\value{kfChoiceCnt}}}]\ignorespaces%
}

% --- 한 줄 5선지 (짧은 선택지용) ---
\newcommand{\kfChoicesInline}[5]{%
  \par\smallskip%
  \noindent\textbf{①}~#1\quad\textbf{②}~#2\quad\textbf{③}~#3\quad\textbf{④}~#4\quad\textbf{⑤}~#5%
  \par\smallskip%
}

% ============================================================
% answer-note.tex
% 답안 작성란 (노트 스타일, 줄친 영역)
% 사용:
%   \kfAnswerLines{5}        % 5줄 답안란
%   \kfAnswerNote{서술형 답안란}{8} % 라벨+8줄
% ============================================================

% 단일 줄 (필요 시 인라인)
\newcommand{\kfAnswerLine}{%
  \par\noindent\textcolor{kfRule}{\rule{\linewidth}{0.4pt}}\par
}

% N줄 답안란 (간단)
\newcommand{\kfAnswerLines}[1]{%
  \par\smallskip%
  \begin{minipage}{\linewidth}%
  \foreach \i in {1,...,#1}{%
    \textcolor{kfRule}{\rule{\linewidth}{0.4pt}}\\[6pt]%
  }%
  \end{minipage}%
  \par\smallskip%
}

% 라벨이 있는 노트형 답안란
\newcommand{\kfAnswerNote}[2]{%
  \par\smallskip\needspace{#2\baselineskip}%
  \begin{tcolorbox}[
    enhanced, breakable=false,
    colback=kfPaper, colframe=kfRule,
    boxrule=0.4pt, arc=2pt,
    left=\kfBoxPad, right=\kfBoxPad, top=\kfBoxPad, bottom=\kfBoxPad,
    title={\footnotesize\bfseries\color{kfMute} #1},
    coltitle=kfMute, colbacktitle=kfPaper,
    attach boxed title to top left={yshift=-1.5mm, xshift=6pt},
    boxed title style={colframe=kfRule, boxrule=0.3pt, arc=1pt}
  ]%
  \foreach \i in {1,...,#2}{%
    \textcolor{kfRule}{\rule{\linewidth}{0.3pt}}\\[6pt]%
  }%
  \end{tcolorbox}\par\smallskip%
}

% 짧은 빈칸 (단답형)
\newcommand{\kfBlank}[1][3cm]{\underline{\hspace{#1}}}
% 넓은 빈칸 (서술 한 줄)
\newcommand{\kfBlankWide}[1][6cm]{\underline{\hspace{#1}}}
% 괄호 빈칸
\newcommand{\kfBlankParen}{(\,\hspace{1cm}\,)}
% O/X 빈칸
\newcommand{\kfBlankOX}{(\,\hspace{1.5cm}\,)}

% ============================================================
% question-block.tex
% 문제 블록 (번호 배지 + 발문 + 보기/조건 + 선택지 + 답안란)
% 모든 구성을 하나의 minipage로 묶어 단/페이지 단절 금지
%
% 사용:
%   \begin{kfQBlock}{1}{객관식}
%     \kfQStem{다음 글에 대한 설명으로 가장 적절한 것은?}
%     \begin{kfEvidence} ... \end{kfEvidence}    % 선택
%     \begin{kfCondition} ... \end{kfCondition}  % 선택
%     \begin{kfChoices}
%       \kfChoice ...
%     \end{kfChoices}
%     \kfAnswerLines{2}                            % 선택
%   \end{kfQBlock}
% ============================================================

% 무단절 minipage로 묶고, 발문/보기/조건/선택지/답안란을 자유롭게 배치
% 사용자 피드백 #4: [객관식]/[서술형] 등 텍스트 라벨 제거 — 번호 배지만 표시
% #2(유형 라벨)는 호환성 인자로만 받고 출력하지 않음.
\newenvironment{kfQBlock}[2]%
  {\par\smallskip\needspace{6\baselineskip}%
   \noindent\begin{minipage}{\linewidth}%
   \samepage%
   \par\noindent
   \kfQNum{#1}\hspace{4pt}%
   \begin{minipage}[t]{\dimexpr\linewidth-30pt\relax}%
  }%
  {\end{minipage}\end{minipage}\par\smallskip}

% 문제 발문
\newcommand{\kfQStem}[1]{%
  \par\noindent#1\par\vspace{2pt}%
}

% 짧은 소문항 라벨 (1), (2), (3)
\newcounter{kfSubQ}
\newcommand{\kfSubQReset}{\setcounter{kfSubQ}{0}}
\newcommand{\kfSubQ}{%
  \stepcounter{kfSubQ}%
  \par\noindent\textbf{(\arabic{kfSubQ})}~\ignorespaces%
}

% ============================================================
% activity-table.tex
% 활동: 정리표 / 내용 확인표 (마크다운 표 → tabularx 변환)
% 사용:
%   % 일반 tabular (직접 컬럼 명시)
%   \begin{kfActTable}{|l|X|X|}
%     \kfTHead{항목} & \kfTHead{설명} & \kfTHead{학생 작성} \\\hline
%     ① & ... & \kfTBlank \\\hline
%   \end{kfActTable}
%
%   % adjustbox 자동 축소 (정해진 폭으로 강제)
%   \kfActTableWrap{
%     ... (\begin{tabular}...\end{tabular} 코드) ...
%   }
% ============================================================

% 사용 시 본문 마지막 행은 \\\hline 으로 끝나야 함.
% v17.1: 흰색 mask로 Canva 배경 본문 침범 차단
\newenvironment{kfActTable}[1]%
  {\par\smallskip\needspace{4\baselineskip}%
   \renewcommand{\arraystretch}{1.3}%
   \arrayrulecolor{kfRule}%
   \begin{tcolorbox}[blanker, colback=white, left=2pt, right=2pt, top=2pt, bottom=2pt]%
   \noindent\begin{adjustbox}{max width=\linewidth}%
   \begin{tabular}{#1}%
   \hline}%
  {\end{tabular}%
   \end{adjustbox}%
   \end{tcolorbox}%
   \arrayrulecolor{black}%
   \par\smallskip}

% 헤더 행 헬퍼
\newcommand{\kfTHead}[1]{\textbf{\color{kfPrimary} #1}}
\newcommand{\kfTHeadRow}[1]{\rowcolor{kfPrimaryLight!50}\textbf{\color{kfPrimary} #1}}

% 빈 셀 (학생 작성용)
\newcommand{\kfTBlank}{\rule[-1.6em]{0pt}{3.4em}}
\newcommand{\kfTBlankSm}{\rule[-1.0em]{0pt}{2.4em}}

% adjustbox 강제 축소 래퍼 — Python에서 변환된 raw tabular 블록을 감쌀 때 사용
\newcommand{\kfActTableWrap}[1]{%
  \par\smallskip\needspace{4\baselineskip}%
  \begin{tcolorbox}[blanker, colback=white, left=2pt, right=2pt, top=2pt, bottom=2pt]%
  \noindent\begin{adjustbox}{max width=\linewidth, center}%
  #1%
  \end{adjustbox}%
  \end{tcolorbox}\par\smallskip%
}

% ============================================================
% activity-vocab.tex
% 어휘 학습 표 (3열: 번호 - 뜻 - 빈칸)
% 사용:
%   \begin{kfVocab}
%     \kfVocabItem{1}{눈으로 보고 마음으로 깨달음}
%     \kfVocabItem{2}{사물의 이치를 따져 헤아림}
%   \end{kfVocab}
% ============================================================

% adjustbox로 폭 강제 + 일반 tabular + p{} 열 사용
\newenvironment{kfVocab}%
  {\par\smallskip\needspace{4\baselineskip}%
   \renewcommand{\arraystretch}{1.35}%
   \arrayrulecolor{kfRule}%
   \noindent\begin{adjustbox}{max width=\linewidth, center}%
   \begin{tabular}{|>{\centering\arraybackslash}m{1.0cm}|m{8.5cm}|>{\centering\arraybackslash}m{4.0cm}|}%
   \hline
   \rowcolor{kfPrimaryLight!50}%
   \multicolumn{1}{|c|}{\textbf{\color{kfPrimary}번호}} &
   \multicolumn{1}{c|}{\textbf{\color{kfPrimary}뜻풀이}} &
   \multicolumn{1}{c|}{\textbf{\color{kfPrimary}어휘 (정답 작성)}} \\\hline
   \ignorespaces}%
  {\end{tabular}%
   \end{adjustbox}%
   \arrayrulecolor{black}%
   \par\smallskip}

\newcommand{\kfVocabItem}[2]{%
  \textbf{#1} & #2 & \rule[-1.0em]{0pt}{2.4em}\\\hline%
}

% ============================================================
% activity-blank.tex
% 빈칸 채우기 활동
% 사용:
%   \begin{kfBlankList}
%     \kfBlankItem 첫 번째 문장에서 (\kfBlank)은(는) ...
%     \kfBlankItem 둘째 문장에서 ...
%   \end{kfBlankList}
% ============================================================

\newenvironment{kfBlankList}%
  {\par\smallskip%
   \begin{enumerate}[leftmargin=20pt, itemsep=4pt, topsep=2pt,
                    label=\textbf{\arabic*.}, labelsep=6pt]}%
  {\end{enumerate}\par\smallskip}

\newcommand{\kfBlankItem}{\item}

% 텍스트 안에 박스형 빈칸 (단어 1개 채우기)
\newcommand{\kfBlankBox}[1][2.4cm]{%
  \tikz[baseline=-0.5ex]{%
    \draw[kfRule, line width=0.4pt] (0,0) -- ++(#1,0);%
  }%
}

% ============================================================
% activity-ox.tex (v15)
% O/X 표기 활동 — 그래픽 디자인
% 사용:
%   \begin{kfOXList}
%     \kfOXItem 글쓴이는 자연을 예찬하고 있다.\kfOX
%     \kfOXItem 1연과 4연은 수미상관 구조이다.\kfOX
%   \end{kfOXList}
%
% \kfOX = 우측에 [O] [X] 두 박스 (학생이 하나 동그라미)
% ============================================================

\newenvironment{kfOXList}%
  {\par\smallskip%
   \begin{enumerate}[leftmargin=20pt, itemsep=5pt, topsep=2pt,
                    label=\textbf{\arabic*.}, labelsep=6pt]}%
  {\end{enumerate}\par\smallskip}

\newcommand{\kfOXItem}{\item}

% v16: O/X 그래픽 박스 키움 (18×14 → 24×20pt) + 영역색 강조
\newcommand{\kfOX}{%
  \hfill\nolinebreak
  \tikz[baseline=-0.9ex]{%
    \node[draw=kfPrimary, fill=kfPrimary!5, line width=0.6pt, rounded corners=3pt,
          minimum width=24pt, minimum height=20pt, inner sep=0pt,
          font=\normalsize\bfseries\color{kfPrimary}] (ob) {O};%
    \node[draw=kfMute, fill=kfMute!5, line width=0.6pt, rounded corners=3pt,
          minimum width=24pt, minimum height=20pt, inner sep=0pt,
          font=\normalsize\bfseries\color{kfMute},
          right=4pt of ob] (xb) {X};%
  }%
}

% ============================================================
% activity-connect.tex
% 좌우 선 긋기 활동 (2단 양식 + 중앙 여백)
% 사용:
%   \begin{kfConnect}
%     \kfPair{사르트르}{실존은 본질에 앞선다}
%     \kfPair{니체}{초인 사상}
%     \kfPair{하이데거}{현존재}
%   \end{kfConnect}
% ============================================================

\newenvironment{kfConnect}%
  {\par\smallskip%
   \renewcommand{\arraystretch}{1.6}%
   \begin{tabular}{@{}>{\raggedleft\arraybackslash}m{4.5cm}
                     @{\hspace{2.5cm}}
                     >{\raggedright\arraybackslash}m{6.5cm}@{}}}%
  {\end{tabular}\par\smallskip}

\newcommand{\kfPair}[2]{%
  \tikz[baseline]{\node[draw=kfRule, fill=kfPrimaryLight!30, rounded corners=2pt,
        inner sep=4pt, font=\small]{#1};} \tikz[baseline]{\node[circle, draw=kfMute, fill=kfPaper, inner sep=1.5pt]{};}
  &
  \tikz[baseline]{\node[circle, draw=kfMute, fill=kfPaper, inner sep=1.5pt]{};} \tikz[baseline]{\node[draw=kfRule, fill=kfNoteBg, rounded corners=2pt,
        inner sep=4pt, font=\small]{#2};} \\
}

% ============================================================
% activity-writing.tex
% 글쓰기 활동 (큰 노트 영역)
% 사용:
%   \kfWriting{독후감}{12}     % 라벨 + 12줄 큰 노트
%   \kfWritingPrompt{프롬프트 텍스트}
% ============================================================

\newcommand{\kfWritingPrompt}[1]{%
  \par\smallskip%
  \begin{tcolorbox}[
    enhanced, breakable=false,
    colback=kfPrimaryLight!30, colframe=kfPrimary,
    boxrule=0.4pt, arc=3pt,
    left=\kfBoxPad, right=\kfBoxPad, top=\kfBoxPad, bottom=\kfBoxPad,
  ]%
  \small\textbf{\color{kfPrimary}글쓰기 안내}\\[2pt]%
  #1
  \end{tcolorbox}\par\smallskip%
}

\newcommand{\kfWriting}[2]{%
  \par\smallskip\needspace{\numexpr#2+2\relax\baselineskip}%
  \begin{tcolorbox}[
    enhanced, breakable=false,
    colback=kfPaper, colframe=kfRule,
    boxrule=0.4pt, arc=2pt,
    left=\kfBoxPad, right=\kfBoxPad, top=\kfBoxPad, bottom=\kfBoxPad,
    title={\footnotesize\bfseries\color{kfMute} #1},
    coltitle=kfMute, colbacktitle=kfPaper,
    attach boxed title to top left={yshift=-1.5mm, xshift=6pt},
    boxed title style={colframe=kfRule, boxrule=0.3pt, arc=1pt}
  ]%
  \foreach \i in {1,...,#2}{%
    \textcolor{kfRule}{\rule{\linewidth}{0.3pt}}\\[7pt]%
  }%
  \end{tcolorbox}\par\smallskip%
}

% ============================================================
% activity-structure.tex
% 구조도 (트리 + 빈칸)
% 사용:
%   \begin{kfStructure}
%     \kfNode{0}{서론}{문제 제기}
%     \kfNode{1}{본론}{(\,\quad\,)}
%     \kfNode{1}{본론}{근거 2}
%     \kfNode{0}{결론}{주장 강화}
%   \end{kfStructure}
% ============================================================

\newenvironment{kfStructure}%
  {\par\smallskip\needspace{6\baselineskip}%
   \begin{tcolorbox}[
     enhanced, breakable=false,
     colback=kfPaper, colframe=kfRule,
     boxrule=0.4pt, arc=2pt,
     left=\kfBoxPad, right=\kfBoxPad, top=\kfBoxPad, bottom=\kfBoxPad,
     title={\footnotesize\bfseries\color{kfMute} 글의 구조},
     coltitle=kfMute, colbacktitle=kfPrimaryLight!40,
     attach boxed title to top left={yshift=-1.5mm, xshift=6pt},
     boxed title style={colframe=kfRule, boxrule=0.3pt, arc=1pt}
   ]%
   \begin{tabbing}
     \hspace{1.4cm}\=\hspace{1.4cm}\=\hspace{8cm}\kill}%
  {\end{tabbing}\end{tcolorbox}\par\smallskip}

% 노드: #1=들여쓰기 단계 (0~3), #2=라벨, #3=내용
\newcommand{\kfNode}[3]{%
  \ifcase#1\relax%
    \>\>\textbf{\color{kfPrimary} ▶ #2}\quad #3\\
  \or
    \>\>\quad\textbf{\color{kfMute} ◦ #2}\quad #3\\
  \or
    \>\>\quad\quad\textbf{\color{kfMute} - #2}\quad #3\\
  \or
    \>\>\quad\quad\quad\textbf{\color{kfMute} \textperiodcentered\ #2}\quad #3\\
  \fi
}

% 빈칸 노드
\newcommand{\kfNodeBlank}[2]{\kfNode{#1}{#2}{(\hspace{2.5cm})}}

% ============================================================
% activity-sentence-reading.tex
% 문장 독해 활동 — 문장 박스 1개 + 그에 묶인 문제 N개를 한 minipage로
% 사용:
%   \begin{kfSentenceUnit}{1}{"바람은 내 귀에 속삭이며..."}
%     \kfSentQ{(1)}{서술어 '속삭이며'의 주어는?}
%     \begin{kfChoices}
%       \kfChoice 바람은
%       ...
%     \end{kfChoices}
%   \end{kfSentenceUnit}
% ============================================================

% kfSentenceBox 매크로는 passage-note.tex에 정의됨

\newenvironment{kfSentenceUnit}[2]%
  {\par\smallskip\needspace{6\baselineskip}%
   \noindent\begin{minipage}{\linewidth}%
   \samepage%
   \kfSentenceBox{문장 #1}{#2}%
   \par\smallskip%
  }%
  {\end{minipage}\par\smallskip}

% 같은 문장에 묶인 소문항 (문장 안내 + 발문)
\newcommand{\kfSentQ}[2]{%
  \par\noindent\textbf{#1}~#2\par\smallskip%
}

% 헤더 없이 문장만 있는 묶음 (sentence_ref가 없는 일반 문항)
\newenvironment{kfSentencelessUnit}%
  {\par\smallskip\noindent\begin{minipage}{\linewidth}\samepage}%
  {\end{minipage}\par\smallskip}

% ============================================================
% chapter-cover.tex (v6)
% 챕터 진입 페이지 (표지) — 하단에 영역 목차 추가
% 사용:
%   \kfChapterCover{1}{근대성과 자아}{Chapter 1}{이번 챕터에서는 ...}
%   \begin{kfChapterAreaList}
%     \kfChapterAreaItem{어휘}{kfAreaVocab}{핵심 어휘 12개}
%     ...
%   \end{kfChapterAreaList}
%
%   영역 목차가 없을 때는 \kfChapterCoverEnd 호출
% ============================================================

\newcommand{\kfChapterCover}[4]{%
  \clearpage%
  \thispagestyle{kfBlank}%
  \kfMarkChapter{Chapter #1. #2}%
  \kfChapterCoverBg%
  % v18.5: 챕터 표지 우상단에 로고
  \begin{tikzpicture}[remember picture, overlay]
    \node[anchor=north east, inner sep=0pt] at ($(current page.north east)+(-18mm,-18mm)$) {\kfLogoMd};
  \end{tikzpicture}%
  \vspace*{20mm}%
  \noindent
  \begin{tikzpicture}[remember picture, overlay]
    % 좌측 액센트 바
    \fill[kfPrimary] (current page.west) ++(12mm,25mm) rectangle ++(5mm, -50mm);
    % 챕터 번호 배지 (이중 원, 약간 작게)
    \node[anchor=north west, inner sep=0pt] at ($(current page.north west)+(24mm,-45mm)$) {%
      \begin{tikzpicture}
        \draw[kfPrimary, line width=1pt] (0,0) circle (10mm);
        \draw[kfPrimary, line width=0.4pt] (0,0) circle (8mm);
        \node[font=\normalsize\bfseries\color{kfPrimary}] at (0,2mm) {CHAPTER};
        \node[font=\LARGE\bfseries\color{kfPrimary}] at (0,-3mm) {#1};
      \end{tikzpicture}%
    };
  \end{tikzpicture}%
  \vspace{42mm}%
  \par\noindent\hspace{32mm}%
  \begin{minipage}{0.65\linewidth}%
    {\footnotesize\color{kfMute} #3}\\[4pt]
    {\LARGE\bfseries\color{kfPrimary} #2}\\[8pt]
    {\color{kfRule}\rule{50mm}{0.5pt}}\\[6pt]
    {\small\color{kfInk} #4}
  \end{minipage}%
  \par\vspace{14mm}%
}

% v6 / v18.5 / v18.6: 챕터 표지의 영역 목차 — 흰색 반투명 박스
% v18.6: 배경 가로선/도형과 안 겹치도록 TikZ overlay로 우하단 절대 위치
\makeatletter
% 영역 항목들을 모아 둘 박스
\newsavebox{\kf@arealistbox}
\newenvironment{kfChapterAreaList}{%
  \begin{lrbox}{\kf@arealistbox}%
  \begin{minipage}{0.62\linewidth}%
    \begin{tcolorbox}[%
      enhanced, colback=white, colframe=white,
      opacityback=0.92, opacityframe=0,
      boxrule=0pt, arc=4pt,
      width=\linewidth,
      top=8pt, bottom=8pt, left=10pt, right=10pt,
    ]%
      {\footnotesize\bfseries\color{kfMute} 이 챕터의 영역}\par\vspace{4pt}%
      {\color{kfRule}\hrule height 0.3pt}\par\vspace{4pt}%
      \begin{itemize}[leftmargin=14pt, itemsep=2pt, topsep=0pt, label={\color{kfPrimary!70}\textbullet}]%
}{%
      \end{itemize}%
    \end{tcolorbox}%
  \end{minipage}%
  \end{lrbox}%
  % v18.7: 배경 상단 가로선 ~ 하단 가로선 사이의 중간 영역에 배치
  % 가로선이 내용을 가리지 않도록 박스 중심이 두 선 사이 중점에 위치
  \begin{tikzpicture}[remember picture, overlay]%
    \node[anchor=center, inner sep=0pt]
      at ($(current page.south)+(0mm,90mm)$)
      {\usebox{\kf@arealistbox}};%
  \end{tikzpicture}%
  \vfill%
  \noindent\hfill\kfSeriesBadge\hspace{12mm}%
  \clearpage%
}
\makeatother

% 영역 항목: \kfChapterAreaItem{영역명}{kfArea색}{부제}
\makeatletter
\newcommand{\kfChapterAreaItem}[3]{%
  \item {\small\bfseries\color{#2} #1}%
  \def\kf@cci@sub{#3}%
  \ifx\kf@cci@sub\@empty\else
    {\small\color{kfMute}\, \textemdash\, #3}%
  \fi
}
\makeatother

% 영역 목록 없이 cover만 (legacy)
\newcommand{\kfChapterCoverEnd}{%
  \vfill%
  \noindent\hfill\kfSeriesBadge\hspace{12mm}%
  \clearpage%
}

% ============================================================
% area-header.tex (v16)
% 영역 시작 표제 — 시리즈별 차별화 라벨 + 큰 영역명 + 부제
% PAGE_BREAK_POLICY.md §1.4
%
% v16 (디자인 고도화 Phase 1):
%   - 시리즈별 라벨 박스 추가:
%     소쉬르 = AREA START (영역색 채움 박스)
%     프레게 = SECTION OPEN (영역색 테두리)
%     러셀   = RUSSELL (영역색 라이트 박스)
%     비트   = WITTGENSTEIN (회색 라이트 박스)
%   - 라벨 위치: 영역명 위 좌측
% ============================================================

\makeatletter

% 시리즈 라벨 텍스트
\newcommand{\kf@serieslabel}{%
  \ifcase\kf@series\relax
    AREA START%
  \or
    AREA START%
  \or
    SECTION OPEN%
  \or
    RUSSELL%
  \or
    WITTGENSTEIN%
  \fi
}

% 시리즈 라벨 박스 스타일 (#1 = 영역색)
\newcommand{\kf@labelbox}[1]{%
  \ifcase\kf@series\relax
    % 0/1 = 소쉬르: 영역색 채움
    \tikz[baseline=-2pt]{\node[fill=#1, text=white,
      rounded corners=3pt, inner xsep=6pt, inner ysep=2pt,
      font=\scriptsize\bfseries]{\kf@serieslabel};}%
  \or
    % 1 = 소쉬르: 영역색 채움
    \tikz[baseline=-2pt]{\node[fill=#1, text=white,
      rounded corners=3pt, inner xsep=6pt, inner ysep=2pt,
      font=\scriptsize\bfseries]{\kf@serieslabel};}%
  \or
    % 2 = 프레게: 영역색 테두리
    \tikz[baseline=-2pt]{\node[draw=#1, line width=0.6pt, text=#1,
      rounded corners=3pt, inner xsep=6pt, inner ysep=2pt,
      font=\scriptsize\bfseries]{\kf@serieslabel};}%
  \or
    % 3 = 러셀: 영역색 라이트 박스
    \tikz[baseline=-2pt]{\node[fill=#1!12, text=#1,
      rounded corners=2pt, inner xsep=5pt, inner ysep=1.5pt,
      font=\scriptsize\bfseries]{\kf@serieslabel};}%
  \or
    % 4 = 비트: 회색 미니 박스
    \tikz[baseline=-2pt]{\node[fill=kfMute!10, text=kfMute,
      rounded corners=2pt, inner xsep=5pt, inner ysep=1.5pt,
      font=\scriptsize\bfseries]{\kf@serieslabel};}%
  \fi
}

% v17: 시리즈+영역별 Canva 배경 자동 매핑
% \kf@hasbg=1로 set되면 \kfAreaStart에서 라벨 박스 출력 생략
\def\kf@areaVocab{어휘}\def\kf@areaGrammar{문법}\def\kf@areaConcept{개념}
\def\kf@areaLit{문학}\def\kf@areaNonlit{비문학}
\newif\ifkf@bgon
% 파일 있으면 배경 적용 + boolean true. 없으면 nothing.
\newcommand{\kfBgSet}[1]{%
  \IfFileExists{#1.pdf}{\kfPageBg{#1}\kf@bgontrue}{}%
}
\newcommand{\kf@trybg}[1]{%
  % v18.4: 영역 배경 PDF 비활성화 — 큰 도형이 본문 영역을 침범해
  % "영역 헤더 후 빈 페이지" 문제를 일으킴. 라벨 박스로만 영역 표시.
  \kf@bgonfalse%
}

% Note: 러셀(rus_*)/비트(wit_*) 배경 PDF 미존재 시 \kfPageBg가 에러 → \IfFileExists로 가드
% (현재는 파일이 모두 존재한다고 가정. 부분 제공 시 cls 수정 필요)

\newcommand{\kfAreaStart}[3]{%
  \clearpage%
  \thispagestyle{fancy}%
  \kfMarkArea{#1}%
  \kf@trybg{#1}%
  \vspace*{1mm}%
  % 배경 없으면 라텍스 라벨 박스 출력 (배경 있으면 일러스트가 라벨 역할)
  \ifkf@bgon
    \vspace*{4mm}%
  \else
    \noindent\kf@labelbox{#2}\par\vspace{4pt}%
  \fi
  % 영역명 + 부제
  \noindent
  \begin{tikzpicture}
    \node[anchor=west, inner sep=0pt] (bar) at (0,0) {\color{#2}\rule{3pt}{22pt}};
    \node[anchor=base west, inner sep=0pt, xshift=8pt, yshift=2pt] (lbl) at (bar.east |- 0,0) {%
      {\LARGE\bfseries\color{#2} #1}%
    };
    \def\kf@areasub{#3}%
    \ifx\kf@areasub\@empty\else
      \node[anchor=base west, inner sep=0pt, xshift=8pt, font=\normalsize\color{kfMute}]
        at (lbl.base east) {\,\textperiodcentered\, #3};%
    \fi
  \end{tikzpicture}%
  \par\vspace{2pt}
  {\color{#2!70}\hrule height 1pt}%
  \par\vspace{1pt}%
  {\color{kfRule}\hrule height 0.3pt}%
  \par\vspace{4pt}%
  \needspace{6\baselineskip}\nopagebreak[4]%
  \global\kf@areaJustOpenedtrue% v18.5: 첫 컨텐츠 needspace 무시
}
\makeatother

% 시리즈/영역 단축 매크로
\newcommand{\kfStartVocab}[1]{\kfAreaStart{어휘}{kfAreaVocab}{#1}}
\newcommand{\kfStartGrammar}[1]{\kfAreaStart{문법}{kfAreaGrammar}{#1}}
\newcommand{\kfStartConcept}[1]{\kfAreaStart{개념}{kfAreaConcept}{#1}}
\newcommand{\kfStartLit}[1]{\kfAreaStart{문학}{kfAreaLiterature}{#1}}
\newcommand{\kfStartNonlit}[1]{\kfAreaStart{비문학}{kfAreaNonlit}{#1}}
\newcommand{\kfStartWeekly}[1]{\kfAreaStart{실력 확인}{kfAreaWeekly}{#1}}

% ============================================================
% section-header.tex
% 섹션 시작 라벨 헤더 — [지문] / [활동] / [문제] 라벨
% 좌측 액센트 바 + 라벨 + 부제 (영역 액센트 색)
%
% 사용:
%   \kfSecHeader{kfAreaLiterature}{지문}{빼앗긴 들에도 봄은 오는가 / 이상화}
%   \kfSecHeader{kfAreaConcept}{활동}{내용 확인}
%   \kfSecHeader{kfAreaNonlit}{문제}{}
%
% 헤더 직후 페이지 분할 금지: \needspace{4\baselineskip} + \nopagebreak
% ============================================================

\providecommand{\kfSecHeader}[3]{%
  \par\needspace{10\baselineskip}\filbreak%
  \vspace{3pt}%
  \noindent
  \begin{tikzpicture}
    \node[anchor=west, inner sep=0pt] (bar) at (0,0) {\color{#1}\rule{3pt}{14pt}};
    \node[anchor=west, inner sep=0pt, xshift=6pt, font=\normalsize\bfseries\color{#1}]
       (lbl) at (bar.east) {[\,#2\,]};
    \node[anchor=west, inner sep=0pt, xshift=4pt, font=\small\color{kfMute}]
       at (lbl.east) {#3};
  \end{tikzpicture}%
  \par\vspace{1pt}%
  {\color{#1!50}\hrule height 0.4pt}%
  \par\vspace{2pt}\nopagebreak%
}

% 영역별 단축 매크로
\newcommand{\kfSecPassageVocab}[1]{\kfSecHeader{kfAreaVocab}{지문}{#1}}
\newcommand{\kfSecPassageGrammar}[1]{\kfSecHeader{kfAreaGrammar}{지문}{#1}}
\newcommand{\kfSecPassageConcept}[1]{\kfSecHeader{kfAreaConcept}{지문}{#1}}
\newcommand{\kfSecPassageLit}[1]{\kfSecHeader{kfAreaLiterature}{지문}{#1}}
\newcommand{\kfSecPassageNonlit}[1]{\kfSecHeader{kfAreaNonlit}{지문}{#1}}
\newcommand{\kfSecPassageWeekly}[1]{\kfSecHeader{kfAreaWeekly}{지문}{#1}}

\newcommand{\kfSecActVocab}[1]{\kfSecHeader{kfAreaVocab}{활동}{#1}}
\newcommand{\kfSecActGrammar}[1]{\kfSecHeader{kfAreaGrammar}{활동}{#1}}
\newcommand{\kfSecActConcept}[1]{\kfSecHeader{kfAreaConcept}{활동}{#1}}
\newcommand{\kfSecActLit}[1]{\kfSecHeader{kfAreaLiterature}{활동}{#1}}
\newcommand{\kfSecActNonlit}[1]{\kfSecHeader{kfAreaNonlit}{활동}{#1}}
\newcommand{\kfSecActWeekly}[1]{\kfSecHeader{kfAreaWeekly}{활동}{#1}}

\newcommand{\kfSecQVocab}[1]{\kfSecHeader{kfAreaVocab}{문제}{#1}}
\newcommand{\kfSecQGrammar}[1]{\kfSecHeader{kfAreaGrammar}{문제}{#1}}
\newcommand{\kfSecQConcept}[1]{\kfSecHeader{kfAreaConcept}{문제}{#1}}
\newcommand{\kfSecQLit}[1]{\kfSecHeader{kfAreaLiterature}{문제}{#1}}
\newcommand{\kfSecQNonlit}[1]{\kfSecHeader{kfAreaNonlit}{문제}{#1}}
\newcommand{\kfSecQWeekly}[1]{\kfSecHeader{kfAreaWeekly}{문제}{#1}}


\begin{document}

\thispagestyle{empty}
\vspace*{40mm}
\begin{center}
{\Huge\bfseries\color{kfPrimary} 비트겐슈타인2}\\[10pt]
{\Large\color{kfMute} 학생용 교재 — 제 1 권}\\[20pt]
{\small\color{kfMute} (챕터 1 \textendash{} 4)}
\end{center}
\vfill
\hfill\kfSeriesBadge\hspace{15mm}
\clearpage
\kfChapterCover{1}{비트겐슈타인2 ch01}{Chapter 1}{이번 챕터: 문법 → 문학 5작품 → 비문학 5지문 → 패턴 워크북.}

\kfStartGrammar{이번 챕터 문법 영역 — 음운 / 음운의 체계 심화}


\kfSecHeader{kfAreaGrammar}{문제}{}

\begin{multicols}{2}\raggedcolumns\setlength{\columnsep}{12pt}
\kfPassageFull{문항 1 지문 (concept)}{%
국어의 단모음 체계는 혀의 높이, 혀의 전후 위치, 입술의 원순성이라는 세 가지 변별적 자질로 구성된다. 현대 국어의 단모음은 표준적으로 10개(ㅏ, ㅓ, ㅗ, ㅜ, ㅡ, ㅣ, ㅐ, ㅔ, ㅚ, ㅟ)를 인정하나, 실제 발음에서 'ㅚ'와 'ㅟ'를 이중 모음으로 발음하는 것이 허용되면서 8모음 체계도 널리 쓰인다. 한편 모음 조화는 양성 모음(ㅏ, ㅗ)과 음성 모음(ㅓ, ㅜ)이 각각 어울리는 현상인데, 중세 국어에서는 이 규칙이 매우 엄격했으나 현대 국어에서는 상당 부분 약화되었다. 특히 'ㅡ'는 양성도 음성도 아닌 중성 모음으로 모음 조화에서 중립적 역할을 한다. 용언의 활용에서 어간 모음이 양성이면 '-아/아서'가, 음성이면 '-어/어서'가 결합하는 것이 원칙이지만, '하다'에는 '-여'가 결합하고, 'ㅡ' 모음 어간은 음성 모음 어미와 결합하는 등 예외도 존재한다.
}
\begin{kfQBlock}{1}{}
\kfQStem{윗글을 바탕으로 <보기>의 ⓐ\textasciitilde{}ⓔ에 대해 탐구한 내용으로 적절하지 않은 것은?

<보기> ⓐ 잡아(잡- + -아)    ⓑ 써서(쓰- + -어서)    ⓒ 하여(하- + -여) ⓓ 놀아(놀- + -아)    ⓔ 담가(담그- + -아)}
\begin{kfChoices}
\kfChoice ⓐ: 어간 '잡-'의 모음 'ㅏ'가 양성 모음이므로 양성 계열 어미 '-아'가 결합한 것이다.
\kfChoice ⓑ: 어간 '쓰-'의 모음 'ㅡ'가 중성 모음이므로 음성 계열 어미 '-어서'와 결합한 것이다.
\kfChoice ⓒ: '하-'는 양성 모음 'ㅏ'를 포함하지만, 특수하게 '-여'가 결합하여 모음 조화의 예외에 해당한다.
\kfChoice ⓓ: 어간 '놀-'의 모음 'ㅗ'가 음성 모음이므로 '-아'가 결합한 것은 모음 조화에 어긋난다.
\kfChoice ⓔ: 어간 '담그-'에서 'ㅡ'가 탈락하고, 잔여 어간 모음 'ㅏ'에 따라 '-아'가 결합한 것이다.
\end{kfChoices}
\end{kfQBlock}

\kfPassageFull{문항 2 지문 (concept)}{%
국어의 자음 19개는 조음 위치(양순음·치조음·경구개음·연구개음·성문음)와 조음 방법(파열음·파찰음·마찰음·비음·유음), 그리고 발성 유형(예사소리·된소리·거센소리)의 세 기준으로 분류된다. 이 중 파열음은 9개(ㅂ, ㅃ, ㅍ, ㄷ, ㄸ, ㅌ, ㄱ, ㄲ, ㅋ)로 가장 많은 수를 차지한다. 파찰음(ㅈ, ㅉ, ㅊ)은 경구개 위치에서만 나타나고, 마찰음(ㅅ, ㅆ, ㅎ) 중 'ㅎ'은 성문 마찰음으로 다른 마찰음과 조음 위치가 다르다. 비음(ㅁ, ㄴ, ㅇ)은 각각 양순·치조·연구개 위치에 하나씩 분포하며, 유음 'ㄹ'은 치조음에 해당한다. 특히 연구개 비음 'ㅇ'은 초성에서 음가가 없고 종성에서만 비음으로 실현되는 독특한 특성을 지닌다.
}
\begin{kfQBlock}{2}{}
\kfQStem{윗글을 바탕으로 <보기>의 밑줄 친 자음들을 분석한 내용으로 적절한 것은?

<보기> 「공장에서 만든 쌀빵이 참 맛있다.」 이 문장의 첫 번째 음절부터 마지막 음절까지에 나타나는 자음을 분석해 보자.}
\begin{kfChoices}
\kfChoice '공'의 초성 'ㄱ'과 종성 'ㅇ'은 조음 위치가 같다.
\kfChoice '장'의 초성 'ㅈ'과 '쌀'의 초성 'ㅆ'은 조음 방법이 같다.
\kfChoice '만'의 초성 'ㅁ'과 '빵'의 초성 'ㅃ'은 발성 유형이 같다.
\kfChoice '참'의 초성 'ㅊ'과 '맛'의 초성 'ㅁ'은 조음 위치가 같다.
\kfChoice '쌀'의 초성 'ㅆ'과 '빵'의 초성 'ㅃ'은 조음 위치가 같다.
\end{kfChoices}
\end{kfQBlock}

\kfPassageFull{문항 3 지문 (concept)}{%
음운은 말의 뜻을 구별해 주는 소리의 최소 단위이다. 두 단어가 오직 하나의 음운만 다르고 나머지는 모두 같을 때, 이 두 단어를 '최소 대립쌍'이라 한다. 예컨대 '불'과 '풀'은 초성 'ㅂ'과 'ㅍ'만 다르므로 최소 대립쌍이다. 이 대립쌍을 통해 'ㅂ'과 'ㅍ'이 별개의 음운임을 증명할 수 있다. 'ㅂ'과 'ㅍ'은 조음 위치(양순)와 조음 방법(파열)이 같고, 오직 발성 유형(예사소리 대 거센소리)만 다르다. 이처럼 최소 대립쌍을 이루는 두 음운이 어떤 자질에서 대립하는지를 분석하면 해당 언어의 음운 체계를 체계적으로 기술할 수 있다.
}
\begin{kfQBlock}{3}{}
\kfQStem{윗글을 바탕으로 할 때, <보기>에서 최소 대립쌍을 이루는 것끼리 짝지은 것으로 적절하지 않은 것은?

<보기> ㄱ. 달 — 탈    ㄴ. 밤 — 밥    ㄷ. 굴 — 귤 ㄹ. 벌 — 벽    ㅁ. 살 — 쌀}
\begin{kfChoices}
\kfChoice ㄱ: '달'과 '탈'은 초성 'ㄷ/ㅌ'만 다르므로 최소 대립쌍이다.
\kfChoice ㄴ: '밤'과 '밥'은 종성 'ㅁ/ㅂ'만 다르므로 최소 대립쌍이다.
\kfChoice ㄷ: '굴'과 '귤'은 모음 'ㅜ/ㅠ'만 다르므로 최소 대립쌍이다.
\kfChoice ㄹ: '벌'과 '벽'은 종성 'ㄹ/ㄱ'만 다르므로 최소 대립쌍이다.
\kfChoice ㅁ: '살'과 '쌀'은 초성 'ㅅ/ㅆ'만 다르므로 최소 대립쌍이다.
\end{kfChoices}
\end{kfQBlock}

\kfPassageFull{문항 4 지문 (concept)}{%
국어의 음절은 '초성(자음) + 중성(모음) + 종성(자음)'의 구조를 가진다. 초성과 종성은 없을 수도 있지만, 중성(모음)은 반드시 있어야 한다. 따라서 가능한 음절 유형은 네 가지이다: ① 모음만(V) — '아', ② 자음+모음(CV) — '나', ③ 모음+자음(VC) — '얼', ④ 자음+모음+자음(CVC) — '달'. 국어에서 하나의 음절에 놓일 수 있는 자음은 초성에 최대 1개, 종성에 최대 1개이다. 겹받침은 표기상 두 자음이지만, 실제 발음에서는 음절의 끝소리 규칙에 의해 하나만 발음된다. 한편 분절 음운의 개수를 셀 때는 이중 모음을 반모음과 단모음의 결합으로 분석하므로, '왕'의 분절 음운은 'ㅇ + ㅗ + ㅏ + ㅇ'의 4개가 된다.
}
\begin{kfQBlock}{4}{}
\kfQStem{윗글을 바탕으로 <보기>의 단어들에 대해 분석한 내용으로 적절한 것은?

<보기> ㄱ. 의사    ㄴ. 원인    ㄷ. 귀환    ㄹ. 뒹굴}
\begin{kfChoices}
\kfChoice ㄱ의 '의'에서 중성은 이중 모음이므로 분절 음운의 수는 3개이다.
\kfChoice ㄴ의 '원'은 CVC 구조로 분절 음운의 수는 4개이다.
\kfChoice ㄷ의 '귀'는 분절 음운의 수가 '환'보다 적다.
\kfChoice ㄹ의 '뒹'은 음절 유형이 VC이다.
\kfChoice ㄴ의 '인'과 ㄹ의 '굴'은 음절 구조가 같고 분절 음운 수도 같다.
\end{kfChoices}
\end{kfQBlock}

\kfPassageFull{문항 5 지문 (concept)}{%
표준어의 단모음 체계는 시대에 따라 변화해 왔다. 중세 국어에는 'ㆍ(아래아)'를 포함하여 7개의 단모음이 있었다. 근대 국어를 거치며 'ㆍ'가 소멸하고, 'ㅐ, ㅔ'가 단모음화하면서 10모음 체계가 성립하였다. 현대 국어에서는 'ㅐ'와 'ㅔ'의 발음 구별이 약화되어 사실상 합류 현상이 나타나고 있으며, 'ㅚ'와 'ㅟ'를 이중 모음으로 발음하는 것이 일반화되면서 실질적으로 7\textasciitilde{}8모음 체계를 사용하는 화자가 많다. 그러나 표준 발음법에서는 'ㅚ, ㅟ'를 단모음으로 발음하는 것을 원칙으로 하면서 이중 모음 발음도 허용하고 있으므로, 규범적으로는 여전히 10단모음 체계를 유지하고 있다. 이처럼 음운 체계는 고정된 것이 아니라 시간에 따라 변화한다.
}
\begin{kfQBlock}{5}{}
\kfQStem{윗글을 바탕으로 할 때, <보기>에 대한 설명으로 적절하지 않은 것은?

<보기> (가) 현대 국어에서 많은 화자가 '개'와 '게'를 같은 발음으로 발음한다. (나) 'ㅚ'를 [외]로도, [웨]로도 발음할 수 있다. (다) 중세 국어의 '㗊'(아래아)은 현대 국어에서 주로 'ㅏ'나 'ㅡ'로 변화하였다.}
\begin{kfChoices}
\kfChoice (가)는 'ㅐ'와 'ㅔ'의 합류 현상을 보여 주는 예이다.
\kfChoice (가)의 현상은 10단모음 체계에서 단모음의 수가 줄어드는 방향의 변화이다.
\kfChoice (나)에서 [외]로 발음하면 단모음, [웨]로 발음하면 이중 모음으로 실현된 것이다.
\kfChoice (나)의 현상이 확산되면 단모음의 수가 늘어나 음운 체계가 복잡해진다.
\kfChoice (다)는 음운 체계가 통시적으로 변화할 수 있음을 보여 주는 예이다.
\end{kfChoices}
\end{kfQBlock}

\end{multicols}


\kfStartLit{이번 챕터 문학 작품 5편}


\kfSecHeader{kfAreaLiterature}{지문}{간 — 윤동주 (현대시) / 윤동주, 『하늘과 바람과 별과 시』(1948)}

\kfPassageNote{간 — 윤동주 (현대시) / 윤동주, 『하늘과 바람과 별과 시』(1948)}{%
바다가 변하야 뽕나무밭이 된다고 \\[4pt]
그 때에 백마 타고 오실 임은 \\[4pt]
어데로 가셨나뇨. \\[14pt]
조개껍데기에 이는 달빛에 \\[4pt]
귀머거리 소경이 되어서 \\[4pt]
이 개울 옆에 주저앉아. \\[14pt]
가을이면 가을이면 \\[4pt]
이 이랑 저 이랑에 \\[4pt]
밀과 콩을 심으리. \\[14pt]
사나이 거운 겨을 풍랑에 \\[4pt]
가슴을 앓고 가슴을 앓고 \\[4pt]
봄 여름 쓸쓸한 농사를 짓노라.
}
\kfSecHeader{kfAreaLiterature}{활동}{작품 정리표 — 간(윤동주)}

\noindent\textit{\small 빈칸에 알맞은 말을 채워 넣으시오. (초성 힌트 참고)}\par\medskip
\begin{kfActTable}{|>{\centering\arraybackslash}m{2.6cm}|m{6cm}|m{4cm}|}
\rowcolor{kfPrimaryLight!50}\textbf{\color{kfPrimary}항목} & \textbf{\color{kfPrimary}내용 (빈칸 채우기)} & \textbf{\color{kfPrimary}답안 작성}\\\hline
갈래 & ㅈㅇㅅ, ㅅㅈㅅ ① & \kfTBlank \\\hline
성격 & ㅅㅊㅈ, ㅅㅈㅈ ② & \kfTBlank \\\hline
제목 의미 & 사이, 간격 → 두 세계 사이의 ㄱㄱ ③ & \kfTBlank \\\hline
1연 고사성어 & 바다가 뽕밭으로 변함 = ㅅㅈㅂㅎ ④ & \kfTBlank \\\hline
임의 상태 & '어데로 가셨나뇨' → 임의 ㅂㅈ ⑤ & \kfTBlank \\\hline
화자의 상태 & ㄱㅁㄱㄹ ㅅㄱ이 되어 ㅈㅈㅇㅇ ⑥ & \kfTBlank \\\hline
대비 소재 1 & 불안정한 세계 = ㅂㄷ, ㅈㄱ ㄲㄷㄱ, ㅍㄹ ⑦ & \kfTBlank \\\hline
대비 소재 2 & 대지에 뿌리 내린 삶 = ㅃㄴㅁㅂ, ㅇㄹ, ㄴㅅ ⑧ & \kfTBlank \\\hline
반복법 & 'ㄱㅅㅇ ㅇㄱ ㄱㅅㅇ ㅇㄱ' → 고통의 지속 ⑨ & \kfTBlank \\\hline
시대 배경 & ㅇㅈ ㄱㅈㄱ 말기(1941년) ⑩ & \kfTBlank \\\hline
주제 1 & 어느 세계에도 속하지 못하는 ㅈㅈㄹㅈ ㅂㅇ ⑪ & \kfTBlank \\\hline
주제 2 & 자기 존재의 자리를 찾으려는 ㅈㅇ ㅌㅅ ⑫ & \kfTBlank \\\hline
\end{kfActTable}
\kfSecHeader{kfAreaLiterature}{문제}{}

\begin{multicols}{2}\raggedcolumns\setlength{\columnsep}{12pt}
\begin{kfQBlock}{1}{}
\kfQStem{윗시에 대한 설명으로 적절하지 않은 것은?}
\begin{kfChoices}
\kfChoice 전설적 시간과 현재 화자의 상황이 대비되어 있다.
\kfChoice '뽕나무밭'과 '바다'의 대비를 통해 세계의 변화를 나타내고 있다.
\kfChoice 화자는 '임'의 부재를 의문형으로 표현하여 그리움을 드러내고 있다.
\kfChoice 화자는 적극적으로 임을 찾아 나서며 의지적 태도를 보이고 있다.
\kfChoice 계절의 순환 구조를 통해 화자의 반복적인 삶을 드러내고 있다.
\end{kfChoices}
\end{kfQBlock}

\begin{kfQBlock}{2}{}
\kfQStem{윗시의 '조개껍데기에 이는 달빛'이 지닌 함축적 의미로 가장 적절한 것은?}
\begin{kfChoices}
\kfChoice 바다의 풍요와 자연의 아름다움을 찬양하는 신비로운 시어
\kfChoice 임이 돌아올 길을 환히 밝혀 주는 희망의 빛으로 작용
\kfChoice 바다가 물러난 흔적에 비치는 아련한 빛, 상실의 잔상
\kfChoice 화자의 경제적 빈곤과 궁핍한 처지를 보여 주는 소재
\kfChoice 화자가 자연에 완전히 동화되어 평안을 얻은 상태
\end{kfChoices}
\end{kfQBlock}

\begin{kfQBlock}{3}{}
\kfQStem{윗시의 마지막 연에서 '쓸쓸한 농사'가 의미하는 바로 가장 적절한 것은?}
\begin{kfChoices}
\kfChoice 임이 돌아올 때를 대비하여 풍성한 수확을 준비하는 행위
\kfChoice 바다를 떠나 육지에 정착하여 안정을 얻은 화자의 만족
\kfChoice 임의 부재와 존재론적 고독 속에서도 이어가야 하는 삶의 지속
\kfChoice 식민 치하에서 농민으로 위장하여 저항 활동을 하는 은유
\kfChoice 자연 속에서 문명을 거부하고 전원생활의 이상을 실현하는 행위
\end{kfChoices}
\end{kfQBlock}

\begin{kfQBlock}{4}{}
\kfQStem{윗시를 바탕으로 <보기>를 감상한 내용으로 적절하지 않은 것은?

<보기> 윤동주의 시에서 '바다'와 '하늘'은 화자가 도달하지 못하는 이상 세계를 상징하는 경우가 많다. 한편 화자는 현실의 '땅' 위에 머물면서 이상과 현실 사이의 괴리를 고뇌한다. 이 시가 쓰인 1941년은 태평양전쟁 직전으로, 조선인의 정체성이 극도로 위협받던 시기이다.}
\begin{kfChoices}
\kfChoice '바다가 변하야 뽕나무밭이 된다'는 이상 세계와 현실 세계의 극단적 거리를 암시하는 것으로 볼 수 있겠군.
\kfChoice 화자가 '주저앉아' 있는 것은 이상 세계에 도달하지 못한 채 현실에 머무는 상태를 보여주는 것이겠군.
\kfChoice '귀머거리 소경이 되어서'는 시대적 억압 속에서 무력해진 지식인의 처지를 형상화한 것으로 볼 수 있겠군.
\kfChoice 화자가 '밀과 콩을 심으리'라고 한 것은 식민 체제를 적극적으로 전복하려는 혁명 의지의 표출이겠군.
\kfChoice '가슴을 앓고'가 반복된 것은 존재론적 불안과 내면의 고통이 지속되고 있음을 보여주는 것이겠군.
\end{kfChoices}
\end{kfQBlock}

\begin{kfQBlock}{5}{}
\kfQStem{'가슴을 앓고 가슴을 앓고'에서 동일한 구절을 반복한 효과로 가장 적절한 것은?}
\begin{kfChoices}
\kfChoice 화자의 내면적 고통이 끊이지 않고 이어지고 있음을 강조한다.
\kfChoice 고통이 이미 사라진 뒤의 안도감과 평온을 부각하여 드러낸다.
\kfChoice 화자가 타인의 고통에 공감하여 함께 아파함을 보여 주고 있다.
\kfChoice 시적 긴장을 해소하여 평온한 결말로 나아가게 유도하고 있다.
\kfChoice 고통의 구체적 원인을 점차 명확하게 밝혀 가고 있음을 보인다.
\end{kfChoices}
\end{kfQBlock}

\end{multicols}

\kfSecHeader{kfAreaLiterature}{지문}{소년 — 윤동주 (현대시) / 윤동주, 『하늘과 바람과 별과 시』(1948)}

\kfPassageNote{소년 — 윤동주 (현대시) / 윤동주, 『하늘과 바람과 별과 시』(1948)}{%
여울 져셔 일어선 무지개. \\[14pt]
소년은 제비같이 나래를 펴고 \\[4pt]
푸른 하늘에 몸을 씻으며 \\[4pt]
어드메 안진 봄새의 노래를 들을까. \\[14pt]
그 무지개 아래에서 소년은 \\[4pt]
맨발로 무엇을 기다리고 있을까. \\[14pt]
저녁 놀이 몰려가면 소년도 따라가고 \\[4pt]
들에 고인 물 속에 \\[4pt]
쓸쓸한 그림자 하나. \\[14pt]
밤이 들면 한 귀퉁이에 웅크리고 \\[4pt]
어른이 될 때까지 울고만 있을까.
}
\kfSecHeader{kfAreaLiterature}{활동}{작품 정리표 — 소년(윤동주)}

\noindent\textit{\small 빈칸에 알맞은 말을 채워 넣으시오. (초성 힌트 참고)}\par\medskip
\begin{kfActTable}{|>{\centering\arraybackslash}m{2.6cm}|m{6cm}|m{4cm}|}
\rowcolor{kfPrimaryLight!50}\textbf{\color{kfPrimary}항목} & \textbf{\color{kfPrimary}내용 (빈칸 채우기)} & \textbf{\color{kfPrimary}답안 작성}\\\hline
갈래 & ㅈㅇㅅ, ㅅㅈㅅ ① & \kfTBlank \\\hline
성격 & ㅅㅊㅈ, ㅅㅈㅈ ② & \kfTBlank \\\hline
핵심 소재 1 & 이상과 희망의 이미지 = ㅁㅈㄱ ③ & \kfTBlank \\\hline
핵심 소재 2 & 고독과 불안의 이미지 = ㅆㅆㅎ ㄱㄹㅈ ④ & \kfTBlank \\\hline
소년의 상태(낮) & ㄴㄹ를 펴고 ㅍㄹ ㅎㄴ에 몸을 씻음 ⑤ & \kfTBlank \\\hline
소년의 상태(저녁) & 저녁 놀에 따라가고 ㅆㅆㅎ ㄱㄹㅈ 하나 ⑥ & \kfTBlank \\\hline
소년의 상태(밤) & ㅎ ㄱㅌㅇ에 ㅇㅋㄹ고 ㅇㄱ만 있을까 ⑦ & \kfTBlank \\\hline
무방비 상징 & 보호 없이 세계 앞에 노출 = ㅁㅂ ⑧ & \kfTBlank \\\hline
시간 구조 & ㄴ → ㅈㄴ → ㅂ으로 점점 어두워짐 ⑨ & \kfTBlank \\\hline
주제 1 & 순수한 자아가 성장 앞에서 겪는 ㅂㅇ ⑩ & \kfTBlank \\\hline
주제 2 & 성장의 문턱에서 느끼는 ㅁㅅㅇ과 ㄷㄹㅇ ⑪ & \kfTBlank \\\hline
출전 & 시집 『ㅎㄴ과 ㅂㄹ과 ㅂ과 ㅅ』 ⑫ & \kfTBlank \\\hline
\end{kfActTable}
\kfSecHeader{kfAreaLiterature}{문제}{}

\begin{multicols}{2}\raggedcolumns\setlength{\columnsep}{12pt}
\begin{kfQBlock}{1}{}
\kfQStem{윗시에 대한 설명으로 적절한 것은?}
\begin{kfChoices}
\kfChoice 소년의 행동이 구체적이고 역동적으로 서술되어 활기찬 분위기를 형성한다.
\kfChoice 시간의 흐름에 따라(낮→저녁→밤) 소년의 내면이 점점 어두워지는 구조를 보인다.
\kfChoice 화자가 소년에게 직접 말을 건네는 대화체를 통해 친밀감을 형성한다.
\kfChoice 자연물의 밝고 화려한 색채가 시 전체를 관통하여 낙관적 전망을 드러낸다.
\kfChoice 과거 회상과 현재의 교차를 통해 소년 시절에 대한 향수를 표현한다.
\end{kfChoices}
\end{kfQBlock}

\begin{kfQBlock}{2}{}
\kfQStem{윗시에서 '무지개'와 '쓸쓸한 그림자'의 관계를 가장 적절하게 설명한 것은?}
\begin{kfChoices}
\kfChoice 무지개는 소년의 이상과 희망을, 그림자는 그것이 사라진 뒤의 고독을 나타낸다.
\kfChoice 무지개와 그림자는 둘 다 소년의 밝은 미래를 예고하는 긍정적 이미지이다.
\kfChoice 무지개는 자연의 위협을, 그림자는 자연에 맞선 인간의 승리를 상징한다.
\kfChoice 무지개는 소년의 과거를, 그림자는 소년의 현재를 시간 순서로 구분한다.
\kfChoice 두 소재 모두 소년이 겪고 있는 물리적 위험의 정도를 함께 보여 준다.
\end{kfChoices}
\end{kfQBlock}

\begin{kfQBlock}{3}{}
\kfQStem{윗시의 '어른이 될 때까지 울고만 있을까'에 담긴 화자의 태도로 가장 적절한 것은?}
\begin{kfChoices}
\kfChoice 성장을 향한 확고한 결의와 낙관적 기대
\kfChoice 어른의 세계를 동경하는 기대와 들뜬 설렘
\kfChoice 성장 앞에 놓인 소년의 불안과 안타까움
\kfChoice 어른이 된 뒤 누릴 행복한 삶에 대한 확신
\kfChoice 소년 시절로 돌아가고픈 퇴행적 욕구의 표출
\end{kfChoices}
\end{kfQBlock}

\begin{kfQBlock}{4}{}
\kfQStem{윗시에서 '맨발로 무엇을 기다리고 있을까'의 '맨발'이 지닌 함축적 의미로 가장 적절한 것은?}
\begin{kfChoices}
\kfChoice 세상의 모진 시련을 견뎌 낸 소년의 강인함
\kfChoice 보호 없이 세계 앞에 노출된 무방비의 순수
\kfChoice 자연 속에서 자유를 누리는 해방감의 표상
\kfChoice 물질적 빈곤으로 인해 겪는 사회적 소외감
\kfChoice 성인 세계의 규범을 거부하는 반항적 태도
\end{kfChoices}
\end{kfQBlock}

\begin{kfQBlock}{5}{}
\kfQStem{윗시를 바탕으로 <보기>를 감상한 내용으로 적절하지 않은 것은?

<보기> 윤동주의 시에서 '소년'은 아직 세계에 의해 오염되지 않은 순수한 자아를 상징하는 경우가 많다. 이 순수한 자아는 성장 과정에서 세계의 폭력성·불확실성과 만나면서 위기에 처하게 된다. 윤동주는 이러한 순수의 위기를 반복적으로 형상화함으로써 식민지 시대 지식인의 내면적 갈등을 보여준다.}
\begin{kfChoices}
\kfChoice 소년이 '나래를 펴고 푸른 하늘에 몸을 씻'는 것은 순수한 자아의 이상적 상태를 형상화한 것이겠군.
\kfChoice '저녁 놀이 몰려가면 소년도 따라가고'에서 소년이 세계의 변화에 수동적으로 끌려가는 모습을 볼 수 있겠군.
\kfChoice 시의 후반부로 갈수록 소년의 상태가 어두워지는 것은 순수한 자아가 위기에 처한 과정을 보여주는 것이겠군.
\kfChoice '울고만 있을까'는 세계의 불확실성 앞에서 순수한 자아가 능동적으로 저항하는 의지를 나타내는 것이겠군.
\kfChoice 소년의 존재론적 불안은 <보기>에서 말한 식민지 시대 지식인의 내면적 갈등과 연결 지어 이해할 수 있겠군.
\end{kfChoices}
\end{kfQBlock}

\end{multicols}

\kfSecHeader{kfAreaLiterature}{지문}{꽃피는 시절 — 이성복 (현대시) / 이성복, 『남해금산』(1986), 수능·학평 기출}

\kfPassageNote{꽃피는 시절 — 이성복 (현대시) / 이성복, 『남해금산』(1986), 수능·학평 기출}{%
꽃이 피면 \\[4pt]
참 참 쓸쓸하다 \\[4pt]
핀 꽃 속에서도 \\[4pt]
피는 꽃 속에서도 \\[4pt]
필 꽃 속에서도 \\[4pt]
외로운 사람아 \\[4pt]
내 그대 잊지 않겠노라 \\[4pt]
길 모르는 그대가 \\[4pt]
꽃과 바람과 어둠 속에서 \\[4pt]
모르는 길을 가고 있으니 \\[4pt]
참 참 쓸쓸하다
}
\kfSecHeader{kfAreaLiterature}{활동}{작품 정리표 — 꽃피는 시절(이성복)}

\noindent\textit{\small 빈칸에 알맞은 말을 채워 넣으시오. (초성 힌트 참고)}\par\medskip
\begin{kfActTable}{|>{\centering\arraybackslash}m{2.6cm}|m{6cm}|m{4cm}|}
\rowcolor{kfPrimaryLight!50}\textbf{\color{kfPrimary}항목} & \textbf{\color{kfPrimary}내용 (빈칸 채우기)} & \textbf{\color{kfPrimary}답안 작성}\\\hline
갈래 & ㅈㅇㅅ, ㅅㅈㅅ ① & \kfTBlank \\\hline
출전 & 시집 『ㄴㅎ ㄱㅅ』(1986) ② & \kfTBlank \\\hline
핵심 역설 & ㄲ피는 아름다운 시절인데 ㅆㅆㅎ다 ③ & \kfTBlank \\\hline
시제 변화 & ㅍ 꽃(과거) / ㅍㄴ 꽃(현재) / ㅍ 꽃(미래) ④ & \kfTBlank \\\hline
시제 변화 효과 & 어떤 ㅅㄱ에도 외로움이 해소되지 않음 ⑤ & \kfTBlank \\\hline
화자의 다짐 & '내 ㄱㄷ ㅇㅈ ㅇㄱㄴㄹ' ⑥ & \kfTBlank \\\hline
그대의 상태 & ㄱ 모르는 그대가 모르는 길을 가고 있음 ⑦ & \kfTBlank \\\hline
반복법 & 'ㅊ ㅊ ㅆㅆㅎㄷ'의 반복 → 정서 강조 ⑧ & \kfTBlank \\\hline
표현 기법 & 밝은 제목 vs 어두운 내용 = ㅇㅅ ⑨ & \kfTBlank \\\hline
주제 1 & 아름다운 시절 속의 ㅈㅈㅇ ㅂㅇ ⑩ & \kfTBlank \\\hline
주제 2 & 자아와 세계의 근원적 ㅂㅎ ⑪ & \kfTBlank \\\hline
시인 특징 & 일상적 사물로 존재의 ㄱㅇ을 포착 ⑫ & \kfTBlank \\\hline
\end{kfActTable}
\kfSecHeader{kfAreaLiterature}{문제}{}

\begin{multicols}{2}\raggedcolumns\setlength{\columnsep}{12pt}
\begin{kfQBlock}{1}{}
\kfQStem{윗시에 대한 설명으로 적절하지 않은 것은?}
\begin{kfChoices}
\kfChoice '핀 꽃', '피는 꽃', '필 꽃'의 시제 변화를 통해 시간의 연속성을 드러내고 있다.
\kfChoice '참 참 쓸쓸하다'의 반복을 통해 화자의 정서를 강조하고 있다.
\kfChoice 꽃이 피는 아름다운 시절을 배경으로 하면서도 쓸쓸한 정서를 표현하는 역설적 구조를 보인다.
\kfChoice 화자는 '그대'에게 직접 말을 건네며 자신의 다짐을 표현하고 있다.
\kfChoice 꽃의 아름다움에 대한 감탄과 자연에 대한 찬미가 시의 중심 정서이다.
\end{kfChoices}
\end{kfQBlock}

\begin{kfQBlock}{2}{}
\kfQStem{윗시에서 '길 모르는 그대가 / 모르는 길을 가고 있으니'가 형상화하는 존재의 상태로 가장 적절한 것은?}
\begin{kfChoices}
\kfChoice 명확한 목표를 향해 묵묵히 나아가는 의지적 존재
\kfChoice 길을 잃었지만 곧 방향을 찾을 것이라는 희망
\kfChoice 방향도 목적지도 없이 세계 속을 헤매는 불안한 존재
\kfChoice 다른 사람들과 함께 길을 찾아가는 공동체적 존재
\kfChoice 과거의 길을 돌아보며 후회하는 존재
\end{kfChoices}
\end{kfQBlock}

\begin{kfQBlock}{3}{}
\kfQStem{'핀 꽃 속에서도 / 피는 꽃 속에서도 / 필 꽃 속에서도 / 외로운 사람아'에 담긴 의미로 가장 적절한 것은?}
\begin{kfChoices}
\kfChoice 꽃의 아름다움이 시간이 갈수록 증가하여 화자의 기쁨도 커진다.
\kfChoice 과거·현재·미래 어느 시간에도 외로움이 해소되지 않는다.
\kfChoice 꽃이 피고 지는 자연의 순환이 인간에게 위안을 준다.
\kfChoice 꽃의 상태 변화에 따라 외로움의 성격이 근본적으로 달라진다.
\kfChoice 꽃이 핀 뒤에야 비로소 외로움이 시작된다.
\end{kfChoices}
\end{kfQBlock}

\begin{kfQBlock}{4}{}
\kfQStem{윗시를 바탕으로 <보기>를 감상한 내용으로 적절한 것은?

<보기> 이성복의 시에서 자아와 세계는 화해하지 못한다. 세계가 아무리 아름다운 모습을 보여주어도, 자아는 그 아름다움 안에서 오히려 자신의 고독과 이질감을 더 강하게 의식한다. 이는 하이데거가 말한 '불안' — 일상적 세계의 의미 연관이 무너지고 존재 자체가 문제가 되는 근원적 기분 — 과 유사한 맥락에서 이해할 수 있다.}
\begin{kfChoices}
\kfChoice '꽃이 피면 참 참 쓸쓸하다'는 세계의 아름다움이 화자의 외로움을 치유하지 못함을 보여주겠군.
\kfChoice 화자의 쓸쓸함은 꽃이 지고 나면 자연스럽게 해소될 일시적인 감정으로 보아야겠군.
\kfChoice '외로운 사람아'의 '그대'는 화자 자신이 아니라 현실에서 성공을 거둔 사람을 가리키겠군.
\kfChoice 이 시의 화자는 세계의 아름다움에 온전히 동화되어 깊은 기쁨을 느끼고 있는 상태이겠군.
\kfChoice '꽃과 바람과 어둠'은 화자가 세계 전체를 이해하고 통제할 수 있음을 상징하는 소재이겠군.
\end{kfChoices}
\end{kfQBlock}

\begin{kfQBlock}{5}{}
\kfQStem{윗시에서 '내 그대 잊지 않겠노라'가 지닌 기능으로 가장 적절한 것은?}
\begin{kfChoices}
\kfChoice 외로운 존재에 대한 연대와 공감의 의지를 드러내어 시에 따뜻한 결을 더한다.
\kfChoice 꽃이 보여 준 아름다움을 영원히 기억하겠다는 자연 예찬의 다짐을 보인다.
\kfChoice 과거의 인연에 대한 원망과 비난의 마음을 우회적으로 표현한 구절이다.
\kfChoice 화자의 외로움이 마침내 완전히 해소되었음을 보여 주는 마무리 진술이다.
\kfChoice 사회적 연대를 통해 구조적 문제를 풀어내겠다는 실천적 의지를 표명한다.
\end{kfChoices}
\end{kfQBlock}

\end{multicols}

\kfSecHeader{kfAreaLiterature}{지문}{겨울 뜸부기 — 오정희 (현대소설) / 오정희, 「겨울 뜸부기」(1978), 수능특강 수록}

\kfPassageNote{겨울 뜸부기 — 오정희 (현대소설) / 오정희, 「겨울 뜸부기」(1978), 수능특강 수록}{%
나는 건너편 집 처마 밑의 고드름을 바라보았다. 그것은 차고 투명한 빛으로 빛나고 있었다. 닿을 수 없는 아름다운 것. 나는 그것을 원했고 원하는 것이 고통이라는 것을 알았다. 아버지는 어디에선가 술을 마시고 돌아왔고 어머니의 배는 날로 불러 갔다. 밤이면 어머니는 몸을 이리저리 뒤척이며 신음 소리를 냈다. 나는 그 소리가 무서웠다. 어둠 속에서 들려오는 신음은 마치 세상 끝에서 들려오는 소리 같았다.

학교에서 돌아오면 나는 개울가에 앉아 얼음장 밑으로 흐르는 물소리를 들었다. 그것은 아무것도 아닌 것 같으면서 동시에 모든 것인 것 같았다. 나는 그 소리를 들으면서 점점 추워졌고, 추위가 내 몸의 일부가 되어 가는 것을 느꼈다. 어쩌면 추위는 나의 본래 모습인지도 모른다고 생각했다.

뜸부기가 우는 소리를 들은 것은 그 무렵이었다. 겨울에 뜸부기라니. 아이들은 그 소리를 듣고 놀라 웃었지만 나는 웃을 수 없었다. 겨울에 우는 뜸부기, 그것은 나와 같은 존재라고 느꼈기 때문이다. 제때에 울지 못하고, 제자리에 있지 못하는 것. 나는 뜸부기 소리가 나는 쪽을 오래도록 바라보았다.
}
\kfSecHeader{kfAreaLiterature}{활동}{작품 정리표 — 겨울 뜸부기(오정희)}

\noindent\textit{\small 빈칸에 알맞은 말을 채워 넣으시오. (초성 힌트 참고)}\par\medskip
\begin{kfActTable}{|>{\centering\arraybackslash}m{2.6cm}|m{6cm}|m{4cm}|}
\rowcolor{kfPrimaryLight!50}\textbf{\color{kfPrimary}항목} & \textbf{\color{kfPrimary}내용 (빈칸 채우기)} & \textbf{\color{kfPrimary}답안 작성}\\\hline
갈래 & ㄷㅍ ㅅㅅ ① & \kfTBlank \\\hline
서술 시점 & 1ㅇㅊ ㅈㅇㄱ ㅅㅈ ② & \kfTBlank \\\hline
서술자 & 사춘기 ㅅㄴ '나' ③ & \kfTBlank \\\hline
배경 & ㄱㅇ의 ㅎㄹㅎ ㅅㅁ ㅁㅇ ④ & \kfTBlank \\\hline
고드름 상징 & ㄷㅇ 수 없는 ㅇㄹㄷㅇ 것 → 갈망과 고통 ⑤ & \kfTBlank \\\hline
추위 의미 & 'ㅊㅇ는 나의 ㅂㄹ ㅁㅅ' → 불안의 내면화 ⑥ & \kfTBlank \\\hline
뜸부기 상징 & 제때에 울지 못하는 ㄱㄱㅈ ㅈㅈ ⑦ & \kfTBlank \\\hline
서술 특징 & ㄱㄱㅈ ㅁㅅ로 내면 심리를 간접 전달 ⑧ & \kfTBlank \\\hline
제한적 이해 & 소녀가 세계의 의미를 ㄱㄱ으로 포착하되 ㄱㄴ으로 파악 못함 ⑨ & \kfTBlank \\\hline
주제 1 & ㅅㅈ의 문턱에서 겪는 ㅈㅈ ㅂㅇ ⑩ & \kfTBlank \\\hline
주제 2 & 세계에 온전히 속하지 못하는 ㄱㄷ ⑪ & \kfTBlank \\\hline
발표 시기 & 1978년 ⑫ & \kfTBlank \\\hline
\end{kfActTable}
\kfSecHeader{kfAreaLiterature}{문제}{}

\begin{multicols}{2}\raggedcolumns\setlength{\columnsep}{12pt}
\begin{kfQBlock}{1}{}
\kfQStem{윗글에 대한 설명으로 적절한 것은?}
\begin{kfChoices}
\kfChoice 1인칭 서술자의 감각적 경험으로 내면 심리를 간접 전달하고 있다.
\kfChoice 전지적 작가 시점에서 인물의 과거와 미래를 모두 폭넓게 조망하고 있다.
\kfChoice 인물 간 대화의 빈번한 교차로 갈등의 양상을 직접적으로 드러내고 있다.
\kfChoice 서술자가 사건의 인과 관계를 논리적 분석으로 차근차근 풀어내고 있다.
\kfChoice 빠른 장면 전환과 짧은 단락 구성으로 긴박한 서사적 긴장감을 조성한다.
\end{kfChoices}
\end{kfQBlock}

\begin{kfQBlock}{2}{}
\kfQStem{윗글에서 '고드름'이 지닌 상징적 의미로 가장 적절한 것은?}
\begin{kfChoices}
\kfChoice 닿을 수 없는 아름다움, 곧 갈망해도 도달 못 할 이상
\kfChoice 겨울의 매서운 위험을 경고하는 자연 그대로의 신호
\kfChoice 아버지의 냉담하고 무뚝뚝한 태도를 비유하는 소재
\kfChoice 어머니의 깊은 고통을 치유할 수 있는 희망의 상징
\kfChoice 소녀가 성장한 뒤 되돌아볼 어린 시절의 따뜻한 추억
\end{kfChoices}
\end{kfQBlock}

\begin{kfQBlock}{3}{}
\kfQStem{소녀가 겨울 뜸부기에 대해 '나와 같은 존재라고 느꼈다'고 한 까닭으로 가장 적절한 것은?}
\begin{kfChoices}
\kfChoice 뜸부기가 소녀와 같이 추운 환경에서 외로이 살아가고 있어서
\kfChoice 제때 울지 못하고 제자리에 있지 못하는 처지가 소녀와 닮아서
\kfChoice 뜸부기의 울음소리가 소녀의 가느다란 목소리와 매우 비슷해서
\kfChoice 뜸부기가 소녀에게 직접 다가와 따뜻하게 교감했기 때문에
\kfChoice 뜸부기를 직접 길러 보는 것이 소녀의 오래된 꿈이어서
\end{kfChoices}
\end{kfQBlock}

\begin{kfQBlock}{4}{}
\kfQStem{윗글을 바탕으로 <보기>를 감상한 내용으로 적절하지 않은 것은?

<보기> 오정희의 소설에서 사춘기 소녀는 세계를 감각적으로 경험하지만 그 의미를 개념적으로 이해하지 못한다. 이 간극이 독특한 서사적 긴장을 만든다. 소녀의 감각 경험은 단순한 배경 묘사가 아니라 존재론적 불안의 신체적 표현이다. 소녀는 하이데거가 말한 '세계 내 존재'로서 세계에 던져져 있지만, 그 세계에 온전히 속하지 못하는 경계적 존재이다.}
\begin{kfChoices}
\kfChoice '추위가 내 몸의 일부가 되어 가는 것을 느꼈다'는 존재론적 불안이 신체적 감각으로 체현되는 것으로 볼 수 있겠군.
\kfChoice 소녀가 고드름을 '닿을 수 없는 아름다운 것'으로 인식하는 것은 세계와의 간극을 감각적으로 경험하는 것이겠군.
\kfChoice 소녀가 어머니의 신음 소리를 '세상 끝에서 들려오는 소리'로 느끼는 것은 감각 경험과 의미 이해 사이의 간극을 보여주는군.
\kfChoice 소녀가 아이들처럼 뜸부기 소리에 웃지 못하는 것은 이미 세계에 완전히 적응한 성인의 태도를 보여주는군.
\kfChoice 겨울 뜸부기가 '제때에 울지 못하는' 존재라는 점은 소녀가 세계에 온전히 속하지 못하는 경계적 존재임을 상징하는군.
\end{kfChoices}
\end{kfQBlock}

\begin{kfQBlock}{5}{}
\kfQStem{윗글에서 '어쩌면 추위는 나의 본래 모습인지도 모른다'라는 서술의 효과로 가장 적절한 것은?}
\begin{kfChoices}
\kfChoice 추위라는 외부 환경이 소녀 존재와 동일시되어 불안의 내면화를 드러낸다.
\kfChoice 소녀가 추위에 완전히 적응하여 더 이상 어떠한 고통도 느끼지 않게 되었음을 보여준다.
\kfChoice 소녀의 체질이 추운 날씨에 유난히 강하다는 생물학적 특성을 객관적으로 설명한다.
\kfChoice 소녀가 차가운 자연 한가운데에서 비로소 평화와 안정을 찾았음을 나타낸다.
\kfChoice 소녀가 매서운 추위를 이겨내기 위한 구체적인 방법을 차분히 모색하고 있음을 드러낸다.
\end{kfChoices}
\end{kfQBlock}

\end{multicols}

\kfSecHeader{kfAreaLiterature}{지문}{이생규장전 — 김시습 (고전소설) / 김시습, 『금오신화』 중 「이생규장전」, 수능·학평 기출 다수}

\kfPassageNote{이생규장전 — 김시습 (고전소설) / 김시습, 『금오신화』 중 「이생규장전」, 수능·학평 기출 다수}{%
이생은 영기(靈氣)가 있는 것처럼 느꼈으나 혹 의심하여 물었다. "그대는 이미 죽은 사람이 아니오? 어찌 다시 인간 세상에 올 수 있었소?" 최랑이 눈물을 흘리며 말했다. "저는 진실로 죽은 사람이옵니다. 그러나 낭군과의 인연이 아직 다하지 않아 이렇게 다시 왔사옵니다. 비록 산 사람과 죽은 사람의 길이 다르오나, 제 정은 살아 있을 때와 다름이 없사옵니다." 이생은 비록 그 말을 듣고 두려움을 느꼈으나, 다시 함께하는 기쁨에 겨워 의심을 거두었다. 이로부터 둘은 예전처럼 함께 살며 시를 짓고 글을 읽으며 날을 보냈다.

수년이 지난 어느 날, 최랑이 말했다. "세 번 가을이 돌아왔으매, 인간 세상의 연분도 다한 듯하옵니다. 이제 가야 할 때가 되었나이다." 이생이 비통하게 울며 붙잡았으나, 최랑은 홀연히 사라져 종적을 알 수 없었다. 이생은 그 뒤로 세상의 모든 일에 뜻을 잃고, 마침내 병이 들어 세상을 떠났다.

뒷사람이 말하기를, "두 사람의 정이 생사를 넘어섰으니, 참으로 기이한 일이로다."
}
\kfSecHeader{kfAreaLiterature}{활동}{작품 정리표 — 이생규장전(김시습)}

\noindent\textit{\small 빈칸에 알맞은 말을 채워 넣으시오. (초성 힌트 참고)}\par\medskip
\begin{kfActTable}{|>{\centering\arraybackslash}m{2.6cm}|m{6cm}|m{4cm}|}
\rowcolor{kfPrimaryLight!50}\textbf{\color{kfPrimary}항목} & \textbf{\color{kfPrimary}내용 (빈칸 채우기)} & \textbf{\color{kfPrimary}답안 작성}\\\hline
갈래 & ㅎㅁ ㅅㅅ (전기소설) ① & \kfTBlank \\\hline
출전 & 김시습, 『ㄱㅇ ㅅㅎ』 ② & \kfTBlank \\\hline
주요 인물 & ㅇㅅ(남편)과 ㅊㄹ(아내) ③ & \kfTBlank \\\hline
비극적 사건 & ㅎㄱㅈ ㅊㅇ으로 최랑이 죽음 ④ & \kfTBlank \\\hline
초자연적 요소 & 최랑의 ㅎㅂ이 이생을 찾아와 함께 삶 ⑤ & \kfTBlank \\\hline
재회의 근거 & '낭군과의 ㅇㅇ이 아직 다하지 않아' ⑥ & \kfTBlank \\\hline
최랑의 고백 & 산 사람과 죽은 사람의 길이 다르나, ㅈ은 같다 ⑦ & \kfTBlank \\\hline
이별의 이유 & 'ㅅ ㅂ ㄱㅇ이 돌아왔으매 ㅇㅂ도 다함' ⑧ & \kfTBlank \\\hline
이생의 최후 & 세상의 모든 일에 ㄷ을 잃고 ㅂ이 들어 죽음 ⑨ & \kfTBlank \\\hline
서사 구조 & 만남-이별-ㅇㅎ-재회-ㅈㅂ ⑩ & \kfTBlank \\\hline
주제 1 & ㅈㅇ을 넘어선 ㅅㄹ ⑪ & \kfTBlank \\\hline
주제 2 & ㅈㅈ와 ㅂㅈ의 경계 ⑫ & \kfTBlank \\\hline
\end{kfActTable}
\kfSecHeader{kfAreaLiterature}{문제}{}

\begin{multicols}{2}\raggedcolumns\setlength{\columnsep}{12pt}
\begin{kfQBlock}{1}{}
\kfQStem{윗글에 대한 설명으로 적절하지 않은 것은?}
\begin{kfChoices}
\kfChoice 초자연적 존재(혼백)가 현실 세계에 등장하는 전기적(傳奇的) 요소가 나타난다.
\kfChoice 대화를 통해 인물의 내면 감정과 상황이 직접적으로 드러나고 있다.
\kfChoice 최랑의 재등장과 최종 이별이라는 서사적 반전이 나타난다.
\kfChoice 서술자가 사건에 직접 참여하여 1인칭 시점으로 이야기를 전달하고 있다.
\kfChoice 작품 말미에 후대의 평가('뒷사람이 말하기를')가 부기되어 있다.
\end{kfChoices}
\end{kfQBlock}

\begin{kfQBlock}{2}{}
\kfQStem{윗글에서 최랑이 '산 사람과 죽은 사람의 길이 다르오나, 제 정은 살아 있을 때와 다름이 없사옵니다'라고 한 말의 함축적 의미로 가장 적절한 것은?}
\begin{kfChoices}
\kfChoice 죽음이 사랑의 감정까지 소멸시키지는 못한다는 초월적 사랑의 표현
\kfChoice 죽은 뒤에도 산 사람과 같은 물리적 능력을 갖추고 있다는 자신감
\kfChoice 삶과 죽음의 구분이 전혀 무의미하다는 허무주의적 인식
\kfChoice 이생에 대한 원한을 표현하기 위해 일부러 한 거짓말
\kfChoice 유교적 예법에 따라 남편에게 복종하겠다는 의지
\end{kfChoices}
\end{kfQBlock}

\begin{kfQBlock}{3}{}
\kfQStem{이생이 최랑의 정체를 알면서도 '의심을 거두'고 함께 살기로 한 까닭으로 가장 적절한 것은?}
\begin{kfChoices}
\kfChoice 최랑의 혼백이 해를 끼치지 않을 것이라는 논리적 판단
\kfChoice 두려움보다 사랑하는 이와 재회한 기쁨이 더 컸기 때문
\kfChoice 최랑이 살아 돌아왔다고 완전히 확신했기 때문
\kfChoice 유교적 의무로서 아내를 돌봐야 한다는 책임감
\kfChoice 최랑의 혼백이 강제로 이생을 억류했기 때문
\end{kfChoices}
\end{kfQBlock}

\begin{kfQBlock}{4}{}
\kfQStem{윗글을 바탕으로 <보기>를 감상한 내용으로 적절하지 않은 것은?

<보기> 김시습의 「이생규장전」은 삶과 죽음의 경계를 사랑이 넘어서는 이야기이다. 이 작품은 존재와 부재의 문제를 다룬다. 최랑은 물리적으로 부재하지만(죽었지만), 사랑이라는 지향에 의해 존재한다. 이는 존재가 물리적 실체에만 의존하지 않음을 시사한다. 그러나 이 초월적 사랑은 영원하지 않으며, '인연이 다하면' 최랑은 떠나야 한다. 이별의 불가피성은 인간 존재의 유한성을 드러낸다.}
\begin{kfChoices}
\kfChoice 최랑이 혼백으로 이생을 찾아온 것은 사랑이라는 지향에 의해 물리적 부재를 넘어 존재할 수 있음을 보여주는군.
\kfChoice '인연이 다하지 않아'라는 최랑의 말은 사랑이 존재의 조건으로 작용하고 있음을 시사하는군.
\kfChoice '세 번 가을이 돌아왔으매' 떠나야 한다는 것은 초월적 사랑도 시간에 의해 제한됨을 보여주는군.
\kfChoice 이생이 최랑의 떠남 이후 세상을 떠난 것은 존재의 의미가 사랑하는 이와의 관계에 의존하고 있음을 암시하는군.
\kfChoice '뒷사람이 말하기를'의 평가는 이 사랑이 현실적으로 모든 부부가 따라야 할 규범임을 선언하는군.
\end{kfChoices}
\end{kfQBlock}

\begin{kfQBlock}{5}{}
\kfQStem{이생이 최랑이 떠난 뒤 '세상의 모든 일에 뜻을 잃고' 죽음에 이른 것이 보여주는 바로 가장 적절한 것은?}
\begin{kfChoices}
\kfChoice 사랑하는 이의 부재가 존재의 의미 전체를 붕괴시킬 수 있음
\kfChoice 유교적 도덕에 따라 아내를 따라 죽는 것이 의무임
\kfChoice 이생이 원래부터 세상에 관심이 없는 은둔형 인물이었음
\kfChoice 혼백과 함께 산 것이 신체적 건강을 해쳐 병이 난 것
\kfChoice 이생이 최랑의 죽음에 대한 복수를 포기하고 체념한 것
\end{kfChoices}
\end{kfQBlock}

\end{multicols}

\kfStartNonlit{이번 챕터 비문학 지문 5편}


\kfSecHeader{kfAreaNonlit}{지문}{하이데거의 존재와 시간 (인문)}

\kfPassageNote{하이데거의 존재와 시간 (인문)}{%
마르틴 하이데거는 서양 형이상학의 역사가 '존재 망각'의 역사라고 진단했다. 그에 따르면, 고대 그리스 이래 철학자들은 '존재자(das Seiende)'에 대해서는 끊임없이 탐구해 왔으나, 정작 '존재(das Sein)' 그 자체, 즉 존재자를 존재하게 하는 것이 무엇인지에 대해서는 묻지 않았다. 플라톤이 이데아를, 아리스토텔레스가 실체를, 데카르트가 주체를 철학의 중심에 놓았을 때, 이들은 모두 특정한 '존재자'를 탐구한 것이지 '존재' 자체의 의미를 해명한 것은 아니었다. 하이데거는 이처럼 존재의 의미에 대한 물음이 망각된 상태를 극복하고자 자신의 주저 『존재와 시간』(1927)을 집필했다.

하이데거가 존재의 의미를 탐구하기 위해 선택한 출발점은 '현존재(Dasein)'이다. 현존재란 자신의 존재를 문제 삼을 수 있는 유일한 존재자, 곧 인간을 가리킨다. 그런데 하이데거가 말하는 인간은 데카르트적 '사유하는 주체(cogito)'와 근본적으로 다르다. 데카르트는 세계와 분리된 고립된 주관에서 출발하여 외부 세계의 존재를 증명하려 했지만, 하이데거는 인간이 처음부터 세계 속에 던져져 있다고 보았다. 이것이 '세계-내-존재(In-der-Welt-sein)'라는 개념이다. 현존재는 세계 바깥에서 세계를 관조하는 주관이 아니라, 언제나 이미 세계 속에서 사물을 사용하고 타인과 관계 맺으며 살아가고 있는 존재이다.

현존재의 일상적인 존재 방식을 하이데거는 '비본래적 존재'라 부른다. 일상 속에서 현존재는 '세인(das Man)', 즉 익명적인 다수의 사람들이 하는 대로 따라 하며 살아간다. 세인이 읽는 것을 읽고, 세인이 판단하는 대로 판단하며, 세인이 분노하는 것에 분노한다. 이 상태에서 현존재는 자기 고유한 가능성을 잊고, 자기 자신으로부터 도피한다. 그러나 하이데거는 이 비본래적 존재를 도덕적으로 비난하지 않는다. 그것은 현존재의 일상적 존재 구조이며, 본래적 존재의 '결여 양태'일 뿐이다.

현존재를 비본래적 존재에서 벗어나게 하는 계기는 '불안(Angst)'이다. 불안은 특정 대상에 대한 두려움과 달리 대상이 없는 근원적 기분이다. 불안 속에서 일상적 세계의 의미 연관이 무너지고, 현존재는 자신이 아무런 근거 없이 세계에 던져져 있다는 사실과 마주한다. 이 '피투성(Geworfenheit)', 즉 던져져 있음의 경험을 통해 현존재는 비로소 자기 고유한 존재 가능성을 직시하게 된다.

하이데거에게 현존재의 가장 고유한 가능성은 '죽음을 향한 존재(Sein zum Tode)'이다. 죽음은 대체 불가능하고 회피할 수 없는 가장 극단적인 가능성이다. 세인의 세계에서 죽음은 '사람은 언젠가 죽는다'는 일반적 사실로 은폐되지만, 본래적 존재는 자신의 죽음을 고유한 것으로 '앞질러 달려감(Vorlaufen)'으로써 받아들인다. 이 선구적 결의를 통해 현존재는 유한한 시간 속에서 자기 자신의 고유한 가능성을 기투하며 살아가는 본래적 존재가 된다. 하이데거에게 존재의 의미는 바로 이 '시간성(Zeitlichkeit)'에서 드러나는 것이다.
}
\kfSecHeader{kfAreaNonlit}{문제}{}

\begin{multicols}{2}\raggedcolumns\setlength{\columnsep}{12pt}
\begin{kfQBlock}{1}{}
\kfQStem{윗글의 내용과 일치하지 않는 것은?}
\begin{kfChoices}
\kfChoice 하이데거는 서양 형이상학이 '존재자'를 탐구하면서 '존재' 자체의 의미를 망각했다고 보았다.
\kfChoice 현존재는 자신의 존재를 문제 삼을 수 있는 존재자로서 인간을 가리킨다.
\kfChoice 세인의 상태에서 현존재는 자기 고유한 가능성을 적극적으로 기투한다.
\kfChoice 불안은 특정 대상에 대한 두려움과 달리 대상이 없는 근원적 기분이다.
\kfChoice 죽음은 현존재의 가장 고유하고 대체 불가능한 가능성이다.
\end{kfChoices}
\end{kfQBlock}

\begin{kfQBlock}{2}{}
\kfQStem{윗글에 따를 때, 하이데거가 데카르트를 비판하는 근거로 가장 적절한 것은?}
\begin{kfChoices}
\kfChoice 데카르트는 존재 의미를 물었으나 충분히 답을 제시하지는 못하였다.
\kfChoice 데카르트는 고립된 주관에서 출발해 존재의 관계적 성격을 놓쳤다.
\kfChoice 데카르트는 인간의 유한성은 인정했으나 시간성 개념을 도입하지 못했다.
\kfChoice 데카르트는 비본래적 존재와 본래적 존재를 명료하게 구분하지 않았다.
\kfChoice 데카르트는 현존재가 지닌 불안 개념을 지나치게 좁게 정의하고 말았다.
\end{kfChoices}
\end{kfQBlock}

\begin{kfQBlock}{3}{}
\kfQStem{윗글을 바탕으로 <보기>를 이해한 내용으로 적절하지 않은 것은?

<보기> A 씨는 오랫동안 직장에서 상사의 지시에 따르고, 동료들이 하는 대로 행동하며 살아왔다. 어느 날 갑작스러운 건강 이상으로 자신의 죽음을 실감하게 된 A 씨는 극심한 불안에 빠졌다. 이후 A 씨는 이전과는 달리 자신이 진정으로 원하는 삶이 무엇인지를 진지하게 고민하기 시작했다.}
\begin{kfChoices}
\kfChoice A 씨가 상사의 지시에 따르고 동료들처럼 행동한 것은 '세인'의 방식으로 살아간 비본래적 존재에 해당하겠군.
\kfChoice A 씨가 겪은 극심한 불안은 특정 대상에 대한 두려움이 아니라 존재론적 '불안'으로 볼 수 있겠군.
\kfChoice A 씨가 자신의 죽음을 실감한 것은 '죽음을 향한 존재'로서 본래적 가능성을 직면한 것이겠군.
\kfChoice A 씨가 진정으로 원하는 삶을 고민하기 시작한 것은 자기 고유한 가능성을 기투하려는 시도로 볼 수 있겠군.
\kfChoice A 씨의 비본래적 존재는 도덕적으로 타락한 상태이므로, 건강 이상이라는 벌을 받은 것이겠군.
\end{kfChoices}
\end{kfQBlock}

\begin{kfQBlock}{4}{}
\kfQStem{윗글의 '앞질러 달려감(Vorlaufen)'이 의미하는 바로 가장 적절한 것은?}
\begin{kfChoices}
\kfChoice 죽음이라는 사건을 물리적으로 앞당겨서 직접 경험하는 행위
\kfChoice 타인의 죽음을 옆에서 목격함으로써 자기 죽음을 미리 이해하는 일
\kfChoice 죽음의 가능성을 고유한 것으로 받아들여 유한한 삶을 본래적으로 사는 것
\kfChoice 죽음에 대한 공포를 극복하려고 내세를 굳게 믿는 종교적 태도를 취함
\kfChoice 죽음을 누구에게나 닥치는 일반적 사실로 인정하며 일상에 충실하려는 태도
\end{kfChoices}
\end{kfQBlock}

\end{multicols}

\kfSecHeader{kfAreaNonlit}{지문}{데카르트 vs 하이데거 (인문)}

\kfPassageNote{데카르트 vs 하이데거 (인문)}{%
근대 철학의 아버지로 불리는 르네 데카르트는 확실한 지식의 토대를 찾기 위해 방법적 회의를 수행했다. 그는 감각 경험, 수학적 진리, 나아가 세계의 존재까지 모두 의심했지만, 의심하는 행위 자체를 수행하는 '나'의 존재만은 의심할 수 없음을 발견했다. 이것이 '나는 생각한다, 그러므로 나는 존재한다(Cogito ergo sum)'라는 유명한 명제이다. 여기서 '나'는 신체와 구분되는 순수한 사유 실체(res cogitans)로서, 물질적 세계(res extensa)와는 본질적으로 다른 영역에 속한다. 데카르트는 이 사유하는 자아를 철학의 절대적 출발점으로 삼았다.

데카르트의 이 접근은 이른바 주관-객관 이분법을 확립했다. 사유하는 주관은 세계라는 객관과 분리되어, 주관 쪽에서 객관을 '표상(repraesentatio)'한다. 인식이란 주관이 객관을 정확하게 표상하는 것이며, 인식론의 핵심 과제는 이 표상이 실제 세계와 일치하는지를 보증하는 일이 된다. 데카르트는 이 보증의 역할을 신(God)에게 부여했다. 완전한 존재인 신이 우리를 기만하지 않으므로, 명석판명하게 인식한 것은 참이라는 논증이다. 이 구도에서 세계는 주관 바깥에 놓인 대상이 되고, 인식은 주관이 대상을 정확히 모사하는 행위로 이해된다.

하이데거는 이러한 데카르트적 구도를 근본적으로 비판한다. 그에 따르면, 데카르트의 코기토는 이미 존재론적 전제를 무비판적으로 받아들이고 있다. '나는 생각한다'에서 '나'란 무엇인가? 데카르트는 이 '나'를 사유 실체, 즉 사물과 동일한 범주의 존재자로 규정함으로써 전통 형이상학의 실체 개념을 답습했다. 하이데거가 보기에 인간은 사물처럼 고정된 본질을 가진 실체가 아니라, 자기 존재를 끊임없이 문제 삼고 가능성을 향해 자신을 던지는 존재이다. 하이데거는 이러한 인간의 존재 방식을 '실존(Existenz)'이라 불렀다.

나아가 하이데거는 데카르트의 주관-객관 이분법 자체가 잘못된 출발점이라고 본다. 인간은 먼저 고립된 주관으로 존재한 뒤 나중에 세계와 관계를 맺는 것이 아니다. 인간은 언제나 이미 세계 속에서 사물을 사용하고 타인과 함께 살아가고 있다. 하이데거의 '세계-내-존재' 개념은 주관과 객관의 분리 이전에 인간이 이미 세계와 불가분하게 얽혀 있음을 드러낸다. 예컨대 우리가 망치를 사용할 때, 망치는 주관 앞에 놓인 물리적 대상이 아니라 '못을 박기 위한 도구'로서 이미 의미의 연관 속에 놓여 있다. 이처럼 사물은 고립된 대상이 아니라 세계의 의미 맥락 안에서 먼저 만나지는 것이다.

이 대비는 인식론적 전통에서 존재론적 전환으로의 이행을 보여준다. 데카르트에게 철학의 근본 물음은 '우리가 세계를 어떻게 정확히 인식하는가?'라는 인식론적 문제였다. 반면 하이데거에게 근본 물음은 '존재란 무엇인가?'라는 존재론적 문제이다. 인식의 가능성을 묻기 전에 인식하는 존재자, 즉 현존재의 존재 방식 자체를 해명해야 한다는 것이다. 하이데거의 전환은 인식론을 폐기하는 것이 아니라, 인식론이 전제하고 있던 존재론적 기반을 드러냄으로써 인식 문제를 보다 근원적인 차원에서 다시 묻는 것이다.
}
\kfSecHeader{kfAreaNonlit}{문제}{}

\begin{multicols}{2}\raggedcolumns\setlength{\columnsep}{12pt}
\begin{kfQBlock}{1}{}
\kfQStem{윗글의 내용과 일치하는 것은?}
\begin{kfChoices}
\kfChoice 데카르트는 감각 경험뿐 아니라 사유하는 자아의 존재까지 의심할 수 있다고 보았다.
\kfChoice 데카르트의 인식론에서 주관은 세계와 분리되어 객관을 표상하는 구도를 취한다.
\kfChoice 하이데거는 인간을 고정된 본질을 가진 사유 실체로 규정했다.
\kfChoice 하이데거에 따르면 인간은 먼저 고립된 주관으로 존재한 뒤 세계와 관계를 맺는다.
\kfChoice 하이데거의 존재론적 전환은 인식론 자체를 폐기하려는 시도이다.
\end{kfChoices}
\end{kfQBlock}

\begin{kfQBlock}{2}{}
\kfQStem{하이데거가 데카르트의 코기토를 비판하는 핵심 논거로 가장 적절한 것은?}
\begin{kfChoices}
\kfChoice 방법적 회의가 충분히 철저하게 수행되지 못했다는 점
\kfChoice 신의 존재 증명이 논리적으로 타당하게 성립하지 않는다는 점
\kfChoice '나'를 사유 실체로 규정해 전통 실체 개념을 답습한 점
\kfChoice 감각 경험의 확실성을 지나치게 낮은 수준으로 평가한 점
\kfChoice 수학적 진리에 대한 의심이 굳이 필요하지 않다고 본 점
\end{kfChoices}
\end{kfQBlock}

\begin{kfQBlock}{3}{}
\kfQStem{윗글의 '망치' 사례가 설명하려는 바로 가장 적절한 것은?}
\begin{kfChoices}
\kfChoice 사물은 주관 앞에 놓인 물리적 대상으로서 인식의 객체가 된다는 점
\kfChoice 사물의 본질은 엄밀한 과학적 분석을 거쳐서만 밝혀질 수 있다는 점
\kfChoice 사물은 고립된 대상이 아니라 세계의 의미 맥락 안에서 만나진다는 점
\kfChoice 주관이 외부 대상을 정확히 표상할 때 비로소 참된 인식이 성립한다.
\kfChoice 도구 사용은 인간을 동물과 구별 짓게 만드는 가장 핵심적인 기준이다.
\end{kfChoices}
\end{kfQBlock}

\begin{kfQBlock}{4}{}
\kfQStem{윗글을 통해 알 수 있는 데카르트와 하이데거의 공통점으로 가장 적절한 것은?}
\begin{kfChoices}
\kfChoice 철학의 근본 물음이 존재의 의미를 해명하는 것이라고 보았다.
\kfChoice 인간의 존재 방식을 철학적 탐구의 중심에 놓았다.
\kfChoice 주관과 객관의 이분법을 철학의 출발점으로 삼았다.
\kfChoice 신의 존재가 인식의 확실성을 보증한다고 주장했다.
\kfChoice 인간을 고정된 본질을 지닌 사유 실체로 규정했다.
\end{kfChoices}
\end{kfQBlock}

\end{multicols}

\kfSecHeader{kfAreaNonlit}{지문}{후설의 현상학적 환원 (인문)}

\kfPassageNote{후설의 현상학적 환원 (인문)}{%
에드문트 후설은 '엄밀한 학으로서의 철학'을 정립하고자 했다. 그는 19세기 말 자연과학의 성공에 자극받아 철학 역시 확실한 토대 위에 세워져야 한다고 보았다. 이 토대를 마련하기 위해 후설이 제시한 것이 현상학(Phanomenologie)이다. 현상학이란 '사태 자체로!'라는 표어가 보여주듯이, 선입견이나 이론적 가정을 배제하고 의식에 주어지는 현상 그 자체를 기술하려는 철학적 방법이다.

후설은 우리가 일상적으로 취하는 태도를 '자연적 태도(natürliche Einstellung)'라 부른다. 자연적 태도에서 우리는 외부 세계가 의식과 독립적으로 존재한다고 암묵적으로 전제한다. 내 앞의 책상이 실재한다, 다른 사람들이 존재한다, 물리 법칙이 객관적으로 성립한다 등의 믿음은 모두 자연적 태도에 속한다. 그런데 후설은 이 자연적 태도가 철학의 출발점이 될 수 없다고 보았다. 왜냐하면 그것은 검증되지 않은 전제 위에 서 있기 때문이다. 세계가 의식과 독립적으로 존재한다는 믿음 자체가 이미 하나의 이론적 가정이라는 것이다.

이 자연적 태도를 중지시키는 방법이 '판단 중지(에포케, Epoché)'이다. 에포케는 외부 세계의 존재에 대한 자연적 태도의 전제를 '괄호 안에 넣는' 것, 즉 그것을 긍정하지도 부정하지도 않고 잠정적으로 보류하는 것이다. 중요한 것은 에포케가 세계의 존재를 부정하는 회의주의가 아니라는 점이다. 후설은 세계가 존재하지 않는다고 주장하는 것이 아니라, 세계의 존재에 대한 판단을 일시적으로 중지함으로써 의식에 주어지는 현상만을 순수하게 기술할 수 있는 조건을 마련하려는 것이다.

에포케를 통해 드러나는 의식의 핵심 구조가 '지향성(Intentionalität)'이다. 후설에 따르면 의식은 항상 '무엇인가에 대한' 의식이다. 지각은 무엇인가를 지각하는 것이고, 기억은 무엇인가를 기억하는 것이며, 상상은 무엇인가를 상상하는 것이다. 의식은 텅 빈 그릇처럼 내용물을 담는 것이 아니라, 본래적으로 대상을 향해 있는 구조를 갖는다. 여기서 의식이 향하는 방식을 '노에시스(noesis)', 의식이 향하는 대상적 의미를 '노에마(noema)'라 한다. 예를 들어 하나의 나무를 지각할 때, 지각 작용(노에시스)은 매 순간 달라지지만(앞에서 보기, 뒤에서 보기, 만져 보기 등), 지각된 나무라는 의미(노에마)는 동일하게 유지된다.

나아가 후설은 개별적 사실을 넘어 '본질(Wesen)'을 직관할 수 있다고 보았다. 이것이 '본질 직관(Wesensschau)'이다. 우리가 빨간 사과, 빨간 장미, 빨간 깃발 등 여러 개별 사례를 경험할 때, 이들에서 공통적으로 직관되는 '빨강'이라는 본질이 있다. 후설에 따르면, 이 본질은 경험에서 귀납적으로 추상된 것이 아니라 의식의 직관 작용을 통해 직접 파악되는 것이다. 현상학적 환원의 최종 목표는 이처럼 에포케와 지향성 분석을 통해 의식에 주어지는 현상의 본질 구조를 해명하고, 이를 통해 모든 학문의 궁극적 토대를 확보하는 것이다.
}
\kfSecHeader{kfAreaNonlit}{문제}{}

\begin{multicols}{2}\raggedcolumns\setlength{\columnsep}{12pt}
\begin{kfQBlock}{1}{}
\kfQStem{윗글의 내용과 일치하지 않는 것은?}
\begin{kfChoices}
\kfChoice 후설은 선입견과 이론적 가정을 배제하고 현상 자체를 기술하려는 방법을 제시했다.
\kfChoice 자연적 태도에서 우리는 외부 세계가 의식과 독립적으로 존재한다고 전제한다.
\kfChoice 에포케는 외부 세계의 존재를 부정하는 회의주의적 태도이다.
\kfChoice 후설에 따르면, 의식은 항상 무엇인가에 대한 의식이다.
\kfChoice 본질 직관은 개별 사례에서 공통적으로 직관되는 본질을 파악하는 것이다.
\end{kfChoices}
\end{kfQBlock}

\begin{kfQBlock}{2}{}
\kfQStem{윗글에 따를 때, 후설이 '자연적 태도'를 철학의 출발점으로 삼을 수 없다고 본 까닭으로 가장 적절한 것은?}
\begin{kfChoices}
\kfChoice 자연적 태도가 과학적 방법론과 모순되기 때문이다.
\kfChoice 자연적 태도가 의식의 지향성을 설명하지 못하기 때문이다.
\kfChoice 자연적 태도가 세계의 독립적 존재라는 이론적 가정 위에 서 있기 때문이다.
\kfChoice 자연적 태도가 본질을 직관하는 능력을 가로막아 불가능하게 만들기 때문이다.
\kfChoice 자연적 태도가 노에시스와 노에마라는 의식 구조를 구분하지 못하기 때문이다.
\end{kfChoices}
\end{kfQBlock}

\begin{kfQBlock}{3}{}
\kfQStem{윗글을 바탕으로 <보기>를 이해한 내용으로 적절한 것은?

<보기> 한 미술관 관람객이 같은 조각품을 정면, 측면, 후면에서 번갈아 보고 있다. 시선의 위치와 각도에 따라 보이는 모습은 매번 달라지지만, 관람객은 그것이 '같은 하나의 조각품'이라는 사실을 알고 있다.}
\begin{kfChoices}
\kfChoice 정면에서 본 조각품만이 진정한 본질이고, 나머지는 왜곡된 현상이다.
\kfChoice 시선의 변화에 따라 달라지는 지각 작용은 노에마에, 동일하게 유지되는 조각품의 의미는 노에시스에 해당한다.
\kfChoice 매번 달라지는 지각 작용(노에시스)에도 불구하고 지각된 대상의 의미(노에마)는 동일하게 유지된다.
\kfChoice 관람객이 조각품의 존재를 전제하는 것은 에포케를 수행하는 것에 해당한다.
\kfChoice 조각품의 각 면은 자연적 태도에, 전체 조각품은 현상학적 태도에 대응한다.
\end{kfChoices}
\end{kfQBlock}

\begin{kfQBlock}{4}{}
\kfQStem{윗글에서 '괄호 안에 넣는다'는 비유가 의미하는 바로 가장 적절한 것은?}
\begin{kfChoices}
\kfChoice 외부 세계의 존재를 완전히 부정하고 의식만을 인정한다.
\kfChoice 외부 세계에 대한 긍정도 부정도 하지 않고 판단을 보류한다.
\kfChoice 외부 세계를 논리적 형식으로 환원하여 수학적으로 분석한다.
\kfChoice 외부 세계의 존재를 일단 인정한 뒤 개별 사실을 분류한다.
\kfChoice 의식의 내용을 외부 세계에 투사하여 객관적으로 확인한다.
\end{kfChoices}
\end{kfQBlock}

\end{multicols}

\kfSecHeader{kfAreaNonlit}{지문}{인격 동일성 논쟁 (인문)}

\kfPassageNote{인격 동일성 논쟁 (인문)}{%
시간이 흐르면서 신체와 정신이 끊임없이 변화하는데, 어제의 '나'와 오늘의 '나'가 같은 사람이라고 할 수 있는 근거는 무엇인가? 이것이 인격 동일성(personal identity) 문제이다. 이 물음은 단순한 호기심이 아니라, 법적 책임, 도덕적 귀속, 사후 생존 등의 실천적 문제와 직결된다. 10년 전에 범죄를 저지른 사람이 현재의 '나'와 같은 사람인지 여부는 처벌의 정당성을 좌우한다.

존 로크는 인격 동일성의 기준을 '기억'에서 찾았다. 그에 따르면, 한 사람이 과거의 특정 경험을 자신의 것으로 기억할 수 있다면, 그 경험을 한 존재와 현재의 자기는 동일한 인격이다. 로크가 영혼이나 신체가 아닌 기억을 기준으로 삼은 이유는 명확하다. 영혼은 관찰 불가능하고, 신체는 물질적으로 끊임없이 변화하기 때문이다. 반면 기억은 과거와 현재를 주관적으로 연결하는 의식의 끈이다. 그러나 로크의 기억 이론은 곧바로 난점에 부딪힌다. 토마스 리드의 유명한 반론이 대표적이다. 소년 시절에 매를 맞은 기억을 가진 장교가 있고, 노년의 장군이 장교 시절의 공적을 기억하지만 소년 시절의 매를 기억하지 못한다면, 장교는 소년과 같고 장군은 장교와 같지만 장군은 소년과 같지 않게 된다. 이는 동일성의 이행성(A=B이고 B=C이면 A=C)에 위배되는 모순이다.

데이비드 흄은 인격 동일성이라는 관념 자체를 해체했다. 그에 따르면, 우리가 내성(introspection)을 통해 발견하는 것은 끊임없이 변화하는 지각들의 다발일 뿐, 그 배후에 지속하는 '자아'라는 실체는 없다. 흄은 이를 '자아는 지각의 다발(bundle of perceptions)에 불과하다'라고 표현했다. 우리가 '나'라는 통일적 존재를 믿는 것은 유사성과 인과성이라는 심리적 습관 때문이다. 비슷한 지각이 반복되고 이전 지각이 이후 지각을 인과적으로 산출하는 것처럼 보이기 때문에, 우리는 그 배후에 동일한 자아가 있다는 환상을 갖게 된다는 것이다.

20세기 영국의 철학자 데릭 파핏은 로크의 기억 기준을 정교하게 발전시켜 '심리적 연속성(psychological continuity)' 이론을 제시했다. 파핏에 따르면, 인격 동일성의 기준은 단일한 기억이 아니라 기억, 성격, 의도, 믿음 등 심리적 상태들의 연속적 연쇄이다. A와 B 사이에 직접적 기억 연결이 없더라도, A→C→B처럼 중간 단계를 거쳐 심리적 연속성이 유지되면 동일 인격으로 볼 수 있다. 이것으로 리드의 반론이 해소된다.

파핏의 가장 도발적인 기여는 '인격 분열' 사고실험이다. 한 사람의 뇌를 반으로 나누어 두 사람의 몸에 각각 이식한다고 가정하자. 이식 후 두 사람 모두 원래 사람의 기억과 성격을 가지고 있다면, 원래 사람은 둘 중 누구와 동일한가? 파핏은 이 사고실험을 통해 인격 동일성이란 '전부 아니면 전무'의 관계가 아니라 정도의 차이가 있는 관계일 수 있음을 보였다. 나아가 파핏은 이러한 결론에서 윤리적 함의를 도출한다. 미래의 '나'가 현재의 '나'와 완전히 동일한 것이 아니라면, 현재의 이기적 관심을 미래 자아에게까지 강하게 투사하는 것은 합리적이지 않을 수 있다. 이는 역설적으로 타인에 대한 관심의 확장을 정당화하는 근거가 될 수 있다고 파핏은 주장했다.
}
\kfSecHeader{kfAreaNonlit}{문제}{}

\begin{multicols}{2}\raggedcolumns\setlength{\columnsep}{12pt}
\begin{kfQBlock}{1}{}
\kfQStem{윗글의 내용과 일치하는 것은?}
\begin{kfChoices}
\kfChoice 로크는 영혼의 동일성이 인격 동일성의 핵심 기준이라고 보았다.
\kfChoice 흄은 내성을 통해 지속하는 자아 실체를 발견할 수 있다고 주장했다.
\kfChoice 리드의 반론은 기억 기준이 동일성의 이행성에 위배될 수 있음을 보여 준다.
\kfChoice 파핏은 인격 동일성이 단일한 기억에 의해 결정된다고 보았다.
\kfChoice 파핏의 인격 분열 사고실험은 인격 동일성이 전부 아니면 전무의 관계임을 입증한다.
\end{kfChoices}
\end{kfQBlock}

\begin{kfQBlock}{2}{}
\kfQStem{윗글에 따를 때, 파핏이 '심리적 연속성' 개념을 도입한 이유로 가장 적절한 것은?}
\begin{kfChoices}
\kfChoice 흄의 자아 다발 이론을 계승하여 자아 실체를 부정하기 위해
\kfChoice 로크의 기억 기준에 대한 리드의 반론을 해소하기 위해
\kfChoice 인격 동일성 문제 자체가 허위 문제임을 증명하기 위해
\kfChoice 신체의 물질적 연속성이 인격 동일성의 기준임을 입증하기 위해
\kfChoice 법적 책임과 도덕적 귀속의 근거를 폐기하기 위해
\end{kfChoices}
\end{kfQBlock}

\begin{kfQBlock}{3}{}
\kfQStem{윗글을 바탕으로 <보기>를 이해한 내용으로 적절하지 않은 것은?

<보기> 갑은 사고로 기억을 모두 잃었지만, 사고 전의 성격, 습관, 가치관은 그대로 유지하고 있다. 을은 갑의 사고 이전 기억을 완벽하게 이식받았지만, 성격과 습관은 을 자신의 것을 유지하고 있다.}
\begin{kfChoices}
\kfChoice 로크의 관점에서 갑은 사고 이전의 갑과 동일한 인격이 아닐 수 있다.
\kfChoice 로크의 관점에서 을은 갑의 기억을 가지므로 사고 이전의 갑과 동일한 인격일 수 있다.
\kfChoice 흄의 관점에서 갑이 동일한 자아라는 믿음은 유사한 지각의 반복에 의한 환상일 수 있다.
\kfChoice 파핏의 관점에서 갑은 성격·습관·가치관이 연속적이므로 사고 이전의 갑과 심리적 연속성을 유지한다.
\kfChoice 파핏의 관점에서 을은 기억만 이식받았으므로 사고 이전의 갑과 완전한 동일성을 갖는다.
\end{kfChoices}
\end{kfQBlock}

\begin{kfQBlock}{4}{}
\kfQStem{파핏이 '인격 분열' 사고실험에서 도출한 윤리적 함의로 가장 적절한 것은?}
\begin{kfChoices}
\kfChoice 미래의 자아와 현재의 자아가 완전히 동일하므로 이기적 관심이 합리적이다.
\kfChoice 인격 동일성이 정도의 문제라면, 타인에 대한 관심의 확장도 정당화될 수 있다.
\kfChoice 인격이 분열될 수 있으므로 모든 도덕적 책임을 면제해야 한다.
\kfChoice 뇌 이식 기술이 발전하면 인격 동일성 문제가 완전히 해소된다.
\kfChoice 자아가 환상이므로 법적 처벌의 근거가 사라진다.
\end{kfChoices}
\end{kfQBlock}

\end{multicols}

\kfSecHeader{kfAreaNonlit}{지문}{메를로퐁티의 몸 현상학 (인문)}

\kfPassageNote{메를로퐁티의 몸 현상학 (인문)}{%
모리스 메를로퐁티는 서양 철학의 오랜 전통인 심신이원론을 극복하고자 했다. 데카르트 이래 서양 철학은 정신(사유 실체)과 신체(연장 실체)를 본질적으로 다른 두 영역으로 분리해 왔다. 이 이원론에서 신체는 정신의 도구이거나 기계적 메커니즘일 뿐, 인식이나 의미 부여의 주체가 아니다. 메를로퐁티는 이러한 이분법이 인간 경험의 실상을 왜곡한다고 비판하며, 몸이야말로 세계와 관계 맺는 근본적 방식이라고 주장했다.

메를로퐁티가 주목한 것은 '체험된 몸(le corps vécu)'이다. 이것은 생리학이 연구하는 객관적 신체도, 의학이 다루는 해부학적 신체도 아닌, 우리가 살아가면서 직접 경험하는 몸이다. 체험된 몸은 '신체도식(schéma corporel)'을 통해 작동한다. 신체도식이란 자신의 몸의 위치, 자세, 움직임 가능성에 대한 비주제적 앎, 즉 명시적으로 의식하지 않아도 작동하는 신체적 자기 인식이다. 예를 들어 우리는 문손잡이의 높이를 계산하지 않고도 자연스럽게 손을 뻗어 문을 연다. 이때 몸은 의식적 판단 없이 세계의 사물과 직접 '대화'하고 있다.

이러한 신체도식의 특성을 메를로퐁티는 '습관적 신체(le corps habituel)'라는 개념으로 발전시킨다. 습관적 신체란 반복된 경험을 통해 몸에 체화된 행동 양식이다. 피아노 연주자는 악보를 보면서 어떤 건반을 눌러야 할지 의식적으로 계산하지 않는다. 손가락이 '알아서' 건반 위를 움직인다. 이는 기억이 뇌에 표상으로 저장되어 있다가 인출되는 것이 아니라, 몸 자체가 과거 경험을 '운동적 의미'로 보유하고 있기 때문이다. 메를로퐁티는 이를 두고 '몸은 과거를 주제적으로 기억하는 것이 아니라 과거를 살고 있다'라고 표현했다.

메를로퐁티의 가장 핵심적인 주장은 몸이 단순한 객체가 아니라 '몸-주체(corps-sujet)'라는 것이다. 전통 철학에서 주체는 의식이었다. 의식이 세계를 인식하고, 판단하고, 의미를 부여하는 주체이며, 신체는 그 의식에 종속된 도구에 불과했다. 그러나 메를로퐁티는 지각 자체가 몸을 통해 이루어짐을 강조한다. 우리가 세계를 지각하는 것은 순수 의식이 아니라 보는 눈, 만지는 손, 걷는 다리를 가진 몸을 통해서이다. 몸은 세계를 수동적으로 반영하는 거울이 아니라, 세계에 능동적으로 참여하면서 의미를 만들어 내는 주체이다.

메를로퐁티의 몸 현상학은 데카르트적 이분법의 양극단, 즉 객관주의와 주관주의를 동시에 넘어서려는 시도이다. 객관주의는 세계를 주관과 무관한 물리적 사실의 총합으로 보고, 주관주의는 세계를 의식의 구성물로 본다. 메를로퐁티에게 세계는 이 둘 모두가 아니라, 체험된 몸이 참여하는 '지각의 장(champ perceptif)'이다. 세계는 몸 바깥에 놓인 객체도 아니고 의식이 투사한 표상도 아니며, 몸과 세계가 상호 침투하는 살아 있는 관계의 영역이다. 이처럼 메를로퐁티는 몸을 철학의 중심에 놓음으로써 인식론과 존재론을 새로운 지평에서 통합하고자 했다.
}
\kfSecHeader{kfAreaNonlit}{문제}{}

\begin{multicols}{2}\raggedcolumns\setlength{\columnsep}{12pt}
\begin{kfQBlock}{1}{}
\kfQStem{윗글의 내용과 일치하지 않는 것은?}
\begin{kfChoices}
\kfChoice 데카르트의 이원론에서 신체는 인식의 주체가 아니라 정신의 도구로 간주된다.
\kfChoice 신체도식은 자신의 몸에 대해 명시적으로 의식하지 않아도 작동하는 신체적 자기 인식이다.
\kfChoice 습관적 신체는 과거 경험이 뇌의 표상으로 저장되어 인출되는 메커니즘이다.
\kfChoice 메를로퐁티에게 몸은 세계에 능동적으로 참여하면서 의미를 만들어 내는 주체이다.
\kfChoice 메를로퐁티는 객관주의와 주관주의를 동시에 넘어서려 했다.
\end{kfChoices}
\end{kfQBlock}

\begin{kfQBlock}{2}{}
\kfQStem{윗글에서 '문손잡이' 사례를 통해 설명하려는 바로 가장 적절한 것은?}
\begin{kfChoices}
\kfChoice 의식이 신체에 명령을 내려 정확한 동작을 수행하도록 통제한다.
\kfChoice 신체도식이 의식의 판단 없이 몸과 세계의 직접적 대화를 가능하게 한다.
\kfChoice 문손잡이의 높이를 정확히 인식하려면 반복적 훈련 과정이 필요하다.
\kfChoice 일상적인 행동은 뇌가 수행하는 무의식적 계산에 의해 이루어진다.
\kfChoice 신체 감각은 의식보다 한발 먼저 외부 자극에 즉각적으로 반응한다.
\end{kfChoices}
\end{kfQBlock}

\begin{kfQBlock}{3}{}
\kfQStem{윗글의 '몸-주체(corps-sujet)' 개념이 기존 철학과 구별되는 점을 가장 적절하게 설명한 것은?}
\begin{kfChoices}
\kfChoice 전통 철학과 달리 의식이 아닌 몸 자체가 세계에 참여하며 의미를 만드는 주체가 된다.
\kfChoice 전통 철학에서 몸이 주체였다는 입장과 달리 의식만이 세계를 인식하는 진정한 주체이다.
\kfChoice 몸이 의식에 종속되지 않고 독립적인 물질적 실체로 존재함을 강조한다.
\kfChoice 몸과 의식이 완전히 분리된 두 주체로서 각각 세계와 관계를 맺는다.
\kfChoice 몸은 세계를 반영하는 수동적 거울이며 의식이 의미를 부여한다.
\end{kfChoices}
\end{kfQBlock}

\begin{kfQBlock}{4}{}
\kfQStem{윗글을 바탕으로 <보기>를 이해한 내용으로 적절하지 않은 것은?

<보기> 숙련된 자전거 선수가 울퉁불퉁한 산길을 고속으로 주행하고 있다. 선수는 매 순간 변하는 노면 상태에 즉각적으로 반응하여 몸의 균형을 유지하면서, 의식적으로는 코너를 돌아갈 전략만 생각하고 있다.}
\begin{kfChoices}
\kfChoice 노면에 대한 즉각적 반응은 신체도식이 의식적 판단 없이 작동하는 사례로 볼 수 있겠군.
\kfChoice 반복된 훈련을 통해 자전거 주행이 몸에 체화된 것은 습관적 신체의 형성으로 볼 수 있겠군.
\kfChoice 선수의 몸이 노면의 변화에 능동적으로 참여하는 것은 몸-주체 개념으로 이해할 수 있겠군.
\kfChoice 선수가 의식적으로 코너 전략만 생각하는 것은 몸과 의식이 완전히 분리되어 작동함을 보여주는군.
\kfChoice 선수의 균형 유지는 뇌에 저장된 표상의 인출이 아니라 몸이 운동적 의미를 보유한 것으로 설명할 수 있겠군.
\end{kfChoices}
\end{kfQBlock}

\end{multicols}

\kfStartWeekly{패턴 워크북 — 총 ?문항}


\kfSecHeader{kfAreaWeekly}{지문}{문항 1 지문 / 섞어치기}

\kfPassageNote{문항 1 지문 / 섞어치기}{%
식물은 동물과 달리 생애 동안 지속적으로 성장하는 개체이며, 이러한 생장은 분열 조직의 활동에 의해 유지된다. 식물의 분열 조직은 크게 측생 분열 조직과 정단 분열 조직으로 나뉘는데, 측생 분열 조직은 줄기와 뿌리의 측면에서 부피 생장을, 정단 분열 조직은 줄기와 뿌리 끝에서 길이 생장을 담당한다.

측생 분열 조직의 대표로는 관다발 형성층과 코르크 형성층이 있다. 관다발 형성층은 물관과 체관 사이에서 새로운 관다발 조직을 만들고, 코르크 형성층은 표피 아래에서 보호 조직을 형성해 수분 손실을 ⓐ막는 역할을 한다. 이들의 지속적 활동은 줄기와 뿌리를 굵고 단단하게 유지하며 나이테 형성과도 관련된다.

(중략)
}
\begin{kfQBlock}{1}{}
\kfQStem{<보기>의 ㄱ\textasciitilde{}ㅁ의 이유를 추론한 내용으로 적절하지 않은 것은?}
\begin{kfEvidence}<보기> ㄱ. 식물의 잎을 제거하면 식물이 정상적으로 자라지 못한다. ㄴ. 뿌리 분열대에 해당하는 부분을 잘라 내면 길이 생장이 일어나지 않는다. ㄷ. 당근을 신장대 자극 용액에 담가 두자 뿌리가 빠르게 길어졌다. ㄹ. 뿌리 표면 한 부분에 잉크 점을 찍고 시간이 지나자 뿌리는 길어졌지만 점 위치는 변하지 않았다. ㅁ. 나무줄기 횡단면에는 가로·세로 무늬 차이가 나타난다.\end{kfEvidence}
\begin{kfChoices}
\kfChoice ㄷ: 신장대는 세포 분열이 가장 활발한 구역이기 때문이겠군.
\kfChoice ㄱ: 슈트인 잎은 생장 유지에 필요한 생리 기능을 수행하기 때문이겠군.
\kfChoice ㄴ: 뿌리 분열대에 뿌리 정단 분열 조직이 위치하기 때문이겠군.
\end{kfChoices}
\end{kfQBlock}

\kfSecHeader{kfAreaWeekly}{지문}{문항 2 지문 / 부정상대어}

\kfPassageNote{문항 2 지문 / 부정상대어}{%
최근까지 일반적인 컴퓨터 환경에서는 운영 체제를 제어하거나 사용자 인터페이스를 처리하고, 문서를 작성하거나 웹을 탐색하는 등 다양한 범용 작업이 주요 연산 대상이었다. 이러한 작업은 명령어의 종류가 많고 그 순서가 수시로 바뀌며, 연산 간 종속성이 크다는 특징이 있다. 컴퓨터의 연산은 주로 여러 개의 코어로 구성된 ㉠중앙 처리 장치(CPU)에서 수행된다. 코어는 컴퓨터의 두뇌 역할을 하는 것으로 명령어를 해석하고 각종 연산을 실행한다. CPU는 여러 개의 고성능 코어로 구성되어 복잡한 제어 흐름을 정밀하게 조율할 수 있도록 설계되었다. 그러나 인공 지능 학습, 고해상도 영상 처리, 과학 기술 관련 계산 등 대규모 데이터를 대상으로 동일한 연산을 반복 수행해야 하는 작업이 증가하면서 기존의 연산 환경에 변화가 생겼다. CPU는 코어 수에 한계가 있고 개별 코어가 복잡한 제어 흐름에 맞추어 다양한 연산을 처리하도록 설계되어 있기 때문에 동일한 연산을 대량으로 반복 수행해야 하는 현대의 

(중략)
}
\begin{kfQBlock}{2}{}
\kfQStem{윗글에 대한 설명으로 가장 적절한 것은?}
\begin{kfChoices}
\kfChoice GPU 특성을 바탕으로 연산 환경 변화에 따른 GPU 활용 확대를 설명한다.
\kfChoice CPU와 GPU의 세부 장치 작동 원리를 대칭적으로 비교 분석하며 동등한 비중으로 다룬다.
\kfChoice 데이터 저장 방식의 역사적 발전 과정을 중심으로 통시적으로 설명한다.
\end{kfChoices}
\end{kfQBlock}

\kfSecHeader{kfAreaWeekly}{지문}{문항 3 지문 / 주체객체}

\kfPassageNote{문항 3 지문 / 주체객체}{%
프랑스의 기호학자이자 문예 비평가인 롤랑 바르트는 대중 매체의 이미지와 텍스트가 단순한 오락이나 소비 대상이 아니라 특정 사회·문화적 의미를 담고 있으며, 그 수용 과정에서 특정 이념이 대중에게 자연스럽게 받아들여진다고 보았다. 그는 이러한 현상을 ‘현대의 신화’라고 명명하고, 일상 속 기호에 숨어 있는 이념과 권력의 작동 방식을 분석하고자 하였다.

(중략)

바르트는 소쉬르의 이론을 분석 도구로 활용하였다. 소쉬르는 언어를 기호로 보았고, 기호는 음성이나 문자 등 감각적으로 지각 가능한 형태인 ‘기표’와 그것이 가리키는 개념적 내용인 ‘기의’라는 두 가지 요소로 이루어진다고 보았다. 예를 들어 ‘나무’라는 말은, 말소리 혹은 문자 형태로 주어지는 기표와 ‘줄기와 가지를 가진 목질 식물’이라는 기의가 결합한 하나의 기호이다. 소쉬르는 이러한 기호의 구조에서 기표와 기의 사이의 관계가 본질적으로 자의적이라는 점을 강조하였다. 즉 특정 기표가 반드시 특정 기의를 가리켜야 할 필연성은 없으며, 다만 

(중략)
}
\begin{kfQBlock}{3}{}
\kfQStem{‘1차 기호 체계’와 ‘2차 기호 체계’의 관계에 대한 설명으로 옳은 것은?}
\begin{kfChoices}
\kfChoice 2차 기호 체계에서는 1차 기호가 새로운 기의로만 작동한다.
\kfChoice 2차 기호 체계에서는 1차 기호 전체가 기표처럼 작동하고 여기에 새로운 기의가 덧붙어 의미가 형성된다.
\end{kfChoices}
\end{kfQBlock}

\kfSecHeader{kfAreaWeekly}{지문}{문항 4 지문 / 순서방향}

\kfPassageNote{문항 4 지문 / 순서방향}{%
일반적으로 혼인은 남자와 여자가 부부가 되는 일을 말한다. 과거에는 특정한 법적 절차 없이도 결혼식과 같은 사회적 관습에 의한 의식을 치르면 혼인이 성립한 것으로 보았다. 하지만 현재 대부분의 국가는 법에서 규정한 절차를 따라야만 혼인이 성립되는 법률혼 제도를 두고 있으며, 사회적 관습보다는 혼인 신고와 같은 법적 절차가 더 중요하다. 혼인은 사회를 구성하는 기본 단위인 가족을 형성하는 일이기 때문에 국가의 근간을 이루는 중요한 과정이고 이에 따라 법을 통해 관계를 보호하는 것이다. 따라서 현대 사회에서 혼인은 법적 절차에 의해 이루어지게 되고 여기에 참여한 남녀 간에는 법적 효력이 발생하게 된다.

(중략)

혼인이 성립하기 위해서는 법에서 규정하고 있는 실질적 요건과 형식적 요건을 모두 만족하여야 한다. 실질적 요건에서 가장 중요한 것은 혼인 의사의 합치이다. 혼인 의사는 부부 관계를 성립하겠다는 의사로 반드시 혼인 당사자 간의 합의가 필요하다. 또 다른 실질적 요건은 연령과 혼인 자격에 대한

(중략)
}
\begin{kfQBlock}{4}{}
\kfQStem{윗글을 읽고 알 수 있는 내용으로 적절하지 않은 것은?}
\begin{kfChoices}
\kfChoice 혼인 당사자들이 서로의 4촌 이내의 혈족 및 그 배우자와 인척 관계가 형성되는 것은 혼인 신고서를 관련 기관에 제출한 때를 기준으로 한다.
\kfChoice 현재 대부분의 국가에서는 사회적 관습에 따른 의식을 치르지 않더라도 법적 절차만으로 혼인이 성립될 수 있다.
\kfChoice 당사자 간의 혼인 의사 합치에 따라 관계가 형성되었더라도 혼인 당시 한쪽이 중혼인 경우에는 혼인 취소의 대상이 될 수 있다.
\end{kfChoices}
\end{kfQBlock}

\kfSecHeader{kfAreaWeekly}{지문}{문항 5 지문 / 인과도치}

\kfPassageNote{문항 5 지문 / 인과도치}{%
지구의 과거 자기를 연구하는 고지자기 연구는 1940년대 말 저명한 실험 물리학자인 블래킷의 가설에 의해 학계의 관심을 끌었다. 블래킷은 지구를 포함해서 모든 회전하는 물체는 그것의 각운동량에 비례하는 자기장을 주변에 형성한다고 주장하였다. 이 주장이 맞다면 지구의 자전 방향이 일정할 경우 지구 자기장의 방향은 바뀔 수가 없다. 지구 자기장의 방향이 바뀌었는지 여부는 암석에 잔류하는 자기, 즉 잔류 지자기의 미약한 흔적을 통해 확인할 수 있었다. 왜냐하면 과거에 용암이나 마그마가 굳어서 형성된 바위에는 용암이나 마그마가 식기 전에 그 안의 자성 물질이 지구 자기장의 방향을 따라 정렬된 채 남아 있기 때문이다. 측정 결과, 지자기의 방향이 바뀐다는 것이 드러나자 지구의 자전 방향이 바뀔 가능성이 없다는 판단에 따라 블래킷의 가설은 반박되었지만 잔류 지자기에 대한 연구자들의 관심은 지속되었다.

(중략)

잔류 지자기의 방향이 지질 시대의 다른 시점에서 조금씩 다르게 나타나는 현상을 설명하기 위해 대

(중략)
}
\begin{kfQBlock}{5}{}
\kfQStem{㉠에 대한 이해로 가장 적절한 것은?}
\begin{kfChoices}
\kfChoice 포타슘 붕괴율은 조건에 따라 크게 변해야 연대 계산이 가능하다.
\kfChoice 시료가 형성 후 외부 교란이 적을수록 포타슘-아르곤 비율 해석의 신뢰도가 높아진다.
\end{kfChoices}
\end{kfQBlock}

\kfSecHeader{kfAreaWeekly}{지문}{문항 6 지문 / 의도/목적}

\kfPassageNote{문항 6 지문 / 의도/목적}{%
프랑스의 기호학자이자 문예 비평가인 롤랑 바르트는 대중 매체의 이미지와 텍스트가 단순한 오락이나 소비 대상이 아니라 특정 사회·문화적 의미를 담고 있으며, 그 수용 과정에서 특정 이념이 대중에게 자연스럽게 받아들여진다고 보았다. 그는 이러한 현상을 ‘현대의 신화’라고 명명하고, 일상 속 기호에 숨어 있는 이념과 권력의 작동 방식을 분석하고자 하였다.

바르트는 소쉬르의 이론을 분석 도구로 활용하였다. 소쉬르는 언어를 기호로 보았고, 기호는 음성이나 문자 등 감각적으로 지각 가능한 형태인 ‘기표’와 그것이 가리키는 개념적 내용인 ‘기의’라는 두 가지 요소로 이루어진다고 보았다. 예를 들어 ‘나무’라는 말은, 말소리 혹은 문자 형태로 주어지는 기표와 ‘줄기와 가지를 가진 목질 식물’이라는 기의가 결합한 하나의 기호이다. 소쉬르는 이러한 기호의 구조에서 기표와 기의 사이의 관계가 본질적으로 자의적이라는 점을 강조하였다. 즉 특정 기표가 반드시 특정 기의를 가리켜야 할 필연성은 없으며, 다…(중략)
}
\begin{kfQBlock}{6}{}
\kfQStem{‘현대의 신화’에 대한 롤랑 바르트의 생각으로 가장 적절한 것은?}
\begin{kfChoices}
\kfChoice 사람들이 일상적으로 접할 수 있는 기호에 특정한 의미를 은밀히 덧붙여서 그것이 자연스러운 사실처럼 받아들여지게 만든다.
\kfChoice 사람들이 일상적으로 접할 수 있는 기호에 특정한 의미를 은밀히 덧붙여서 그것이 자연스러운 사실처럼 받아들여지게 만든다.
\kfChoice 콘텐츠에 담긴 사회적 맥락, 역사적 조건, 이념적 작용 등을 분석하고 해석할 수 있는 기준이 된다.
\kfChoice 현실의 모습을 있는 그대로 묘사하며 대중들에게 특정한 사회적·이념적 의미를 강압적으로 주입한다.
\end{kfChoices}
\end{kfQBlock}

\kfSecHeader{kfAreaWeekly}{지문}{문항 7 지문 / 상관관계 반전}

\kfPassageNote{문항 7 지문 / 상관관계 반전}{%
(가) 역사적으로 예술 형식을 구현하는 수단인 매체의 자리를 두고 경쟁한 문자와 이미지의 관계는 매체 연구의 주된 주제였다. 레싱은 시와 회화의 관계를 논하면서 당시 지배적이던 호라티우스의 경구 '시는 그림과 같이'를 대체할 새로운 관점을 제시했다. 호라티우스의 경구는 시와 회화의 공통 미메시스를 설명하는 말로 읽히기도 하고, 시가 회화를 모범으로 삼아야 한다는 말로 해석되기도 했다. 레싱은 회화의 원칙을 시에 확장 적용해 두 예술 형식의 경계를 모호하게 만드는 고전적 태도를 비판했다. 그는 각 예술 형식에서 매체와 매체 효과를 구분해 보아야 한다고 주장했다. 레싱에 따르면 시와 회화는 정보를 처리하는 방식이 달라 근본적 차이가 발생한다. 그는 미메시스라 통칭되던 공통 기능을 회화적 묘사에서는 표현으로, 서술적 묘사에서는 재현으로 나누어 매체 차별성을 강조했다. 레싱은 「라오콘 군상」을 사례로 들어, 조각은 고통을 표정으로 직접 제시하지만 시는 독자가 장면을 상상해 구성하게 만든다고…(중략)
}
\begin{kfQBlock}{7}{}
\kfQStem{㉠의 이유로 가장 적절하지 않은 것은?}
\begin{kfChoices}
\kfChoice 문자 예술이 이미지 예술보다 수용자의 상상력을 더 자극하기 때문에
\kfChoice 문자 예술이 이미지 예술보다 수용자의 상상력을 더 자극하기 때문에
\kfChoice 문자 예술은 이미지 예술보다 동시적 묘사에 탁월하기 때문에
\kfChoice 문자 예술의 정보 처리 방식이 이미지 예술에 적용되기 때문에
\end{kfChoices}
\end{kfQBlock}

\kfSecHeader{kfAreaWeekly}{지문}{문항 8 지문 / 필요조건}

\kfPassageNote{문항 8 지문 / 필요조건}{%
우리는 일반적으로 마음속 상태를 겉으로 분명하게 드러내거나 나타낼 때 ‘\uline{표현}’이라는 말을 사용한다. 그러나 표현이라는 말이 단지 자신의 심리적 상태를 드러내는 데에만 쓰이는 것은 아니다. 때로는 자연을 대하는 우리의 태도나 반응에도 표현이라는 말이 적용된다. 예를 들어 축 늘어진 나뭇가지를 보고 “애처롭다.”라고 말하는 경우가 있다. 이는 나뭇가지 자체가 감정을 지닌 것은 아니지만, 우리가 그것을 바라보며 내면의 정서를 \uline{투영}하고 외부로 드러내는 행위를 통해 표현이라는 말을 사용하게 되는 것이다. 하지만 표현이라는 말이 가장 널리 사용되는 분야는 다름 아닌 예술이며, 예술과 표현의 관계에 대해서는 오래전부터 다양한 견해가 논의되어 왔다.

첫째로, \uline{㉠표현}을 예술가가 창작 과정 속에서 행하는 내적 활동으로 보는 견해가 있다. 이 입장에서 표현은 예술가의 마음속에서 일어나는 정신적 과정, 즉 내적인 작용으로 이해된다. 이러한 입장을 대표하는 사상가가 …(중략)
}
\begin{kfQBlock}{8}{}
\kfQStem{윗글을 읽고 <보기>를 이해한 내용으로 적절하지 않은 것은?}
\begin{kfEvidence}• 화가 A는 폭풍우 치는 바다를 보며 두려움과 경외가 뒤섞인 내적 이미지를 형성했다. A는 이를 캔버스에 옮기기 위해 추상적인 선과 색채로 화면을 구성했다. • 작곡가 B는 불안한 심리를 점점 빨라지는 리듬과 불협화음의 반복으로 표현했다. 공연장에서 이 곡을 들은 감상자 C는 불안과 긴장을 알아보았다.\end{kfEvidence}
\begin{kfChoices}
\kfChoice 콜링우드는 화가 A의 캔버스에 구현된 추상적 선과 색채 그 자체를 예술의 표현으로 보겠군.
\kfChoice 콜링우드는 화가 A의 캔버스에 구현된 추상적 선과 색채 그 자체를 예술의 표현으로 보겠군.
\kfChoice 화가 A가 내적 이미지로 형성한 두려움과 경외를 적절한 선과 색채로 찾아 캔버스에 구현했더라도, 크로체는 이를 표현 그 자체가 아니라 내적 직관을 전달하기 위한 외적 수단으로 보겠군.
\kfChoice 작곡가 B의 곡을 들은 감상자 C가 불안과 긴장을 알아본 것에 대해, 톨스토이는 작곡가 B가 감정 전달에 성공했다고 보겠군.
\end{kfChoices}
\end{kfQBlock}

\kfSecHeader{kfAreaWeekly}{지문}{문항 9 지문 / 결합/분리}

\kfPassageNote{문항 9 지문 / 결합/분리}{%
일반적으로 [법적 의제]란 실제로는 사실이 아닌 어떤 것을 법적으로는 마치 사실인 것처럼 취급하고 반대되는 증거가 있다 하더라도 그와 같은 취급을 고수하려는 것을 의미한다. 가령 일정 기간 이상 생사가 확인되지 않는 사람에 대해 법원이 실종 선고를 내리면 그 사람은 사망한 것으로 간주되어, 법원이 다시 실종 선고를 취소하기 전에는 설령 그가 살아 있다 하더라도 법적으로 여전히 사망한 것으로 취급되는데, 이러한 결과가 바로 법적 의제에 따른 것이라 할 수 있다.

법적 의제의 활용에 대해서는 다양한 평가가 존재해 왔다. 벤담은 이를 가장 해롭고 비열한 거짓말이라 비판하는데, 그에 의하면 정의를 실현하는 일에 의제를 동원하는 것은 거래를 한다면서 사기나 치는 것과 같다. 반면에 블랙스톤은 법적 의제가 바람직한 목표의 제도적 달성에 어색하게나마 기여한다는 점에서 매우 유용한 수단이라며 옹호한 바 있다.

메인도 일견 법적 의제의 유용성을 인정하는 입장을 제시한다. 그는 정체된 사회와 진보…(중략)
}
\begin{kfQBlock}{9}{}
\kfQStem{‘메인’의 입장에 대한 설명으로 가장 적절한 것은?}
\begin{kfChoices}
\kfChoice 메인은 의제의 과거 유용성은 인정하되 성숙한 법체계에서의 의존 확대에는 반대했다.
\kfChoice 메인은 의제의 과거 유용성은 인정하되 성숙한 법체계에서의 의존 확대에는 반대했다.
\kfChoice 메인은 벤담의 비판을 그대로 수용해 의제를 전면 부정했다.
\kfChoice 메인은 의제의 역사적 유용성을 부정하고 블랙스톤만 비판했다.
\end{kfChoices}
\end{kfQBlock}

\kfSecHeader{kfAreaWeekly}{지문}{문항 10 지문 / 조건 왜곡}

\kfPassageNote{문항 10 지문 / 조건 왜곡}{%
현대 사회에서 디지털 금융 기술의 발달은 금융 거래의 속도와 편의성을 극적으로 향상시켰다. 모바일 뱅킹 등을 통해 몇 번의 터치만으로 송금이 이루어지는 시대가 도래했지만, 이러한 기술 발전은 예상치 못한 문제도 야기했다. 대표적인 사례가 바로 착오 송금이다. 착오 송금이란 송금인이 실수로 의도하지 않은 계좌에 돈을 보내는 경우를 말하며, 입력 오류, 수취인 착오 선택 등 다양한 원인에 의해 발생한다.

착오 송금이 발생했을 때, 송금인은 즉시 해당 은행에 이를 신고하고 반환을 요청하며, 은행은 중개자로서 수취인에게 연락하여 송금 실수임을 알리고 자발적 반환을 유도한다. 이때 수취인이 동의하면 간단하게 환급 절차가 이루어지지만, 반환을 거부하거나 연락이 닿지 않는 경우에는 상황이 복잡해진다. 송금인이 직접 법적 절차를 진행해야 하기 때문이다. 착오 송금인, 즉 실수로 돈을 잘못 보낸 사람은 그 실수를 바로잡기 위해 직접 회수를 시도하지만, 법적 대응에 들어가는 비용과 시간, 절차적 복…(중략)
}
\begin{kfQBlock}{10}{}
\kfQStem{윗글을 통해 알 수 있는 내용으로 적절한 것은?}
\begin{kfChoices}
\kfChoice 착오 송금 반환지원 제도에서는 회수 비용을 공제한 뒤 잔액을 반환할 수 있다.
\kfChoice 착오 송금 반환지원 제도에서는 회수 비용을 공제한 뒤 잔액을 반환할 수 있다.
\kfChoice 수취인이 반환을 거부하면 은행은 수취인 주소를 송금인에게 반드시 제공해야 한다.
\kfChoice 민사 소송은 착오 송금 금액이 5만 원 미만이면 법적으로 금지된다.
\end{kfChoices}
\end{kfQBlock}

\kfSecHeader{kfAreaWeekly}{지문}{문항 11 지문 / 긍부정 반전}

\kfPassageNote{문항 11 지문 / 긍부정 반전}{%
\#\#\# [22～27] 다음 글을 읽고 물음에 답하시오.

(가)   첩첩산중에도 없는 마을이 여긴 있습니다. 잎 진 사잇길 저 모랫둑, 그 너머 강기슭에서도 보이진 않습니다. 허방다리* 들어내면 보이는 마을.   갱 속 같은 마을. ㉠꼴깍, 해가, 노루꼬리 해가 지면 집집마다 봉당에 불을 켜지요. 콩깍지, 콩깍지처럼 후미진 외딴집, 외딴집에도 불빛은 앉아 이슥토록 창문은 모과빛입니다.   기인 밤입니다. 외딴집 노인은 홀로 잠이 깨어 출출한 나머지 무우를 깎기도 하고 고구마를 깎다, 문득 바람도 없는데 시나브로 풀려 풀려 내리는 짚단, 짚오라기의 설레임을 듣습니다. 귀를 모으고 듣지요. ㉡후루룩 후루룩 처마 깃에 나래 묻는 이름 모를 새, 새들의 온기를 생각합니다. 숨을 죽이고 생각하지요.   참 오래오래, 노인의 자리맡에 밭은기침 소리도 없을 양이면 벽 속에서 겨울 귀뚜라미는 울지요. 떼를 지어 웁니다, 벽이 무너지라고 웁니다.   어느덧 밖에는 눈발이라도 치는지, 펄펄 함박눈이라…(중략)
}
\begin{kfQBlock}{11}{}
\kfQStem{<보기>를 바탕으로 (가), (다)를 이해한 내용으로 가장 적절한 것은? [3점]}
\begin{kfChoices}
\kfChoice (다)의 ‘파도’와 ‘깊은 물’은 바다의 형상이라는 유사성으로 관계를 맺으며 물에 사는 사람이 살면서 만나게 되는 환경이라는 의미를 생성하고 있군.
\kfChoice (다)의 ‘파도’와 ‘깊은 물’은 바다의 형상이라는 유사성으로 관계를 맺으며 물에 사는 사람이 살면서 만나게 되는 환경이라는 의미를 생성하고 있군.
\kfChoice (가)의 ‘허방다리 들어내면 보이는 마을’, ‘갱 속 같은 마을’은 얕음과 깊음의 대비를 이루어 숨어 있는 두 공간의 차이를 부각하고 있군.
\kfChoice (가)의 ‘무우’와 ‘고구마’는 차가움과 따뜻함의 대비를 이루어 밤에 출출함을 달래기 위해 먹는 다양한 음식의 속성을 부각하고 있군.
\end{kfChoices}
\end{kfQBlock}

\kfSecHeader{kfAreaWeekly}{지문}{문항 12 지문 / 비교대상 뒤바꿈}

\kfPassageNote{문항 12 지문 / 비교대상 뒤바꿈}{%
[22 \textasciitilde{} 26] 다음 글을 읽고 물음에 답하시오.

(가)   목숨이란 마치 깨어진 배 조각   여기저기 흩어져 마을이 구죽죽한 어촌보담 어설프고   삶의 티끌만 오래 묵은 포범(布帆)처럼 달아 매었다

남들은 기뻤다는 젊은 날이었건만   밤마다 내 꿈은 서해를 밀항하는 쩡크*와 같아   소금에 절고 조수(潮水)에 부풀어 올랐다

항상 흐렷한 밤 암초를 벗어나면 태풍과 싸워 가고   전설에 읽어 본 산호도(珊瑚島)는 구경도 못 하는   그곳은 남십자성이 비쳐 주도 않았다

쫓기는 마음 지친 몸이길래   그리운 지평선을 한숨에 기오르면   시궁치*는 열대 식물처럼 발목을 오여쌌다

새벽 밀물에 밀려온 거미이냐   다 삭아 빠진 소라 껍질에 나는 붙어 왔다   먼 항구의 노정(路程)*에 흘러간 생활을 들여다보며

* 이육사, ｢노정기｣ -   * 쩡크: 정크(Junk). 중국 연해나 하천에서 사람과 짐을 실어 나르는 배.   * 시궁치: 더러운 물이 잘 빠지지 않고 썩어서 질척질척하게…(중략)
}
\begin{kfQBlock}{12}{}
\kfQStem{(가) \textasciitilde{} (다)에 대한 설명으로 가장 적절한 것은?}
\begin{kfChoices}
\kfChoice (나)와 (다) 모두 설의적 표현을 활용하여 의미를 부각하고 있다.
\kfChoice (나)와 (다) 모두 설의적 표현을 활용하여 의미를 부각하고 있다.
\kfChoice (가)와 (나) 모두 청유형 어미를 활용하여 친근감을 드러내고 있다.
\kfChoice (가)와 (다) 모두 반어적 표현을 활용하여 현실을 비판하고 있다.
\end{kfChoices}
\end{kfQBlock}

\kfSecHeader{kfAreaWeekly}{지문}{문항 13 지문 / 주체/객체 혼동}

\kfPassageNote{문항 13 지문 / 주체/객체 혼동}{%
\#\#\# [39\textasciitilde{}42] 다음 글을 읽고 물음에 답하시오.

(가)   苦忘亂抽書 잊기를 자주 하여 어지러이 뽑아 놓은 책들   散漫還復整 흩어진 걸 다시 또 정리하자니   曜靈忽西頹 해는 문득 서쪽으로 기울고   江光搖林影 강 위에 숲 그림자 흔들린다.   扶筇下中庭 막대 짚고 마당 가운데 내려서서   矯首望雲嶺 고개 들어 구름 낀 고개 바라보니   漠漠炊烟生 아득히 밥 짓는 연기가 피어나고   蕭蕭原野冷 쓸쓸히 들판은 서늘하구나.   田家近秋穫 농삿집 가을걷이 가까워지니   喜色動臼井 절구질 우물가에 기쁜 빛 돌아   鴉還天機熟 갈까마귀 돌아오니 절기가 무르익고   鷺立風標逈 해오라기 서 있는 모습 우뚝하고 훤하다.   我生獨何爲 내 인생은 홀로 무얼 하는 것인지   宿願久相梗 숙원이 오래도록 풀리질 않네.   無人語此懷 이 회포 털어놓을 사람 아무도 없어   搖琴彈夜靜 거문고만 둥둥 탄다, 고요한 밤에.

* 이황, ｢만보(晩步)｣ -

(나)   밤이다. ㉠ 하늘은 푸르다 못해 …(중략)
}
\begin{kfQBlock}{13}{}
\kfQStem{[A]와 [B]에 대한 설명으로 가장 적절한 것은?}
\begin{kfChoices}
\kfChoice [B]는 [A]와 달리 자연물에 대한 변화된 인식을 제시하고 있다.
\kfChoice [B]는 [A]와 달리 자연물에 대한 변화된 인식을 제시하고 있다.
\kfChoice [A]는 [B]와 달리 자연물에 감정을 이입하여 심리적 변화를 우회적으로 드러내고 있다.
\kfChoice [A]와 [B]는 모두 계절감을 드러내는 자연물을 통해 결실에 대한 기쁨을 나타내고 있다.
\end{kfChoices}
\end{kfQBlock}

\kfSecHeader{kfAreaWeekly}{지문}{문항 14 지문 / 내용과 방법 혼동}

\kfPassageNote{문항 14 지문 / 내용과 방법 혼동}{%
[27 \textasciitilde{} 30] 다음 글을 읽고 물음에 답하시오.

2424 혹은 5454번의 전화번호를 보디에 커다랗게 써 붙인 삼륜차 또는 픽업이 대충 비슷비슷한 내용물들을 실은 채 속속들이 닿고 있었고, 감색 유니폼의 관리인들이 요소요소마다 늘어선 채 똑같은 말들을 외쳐 대고 있었다. 일테면,   “차는 현관 옆으로 바짝 붙여 주십시오!”   “호실 키는 임시 관리 사무소에서 입주증과 교환해 드리고 있습니다. 관리 사무소는 217동과 219동 사이에 위치하고 있습니다…….”   “계단이 혼잡하오니 도착순대로 짐을 올리시고, 화장실 및 주방의 부착물은 248동과 249동 간에 위치하고 있는…….”   삼륜차 위에서 나는 한동안 멍청하게 흔들리고만 있었다. 수백 수천의 똑같은 5층짜리 콘크리트 건물군과 그리고 그 협곡 사이사이마다 출렁이고 있는 입주자들의 행렬……. 그것은 실로 기이한 대조였다. 나는 무거운 압박감과 마음 붙일 곳 없는 황량함을 동시에 의식하지 않을 수 없었다. 차가움, 견고함,…(중략)
}
\begin{kfQBlock}{14}{}
\kfQStem{[A]와 [B]의 서술상 특징으로 가장 적절한 것은?}
\begin{kfChoices}
\kfChoice [A]는 장면에 대한 관찰을 중심으로, [B]는 인물의 복잡한 내면을 중심으로 서술하고 있다.
\kfChoice [A]는 장면에 대한 관찰을 중심으로, [B]는 인물의 복잡한 내면을 중심으로 서술하고 있다.
\kfChoice [A]는 사건 해결의 실마리를 중심으로, [B]는 인물의 행위에 담긴 의미를 중심으로 서술하고 있다.
\kfChoice [A]는 인물들 간에 심화되는 갈등을 중심으로, [B]는 인물이 겪는 내적 갈등을 중심으로 서술하고 있다.
\end{kfChoices}
\end{kfQBlock}

\kfSecHeader{kfAreaWeekly}{지문}{문항 15 지문 / 내용 왜곡}

\kfPassageNote{문항 15 지문 / 내용 왜곡}{%
[26\textasciitilde{}31] 다음 글을 읽고 물음에 답하시오.

(가)   십 년 종사 후에 고향으로 도라오니   산천 의구하되 인사(人事)는 달라졌구나   아마도 세간의 존멸을 못내 슬허 하노라 <제1수>

산화(山花)는 믈의 픠고 물새는 산의 운다   일신이 한가하야 산수간의 누어시니   세상의 어즈러은 긔별을 나는 몰라 하노라 <제4수>

거믄고 빗기 들고 산수를 희롱하니   청풍은 건듯 블고 명월도 도라온다   하물며 유신(有信)한 갈매기는 오명 가명 하나니 <제5수>

거믄고 흥진(興盡)커던 조대(釣臺)로 내려가니   도화 뜬 말근 믈 뛰노나니 고기로다   아이야 밋기 다지 마라 취적(取適)*이나 하오리라 <제7수>

* 신교, ｢귀산음(歸山吟)｣ -   * 취적 : 낚시질의 참뜻이 세상 생각을 잊고자 하는 데 있음.

(나)   백수(白首)에 산수 구경 늦은 줄 알지마는   평생 품은 뜻을 이루고야 말리라 여겨   병자년 봄에 봄옷을 새로 입고   죽장망혜(竹杖芒鞋)로 노계 깊은 골에 …(중략)
}
\begin{kfQBlock}{15}{}
\kfQStem{(가)에 대한 이해로 적절하지 않은 것은?}
\begin{kfChoices}
\kfChoice <제7수>에서는 말을 건네는 방식을 사용하여 상대와의 동질감을 표현한다.
\kfChoice <제7수>에서는 말을 건네는 방식을 사용하여 상대와의 동질감을 표현한다.
\kfChoice <제1수>에서는 영탄적 표현을 통해 화자의 정서를 드러낸다.
\kfChoice <제4수>에서는 대구의 방식을 활용하여 시적 상황을 표현한다.
\end{kfChoices}
\end{kfQBlock}

\kfSecHeader{kfAreaWeekly}{지문}{문항 16 지문 / 과잉 해석}

\kfPassageNote{문항 16 지문 / 과잉 해석}{%
\# 
}
\begin{kfQBlock}{16}{}
\kfQStem{<보기>를 참고하여 (가), (나)를 감상한 내용으로 적절하지 않은 것은? [3점]}
\begin{kfChoices}
\kfChoice (가)의 ‘항산’이 있어야 ‘선심’이 나온다는 것에서 생업을 도덕적 삶의 실천을 위한 전제로 보는, (나)의 ‘체를 만드는 것과 글을 짓는 것’은 ‘재료’와 ‘소득’의 ‘차이만 있’다는 것에서 문학을 생업을 비판하기 위한 도구로 보는 사대부의 현실 인식을 엿볼 수 있군.
\kfChoice (가)의 ‘항산’이 있어야 ‘선심’이 나온다는 것에서 생업을 도덕적 삶의 실천을 위한 전제로 보는, (나)의 ‘체를 만드는 것과 글을 짓는 것’은 ‘재료’와 ‘소득’의 ‘차이만 있’다는 것에서 문학을 생업을 비판하기 위한 도구로 보는 사대부의 현실 인식을 엿볼 수 있군.
\kfChoice (가)의 ‘인명’은 ‘하늘이 만’드니 ‘천민이 되어 나서 본업’을 해야 한다는 것에서 운명론적 관점으로 생업을 바라보는 사대부의 모습을 엿볼 수 있군.
\kfChoice (나)의 ‘벼슬한 자들’이 ‘스스로’에 대한 평가와는 달리 ‘남을 잡스럽다고’ 한다는 것에서 당대의 현실을 부조리한 사회로 바라보는 사대부의 관점을 엿볼 수 있군.
\end{kfChoices}
\end{kfQBlock}

\kfSecHeader{kfAreaWeekly}{지문}{문항 17 지문 / 범위 오류}

\kfPassageNote{문항 17 지문 / 범위 오류}{%
\#\#\# [25\textasciitilde{}27] 다음 글을 읽고 물음에 답하시오.

[앞부분의 줄거리] 영웅적 면모를 지녔으나 먹고 자기만 하며 허송세월을 보내던 이생은 장인이 죽자 장모에 의해 쫓겨나고, 이생의 부인 양 소저는 시부모의 제사를 지내면서 이생을 기다린다. 한편 집을 나온 이생은 곤경에 빠진 사람에게 자신이 가진 삼백 냥의 돈을 다 내어 준 후 어떤 노인의 도움을 받게 된다.

“그대 오늘 큰 적선을 하였으니 깊이 감격하노라.”   이생이 이렇듯 노인의 신기함을 보고 평범한 인물이 아니리라 생각하며 의아해 마지않았다.   “존옹의 물으심이 무슨 일이니이꼬? 저는 적선한 일이 없소이다.”   노인 말하기를,   “대인은 사람 속이기를 아니하나니라.”   이에 또 말하기를,   “그대 저리 먹는 양에 양식도 없이 어찌 살아가려 하느뇨?”   이생이 답하기를,   “이처럼 얻어먹으면 아니 살아가리이까?”   노인이 크게 웃으며 말하였다.   “㉠ 소년의 말이 사리에 어둡도다. 그러나 나는 그대를 알거…(중략)
}
\begin{kfQBlock}{17}{}
\kfQStem{㉠∼㉣에 대한 설명으로 가장 적절하지 않은 것은?}
\begin{kfChoices}
\kfChoice ㉢에서는 과거의 일을 근거로 자책하고 있고, ㉣에서는 미래의 일을 예측하며 상대를 만류하고 있다.
\kfChoice ㉢에서는 과거의 일을 근거로 자책하고 있고, ㉣에서는 미래의 일을 예측하며 상대를 만류하고 있다.
\kfChoice ㉠에서는 당위를 내세워 상대를 설득하고, ㉡에서는 감정에 호소하여 상대의 동의를 이끌어 내고 있다.
\kfChoice ㉠에서는 상대보다 우월한 신분을 통해, ㉣에서는 상대에 대한 배려를 통해 상대의 태도 변화를 촉구하고 있다.
\end{kfChoices}
\end{kfQBlock}

\kfSecHeader{kfAreaWeekly}{지문}{문항 18 지문 / 시점/화자 혼동}

\kfPassageNote{문항 18 지문 / 시점/화자 혼동}{%
[26\textasciitilde{}31] 다음 글을 읽고 물음에 답하시오.

(가)   십 년 종사 후에 고향으로 도라오니   산천 의구하되 인사(人事)는 달라졌구나   아마도 세간의 존멸을 못내 슬허 하노라 <제1수>

산화(山花)는 믈의 픠고 물새는 산의 운다   일신이 한가하야 산수간의 누어시니   세상의 어즈러은 긔별을 나는 몰라 하노라 <제4수>

거믄고 빗기 들고 산수를 희롱하니   청풍은 건듯 블고 명월도 도라온다   하물며 유신(有信)한 갈매기는 오명 가명 하나니 <제5수>

거믄고 흥진(興盡)커던 조대(釣臺)로 내려가니   도화 뜬 말근 믈 뛰노나니 고기로다   아이야 밋기 다지 마라 취적(取適)*이나 하오리라 <제7수>

* 신교, ｢귀산음(歸山吟)｣ -   * 취적 : 낚시질의 참뜻이 세상 생각을 잊고자 하는 데 있음.

(나)   백수(白首)에 산수 구경 늦은 줄 알지마는   평생 품은 뜻을 이루고야 말리라 여겨   병자년 봄에 봄옷을 새로 입고   죽장망혜(竹杖芒鞋)로 노계 깊은 골에 …(중략)
}
\begin{kfQBlock}{18}{}
\kfQStem{<보기>를 참고하여 (가), (나)를 감상한 내용으로 적절하지 않은 것은? [3점]}
\begin{kfChoices}
\kfChoice (가)의 ‘십 년’, (나)의 ‘백수’는 자신이 바라던 생활을 누릴 수 있는 공간을 찾기 위해 화자가 노력한 세월로 볼 수 있군.
\kfChoice (가)의 ‘십 년’, (나)의 ‘백수’는 자신이 바라던 생활을 누릴 수 있는 공간을 찾기 위해 화자가 노력한 세월로 볼 수 있군.
\kfChoice (가)의 자연은 화자가 ‘고향’의 ‘산천’이 ‘의구하’다고 말하는 것으로 보아 화자에게 익숙한 곳으로 볼 수 있군.
\kfChoice (나)의 자연은 ‘임자 없이’ 감춰져 있던 곳이라는 점에서 사람들이 쉽게 찾지 못했던 곳으로 볼 수 있군.
\end{kfChoices}
\end{kfQBlock}



\clearpage

\kfChapterCover{2}{비트겐슈타인2 ch02}{Chapter 2}{이번 챕터: 문법 → 문학 5작품 → 비문학 5지문 → 패턴 워크북.}

\kfStartGrammar{이번 챕터 문법 영역 — 음운 / 음절의 끝소리 규칙 + 연음 심화}


\kfSecHeader{kfAreaGrammar}{문제}{}

\begin{multicols}{2}\raggedcolumns\setlength{\columnsep}{12pt}
\kfPassageFull{문항 1 지문 (concept)}{%
국어에서 음절 끝에 발음될 수 있는 자음은 'ㄱ, ㄴ, ㄷ, ㄹ, ㅁ, ㅂ, ㅇ'의 7개뿐이다. 이 7개 이외의 자음이 받침에 오면 이 7개 중 하나로 바뀌어 발음되는데, 이를 '음절의 끝소리 규칙'이라 한다. 'ㅋ, ㄲ'은 [ㄱ]으로, 'ㅅ, ㅆ, ㅈ, ㅊ, ㅌ'은 [ㄷ]으로, 'ㅍ'은 [ㅂ]으로 교체된다. 겹받침의 경우 두 자음 중 하나만 발음되는데, 'ㄳ, ㄵ, ㄼ, ㄽ, ㄾ, ㅀ, ㅄ'은 앞 자음이, 'ㄺ, ㄻ, ㄿ'은 뒤 자음이 발음된다. 다만 예외로 '밟다'의 'ㄼ'은 [ㅂ]으로, '넓적하다·넓둥글다'의 'ㄼ'도 [ㅂ]으로 발음된다.
}
\begin{kfQBlock}{1}{}
\kfQStem{윗글을 바탕으로 <보기>의 발음으로 적절하지 않은 것은?

<보기> ⓐ 부엌[부억]  ⓑ 낫[낟]  ⓒ 삶[삼]  ⓓ 밟다[밥따]  ⓔ 읽다[일따]}
\begin{kfChoices}
\kfChoice ⓐ: 'ㅋ'이 [ㄱ]으로 교체되어 [부억]이 된다.
\kfChoice ⓑ: 'ㅅ'이 [ㄷ]으로 교체되어 [낟]이 된다.
\kfChoice ⓒ: 겹받침 'ㄻ'에서 뒤 자음 [ㅁ]이 발음되어 [삼]이 된다.
\kfChoice ⓓ: 'ㄼ'에서 예외적으로 뒤 자음 [ㅂ]이 발음되고 된소리되기가 적용되어 [밥따]가 된다.
\kfChoice ⓔ: 겹받침 'ㄺ'에서 앞 자음 [ㄹ]이 발음되어 [일따]가 된다.
\end{kfChoices}
\end{kfQBlock}

\kfPassageFull{문항 2 지문 (concept)}{%
연음(이어 나기)이란 받침 자음이 뒤에 오는 모음 음절의 초성으로 옮겨 발음되는 현상이다. 체언과 조사, 어간과 어미 등 같은 말 안에서 받침 뒤에 모음이 오면 끝소리 규칙이 적용되지 않고 원래의 받침 자음이 그대로 연음된다. '옷이'에서 'ㅅ'은 대표음 [ㄷ]으로 바뀌지 않고 그대로 연음되어 [오시]가 된다. 겹받침의 연음에서는 뒤 자음이 다음 음절 초성으로 옮겨가고 앞 자음이 받침에 남는다. '읽어'는 'ㄹ'이 남고 'ㄱ'이 연음되어 [일거]가 된다. 반면 합성어 내부나 단어 경계에서는 끝소리 규칙이 먼저 적용된 후 연음되기도 한다. '꽃 위'는 'ㅅ'→[ㄷ](끝소리 규칙) → 연음 → [꼬뒤]가 된다.
}
\begin{kfQBlock}{2}{}
\kfQStem{윗글을 바탕으로 <보기>의 발음을 분석한 것으로 적절한 것은?

<보기> ⓐ 낮이[나지]    ⓑ 삶을[살믈]    ⓒ 밖에[바게] ⓓ 꽃 아래[꼬다래]    ⓔ 없어[업서]}
\begin{kfChoices}
\kfChoice ⓐ에서 'ㅈ'은 끝소리 규칙으로 [ㄷ]이 된 후 연음되어 [나디]가 된다.
\kfChoice ⓑ에서 겹받침 'ㄻ'의 앞 자음 'ㄹ'이 연음되어 [살믈]이 된다.
\kfChoice ⓒ에서 'ㄲ'은 끝소리 규칙 없이 바로 연음되어 [바께]가 된다.
\kfChoice ⓓ에서 'ㅅ'이 바로 연음되어 [꼬사래]가 된다.
\kfChoice ⓔ에서 겹받침 'ㅄ'의 뒤 자음 'ㅅ'이 연음되어 [업서]가 되는데, 이때 'ㅅ'에 끝소리 규칙이 적용된다.
\end{kfChoices}
\end{kfQBlock}

\kfPassageFull{문항 3 지문 (concept)}{%
겹받침의 발음에는 체계적인 원칙이 있지만, 동시에 주의해야 할 예외도 존재한다. 'ㄼ' 받침은 원칙적으로 앞 자음 [ㄹ]이 발음되어 '넓다'는 [널따]로 발음된다. 그러나 '밟다'는 예외적으로 뒤 자음 [ㅂ]이 발음되어 [밥따]가 된다. 또한 '넓적하다, 넓둥글다'에서도 [ㅂ]이 발음되어 각각 [넙쩌카다], [넙뚱글다]가 된다. 'ㄺ' 받침은 원칙적으로 뒤 자음 [ㄱ]이 발음되지만, '맑게'에서는 [막게]→[막게]? 표준 발음법 제10항에 따르면 'ㄺ'은 [ㄱ]으로 발음하므로 '맑다'[막따], '맑고'[막꼬]가 된다.
}
\begin{kfQBlock}{3}{}
\kfQStem{윗글을 바탕으로 <보기>의 발음이 표준 발음법에 맞는지 판단한 것으로 적절하지 않은 것은?

<보기> ⓐ 넓다[널따]  ⓑ 밟고[밥꼬]  ⓒ 넓적[넙쩍]  ⓓ 맑다[말따]  ⓔ 삶다[삼따]}
\begin{kfChoices}
\kfChoice ⓐ는 'ㄼ'의 원칙적 발음으로 적절하다.
\kfChoice ⓑ는 예외적으로 [ㅂ]이 발음된 후 된소리되기가 적용되어 적절하다.
\kfChoice ⓒ는 '넓적'이 예외적으로 [ㅂ]이 발음되어 적절하다.
\kfChoice ⓓ는 'ㄺ'에서 앞 자음 [ㄹ]이 발음되어 적절하다.
\kfChoice ⓔ는 'ㄻ'에서 뒤 자음 [ㅁ]이 발음되어 적절하다.
\end{kfChoices}
\end{kfQBlock}

\kfPassageFull{문항 4 지문 (concept)}{%
음절의 끝소리 규칙은 단독으로 적용되기도 하지만, 다른 음운 변동의 전 단계가 되기도 한다. '꽃나무'를 예로 들면, 먼저 'ㅅ'이 끝소리 규칙에 의해 [ㄷ]으로 교체된다. 그 다음 [ㄷ]이 비음 'ㄴ' 앞에서 비음화하여 [ㄴ]이 되므로 최종 발음은 [꼰나무]이다. 이처럼 한 단어에 두 가지 이상의 음운 변동이 순차적으로 적용되는 것을 '복합 음운 변동'이라 한다. 또한 '있고'에서는 'ㅆ'→[ㄷ](끝소리 규칙) 후, [ㄷ] 뒤에서 'ㄱ'이 된소리가 되어 [익꼬]가 된다. 변동의 적용 순서는 일반적으로 끝소리 규칙 → 비음화/유음화 → 된소리되기 순이다.
}
\begin{kfQBlock}{4}{}
\kfQStem{윗글을 바탕으로 <보기>의 단어에 적용되는 음운 변동 과정을 바르게 분석한 것은?

<보기> '젖먹이'의 발음 과정을 단계별로 분석하시오.}
\begin{kfChoices}
\kfChoice 'ㅈ'→[ㄷ](끝소리 규칙) → [ㄴ](비음화)만 적용되어 [전머기]가 된다.
\kfChoice 'ㅈ'→[ㄷ](끝소리 규칙) → [ㄴ](비음화) 후 [젼머기]가 된다.
\kfChoice 'ㅈ'→[ㄷ](끝소리) → [ㄴ](비음화) → 'ㄱ' 연음 → [전머기]
\kfChoice 'ㅈ'→[ㄷ](끝소리 규칙) → 된소리되기가 적용되어 [젇떠기]가 된다.
\kfChoice 끝소리 규칙 없이 'ㅈ'이 바로 비음화되어 [전머기]가 된다.
\end{kfChoices}
\end{kfQBlock}

\kfPassageFull{문항 5 지문 (concept)}{%
연음이 적용되는 환경은 크게 두 가지로 나뉜다. 첫째, 체언+조사, 어간+어미처럼 같은 말 안에서 받침 뒤에 모음이 오는 경우에는 원래 받침 자음이 그대로 다음 음절 초성으로 옮겨간다. '옷을'[오슬], '꽃이'[꼬치], '밖에'[바께]가 이에 해당한다. 둘째, 합성어 내부나 단어와 단어 사이(단어 경계)에서는 끝소리 규칙이 먼저 적용된 후 그 결과가 연음된다. '맛없다'에서 'ㅅ'→[ㄷ](끝소리 규칙) → 연음 → [마덥따]가 된다. 다만 '맛있다'는 [마딧따]와 [마싣따] 두 발음이 모두 허용되는 특수한 경우이다.
}
\begin{kfQBlock}{5}{}
\kfQStem{윗글을 바탕으로 <보기>의 밑줄 친 부분의 발음과 그 이유를 분석한 것으로 적절하지 않은 것은?

<보기> ⓐ 닭을 [달글]    ⓑ 꽃 위 [꼬뒤] ⓒ 헛일 [헌닐]    ⓓ 값어치 [가버치]}
\begin{kfChoices}
\kfChoice ⓐ: 체언+조사이므로 겹받침 'ㄺ'에서 'ㄱ'이 그대로 연음되어 [달글]이 된다.
\kfChoice ⓑ: 단어 경계이므로 끝소리 규칙 'ㅅ'→[ㄷ]이 먼저 적용되고 연음되어 [꼬뒤]가 된다.
\kfChoice ⓒ: 합성어이므로 끝소리 규칙 'ㅅ'→[ㄷ]이 적용되고 'ㄴ' 첨가와 비음화를 거쳐 [헌닐]이 된다.
\kfChoice ⓓ: '값'의 겹받침 'ㅄ'에서 끝소리 규칙으로 [ㄱ]이 남고 연음되어 [가버치]가 된다.
\kfChoice ⓐ와 ⓑ는 연음이 적용되는 환경이 다르다.
\end{kfChoices}
\end{kfQBlock}

\end{multicols}


\kfStartLit{이번 챕터 문학 작품 5편}


\kfSecHeader{kfAreaLiterature}{지문}{장마 — 윤흥길 (현대소설) / 윤흥길, 「장마」(1973), 수능·학평 기출 다수}

\kfPassageNote{장마 — 윤흥길 (현대소설) / 윤흥길, 「장마」(1973), 수능·학평 기출 다수}{%
외할머니가 다시 친할머니에게 말을 건넸다. "삼 형제 중에 하나가 그러고 나면 남은 식구들이야 오죽하겠소." 친할머니는 대꾸하지 않았다. 한동안 침묵이 흘렀다.

그러다 느닷없이 마당에 큰 구렁이 한 마리가 나타났다. 식구들이 놀라 소리를 질렀지만, 구렁이는 다른 사람들에게는 눈길도 주지 않고 곧장 외할머니 앞으로 다가갔다. 외할머니는 구렁이를 바라보며 한참 동안 눈물을 흘렸다. "너 정녕 내 아들이냐. 네가 왔구나." 그리고 구렁이를 향해 말했다. "이제 한이 풀렸으면 편히 가거라."

나는 무서웠지만 그 자리에서 움직일 수가 없었다. 친할머니도 처음에는 돌처럼 굳어 있었다. 그러나 구렁이가 마당을 가로질러 천천히 나가는 것을 보며 친할머니가 낮은 목소리로 말했다. "가거라. 가."

그 순간 나는 친할머니의 목소리에서 미움도, 분노도, 슬픔도 아닌 무엇인가를 느꼈다. 그것은 내가 어려서 알 수 없었지만, 어쩌면 용서라는 것이었을지도 모른다.

장마가 개었다. 햇빛이 마당에 쏟아졌고, 빨래가 바람에 펄럭이고 있었다.
}
\kfSecHeader{kfAreaLiterature}{활동}{작품 정리표 — 장마(윤흥길)}

\noindent\textit{\small 빈칸에 알맞은 말을 채워 넣으시오. (초성 힌트 참고)}\par\medskip
\begin{kfActTable}{|>{\centering\arraybackslash}m{2.6cm}|m{6cm}|m{4cm}|}
\rowcolor{kfPrimaryLight!50}\textbf{\color{kfPrimary}항목} & \textbf{\color{kfPrimary}내용 (빈칸 채우기)} & \textbf{\color{kfPrimary}답안 작성}\\\hline
갈래 & ㅈㅍ ㅅㅅ ① & \kfTBlank \\\hline
시대 배경 & 6·25 ㅈㅈ 중 ② & \kfTBlank \\\hline
서술 시점 & 1ㅇㅊ ㄱㅊㅈ ㅅㅈ (어린 '나') ③ & \kfTBlank \\\hline
갈등 축 1 & 친할머니 아들 = ㄱㄱ ④ & \kfTBlank \\\hline
갈등 축 2 & 외할머니 아들 = ㅇㅁㄱ ⑤ & \kfTBlank \\\hline
갈등 원인 & ㅇㄴ 대립이 ㅎㅇ관계와 충돌 ⑥ & \kfTBlank \\\hline
화해 계기 & ㄱㄹㅇ의 등장 → 외할머니: '내 아들이냐' ⑦ & \kfTBlank \\\hline
화해 장면 & 친할머니: 'ㄱㄱㄹ. ㄱ.' ⑧ & \kfTBlank \\\hline
서술자 해석 & 어쩌면 ㅇㅅ라는 것이었을지도 ⑨ & \kfTBlank \\\hline
제목 상징 1 & ㅈㅁ = 전쟁의 ㅂㄱ과 갈등 ⑩ & \kfTBlank \\\hline
제목 상징 2 & 장마가 ㄱㅇ다 = ㅎㅎ의 가능성 ⑪ & \kfTBlank \\\hline
주제 & ㅇㄴ 갈등의 ㅊㅇ과 인간적 ㅎㅎ ⑫ & \kfTBlank \\\hline
\end{kfActTable}
\kfSecHeader{kfAreaLiterature}{문제}{}

\begin{multicols}{2}\raggedcolumns\setlength{\columnsep}{12pt}
\begin{kfQBlock}{1}{}
\kfQStem{윗글에 대한 설명으로 적절한 것은?}
\begin{kfChoices}
\kfChoice 어린 서술자의 시선이 이념 갈등의 의미를 완전히 파악하여 분석적으로 전달한다.
\kfChoice 초자연적 존재(구렁이)의 등장이 인물 간 갈등 해소의 계기로 기능한다.
\kfChoice 대화를 통해 두 할머니가 이념적 입장을 논리적으로 공방한다.
\kfChoice 서술자가 전지적 시점에서 모든 인물의 내면을 직접 서술한다.
\kfChoice 과거와 현재를 교차하는 서술 구조를 통해 전쟁의 원인을 규명한다.
\end{kfChoices}
\end{kfQBlock}

\begin{kfQBlock}{2}{}
\kfQStem{윗글에서 '장마가 개었다'가 지닌 상징적 의미로 가장 적절한 것은?}
\begin{kfChoices}
\kfChoice 전쟁이 다시 일어나 더 큰 비극이 시작될 것이라는 어두운 암시
\kfChoice 이념 갈등으로 인한 고통이 해소되고 화해의 가능성이 열리는 것
\kfChoice 여름이 지나가는 계절의 자연스러운 변화를 사실적으로 묘사한 것
\kfChoice 인물들이 전쟁의 아픈 기억을 완전히 잊어버렸음을 나타낸 것
\kfChoice 친할머니의 오랜 복수심이 마침내 실현되었음을 상징한 장면
\end{kfChoices}
\end{kfQBlock}

\begin{kfQBlock}{3}{}
\kfQStem{친할머니가 구렁이에게 '가거라. 가.'라고 말한 행위에 담긴 의미로 가장 적절한 것은?}
\begin{kfChoices}
\kfChoice 구렁이에 대한 두려움 때문에 황급히 물리치려는 것
\kfChoice 외할머니 아들의 혼령을 인정하며 적대감을 내려놓는 화해 표현
\kfChoice 외할머니의 미신적인 행동을 속으로 비웃고 무시하는 것
\kfChoice 자기 아들의 죽음에 대한 깊은 원한을 다시 확인하려는 것
\kfChoice 구렁이가 어린 식구들에게 해를 끼칠까 봐 쫓아내는 것
\end{kfChoices}
\end{kfQBlock}

\begin{kfQBlock}{4}{}
\kfQStem{윗글을 바탕으로 <보기>를 감상한 내용으로 적절하지 않은 것은?

<보기> 윤흥길의 「장마」는 6·25 전쟁이 한 가정에 남긴 이념 갈등을 다룬다. 이 작품에서 갈등의 핵심은 국군과 인민군이라는 적대적 이념이 혈연관계와 충돌하는 데 있다. 어린 서술자는 이 갈등의 의미를 완전히 이해하지 못하지만, 감각적으로 포착함으로써 독자에게 전달한다. 결말의 화해는 이념의 논리가 아닌 인간적 연민에 의해 가능해진다.}
\begin{kfChoices}
\kfChoice 외할머니가 구렁이를 아들로 인식하는 것은 혈연의 정이 이념의 논리를 넘어서는 장면이겠군.
\kfChoice 친할머니의 '가거라, 가'는 적군의 혼령을 받아들이는 인간적 연민의 발로이겠군.
\kfChoice '나'가 '용서라는 것이었을지도 모른다'고 한 것은 어린 서술자가 갈등의 의미를 감각적으로 포착한 것이겠군.
\kfChoice 두 할머니의 화해는 이념적 논쟁을 통해 어느 한쪽이 승리하여 이루어진 것이겠군.
\kfChoice '장마가 개었다'는 갈등의 해소와 화해의 가능성이 열렸음을 상징하겠군.
\end{kfChoices}
\end{kfQBlock}

\begin{kfQBlock}{5}{}
\kfQStem{윗글 결말에서 '햇빛이 마당에 쏟아졌고, 빨래가 바람에 펄럭이고 있었다'라는 서술의 효과로 가장 적절한 것은?}
\begin{kfChoices}
\kfChoice 일상의 평온한 이미지로 화해 뒤 회복의 분위기를 형상화한다.
\kfChoice 날씨의 급격한 변화를 사실 그대로 보고하는 객관적 서술이다.
\kfChoice 인물들의 갈등이 해소되지 않았음을 역설적으로 보여준다.
\kfChoice 전쟁이 곧 재개될 것이라는 불길한 전조이다.
\kfChoice 서술자가 성인이 된 뒤의 시점으로 전환되었음을 나타낸다.
\end{kfChoices}
\end{kfQBlock}

\end{multicols}

\kfSecHeader{kfAreaLiterature}{지문}{완장 — 윤흥길 (현대소설) / 윤흥길, 「완장」(1978), 수능특강 수록}

\kfPassageNote{완장 — 윤흥길 (현대소설) / 윤흥길, 「완장」(1978), 수능특강 수록}{%
완장을 팔뚝에 두른 순간부터 그는 딴사람이 되었다. 전에는 아무한테서나 구박을 받던 사람이 이젠 말 한마디에 사람들이 고개를 숙이는 것이었다. 그것은 신기한 경험이었고, 또 무섭도록 달콤한 것이었다.

저수지에 낚시를 하러 온 사람이 있었다. 허가증이 없었다. 그는 큰 소리로 호통을 쳤다. "이 사람이, 여기가 어디라고 함부로 들어와. 당장 나가." 낚시꾼이 사정을 했지만 그는 들은 체도 하지 않았다. 허가증 없이 낚시하는 것은 규칙 위반이니까. 그러나 그가 정말로 지키고 싶었던 것이 규칙이었을까.

그날 밤 그는 거울 앞에 서서 완장을 바라보았다. 팔뚝 위의 완장은 이상하게 빛나 보였다. 예전에 자신을 무시하던 놈들을 떠올리며 그는 비릿한 미소를 지었다. 내일은 그놈들에게도 한번 보여줘야지. 나도 이제 사람 구실을 하는 거야.

그러나 어느 날 한 노인이 찾아왔다. 노인의 눈에서 그는 자기 아버지의 모습을 보았다. 평생을 남의 눈치를 보며 살다 간 아버지. 완장 하나에 기고만장해진 자신이 갑자기 부끄러워졌다. 그렇지만 그 부끄러움은 오래가지 않았다. 다음 날 아침 완장을 두르자 어제의 감정은 안개처럼 사라져 버렸다.
}
\kfSecHeader{kfAreaLiterature}{활동}{작품 정리표 — 완장(윤흥길)}

\noindent\textit{\small 빈칸에 알맞은 말을 채워 넣으시오. (초성 힌트 참고)}\par\medskip
\begin{kfActTable}{|>{\centering\arraybackslash}m{2.6cm}|m{6cm}|m{4cm}|}
\rowcolor{kfPrimaryLight!50}\textbf{\color{kfPrimary}항목} & \textbf{\color{kfPrimary}내용 (빈칸 채우기)} & \textbf{\color{kfPrimary}답안 작성}\\\hline
갈래 & ㄷㅍ ㅅㅅ ① & \kfTBlank \\\hline
핵심 소재 & ㅇㅈ = 소규모 ㄱㄹ의 상징 ② & \kfTBlank \\\hline
주인공 변화 & 완장을 두르자 'ㄸㅅㄹ'이 됨 ③ & \kfTBlank \\\hline
권력의 느낌 & '무섭도록 ㄷㅋㅎ 것' ④ & \kfTBlank \\\hline
권력 남용 사례 & ㅎㄱㅈ 없는 낚시꾼에게 ㅎㅌ ⑤ & \kfTBlank \\\hline
서술자 의문 & '정말 지키고 싶었던 것이 ㄱㅊ이었을까' ⑥ & \kfTBlank \\\hline
복수심 & 무시하던 놈들에게 ㅂㄹㅎ ㅁㅅ ⑦ & \kfTBlank \\\hline
자기 반성 계기 & ㄴㅇ에게서 ㅇㅂㅈ의 모습을 봄 ⑧ & \kfTBlank \\\hline
반성의 한계 & 부끄러움이 ㅇㄹㄱㅈ 않았다 ⑨ & \kfTBlank \\\hline
아이러니 & ㅁㅁㅎ 권력에도 ㅂㅈ → ㅂㅍㅈ 위험 ⑩ & \kfTBlank \\\hline
주제 1 & ㄱㄹㅇ ㅇㅎ과 ㄷㄷㅈ ㅈㅊ ㅂㄱ ⑪ & \kfTBlank \\\hline
발표 시기 & 1978년 ⑫ & \kfTBlank \\\hline
\end{kfActTable}
\kfSecHeader{kfAreaLiterature}{문제}{}

\begin{multicols}{2}\raggedcolumns\setlength{\columnsep}{12pt}
\begin{kfQBlock}{1}{}
\kfQStem{윗글에 대한 설명으로 적절한 것은?}
\begin{kfChoices}
\kfChoice 서술자가 인물의 외면 행동에 더하여 내면 심리도 함께 서술하고 있다.
\kfChoice 1인칭 서술자가 자신의 지난 경험을 회고하는 형식으로 서술한다.
\kfChoice 인물 간 대화를 중심으로 갈등이 전개되고 최종적으로 해결된다.
\kfChoice 자연 풍경에 대한 묘사가 서사의 중심을 이루고 있다.
\kfChoice 사건의 시간 순서를 의도적으로 뒤섞어 서술하고 있다.
\end{kfChoices}
\end{kfQBlock}

\begin{kfQBlock}{2}{}
\kfQStem{윗글에서 '완장'이 상징하는 바로 가장 적절한 것은?}
\begin{kfChoices}
\kfChoice 사회적 약자를 안전하게 보호하기 위한 제도적 장치
\kfChoice 작은 권력이 개인의 도덕적 주체를 변질시키는 매개
\kfChoice 정당한 노력으로 얻어 낸 사회적 인정과 정당한 보상
\kfChoice 전통적 공동체의 질서를 굳건히 유지하는 권위의 상징
\kfChoice 인물이 자신의 과거를 깊이 반성하고 성장하게 되는 계기
\end{kfChoices}
\end{kfQBlock}

\begin{kfQBlock}{3}{}
\kfQStem{주인공이 노인에게서 '자기 아버지의 모습'을 본 뒤 '부끄러워졌다'가 오래가지 않은 이유로 가장 적절한 것은?}
\begin{kfChoices}
\kfChoice 노인이 사과했기 때문에 부끄러움이 해소되었다.
\kfChoice 완장이 주는 권력의 유혹이 자기 반성보다 더 강했기 때문이다.
\kfChoice 아버지에 대한 기억이 흐릿해져 더 이상 떠오르지 않았기 때문이다.
\kfChoice 주인공이 도덕적 성찰을 통해 성장하여 부끄러움을 극복했기 때문이다.
\kfChoice 동료들이 주인공의 행동이 정당하다고 격려했기 때문이다.
\end{kfChoices}
\end{kfQBlock}

\begin{kfQBlock}{4}{}
\kfQStem{윗글을 바탕으로 <보기>를 감상한 내용으로 적절하지 않은 것은?

<보기> 윤흥길의 「완장」은 권력이 개인의 도덕적 주체를 어떻게 붕괴시키는지를 보여준다. 특히 작품의 아이러니는 이 권력이 극히 미미한 것이라는 데 있다. 저수지 관리원이라는 사소한 지위조차 인간을 변질시킬 수 있다면, 이는 권력의 유혹이 인간 본성에 깊이 뿌리 내리고 있음을 시사한다.}
\begin{kfChoices}
\kfChoice 주인공이 '딴사람'이 된 것은 권력이 도덕적 주체를 붕괴시키는 과정을 보여주는 것이겠군.
\kfChoice 주인공이 낚시꾼에게 호통을 친 것은 규칙 준수가 아니라 권력 남용의 사례이겠군.
\kfChoice '무섭도록 달콤한 것'이라는 표현은 권력의 유혹이 갖는 위험한 매력을 형상화한 것이겠군.
\kfChoice 완장이 미미한 권력의 상징이라는 점에서, 주인공의 변질은 큰 권력에서만 나타나는 예외적 현상이겠군.
\kfChoice 부끄러움이 완장을 두르자 사라진 것은 권력의 유혹이 자기 반성보다 강함을 보여주는 것이겠군.
\end{kfChoices}
\end{kfQBlock}

\begin{kfQBlock}{5}{}
\kfQStem{'비릿한 미소'라는 표현이 드러내는 주인공의 심리로 가장 적절한 것은?}
\begin{kfChoices}
\kfChoice 자신을 무시했던 사람들에 대한 복수심과 뒤틀린 쾌감
\kfChoice 변해 가는 자신의 모습에 대한 자기 비하와 수치심
\kfChoice 완장을 통해 사회적으로 인정받게 된 순수한 기쁨
\kfChoice 낚시꾼에 대한 동정과 미안함
\kfChoice 아버지의 삶을 이해하게 된 감동
\end{kfChoices}
\end{kfQBlock}

\end{multicols}

\kfSecHeader{kfAreaLiterature}{지문}{위로 — 윤동주 (현대시(산문시)) / 윤동주, 「위로」(1940), 학평·모평 기출 다수}

\kfPassageNote{위로 — 윤동주 (현대시(산문시)) / 윤동주, 「위로」(1940), 학평·모평 기출 다수}{%
거미란 놈이 흉한 심보로 병원 뒤뜰 난간과 꽃밭 사이 사람 발이 잘 닿지 않는 곳에 그물을 쳐 놓았다. 옥외 요양을 받는 젊은 사나이가 누워서 치어다보기 바르게―

나비가 한 마리 꽃밭에 날아들다 그물에 걸리었다. 노오란 날개를 파득거려도 파득거려도 나비는 자꾸 감기우기만 한다. 거미가 쏜살같이 가더니 끝없는 끝없는 실을 뽑아 나비의 온몸을 감아 버린다. 사나이는 긴 한숨을 쉬었다.

나이보담 무수한 고생 끝에 때를 잃고 병을 얻은 이 사나이를 위로할 말이―거미줄을 헝클어 버리는 것밖에 위로의 말이 없었다.
}
\kfSecHeader{kfAreaLiterature}{활동}{작품 정리표 — 위로(윤동주)}

\noindent\textit{\small 빈칸에 알맞은 말을 채워 넣으시오. (초성 힌트 참고)}\par\medskip
\begin{kfActTable}{|>{\centering\arraybackslash}m{2.6cm}|m{6cm}|m{4cm}|}
\rowcolor{kfPrimaryLight!50}\textbf{\color{kfPrimary}항목} & \textbf{\color{kfPrimary}내용 (빈칸 채우기)} & \textbf{\color{kfPrimary}답안 작성}\\\hline
갈래 & ㅈㅇㅅ, ㅅㅈㅅ(ㅅㅁㅅ) ① & \kfTBlank \\\hline
성격 & ㅅㄱㅈ, ㅅㅅㅈ, ㅅㅁㅈ ② & \kfTBlank \\\hline
공간 & ㅂㅇ ㄷㄸ ③ & \kfTBlank \\\hline
거미의 행위 & ㅎㅎ ㅅㅂ로 ㄱㅁ을 침 ④ & \kfTBlank \\\hline
나비의 처지 & ㄴㅇㄹ ㄴㄱ를 ㅍㄷㄱㄹ도 자꾸 감김 ⑤ & \kfTBlank \\\hline
거미의 결정타 & ㅆㅅㄱㅇ 가더니 끝없는 실로 감음 ⑥ & \kfTBlank \\\hline
사나이의 반응 & ㄱ ㅎㅅ을 쉼 ⑦ & \kfTBlank \\\hline
수법 & ㅇㅎㅈ 수법, ㅇㅇㅎ ⑧ & \kfTBlank \\\hline
상징 1 & 거미 = ㅇㅈ ⑨ & \kfTBlank \\\hline
상징 2 & 나비 = ㅇㄹ ㅁㅈ ⑩ & \kfTBlank \\\hline
상징 3 & 젊은 사나이 = ㅁㄱㄹㅎ ㄷㅅㄷㅇ ⑪ & \kfTBlank \\\hline
주제 & 암담한 현실에 처한 존재에 대한 ㅇㄹ와 ㅁㄱㄹ에 대한 ㅇㅌㄲㅇ ⑫ & \kfTBlank \\\hline
\end{kfActTable}
\kfSecHeader{kfAreaLiterature}{문제}{}

\begin{multicols}{2}\raggedcolumns\setlength{\columnsep}{12pt}
\begin{kfQBlock}{1}{}
\kfQStem{윗시에 대한 설명으로 적절하지 않은 것은?}
\begin{kfChoices}
\kfChoice 산문적 어조로 서술하면서도 ‘\textasciitilde{}다’의 종결 표현을 반복하여 운율감을 형성하고 있다.
\kfChoice ‘거미’와 ‘나비’를 의인화하는 우화적 수법으로 상황을 형상화하고 있다.
\kfChoice 줄표(―)를 사용하여 특정 장면이나 상황을 강조하고 있다.
\kfChoice 시각적 심상을 주로 활용하여 사건을 구체적으로 묘사하고 있다.
\kfChoice 공간의 빈번한 이동을 통해 화자의 정서 변화를 단계적으로 보여 주고 있다.
\end{kfChoices}
\end{kfQBlock}

\begin{kfQBlock}{2}{}
\kfQStem{윗시의 ‘젊은 사나이’에 대한 이해로 가장 적절한 것은?}
\begin{kfChoices}
\kfChoice 거미줄을 끊어 나비를 직접 구해 내는 적극적 인물이다.
\kfChoice 병원에서 옥외 요양을 받으며 거미와 나비의 사건을 무력하게 지켜보는 인물이다.
\kfChoice 거미를 직접 잡아 죽임으로써 부정적 현실을 단번에 근본적으로 해결하는 인물이다.
\kfChoice 꽃밭을 가꾸는 정원사로서 자연의 질서를 통제하는 인물이다.
\kfChoice 사건과 무관한 외부 관찰자로서 어떤 감정도 드러내지 않는 인물이다.
\end{kfChoices}
\end{kfQBlock}

\begin{kfQBlock}{3}{}
\kfQStem{윗시에서 ‘거미줄을 헝클어 버리는 것밖에 위로의 말이 없었다’가 함축하는 의미로 가장 적절한 것은?}
\begin{kfChoices}
\kfChoice 거미줄을 헝클어 버림으로써 사나이의 모든 문제가 완전히 해결되었음을 선언한다.
\kfChoice 위로조차 일시적이며 부정적 현실을 근본적으로 극복할 수 없다는 무력감을 드러낸다.
\kfChoice 사나이의 병이 곧 완치될 것이라는 화자의 낙관적 전망을 보여 준다.
\kfChoice 거미를 직접 잡아 죽이지 못한 화자의 분노가 폭력적으로 표출된다.
\kfChoice 꽃밭을 가꾸는 새로운 일과를 시작하려는 화자의 의지를 드러낸다.
\end{kfChoices}
\end{kfQBlock}

\begin{kfQBlock}{4}{}
\kfQStem{<보기>를 바탕으로 윗시를 감상한 내용으로 적절하지 않은 것은? [3점]

<보기> 윤동주의 「위로」는 일제 강점기의 부정적 현실을 배경으로, ‘거미–나비–사나이’의 우화적 장면을 통해 우리 민족의 무기력한 처지를 형상화한 산문시이다. ‘거미’는 일제를, ‘나비’는 우리 민족을, ‘젊은 사나이’는 식민지 현실에 무기력하게 머무는 시인 자신과 동시대인을 상징한다. 화자는 거미줄을 헝클어 버리는 행위를 통해 사나이를 위로하지만, 이 행위가 근본적 해결이 아님을 자각하면서 안타까움도 함께 드러낸다.}
\begin{kfChoices}
\kfChoice ‘흉한 심보’로 그물을 쳐 놓는 ‘거미’는 우리 민족을 억압하는 일제를 상징한다고 볼 수 있겠군.
\kfChoice ‘날개를 파득거려도 파득거려도’ 그물에 감기우기만 하는 ‘나비’는 일제 강점기 우리 민족의 무기력한 처지를 보여 준다고 할 수 있겠군.
\kfChoice ‘긴 한숨’만 쉴 뿐 어떤 행동도 하지 못하는 ‘젊은 사나이’는 식민지 현실에 무기력하게 머무는 동시대인을 상징한다고 볼 수 있겠군.
\kfChoice ‘거미줄을 헝클어 버리는’ 화자의 행위에서 일시적 위로를 건네지만 근본적 해결에 이르지 못하는 안타까움을 함께 읽을 수 있겠군.
\kfChoice ‘위로의 말이 없었다’는 진술은 사나이의 고통이 이미 모두 해결된 뒤의 평온을 보여 주는 표현이라고 할 수 있겠군.
\end{kfChoices}
\end{kfQBlock}

\begin{kfQBlock}{5}{}
\kfQStem{윗시에서 ‘노오란 날개를 파득거려도 파득거려도 나비는 자꾸 감기우기만 한다’의 표현 효과로 가장 적절한 것은?}
\begin{kfChoices}
\kfChoice 동일한 시구의 반복으로 나비의 안간힘과 그럼에도 벗어나지 못하는 상황을 강조한다.
\kfChoice 의인법을 활용하여 거미의 잔인함을 직접 비판하며 화자의 격앙된 분노를 토로한다.
\kfChoice 공감각적 심상을 통해 봄날 꽃밭의 화사한 분위기를 부각한다.
\kfChoice 수미상관 구조로 시 전체의 통일성을 형성하고 주제를 강조한다.
\kfChoice 반어적 표현을 사용하여 나비의 자유를 역설적으로 드러낸다.
\end{kfChoices}
\end{kfQBlock}

\end{multicols}

\kfSecHeader{kfAreaLiterature}{지문}{원미동 시인 — 양귀자 (현대소설) / 양귀자, 『원미동 사람들』(1987) 중 「원미동 시인」, 수능특강 수록}

\kfPassageNote{원미동 시인 — 양귀자 (현대소설) / 양귀자, 『원미동 사람들』(1987) 중 「원미동 시인」, 수능특강 수록}{%
김 시인은 구멍가게의 셔터를 올리며 생각했다. 오늘도 외상 장부에 빨간 줄이 늘어날 것이다. 외상값을 갚지 못하는 사람들을 탓할 수도 없는 노릇이었다. 그들도 다 사정이 있는 것이니까.

한 달 전, 옆집 최씨 아줌마가 울면서 찾아왔다. 남편이 실직하여 아이들 학원비를 낼 수 없다는 것이었다. 김 시인은 잠시 망설였다. 자기네도 넉넉한 형편이 아니었다. 그러나 그는 서랍에서 봉투를 꺼내 최씨 아줌마에게 내밀었다. "이거 갚을 생각 하지 마시오. 그냥 받으시오." 그것은 아내 몰래 모아 둔 원고료였다. 시 한 편에 만 원도 안 되는 원고료. 그 돈이면 아이들 운동화를 살 수 있었는데.

아내가 알면 난리가 날 것이다. 하지만 김 시인은 이상하게도 마음이 가벼웠다. 그날 밤 그는 노트를 펴고 시를 한 편 썼다. '가난해도 가난에 구차하지 않은 사람이 되고 싶다'는 내용이었다.

구멍가게 앞을 지나가던 최씨 아줌마의 아이가 고개를 숙여 인사했다. 김 시인은 손을 흔들며 생각했다. 시를 쓴다는 것은 이런 것이 아닐까. 삶에 한 줄의 품위를 더하는 것.
}
\kfSecHeader{kfAreaLiterature}{활동}{작품 정리표 — 원미동 시인(양귀자)}

\noindent\textit{\small 빈칸에 알맞은 말을 채워 넣으시오. (초성 힌트 참고)}\par\medskip
\begin{kfActTable}{|>{\centering\arraybackslash}m{2.6cm}|m{6cm}|m{4cm}|}
\rowcolor{kfPrimaryLight!50}\textbf{\color{kfPrimary}항목} & \textbf{\color{kfPrimary}내용 (빈칸 채우기)} & \textbf{\color{kfPrimary}답안 작성}\\\hline
갈래 & ㄷㅍ ㅅㅅ ① & \kfTBlank \\\hline
출전 & 『ㅇㅁㄷ ㅅㄹㄷ』(1987) ② & \kfTBlank \\\hline
주인공 & ㄱㅁㄱㄱ 주인이자 ㅅㅇ ③ & \kfTBlank \\\hline
윤리적 선택 & ㅇㄱㄹ를 최씨 아줌마에게 줌 ④ & \kfTBlank \\\hline
선택의 조건 & 자기네도 ㄴㄴㅎ 형편이 아님 ⑤ & \kfTBlank \\\hline
시의 내용 & 가난해도 ㄱㅊ하지 않은 사람 ⑥ & \kfTBlank \\\hline
시 쓰기의 의미 & 삶에 한 줄의 ㅍㅇ를 더하는 것 ⑦ & \kfTBlank \\\hline
공동체 교류 & 최씨 아줌마 아이의 ㅇㅅ ⑧ & \kfTBlank \\\hline
서술 특징 & 일상적이고 ㄷㄷㅎ 서술 ⑨ & \kfTBlank \\\hline
주제 1 & 서민 삶 속 ㅇㄹㅈ ㅅㅌ과 ㅇㄷ ⑩ & \kfTBlank \\\hline
주제 2 & 가난 속에서도 ㅇㄱㅈ ㅈㅇ 유지 ⑪ & \kfTBlank \\\hline
발표 시기 & 1987년 ⑫ & \kfTBlank \\\hline
\end{kfActTable}
\kfSecHeader{kfAreaLiterature}{문제}{}

\begin{multicols}{2}\raggedcolumns\setlength{\columnsep}{12pt}
\begin{kfQBlock}{1}{}
\kfQStem{윗글에 대한 설명으로 적절한 것은?}
\begin{kfChoices}
\kfChoice 내면 심리와 행동이 함께 서술되어 인물 성격이 입체적으로 드러난다.
\kfChoice 급격한 사건의 전개를 통해 긴박한 서사적 긴장감을 강하게 조성한다.
\kfChoice 화려한 수사와 다양한 비유를 통해 도시의 화려한 모습을 묘사한다.
\kfChoice 1인칭 서술자의 주관적인 고백을 통해 일련의 사건을 전달하고 있다.
\kfChoice 역사적인 사건을 중심에 두어 시대적 전환점을 분명히 형상화한다.
\end{kfChoices}
\end{kfQBlock}

\begin{kfQBlock}{2}{}
\kfQStem{윗글에서 김 시인이 원고료를 내주면서 '갚을 생각 하지 마시오'라고 한 행위에 담긴 의미로 가장 적절한 것은?}
\begin{kfChoices}
\kfChoice 자기 이익을 접고서 타인의 고통을 외면하지 않는 윤리적 연대
\kfChoice 경제적으로 여유로운 사람이 가난한 이웃에게 베푸는 일방적 자선
\kfChoice 상대방에게 빚을 지움으로써 관계에서 우위를 점하려는 전략
\kfChoice 아내에 대한 불만을 돈을 통해 간접적으로 표출하는 행위
\kfChoice 시인으로서의 사회적 명성을 높이기 위한 계산된 행위
\end{kfChoices}
\end{kfQBlock}

\begin{kfQBlock}{3}{}
\kfQStem{김 시인이 '가난해도 가난에 구차하지 않은 사람이 되고 싶다'는 시를 쓴 것이 드러내는 인물의 가치관으로 가장 적절한 것은?}
\begin{kfChoices}
\kfChoice 물질적 가난을 벗어나기 위해 더 열심히 돈을 벌어야 한다는 현실 인식
\kfChoice 경제적 빈곤 속에서도 인간적 품위와 도덕적 존엄을 지키려는 의지
\kfChoice 가난한 사람들이 사회에 저항해야 한다는 계급 의식
\kfChoice 가난을 미덕으로 찬양하는 반물질주의적 태도
\kfChoice 시를 쓰면 경제적 보상이 따를 것이라는 기대
\end{kfChoices}
\end{kfQBlock}

\begin{kfQBlock}{4}{}
\kfQStem{윗글을 바탕으로 <보기>를 감상한 내용으로 적절하지 않은 것은?

<보기> 양귀자의 「원미동 시인」은 서민의 일상 속에서 도덕적 선택의 의미를 묻는다. 이 작품에서 윤리란 거창한 원리가 아니라, 이웃의 고통에 반응하는 작은 행위에서 나타난다. 김 시인의 선택은 합리적 이익 계산을 넘어서는 것이며, 시를 쓴다는 행위 자체가 삶에 품위를 더하는 윤리적 실천이다.}
\begin{kfChoices}
\kfChoice 김 시인이 넉넉하지 않으면서도 원고료를 내준 것은 합리적 이익 계산을 넘어서는 윤리적 선택이겠군.
\kfChoice '외상값을 갚지 못하는 사람들을 탓할 수도 없다'는 김 시인의 태도는 이웃의 사정에 대한 공감을 보여주는군.
\kfChoice 김 시인이 시를 통해 '삶에 한 줄의 품위를 더하는 것'이라고 생각한 것은 시 쓰기의 윤리적 의미를 드러내는군.
\kfChoice 김 시인의 행위는 경제적 보상을 기대한 투자이므로, 이타적 동기가 아닌 이기적 동기에 의한 것이겠군.
\kfChoice 최씨 아줌마의 아이가 인사하는 장면은 작은 선의가 공동체 안에서 교류되는 모습을 보여주는군.
\end{kfChoices}
\end{kfQBlock}

\begin{kfQBlock}{5}{}
\kfQStem{김 시인이 돈을 내준 뒤 '마음이 가벼웠다'고 느낀 것이 보여주는 바로 가장 적절한 것은?}
\begin{kfChoices}
\kfChoice 경제적 부담에서 비로소 완전히 벗어났다는 큰 해방감
\kfChoice 도덕적 선택에 대한 내적 충족감과 품위의 자부심
\kfChoice 아내에게 자기 행동의 비밀을 감출 수 있다는 안도감
\kfChoice 최씨 아줌마가 빌린 돈을 반드시 갚을 것이라는 확신
\kfChoice 더 이상 그 일을 두고 고민할 필요가 없다는 무관심
\end{kfChoices}
\end{kfQBlock}

\end{multicols}

\kfSecHeader{kfAreaLiterature}{지문}{수궁가 — 작자 미상 (판소리) / 판소리 「수궁가」, 수능·학평 기출 다수}

\kfPassageNote{수궁가 — 작자 미상 (판소리) / 판소리 「수궁가」, 수능·학평 기출 다수}{%
별주부 수궁 나서며 탄식하되, "내 비록 미물이나 충성이 뼈에 사무쳤으니, 어찌 주상의 환후를 모른 체하리오. 토끼의 간이라면 천리를 밟아서라도 구해 오리라." 하고 육지로 나아가니라.

토끼를 만나 별주부 가로되, "수궁은 가히 극락세계라. 산해진미 진수성찬이 날마다 벌여 있고, 비단 보배가 방마다 가득하니, 그대 같은 호걸을 기다린 지 오래니라." 토끼 듣고 제 마음에 혹하여, "과연 그러한가. 내 한번 가 보리라."

수궁에 이르러 용왕 앞에 꿇려 놓으매, 용왕 가로되, "토끼의 간을 내어 약을 지어라." 토끼 크게 놀라 죽을 판이라. 그러나 토끼 급히 꾀를 내어 아뢰되, "소신이 간을 빼어 깨끗이 씻어 바위 위에 말려 놓고 왔사오니, 다시 육지에 보내 주시면 가져오겠나이다." 용왕이 그 말을 곧이듣고 토끼를 다시 육지로 보내니, 토끼 물 밖에 나서자 깡충깡충 뛰며 돌아보고 하는 말이, "에라, 이 못난 자라야. 세상에 간을 빼어 놓고 사는 놈이 어디 있더냐." 하고 산으로 달아나니라.
}
\kfSecHeader{kfAreaLiterature}{활동}{작품 정리표 — 수궁가(작자 미상)}

\noindent\textit{\small 빈칸에 알맞은 말을 채워 넣으시오. (초성 힌트 참고)}\par\medskip
\begin{kfActTable}{|>{\centering\arraybackslash}m{2.6cm}|m{6cm}|m{4cm}|}
\rowcolor{kfPrimaryLight!50}\textbf{\color{kfPrimary}항목} & \textbf{\color{kfPrimary}내용 (빈칸 채우기)} & \textbf{\color{kfPrimary}답안 작성}\\\hline
갈래 & ㅍㅅㄹ (5마당 중 하나) ① & \kfTBlank \\\hline
원전 & ㅌㄲㅈ ② & \kfTBlank \\\hline
별주부 동기 & 용왕의 ㅎㅎ를 고치기 위한 ㅊㅅ ③ & \kfTBlank \\\hline
유인 수법 & 수궁을 ㄱㄹ ㅅㄱ로 과장 = ㄱㅇㅇㅅ ④ & \kfTBlank \\\hline
토끼의 약점 & 별주부의 말에 ㅎㅎㅇ = ㅇㅁㅎ ⑤ & \kfTBlank \\\hline
토끼의 기지 & ㄱ을 빼어 바위에 ㅁㄹ 놓고 왔다는 ㄱㅈㅁ ⑥ & \kfTBlank \\\hline
풍자 대상 1 & ㅇㅇ의 어리석음 (곧이듣고 보내줌) ⑦ & \kfTBlank \\\hline
풍자 대상 2 & ㅂㅈㅂ의 맹목적 ㅊㅅ ⑧ & \kfTBlank \\\hline
해학적 장면 & 'ㅇㄹ, 이 못난 ㅈㄹ야' ⑨ & \kfTBlank \\\hline
윤리적 쟁점 1 & ㅁㅈ이 ㅅㄷ을 정당화하는가 ⑩ & \kfTBlank \\\hline
윤리적 쟁점 2 & ㅅㅈ을 위한 ㄱㅈㅁ의 정당성 ⑪ & \kfTBlank \\\hline
주제 & ㅊㅇ와 ㅈㅎ의 갈등, ㄱㄹㅈ ㅂㄷ과 ㅇㅈ ㄱㅈ ⑫ & \kfTBlank \\\hline
\end{kfActTable}
\kfSecHeader{kfAreaLiterature}{문제}{}

\begin{multicols}{2}\raggedcolumns\setlength{\columnsep}{12pt}
\begin{kfQBlock}{1}{}
\kfQStem{윗글에 대한 설명으로 적절한 것은?}
\begin{kfChoices}
\kfChoice 판소리 특유의 운문체와 구어체가 혼합된 서술 방식이 보인다.
\kfChoice 인물 간의 심리적 갈등이 인물의 내면 독백으로만 전달된다.
\kfChoice 사건이 시간의 역순으로 배치되어 강한 긴장감을 조성한다.
\kfChoice 배경에 대한 자세하고 상세한 묘사가 서사의 중심을 이룬다.
\kfChoice 인물들 모두가 진실만을 말하며 정직한 자세로 행동하고 있다.
\end{kfChoices}
\end{kfQBlock}

\begin{kfQBlock}{2}{}
\kfQStem{별주부가 토끼에게 '수궁은 가히 극락세계라'라고 한 것의 서사적 기능으로 가장 적절한 것은?}
\begin{kfChoices}
\kfChoice 수궁의 실제 모습을 사실적으로 전달하는 정보 제공
\kfChoice 토끼의 욕심을 자극하여 수궁으로 유인하기 위한 감언이설
\kfChoice 별주부가 토끼에게 진심으로 호의를 베푸는 장면
\kfChoice 용왕의 위엄을 강조하여 토끼에게 경외심을 심어주는 장면
\kfChoice 수궁의 환경 문제를 비판하는 풍자적 표현
\end{kfChoices}
\end{kfQBlock}

\begin{kfQBlock}{3}{}
\kfQStem{토끼가 '간을 빼어 바위 위에 말려 놓고 왔다'는 거짓말의 윤리적 성격에 대한 평가로 가장 적절한 것은?}
\begin{kfChoices}
\kfChoice 생존을 위한 불가피한 거짓말이지만, 거짓말의 윤리적 정당성 문제를 제기한다.
\kfChoice 토끼가 별주부를 돕기 위해 한 선의의 거짓말이다.
\kfChoice 토끼의 거짓말은 도덕적으로 완전히 정당하며 비판의 여지가 없다.
\kfChoice 토끼의 거짓말은 용왕의 명을 어기는 불충이므로 도덕적으로 잘못된 것이다.
\kfChoice 토끼는 거짓말을 한 것이 아니라 진실을 말한 것이다.
\end{kfChoices}
\end{kfQBlock}

\begin{kfQBlock}{4}{}
\kfQStem{윗글을 바탕으로 <보기>를 감상한 내용으로 적절하지 않은 것은?

<보기> 「수궁가」는 충의와 지혜의 갈등, 목적과 수단의 윤리적 문제를 다룬다. 별주부의 충성은 유교적 덕목이지만, 그것을 실현하기 위해 무고한 존재를 속이고 해치는 것은 수단의 정당성 문제를 제기한다. 한편 토끼의 기지는 약자의 생존 전략이면서도 거짓말이라는 윤리적 문제를 안고 있다.}
\begin{kfChoices}
\kfChoice 별주부가 토끼를 속여 수궁으로 데려간 것은 충성이라는 목적을 위해 속임이라는 수단을 사용한 것이겠군.
\kfChoice 별주부의 충성심은 유교적 덕목이지만, 무고한 토끼를 해치는 수단이 정당한지는 별개의 문제이겠군.
\kfChoice 토끼의 기지는 약자의 생존 전략이면서도 거짓말이라는 윤리적 한계를 지니는군.
\kfChoice 용왕이 토끼의 거짓말을 '곧이듣'는 것은 권력자의 어리석음을 풍자하는 것이겠군.
\kfChoice 별주부의 충성은 아무런 윤리적 문제가 없는 완벽한 도덕적 행위이겠군.
\end{kfChoices}
\end{kfQBlock}

\begin{kfQBlock}{5}{}
\kfQStem{토끼가 물 밖에 나서 '에라, 이 못난 자라야. 세상에 간을 빼어 놓고 사는 놈이 어디 있더냐'라고 한 것의 효과로 가장 적절한 것은?}
\begin{kfChoices}
\kfChoice 해학적 어조로 긴장을 풀면서 별주부의 어리석음을 풍자한다.
\kfChoice 토끼가 별주부에게 자신의 잘못을 진심으로 사과하는 장면이다.
\kfChoice 토끼가 머지않아 다시 수궁으로 돌아갈 의사가 있음을 강하게 암시한다.
\kfChoice 토끼와 별주부 사이의 깊은 우정과 신뢰를 새삼 확인하는 대사이다.
\kfChoice 토끼가 자신의 거짓말과 행동에 대해 깊이 반성하는 독백이다.
\end{kfChoices}
\end{kfQBlock}

\end{multicols}

\kfStartNonlit{이번 챕터 비문학 지문 5편}


\kfSecHeader{kfAreaNonlit}{지문}{공리주의의 세 갈래 (인문)}

\kfPassageNote{공리주의의 세 갈래 (인문)}{%
공리주의(utilitarianism)는 행위의 도덕적 옳고 그름을 그 행위가 산출하는 결과, 특히 쾌락(행복)과 고통의 총량에 의해 판단하는 윤리 이론이다. 공리주의의 핵심 원리는 '최대 다수의 최대 행복'으로 요약된다. 이 입장에 따르면, 도덕적으로 옳은 행위란 관련된 모든 사람의 행복을 극대화하는 행위이다. 공리주의는 결과를 기준으로 판단하므로 결과주의(consequentialism)의 대표적 형태이며, 행위의 동기나 의무를 우선시하는 의무론과 대조된다.

제러미 벤담은 근대 공리주의를 체계화한 인물이다. 벤담은 쾌락과 고통이 인간 행동의 두 가지 지배자라고 보았으며, 쾌락의 양을 계산하기 위해 '쾌락 계산법(felicific calculus)'을 제시했다. 이 계산법은 쾌락의 강도, 지속성, 확실성, 원근성, 다산성(추가 쾌락을 낳는 정도), 순수성(고통이 섞이지 않는 정도), 범위(영향받는 사람의 수) 등 일곱 가지 기준으로 쾌락의 양을 측정한다. 벤담에게 쾌락은 질적으로 동일하며, 오직 양적으로만 비교된다. 따라서 그는 '편견을 배제하면, 핀볼 놀이의 쾌락이 시 읽기의 쾌락과 같은 양이라면 둘의 가치는 동등하다'라고 주장했다.

존 스튜어트 밀은 벤담의 양적 공리주의를 수정하여 '질적 공리주의'를 제시했다. 밀은 쾌락 사이에 질적 차이가 있다고 주장했다. 그에 따르면, 지적 쾌락이나 도덕적 감정에서 오는 쾌락은 단순한 신체적 쾌락보다 질적으로 우월하다. 밀은 이를 '배부른 돼지보다 배고픈 소크라테스가 낫다'라는 유명한 표현으로 요약했다. 두 종류의 쾌락을 모두 경험해 본 '유능한 판정자'만이 어떤 쾌락이 질적으로 우월한지 판단할 수 있다고 밀은 보았다. 이 수정은 벤담의 이론이 저급한 쾌락을 지나치게 옹호한다는 비판에 대응한 것이다.

그러나 벤담과 밀의 공리주의는 공통적으로 '행위 공리주의(act utilitarianism)'에 속한다. 행위 공리주의는 개별 행위마다 그 결과를 계산하여 최대 행복을 산출하는 행위를 선택해야 한다고 본다. 이에 대해 '규칙 공리주의(rule utilitarianism)'는 다른 접근을 취한다. 규칙 공리주의에 따르면, 도덕적 판단의 기준은 개별 행위가 아니라 규칙이다. 어떤 규칙을 일반적으로 따르는 것이 최대 행복을 산출하는가를 물어야 한다. 예를 들어, 약속을 어기는 개별 행위가 일시적으로 행복을 증가시키더라도, '약속을 지키라'는 규칙을 일반적으로 따르는 것이 장기적으로 더 큰 사회적 효용을 산출한다면, 약속을 지키는 것이 도덕적으로 옳다.

규칙 공리주의의 강점은 행위 공리주의가 부딪히는 난점들을 해소한다는 데 있다. 행위 공리주의는 극단적 상황에서 무고한 사람을 희생시키는 것도 정당화할 수 있다는 비판을 받는다. 규칙 공리주의는 '무고한 사람을 희생시키지 말라'는 규칙이 장기적으로 최대 효용을 산출하므로 그러한 행위를 금지한다. 그러나 규칙 공리주의 역시 비판에서 자유롭지 못하다. 규칙을 따르는 것이 특정 상황에서 명백히 나쁜 결과를 낳는 경우에도 규칙을 고수해야 하는지의 문제가 남기 때문이다. 이 경우 규칙 공리주의는 결국 행위 공리주의로 회귀하거나, 규칙 자체가 유연해져야 한다는 딜레마에 빠진다.
}
\kfSecHeader{kfAreaNonlit}{문제}{}

\begin{multicols}{2}\raggedcolumns\setlength{\columnsep}{12pt}
\begin{kfQBlock}{1}{}
\kfQStem{윗글의 내용과 일치하지 않는 것은?}
\begin{kfChoices}
\kfChoice 벤담은 쾌락의 양을 일곱 가지 기준으로 측정하는 쾌락 계산법을 제시했다.
\kfChoice 밀은 지적 쾌락이 신체적 쾌락보다 질적으로 우월하다고 보았다.
\kfChoice 벤담과 밀의 공리주의는 모두 규칙 공리주의에 속한다.
\kfChoice 규칙 공리주의는 개별 행위가 아니라 규칙의 일반적 준수를 기준으로 판단한다.
\kfChoice 행위 공리주의는 극단적 상황에서 무고한 사람의 희생을 정당화할 수 있다는 비판을 받는다.
\end{kfChoices}
\end{kfQBlock}

\begin{kfQBlock}{2}{}
\kfQStem{밀이 벤담의 공리주의를 수정한 핵심 이유로 가장 적절한 것은?}
\begin{kfChoices}
\kfChoice 벤담의 이론이 쾌락의 양을 측정할 방법을 제시하지 못했기 때문
\kfChoice 벤담의 이론이 결과가 아닌 동기를 기준으로 판단했기 때문
\kfChoice 벤담의 이론이 저급한 쾌락을 옹호한다는 비판에 대응하기 위함
\kfChoice 벤담의 이론이 규칙의 일반적 준수를 고려하지 않았기 때문
\kfChoice 벤담의 이론이 고통의 최소화만을 목표로 했기 때문
\end{kfChoices}
\end{kfQBlock}

\begin{kfQBlock}{3}{}
\kfQStem{윗글을 바탕으로 <보기>를 이해한 내용으로 적절한 것은?

<보기> 한 의사가 건강한 사람 한 명의 장기를 적출하여 다섯 명의 생명을 구할 수 있는 상황에 처해 있다.}
\begin{kfChoices}
\kfChoice 벤담의 입장에서 장기 적출은 언제나 부도덕하다.
\kfChoice 밀의 입장에서 장기 적출은 질적으로 고급한 쾌락을 산출하므로 정당화된다.
\kfChoice 행위 공리주의에 따르면 다섯 명의 생명을 구하는 것이 총 행복량이 크므로 장기 적출이 정당화될 수 있다.
\kfChoice 규칙 공리주의에 따르면 '무고한 사람을 희생시키지 말라'는 규칙이 장기적으로 효용이 크므로 장기 적출을 금지한다.
\kfChoice 규칙 공리주의는 이 상황에서 규칙을 무시하고 행위 공리주의와 동일한 결론을 내린다.
\end{kfChoices}
\end{kfQBlock}

\begin{kfQBlock}{4}{}
\kfQStem{윗글에서 규칙 공리주의가 직면하는 딜레마로 가장 적절한 것은?}
\begin{kfChoices}
\kfChoice 쾌락의 질적 차이를 인정할 경우 양적 계산이 불가능해지는 문제
\kfChoice 규칙을 따르는 것이 특정 상황에서 나쁜 결과를 낳을 때 규칙을 고수할 것인가의 문제
\kfChoice 최대 다수의 행복이 소수의 행복과 충돌할 때 어느 쪽을 선택할 것인가의 문제
\kfChoice 행위의 동기와 결과 중 어느 것을 도덕 판단의 기준으로 삼을 것인가의 문제
\kfChoice 유능한 판정자를 누가 선정할 것인가의 실천적 문제
\end{kfChoices}
\end{kfQBlock}

\end{multicols}

\kfSecHeader{kfAreaNonlit}{지문}{샤프츠베리와 듀이의 도덕감 (인문)}

\kfPassageNote{샤프츠베리와 듀이의 도덕감 (인문)}{%
도덕적 판단의 근거는 무엇인가? 이성인가, 감정인가, 아니면 경험인가? 이 물음에 대해 서로 다른 시대와 전통에서 독특한 답을 내놓은 두 사상가가 있다. 18세기 영국의 샤프츠베리(3대 백작)와 20세기 미국의 존 듀이이다.

샤프츠베리는 인간에게 선천적인 '도덕 감각(moral sense)'이 있다고 주장했다. 이 도덕 감각은 외부 세계를 지각하는 시각이나 청각과 마찬가지로, 도덕적 선악을 직접 감지하는 내적 능력이다. 샤프츠베리에 따르면, 우리가 어떤 행위를 보고 즉각적으로 '아름답다' 또는 '추하다'고 느끼는 것은 이 도덕 감각의 작용이다. 여기서 주목할 점은 샤프츠베리가 도덕적 판단을 미적 판단에 비유했다는 것이다. 아름다운 음악을 듣고 감동받듯이, 선한 행위를 보고 내적으로 승인하는 감정이 자연스럽게 일어난다는 것이다. 이 때문에 샤프츠베리의 윤리학은 '미적 도덕론'이라고도 불린다.

샤프츠베리의 도덕 감각론은 당시 홉스의 이기주의 인간관에 대한 반론으로 제시되었다. 홉스는 인간이 본질적으로 이기적이어서 사회 계약 없이는 도덕적 질서가 불가능하다고 보았다. 샤프츠베리는 이에 맞서, 인간에게는 타인의 이익을 자연스럽게 배려하고 그에 기뻐하는 선천적 감정이 있다고 주장했다. 도덕은 외부에서 강제되는 것이 아니라 인간 본성 안에 이미 내재해 있다는 것이다. 다만 이 도덕 감각은 교육과 문화를 통해 세련되어야 하며, 왜곡된 환경에서는 둔화될 수 있다.

존 듀이는 샤프츠베리와 전혀 다른 전통, 즉 미국 프래그머티즘(실용주의)에서 도덕의 문제에 접근했다. 듀이에게 도덕은 선천적 감각이나 고정된 원리에 기반하는 것이 아니라, 구체적인 문제 상황 속에서 경험적으로 성장해 가는 것이다. 듀이는 전통 윤리학이 절대적이고 영원한 도덕 원리를 추구하는 것을 비판하며, 도덕적 판단은 과학적 탐구와 마찬가지로 가설을 세우고 실험하고 결과를 관찰하는 과정이라고 보았다. 이것이 듀이의 '도구주의(instrumentalism)' 윤리학이다.

듀이에게 '도덕감'이란 샤프츠베리처럼 선천적으로 주어진 감각이 아니라, 경험을 통해 형성되는 습관적 감수성이다. 한 사람이 타인의 고통에 민감하게 반응하는 것은 타고난 능력이 아니라, 다양한 사회적 상호작용과 반성적 경험을 통해 형성된 것이다. 듀이는 이러한 도덕적 감수성이 고정된 것이 아니라 끊임없이 성장할 수 있다고 보았다. 그에게 도덕적 성장은 새로운 상황에 대처하는 능력의 확장이며, 이는 민주적 공동체 안에서 다양한 관점을 접하고 소통할 때 가장 잘 이루어진다. 결국 샤프츠베리가 도덕감을 인간 본성의 선천적 장비로 보았다면, 듀이는 경험적 성장의 산물로 본 셈이다.
}
\kfSecHeader{kfAreaNonlit}{문제}{}

\begin{multicols}{2}\raggedcolumns\setlength{\columnsep}{12pt}
\begin{kfQBlock}{1}{}
\kfQStem{윗글의 내용과 일치하는 것은?}
\begin{kfChoices}
\kfChoice 샤프츠베리는 도덕 감각이 교육 없이도 항상 완벽하게 작동한다고 보았다.
\kfChoice 샤프츠베리의 도덕 감각론은 홉스의 이기주의 인간관에 대한 반론으로 제시되었다.
\kfChoice 듀이는 절대적이고 영원한 도덕 원리의 존재를 긍정했다.
\kfChoice 듀이에게 도덕감은 선천적으로 주어진 능력이다.
\kfChoice 샤프츠베리와 듀이 모두 도덕이 외부에서 강제된다고 보았다.
\end{kfChoices}
\end{kfQBlock}

\begin{kfQBlock}{2}{}
\kfQStem{윗글에서 샤프츠베리와 듀이의 공통점과 차이점을 가장 적절하게 설명한 것은?}
\begin{kfChoices}
\kfChoice 둘 다 도덕감을 인정하나, 샤프츠베리는 선천적, 듀이는 경험적 성장으로 본다.
\kfChoice 두 사람 모두 도덕감의 존재를 부정하며, 이성만을 도덕 판단의 근거로 삼는다.
\kfChoice 샤프츠베리는 경험적 성장을, 듀이는 선천적 감각을 도덕의 핵심 근거로 본다.
\kfChoice 두 사람 모두 홉스의 이기주의를 계승하여 도덕의 외부적 강제를 주장한다.
\kfChoice 두 사람 모두 절대적 도덕 원리를 추구하며 그 내용에서만 차이를 보인다.
\end{kfChoices}
\end{kfQBlock}

\begin{kfQBlock}{3}{}
\kfQStem{윗글을 바탕으로 <보기>를 이해한 내용으로 적절하지 않은 것은?

<보기> A는 어려서부터 타인의 고통에 무관심했지만, 자원봉사 활동을 통해 점차 타인의 고통에 민감하게 반응하게 되었다.}
\begin{kfChoices}
\kfChoice 듀이는 A의 변화를 도덕적 감수성이 경험을 통해 성장한 사례로 볼 것이다.
\kfChoice 듀이는 A의 자원봉사가 새로운 상황에 대처하는 능력의 확장에 해당한다고 볼 것이다.
\kfChoice 샤프츠베리는 A가 처음에 무관심했던 것을 도덕 감각이 왜곡된 환경에서 둔화된 것으로 설명할 수 있을 것이다.
\kfChoice 샤프츠베리는 A의 도덕 감각이 자원봉사를 통해 처음으로 생겨난 것이라고 볼 것이다.
\kfChoice 듀이는 A의 변화가 민주적 공동체 안에서의 소통과 관련될 수 있다고 볼 것이다.
\end{kfChoices}
\end{kfQBlock}

\begin{kfQBlock}{4}{}
\kfQStem{윗글에서 듀이의 '도구주의 윤리학'이 의미하는 바로 가장 적절한 것은?}
\begin{kfChoices}
\kfChoice 도덕은 사회적 통제를 위한 강제적 도구에 불과하다고 보는 냉소적 관점
\kfChoice 도덕 판단도 과학처럼 가설·실험·관찰을 거쳐 경험적으로 형성된다는 관점
\kfChoice 도덕 원리는 영원불변하므로 실천의 도구로 늘 삼아야 한다는 관점
\kfChoice 물리적 도구를 잘 만들어 활용하는 것이 윤리적으로 가장 중요하다는 관점
\kfChoice 도덕 감각은 본래 선천적으로 주어진 내적 도구라고 규정하는 관점
\end{kfChoices}
\end{kfQBlock}

\end{multicols}

\kfSecHeader{kfAreaNonlit}{지문}{다산 정약용의 윤리학 (인문)}

\kfPassageNote{다산 정약용의 윤리학 (인문)}{%
조선 후기의 실학자 정약용(다산, 1762\textasciitilde{}1836)은 당대 주류였던 주자학의 인성론을 비판하면서 독자적인 윤리 체계를 구축했다. 주자학에서 인간의 본성(性)은 하늘로부터 부여받은 '리(理)'로서, 선한 것이 본래부터 내재해 있다고 본다. 주자학의 이러한 인성론에 따르면, 도덕적 수양이란 이미 갖추어진 본성을 회복하는 것, 즉 기질의 가림을 걷어내어 본래의 선한 성을 드러내는 과정이다.

정약용은 이러한 주자학적 인성론에 근본적인 의문을 제기했다. 그의 핵심 주장이 '성기호설(性嗜好說)'이다. 정약용에 따르면, 인간의 본성은 '리'와 같은 추상적 원리가 아니라 구체적인 '기호(嗜好)', 즉 좋아하고 싫어하는 경향성이다. 인간은 선한 것을 좋아하고 악한 것을 싫어하는 기호를 가지고 있다. 그런데 이 기호는 고정된 실체가 아니라 '경향성'이다. 인간은 선을 좋아하는 경향을 가지지만, 동시에 악을 행할 수도 있다. 이 점에서 정약용의 인성론은 주자학의 '성즉리(性卽理)'와 근본적으로 다르다.

정약용 윤리학의 또 다른 핵심은 '자주지권(自主之權)'이다. 자주지권이란 인간이 선과 악 사이에서 스스로 선택할 수 있는 자유 의지를 말한다. 주자학에서 선한 본성이 이미 주어져 있다면, 악의 발생은 기질에 의한 가림으로 설명된다. 그러나 정약용에게 선과 악은 인간의 자유로운 선택의 결과이다. 인간은 선을 좋아하는 기호를 갖되, 실제로 선을 행할지 악을 행할지는 자주지권에 의해 결정된다. 이 자유로운 선택이야말로 인간을 도덕적 존재로 만드는 근거이다. 만약 선이 본성에 이미 내재해 있어 자동으로 발현된다면, 도덕적 노력의 의미는 퇴색될 것이다.

정약용은 이러한 기초 위에서 실천 윤리를 강조했다. 그에게 도덕은 내면의 본성을 관조하는 것이 아니라, 구체적인 인간관계 속에서 '행사(行事)'하는 것이다. 정약용은 효(孝), 제(悌), 자(慈) 등의 덕목을 관념이 아닌 실천적 행위로 재정의했다. 효는 부모를 봉양하는 구체적 행위이고, 자는 자녀를 사랑하는 구체적 행위이다. 마음속으로 선한 생각만 하고 행동하지 않는다면 그것은 진정한 도덕이 아니다. 이 점에서 정약용은 주자학이 내면의 수양에 치우쳐 실천을 소홀히 한다고 비판했다.

정약용의 윤리학은 동아시아 유학 전통 내에서 독특한 위치를 차지한다. 그는 유학의 핵심 가치를 계승하면서도, 주자학의 형이상학적 틀을 거부하고 경험적·실천적 방향으로 윤리학을 재구성했다. 성기호설은 인성을 추상적 원리에서 구체적 경향성으로, 자주지권은 도덕의 근거를 타고난 본성에서 자유 의지로, 행사의 강조는 수양의 초점을 관조에서 실천으로 전환한 것이다. 이러한 전환은 근대적 자율 윤리학과도 상통하는 면이 있다.
}
\kfSecHeader{kfAreaNonlit}{문제}{}

\begin{multicols}{2}\raggedcolumns\setlength{\columnsep}{12pt}
\begin{kfQBlock}{1}{}
\kfQStem{윗글의 내용과 일치하지 않는 것은?}
\begin{kfChoices}
\kfChoice 주자학에서 인간의 본성은 선한 것이 본래부터 내재해 있다.
\kfChoice 정약용의 성기호설에 따르면, 본성은 추상적 원리가 아니라 좋아하고 싫어하는 경향성이다.
\kfChoice 정약용은 인간이 선과 악 사이에서 자유롭게 선택할 수 있다고 보았다.
\kfChoice 정약용은 주자학의 형이상학적 틀을 계승하면서 덕목의 의미를 확장했다.
\kfChoice 정약용에게 효, 제, 자 등의 덕목은 관념이 아니라 실천적 행위이다.
\end{kfChoices}
\end{kfQBlock}

\begin{kfQBlock}{2}{}
\kfQStem{정약용이 '자주지권'을 강조한 이유로 가장 적절한 것은?}
\begin{kfChoices}
\kfChoice 선한 본성이 자동으로 발현되지 않아 자유 의지에 의한 도덕적 선택이 필요하기 때문
\kfChoice 인간의 기질이 본성을 가리고 있어서 기질을 변화시키는 도덕적 수양이 필요하기 때문
\kfChoice 인간은 본래 악한 존재이므로 사회 계약을 통해 도덕을 강제할 필요가 있기 때문
\kfChoice 선한 본성이 이미 인간 안에 완비되어 있으므로 그것을 회복하기만 하면 충분하기 때문
\kfChoice 인간의 도덕적 판단은 반드시 외부의 권위에 기대어야 정확해질 수 있기 때문
\end{kfChoices}
\end{kfQBlock}

\begin{kfQBlock}{3}{}
\kfQStem{정약용이 주자학의 인성론을 비판하는 핵심 논점으로 가장 적절한 것은?}
\begin{kfChoices}
\kfChoice 주자학이 효, 제, 자의 덕목을 도덕의 핵심으로 인정하지 않았다는 점
\kfChoice 본성을 추상적 원리(리)로 규정해 도덕적 노력의 의미를 약화한다는 점
\kfChoice 주자학이 인간의 악한 본성을 끝내 인정하지 않고 외면해 왔다는 점
\kfChoice 경험적 탐구를 통해 도덕 원리를 직접 검증해야 한다고 주장한 점
\kfChoice 주자학이 인간에게 자유 의지가 있다고 적극적으로 주장하고 있다는 점
\end{kfChoices}
\end{kfQBlock}

\begin{kfQBlock}{4}{}
\kfQStem{윗글을 바탕으로 <보기>를 이해한 내용으로 적절하지 않은 것은?

<보기> 갑은 부모가 아프다는 소식을 듣고 마음속으로 걱정은 했지만, 바쁜 업무를 핑계로 방문하지 않았다. 을은 같은 상황에서 즉시 업무를 정리하고 부모를 찾아가 간병했다.}
\begin{kfChoices}
\kfChoice 정약용의 관점에서 갑은 마음속으로 선한 생각을 했으나 행사(실천)가 없으므로 진정한 효를 행한 것이 아니다.
\kfChoice 정약용의 관점에서 을은 자주지권을 발휘하여 선을 실천한 사례이다.
\kfChoice 주자학의 관점에서 갑은 기질이 본성을 가리고 있는 상태로 설명될 수 있다.
\kfChoice 정약용의 관점에서 갑과 을 모두 선을 좋아하는 기호를 가지고 있으나, 실천 여부에서 차이가 난다.
\kfChoice 주자학의 관점에서 을의 행위는 자주지권에 의한 자유로운 선택의 결과이다.
\end{kfChoices}
\end{kfQBlock}

\end{multicols}

\kfSecHeader{kfAreaNonlit}{지문}{롤스의 정의론 (인문)}

\kfPassageNote{롤스의 정의론 (인문)}{%
존 롤스는 1971년 발표한 『정의론(A Theory of Justice)』에서 사회 정의의 원칙을 도출하기 위한 독창적인 사고실험을 제시했다. 롤스의 핵심 물음은 이것이다. 합리적인 사람들이 자유롭고 평등한 조건에서 사회의 기본 구조를 결정한다면, 어떤 정의 원칙에 합의할 것인가?

이 물음에 답하기 위해 롤스가 설정한 가상의 상황이 '원초적 입장(original position)'이다. 원초적 입장에서 사람들은 사회의 기본 규칙을 선택하는 계약 당사자이다. 그런데 이 당사자들에게는 특별한 제약이 부과된다. 바로 '무지의 베일(veil of ignorance)'이다. 무지의 베일 뒤에서 당사자들은 자신의 사회적 지위, 경제적 능력, 재능, 성별, 인종은 물론이고 자신이 어떤 가치관이나 인생 계획을 가지고 있는지도 모른다. 다만 당사자들은 사회와 인간 심리에 대한 일반적 지식은 알고 있으며, 합리적으로 자기 이익을 추구한다.

무지의 베일의 핵심 기능은 공정성을 확보하는 것이다. 자신이 부자인지 가난한 사람인지, 재능이 많은지 적은지 모르는 상태에서 선택하면, 누구도 자신에게 유리한 규칙을 관철시킬 수 없다. 롤스는 이 상황에서 합리적 당사자들이 두 가지 정의 원칙에 합의할 것이라고 논증했다.

제1원칙은 '평등한 자유의 원칙'이다. 모든 사람은 다른 사람의 동일한 자유와 양립 가능한 한에서, 평등한 기본적 자유에 대한 권리를 갖는다. 여기서 기본적 자유에는 양심의 자유, 표현의 자유, 집회·결사의 자유, 참정권 등이 포함된다. 이 원칙은 어떤 경우에도 양보될 수 없으며, 경제적 이익을 위해 기본적 자유를 제한하는 것은 허용되지 않는다.

제2원칙은 두 부분으로 구성된다. 먼저 '기회 균등의 원칙'으로, 사회적·경제적 불평등이 연결된 직위와 직책은 공정한 기회 균등의 조건하에서 모든 사람에게 개방되어야 한다. 다음으로 '차등 원칙(difference principle)'이 있다. 사회적·경제적 불평등은 오직 사회에서 가장 불리한 처지에 있는 구성원에게 최대 이익이 돌아가도록 편성된 경우에만 허용된다. 즉 불평등 자체를 금지하는 것이 아니라, 그 불평등이 최소 수혜자에게 이익이 되는 경우에만 정당화된다. 롤스는 무지의 베일 뒤의 합리적 당사자가 자신이 최소 수혜자가 될 가능성을 고려하여 최악의 경우를 최선으로 만드는 '맥시민(maximin)' 전략을 택할 것이라고 보았다. 이 두 원칙 사이에는 '우선 순위 규칙'이 적용되어, 제1원칙(자유)이 제2원칙(평등)에 우선하고, 기회 균등이 차등 원칙에 우선한다.
}
\kfSecHeader{kfAreaNonlit}{문제}{}

\begin{multicols}{2}\raggedcolumns\setlength{\columnsep}{12pt}
\begin{kfQBlock}{1}{}
\kfQStem{윗글의 내용과 일치하지 않는 것은?}
\begin{kfChoices}
\kfChoice 무지의 베일 뒤에서 당사자들은 자신의 사회적 지위와 재능을 알지 못한다.
\kfChoice 제1원칙에 따르면 경제적 이익을 위해 기본적 자유를 제한할 수 있다.
\kfChoice 차등 원칙은 가장 불리한 처지의 구성원에게 최대 이익이 돌아가는 불평등만 허용한다.
\kfChoice 맥시민 전략은 최악의 경우를 최선으로 만드는 선택 방식이다.
\kfChoice 제1원칙(자유)이 제2원칙(평등)보다 우선한다.
\end{kfChoices}
\end{kfQBlock}

\begin{kfQBlock}{2}{}
\kfQStem{롤스가 '무지의 베일'을 설정한 목적으로 가장 적절한 것은?}
\begin{kfChoices}
\kfChoice 당사자들이 자기 이익을 포기하고 이타적으로 행동하게 하기 위해
\kfChoice 당사자들이 자신에게 유리한 규칙을 관철시킬 수 없게 하여 공정성을 확보하기 위해
\kfChoice 당사자들이 사회에 대한 모든 지식을 갖추고 최적의 결정을 내리게 하기 위해
\kfChoice 당사자들이 기존 사회 체제를 그대로 유지하도록 유도하기 위해
\kfChoice 당사자들 간의 자유로운 토론을 통해 다수결로 원칙을 결정하기 위해
\end{kfChoices}
\end{kfQBlock}

\begin{kfQBlock}{3}{}
\kfQStem{윗글을 바탕으로 <보기>를 이해한 내용으로 적절한 것은?

<보기> 한 사회에서 의사의 소득이 일반 노동자의 5배이다. 이 불평등이 정당화될 수 있는가?}
\begin{kfChoices}
\kfChoice 제1원칙에 따르면, 소득 불평등 자체가 기본적 자유를 침해하므로 부당하다.
\kfChoice 차등 원칙에 따르면, 소득 불평등은 어떤 경우에도 허용되지 않는다.
\kfChoice 차등 원칙에 따르면, 의사의 높은 소득이 최소 수혜자에게 최대 이익이 되는 경우에만 정당화된다.
\kfChoice 무지의 베일에 따르면, 당사자들은 자신이 의사가 될 것을 알고 있으므로 불평등에 동의할 것이다.
\kfChoice 맥시민 전략에 따르면, 가장 유리한 사람의 이익을 극대화해야 한다.
\end{kfChoices}
\end{kfQBlock}

\begin{kfQBlock}{4}{}
\kfQStem{윗글에서 '맥시민(maximin) 전략'이 의미하는 바로 가장 적절한 것은?}
\begin{kfChoices}
\kfChoice 최선의 결과가 나올 수 있는 가능성을 극대화하는 낙관적 전략
\kfChoice 가능한 결과 중 최악의 경우가 가장 나은 선택지를 택하는 전략
\kfChoice 공동체 모든 구성원이 누릴 평균 이익을 극대화하는 전략
\kfChoice 다수의 이익을 위해 소수가 겪을 희생을 감수하게 하는 전략
\kfChoice 사회 전체가 보유한 총 부(富)를 극대화하려는 경제적 전략
\end{kfChoices}
\end{kfQBlock}

\end{multicols}

\kfSecHeader{kfAreaNonlit}{지문}{윤리적 이기주의 vs 이타주의 (인문)}

\kfPassageNote{윤리적 이기주의 vs 이타주의 (인문)}{%
인간이 도덕적으로 행동하는 궁극적 동기는 무엇인가? 이 물음에 대해 이기주의와 이타주의는 정반대의 답을 제시한다. 그러나 이 두 입장은 단일한 이론이 아니라 다양한 갈래를 포함하고 있으며, 각각의 갈래가 도덕적 동기 논쟁에서 고유한 위치를 차지한다.

토마스 홉스의 '심리적 이기주의(psychological egoism)'는 인간의 모든 행위가 궁극적으로 자기 이익에 의해 동기 부여된다는 경험적 주장이다. 심리적 이기주의에 따르면, 심지어 자선 행위나 자기희생마저도 자기 만족, 사회적 평판, 죄책감 회피 등 자기 이익의 추구로 환원된다. 홉스는 인간이 자연 상태에서 '만인의 만인에 대한 투쟁'을 벌인다고 보았으며, 사회 질서는 이기적 개인들 사이의 사회 계약으로만 가능하다고 주장했다. 그러나 심리적 이기주의는 당위(어떻게 행동해야 하는가)가 아니라 사실(인간이 실제로 어떻게 동기 부여되는가)에 대한 주장이므로, 그 자체로는 윤리 이론이 아니다.

에인 랜드는 이기주의를 도덕적 당위의 차원으로 끌어올렸다. 그녀의 '윤리적 이기주의(ethical egoism)'는 각 개인이 자기 이익을 추구하는 것이 도덕적으로 올바르다고 주장한다. 랜드에 따르면, 이타주의는 인간의 합리적 자기 이익을 부정하는 자기 파괴적 윤리이다. 랜드는 자신의 소설 『아틀라스 어깨를 으쓱하다』에서 생산적 개인의 이기적 동기가 사회 전체를 발전시킨다고 역설했다. 그러나 윤리적 이기주의는 심각한 비판에 직면한다. 만약 모든 사람이 자기 이익만을 추구하는 것이 도덕적으로 올바르다면, 타인의 이익을 체계적으로 침해하는 행위도 정당화될 수 있고, 도덕 자체의 존립 근거가 흔들릴 수 있다.

이기주의에 대한 대안으로 이타주의(altruism)가 있다. 이타주의는 타인의 이익이나 복지를 도덕적 행위의 핵심 동기 또는 목표로 삼는다. 최근 '효과적 이타주의(effective altruism)' 운동은 이타적 동기를 실증적 증거와 결합한다. 효과적 이타주의는 단순히 선한 의도를 가지는 것이 아니라, 한정된 자원으로 가능한 한 많은 선을 실현하기 위해 과학적 증거와 합리적 추론에 기반하여 가장 효과적인 방법을 찾아야 한다고 주장한다. 예를 들어, 감정적으로 호소력이 큰 곳에 기부하는 것보다, 기부금 대비 가장 많은 생명을 구할 수 있는 곳에 기부하는 것이 도덕적으로 우월하다고 본다.

이기주의와 이타주의의 논쟁은 도덕적 동기의 본질에 대한 근본적 물음으로 이어진다. 심리적 이기주의는 이타적 동기의 존재 자체를 부정하고, 윤리적 이기주의는 이기적 동기를 도덕적으로 정당화하며, 이타주의는 타인의 이익을 위한 동기가 도덕의 핵심이라고 본다. 그러나 이 대립은 반드시 양자택일일 필요는 없다. '합리적 이타주의'의 관점에서 보면, 장기적으로 타인의 이익을 증진하는 것이 자기 이익과도 합치할 수 있다. 협력과 신뢰에 기반한 사회가 이기적 경쟁에 기반한 사회보다 모든 구성원에게 더 나은 결과를 가져올 수 있기 때문이다.
}
\kfSecHeader{kfAreaNonlit}{문제}{}

\begin{multicols}{2}\raggedcolumns\setlength{\columnsep}{12pt}
\begin{kfQBlock}{1}{}
\kfQStem{윗글의 내용과 일치하는 것은?}
\begin{kfChoices}
\kfChoice 심리적 이기주의는 인간이 이타적으로 행동해야 한다는 도덕적 당위를 주장한다.
\kfChoice 윤리적 이기주의는 자기 이익 추구가 도덕적으로 올바르다고 주장한다.
\kfChoice 효과적 이타주의는 선한 의도만으로 충분하며 결과는 중요하지 않다고 본다.
\kfChoice 홉스는 이타주의에 기반한 사회 질서를 주장했다.
\kfChoice 랜드는 이타주의가 사회를 발전시키는 원동력이라고 보았다.
\end{kfChoices}
\end{kfQBlock}

\begin{kfQBlock}{2}{}
\kfQStem{윗글에 따를 때, 심리적 이기주의와 윤리적 이기주의의 가장 핵심적인 차이는?}
\begin{kfChoices}
\kfChoice 심리적 이기주의는 이기적 동기를 주장하고, 윤리적 이기주의는 이타적 동기를 주장한다는 점
\kfChoice 심리적 이기주의는 인간 동기에 대한 사실 주장이고, 윤리적 이기주의는 도덕적 당위 주장이다.
\kfChoice 심리적 이기주의는 개인에게 적용되는 이론이고, 윤리적 이기주의는 사회 차원에 적용된다.
\kfChoice 심리적 이기주의는 랜드가 대표하고, 윤리적 이기주의는 홉스가 대표한다는 차이가 있다.
\kfChoice 심리적 이기주의는 사회 계약을 부정하고, 윤리적 이기주의는 사회 계약을 긍정한다는 점
\end{kfChoices}
\end{kfQBlock}

\begin{kfQBlock}{3}{}
\kfQStem{윗글을 바탕으로 <보기>를 이해한 내용으로 적절하지 않은 것은?

<보기> 갑은 아프리카 기아 지역에 대한 감동적인 영상을 보고 즉시 그 단체에 100만 원을 기부했다. 을은 같은 100만 원으로 어떤 단체에 기부하는 것이 가장 많은 생명을 구할 수 있는지 비교 분석한 뒤 기부했다.}
\begin{kfChoices}
\kfChoice 효과적 이타주의의 관점에서 을의 기부 방식이 갑보다 도덕적으로 우월하다고 평가될 수 있다.
\kfChoice 심리적 이기주의의 관점에서 갑의 기부도 자기 만족이라는 이기적 동기로 환원될 수 있다.
\kfChoice 효과적 이타주의는 갑과 을 모두 선한 의도를 가졌으므로 동등하게 평가할 것이다.
\kfChoice 윤리적 이기주의의 관점에서 기부 자체가 자기 이익에 반하므로 도덕적으로 바람직하지 않을 수 있다.
\kfChoice 합리적 이타주의의 관점에서 을의 기부는 타인의 이익 증진이 장기적으로 자기 이익과도 합치하는 사례가 될 수 있다.
\end{kfChoices}
\end{kfQBlock}

\begin{kfQBlock}{4}{}
\kfQStem{윗글의 마지막 문단에서 이기주의와 이타주의의 대립을 넘어서려는 관점으로 가장 적절한 것은?}
\begin{kfChoices}
\kfChoice 장기적으로 타인의 이익 증진이 자기 이익과 합치할 수 있다는 합리적 이타주의
\kfChoice 모든 이기적 행위를 일체 금지하고 이타적 행위만을 강제하는 윤리적 이타주의
\kfChoice 심리적 이기주의의 관점에서 이타적 동기를 완전히 부정하는 것
\kfChoice 사회 계약을 통해 이기적 개인들을 강제로 협력하게 하는 홉스적 방식
\kfChoice 개인의 생산적 이기심이 사회를 발전시킨다는 랜드의 입장
\end{kfChoices}
\end{kfQBlock}

\end{multicols}

\kfStartWeekly{패턴 워크북 — 총 ?문항}


\kfSecHeader{kfAreaWeekly}{지문}{문항 1 지문 / 섞어치기}

\kfPassageNote{문항 1 지문 / 섞어치기}{%
식물은 동물과 달리 생애 동안 지속적으로 성장하는 개체이며, 이러한 생장은 분열 조직의 활동에 의해 유지된다. 식물의 분열 조직은 크게 측생 분열 조직과 정단 분열 조직으로 나뉘는데, 측생 분열 조직은 줄기와 뿌리의 측면에서 부피 생장을, 정단 분열 조직은 줄기와 뿌리 끝에서 길이 생장을 담당한다.

측생 분열 조직의 대표로는 관다발 형성층과 코르크 형성층이 있다. 관다발 형성층은 물관과 체관 사이에서 새로운 관다발 조직을 만들고, 코르크 형성층은 표피 아래에서 보호 조직을 형성해 수분 손실을 ⓐ막는 역할을 한다. 이들의 지속적 활동은 줄기와 뿌리를 굵고 단단하게 유지하며 나이테 형성과도 관련된다.

(중략)
}
\begin{kfQBlock}{1}{}
\kfQStem{윗글의 내용과 일치하지 않는 것은?}
\begin{kfChoices}
\kfChoice 코르크 형성층은 수분 손실 억제와 무관하다.
\kfChoice 측생 분열 조직은 부피 생장과 연관된다.
\kfChoice 정단 분열 조직은 길이 생장과 연관된다.
\end{kfChoices}
\end{kfQBlock}

\kfSecHeader{kfAreaWeekly}{지문}{문항 2 지문 / 부정상대어}

\kfPassageNote{문항 2 지문 / 부정상대어}{%
최근까지 일반적인 컴퓨터 환경에서는 운영 체제를 제어하거나 사용자 인터페이스를 처리하고, 문서를 작성하거나 웹을 탐색하는 등 다양한 범용 작업이 주요 연산 대상이었다. 이러한 작업은 명령어의 종류가 많고 그 순서가 수시로 바뀌며, 연산 간 종속성이 크다는 특징이 있다. 컴퓨터의 연산은 주로 여러 개의 코어로 구성된 ㉠중앙 처리 장치(CPU)에서 수행된다. 코어는 컴퓨터의 두뇌 역할을 하는 것으로 명령어를 해석하고 각종 연산을 실행한다. CPU는 여러 개의 고성능 코어로 구성되어 복잡한 제어 흐름을 정밀하게 조율할 수 있도록 설계되었다. 그러나 인공 지능 학습, 고해상도 영상 처리, 과학 기술 관련 계산 등 대규모 데이터를 대상으로 동일한 연산을 반복 수행해야 하는 작업이 증가하면서 기존의 연산 환경에 변화가 생겼다. CPU는 코어 수에 한계가 있고 개별 코어가 복잡한 제어 흐름에 맞추어 다양한 연산을 처리하도록 설계되어 있기 때문에 동일한 연산을 대량으로 반복 수행해야 하는 현대의 

(중략)
}
\begin{kfQBlock}{2}{}
\kfQStem{윗글에 대한 설명으로 가장 적절한 것은?}
\begin{kfChoices}
\kfChoice 이 글은 연산 환경 변화 속 GPU 활용 확대와 CPU-GPU 역할 분담의 필요성을 설명한다.
\kfChoice 이 글은 GPU 내부 레지스터 규격만을 독립적으로 설명한다.
\kfChoice 이 글은 CPU를 완전히 대체하는 단일 장치 전략만을 주장한다.
\end{kfChoices}
\end{kfQBlock}

\kfSecHeader{kfAreaWeekly}{지문}{문항 3 지문 / 주체객체}

\kfPassageNote{문항 3 지문 / 주체객체}{%
(가) 역사적으로 예술 형식을 구현하는 수단인 매체의 자리를 두고 경쟁한 문자와 이미지의 관계는 매체 연구의 주된 주제였다. 레싱은 시와 회화의 관계를 논하면서 당시 지배적이던 호라티우스의 경구 '시는 그림과 같이'를 대체할 새로운 관점을 제시했다. 호라티우스의 경구는 시와 회화의 공통 미메시스를 설명하는 말로 읽히기도 하고, 시가 회화를 모범으로 삼아야 한다는 말로 해석되기도 했다. 레싱은 회화의 원칙을 시에 확장 적용해 두 예술 형식의 경계를 모호하게 만드는 고전적 태도를 비판했다. 그는 각 예술 형식에서 매체와 매체 효과를 구분해 보아야 한다고 주장했다. 레싱에 따르면 시와 회화는 정보를 처리하는 방식이 달라 근본적 차이가 발생한다. 그는 미메시스라 통칭되던 공통 기능을 회화적 묘사에서는 표현으로, 서술적 묘사에서는 재현으로 나누어 매체 차별성을 강조했다. 레싱은 「라오콘 군상」을 사례로 들어, 조각은 고통을 표정으로 직접 제시하지만 시는 독자가 장면을 상상해 구성하게 만든다고 (중략)
}
\begin{kfQBlock}{3}{}
\kfQStem{(가)와 (나)를 바탕으로 <보기>의 ⓐ와 ⓑ에 대해 이해한 내용으로 적절하지 않은 것은?}
\begin{kfEvidence}△△△ 작가의 작품 ⓐ는 루벤스의 그림 ⓑ「사슬로 묶인 프로메테우스」를 움직이는 영상으로 재구성하고 음악을 삽입한 작업이다.\end{kfEvidence}
\begin{kfChoices}
\kfChoice 지마는 ⓐ가 다른 예술 형식을 통합했다는 점에서 ⓑ보다 예술의 이상에 가깝다고 보겠군.
\kfChoice 레싱은 프로메테우스를 주제로 한 ⓐ의 영상과 ⓑ의 그림을 모두 재현으로 보겠군.
\end{kfChoices}
\end{kfQBlock}

\kfSecHeader{kfAreaWeekly}{지문}{문항 4 지문 / 순서방향}

\kfPassageNote{문항 4 지문 / 순서방향}{%
일반적으로 [법적 의제]란 실제로는 사실이 아닌 어떤 것을 법적으로는 마치 사실인 것처럼 취급하고 반대되는 증거가 있다 하더라도 그와 같은 취급을 고수하려는 것을 의미한다. 가령 일정 기간 이상 생사가 확인되지 않는 사람에 대해 법원이 실종 선고를 내리면 그 사람은 사망한 것으로 간주되어, 법원이 다시 실종 선고를 취소하기 전에는 설령 그가 살아 있다 하더라도 법적으로 여전히 사망한 것으로 취급되는데, 이러한 결과가 바로 법적 의제에 따른 것이라 할 수 있다.

법적 의제의 활용에 대해서는 다양한 평가가 존재해 왔다. 벤담은 이를 가장 해롭고 비열한 거짓말이라 비판하는데, 그에 의하면 정의를 실현하는 일에 의제를 동원하는 것은 거래를 한다면서 사기나 치는 것과 같다. 반면에 블랙스톤은 법적 의제가 바람직한 목표의 제도적 달성에 어색하게나마 기여한다는 점에서 매우 유용한 수단이라며 옹호한 바 있다.

(중략)
}
\begin{kfQBlock}{4}{}
\kfQStem{‘메인’의 입장에서 ㉮에 대해 추론한 내용으로 가장 적절한 것은?}
\begin{kfChoices}
\kfChoice 유년기 사회에서는 법 변화 거부감 때문에 의제 장치의 필요가 커질 수 있다.
\kfChoice 정체 사회일수록 법이 변화에 빠르게 적응한다고 볼 수 있다.
\kfChoice 진보 사회에서는 사회와 법의 간극이 발생하지 않으므로 의제가 불필요하다.
\end{kfChoices}
\end{kfQBlock}

\kfSecHeader{kfAreaWeekly}{지문}{문항 5 지문 / 인과도치}

\kfPassageNote{문항 5 지문 / 인과도치}{%
개인에게 부과되는 일반적 조세와 달리, 법인세는 법적으로만 그 실체가 인정되는 법인에 부과되는 조세이다. 법인세의 과세 대상은 일정 기간 얻은 수입에서 사업에 소요된 제반 비용을 뺀 것으로 정의된다. 세계 여러 나라의 법인세 세율 구조를 보면 단순한 세율 구조를 유지하고 있으며 세율에 차이가 있더라도 누진성과는 관계가 없는 것이 일반적인 추세이다. 이는 법인세가 납세자의 능력에 따라 조세 부담을 지게 만든다는 원칙을 적용하기 힘들기 때문이다. 이 외에도 법인세는 과세 대상이나 감가상각*과 관련하여 살펴봐야 할 점들이 여러 가지가 있다.

(중략)

법인세는 직접세의 일종이지만 실제로 누가 그 부담을 지게 될 것인지 분명하지 않다. 법인세의 부담이 실제로 누구에게 귀착되느냐에 대해서는 ㉠법인 부문에 투자된 자본에 대한 과세라는 견해와 ㉡경제적 이윤에 관한 과세라는 견해가 있다. 전자의 입장에서는 법인세의 과세 대상이 되는 소득을 정의할 때 주주가 공급한 자본의 기회비용을 공제하는 것이 허용되지 않기

(중략)
}
\begin{kfQBlock}{5}{}
\kfQStem{ⓐ의 정책 효과에 대한 설명으로 적절하지 않은 것은?}
\begin{kfChoices}
\kfChoice 가속 상각은 기업의 투자 유인 강화를 정책 목적으로 가질 수 있다.
\kfChoice 가속 상각은 감가상각의 비용 처리 기능을 없애 세 부담을 자동 증가시킨다.
\end{kfChoices}
\end{kfQBlock}

\kfSecHeader{kfAreaWeekly}{지문}{문항 6 지문 / 의도/목적}

\kfPassageNote{문항 6 지문 / 의도/목적}{%
조선에서는 16세기 이후 이황과 이이로 대표되는 사림 계열의 성리학 연구가 심화되었다. [조선의 성리학자들]은 주희의 행적을 이상적인 학자의 모습이라고 인식했다. 성리학을 \uline{ⓐ집대성}한 주희는 무이산에 은거하며 무이정사를 짓고 학문을 닦고 후학 양성에 힘을 쏟았다. 중국에서 무이산은 일찍부터 도교의 성지로 알려졌으나 주희가 은거하면서부터 더욱 유명해졌다. 주희는 친구들과 배를 타고 무이산의 아홉 굽이 계류를 따라 펼쳐진 아름다운 경치를 보고 구곡의 실경을 묘사한 「무이도가」를 지었다. 「무이도가」는 자연 속에서 심신을 수양한다는 주희의 성리학적 자연관을 보여 주는 대표적인 문학 작품으로, 차운시(남이 지은 시의 운자를 따서 지은 시)가 다수 지어졌고 무이구곡이 그림으로도 그려졌다. 조선에서는 주희의 행적을 따라 성리학을 기반으로 심신을 수양하고 학문에 정진하는 것이 사대부 사회의 문화적 유행이 되었다. 다만 그 수용 양상은 이황의 학통을 따랐던 사대부들과 이이의 학통을 따랐…(중략)
}
\begin{kfQBlock}{6}{}
\kfQStem{이황 학통의 사대부들이 「무이구곡도」의 모사에 치중한 이유로 가장 적절한 것은?}
\begin{kfChoices}
\kfChoice 「무이구곡도」가 무이산의 ‘상징적 공간’을 대신한다고 여겨, 빠짐없이 자세히 베껴 그릴 필요가 있다고 보았기 때문에
\kfChoice 「무이구곡도」가 무이산의 ‘상징적 공간’을 대신한다고 여겨, 빠짐없이 자세히 베껴 그릴 필요가 있다고 보았기 때문에
\kfChoice 「무이구곡도」가 도교 교리를 전파하는 경전이기 때문에
\kfChoice 「무이구곡도」가 조선식 구곡 경영의 현장 답사 지도이기 때문에
\end{kfChoices}
\end{kfQBlock}

\kfSecHeader{kfAreaWeekly}{지문}{문항 7 지문 / 상관관계 반전}

\kfPassageNote{문항 7 지문 / 상관관계 반전}{%
(가) 역사적으로 예술 형식을 구현하는 수단인 매체의 자리를 두고 경쟁한 문자와 이미지의 관계는 매체 연구의 주된 주제였다. 레싱은 시와 회화의 관계를 논하면서 당시 지배적이던 호라티우스의 경구 '시는 그림과 같이'를 대체할 새로운 관점을 제시했다. 호라티우스의 경구는 시와 회화의 공통 미메시스를 설명하는 말로 읽히기도 하고, 시가 회화를 모범으로 삼아야 한다는 말로 해석되기도 했다. 레싱은 회화의 원칙을 시에 확장 적용해 두 예술 형식의 경계를 모호하게 만드는 고전적 태도를 비판했다. 그는 각 예술 형식에서 매체와 매체 효과를 구분해 보아야 한다고 주장했다. 레싱에 따르면 시와 회화는 정보를 처리하는 방식이 달라 근본적 차이가 발생한다. 그는 미메시스라 통칭되던 공통 기능을 회화적 묘사에서는 표현으로, 서술적 묘사에서는 재현으로 나누어 매체 차별성을 강조했다. 레싱은 「라오콘 군상」을 사례로 들어, 조각은 고통을 표정으로 직접 제시하지만 시는 독자가 장면을 상상해 구성하게 만든다고…(중략)
}
\begin{kfQBlock}{7}{}
\kfQStem{㉠의 근거를 추론한 내용으로 가장 적절하지 않은 것은?}
\begin{kfChoices}
\kfChoice 문자 예술은 시간적 전개를 통해 수용자의 상상 참여를 촉진할 수 있다.
\kfChoice 문자 예술은 시간적 전개를 통해 수용자의 상상 참여를 촉진할 수 있다.
\kfChoice 이미지 예술은 연속적 서사 전개에서 문자 예술보다 항상 우월하다.
\kfChoice 레싱의 우월성 판단은 미메시스 유무 차이에서만 나온다.
\end{kfChoices}
\end{kfQBlock}

\kfSecHeader{kfAreaWeekly}{지문}{문항 8 지문 / 필요조건}

\kfPassageNote{문항 8 지문 / 필요조건}{%
어느 시대에서나 기적에 대한 보고는 꾸준히 있었고, 오늘날에도 기적을 목격했다는 사람은 부지기수로 많다. 죽은 사람이 살아난다든가 바위가 스스로 움직여 떠다니는 것을 보았다는 증언이 그런 예이다. 어떤 사건을 \uline{㉠기적적 사건}이라고 해야 할까?

기적이 특이한 사건이라야 한다는 점에는 동의할 수 있다. 그러나 일어나기 힘든 일이긴 하지만 자연법칙으로 설명될 수 있는 한 그것은 기적적 사건으로 간주될 수 없다. 이순신 장군의 명량 대첩처럼 13척의 배로 133척의 왜적을 물리친 일은 일어나기에 아주 힘든 일일 뿐이지 자연법칙에 어긋나지는 않는다. 그렇다면 우리에게 알려진 어떤 자연법칙으로도 설명할 수 없는 사건은 기적일까? 사진 감광판을 완전히 어두운 곳에 줄곧 보관했는데도 감광판이 빛에 노출된다는 사실을 맨 처음 알게 되었을 때, 이 사건은 당시에 알려진 어떤 자연법칙으로도 설명할 수 없었다. 그러나 머지않아 방사능에 관한 과학이 정립되어 이 진기한 현상에 대한 과학적 설명…(중략)
}
\begin{kfQBlock}{8}{}
\kfQStem{윗글에 대한 이해로 가장 적절한 것은?}
\begin{kfChoices}
\kfChoice 초자연적 사건이 실제로 일어나지 않았다면 그것에 대한 증언은 거짓이다.
\kfChoice 초자연적 사건이 실제로 일어나지 않았다면 그것에 대한 증언은 거짓이다.
\kfChoice 특이한 사건이 일어나는 것은 초자연적 힘 때문이다.
\kfChoice 현재 알려진 자연법칙의 사례가 아닌 사건은 기적적 사건이다.
\end{kfChoices}
\end{kfQBlock}

\kfSecHeader{kfAreaWeekly}{지문}{문항 9 지문 / 결합/분리}

\kfPassageNote{문항 9 지문 / 결합/분리}{%
소비란 가계가 재화와 서비스를 구입해 사용하는 행위이며, 얼마나 소비할 것인지를 결정하는 요인 중에서 가장 중요한 것은 소득의 크기이다. 특히 개인이 마음대로 쓸 수 있는 가처분 소득의 크기와 밀접한 관계를 맺고 있다. 소비와 가처분 소득 간에는 양(+)의 상관관계가 있어 가처분 소득이 증가하면 소비도 증가한다. 그리고 가처분 소득이 늘어나면서 커지는 소비의 증가분은 가처분 소득의 증가분보다 작다. 소비자들은 증가된 소득을 전부 소비하는 것이 아니라 그중 일부를 저축하기 때문이다.

어떤 국민 경제의 가처분 소득이 Yd, 그리고 소비가 C라고 할 때, \uline{㉠둘 사이의 관계는 ‘C=a+bYd (단, a>0, 0<b<1)’}와 같이 나타낼 수 있다. 이때 ａ는 생존을 위해 필요한 최소한의 소비이고, ｂ는 한계 소비 성향이다. 이 식처럼 소비와 가처분 소득 사이의 관계를 단순한 1차 함수로 표현한 것을 케인스의 소비 함수라고 한다. 이 함수에 나타난 소비와 가처분 소득 사이의 관계…(중략)
}
\begin{kfQBlock}{9}{}
\kfQStem{<보기>를 바탕으로 소비 함수를 이해한 내용으로 적절하지 않은 것은?}
\begin{kfEvidence}<보기> 어떤 나라의 가처분 소득이 2026년에 200억 원, 소비가 150억 원이었다. 2027년에는 가처분 소득이 210억 원, 소비가 156억 원이 되었다. 이 나라의 소비 함수는 C=a+bYd 형태라고 가정한다.\end{kfEvidence}
\begin{kfChoices}
\kfChoice 한계 소비 성향이 0.6이므로 늘어난 소득은 모두 소비되었고 저축 증가분은 0이다.
\kfChoice 한계 소비 성향이 0.6이므로 늘어난 소득은 모두 소비되었고 저축 증가분은 0이다.
\kfChoice 2026년의 평균 소비 성향은 0.75이다.
\kfChoice 2026년에서 2027년으로 갈 때 한계 소비 성향은 0.6이다.
\end{kfChoices}
\end{kfQBlock}

\kfSecHeader{kfAreaWeekly}{지문}{문항 10 지문 / 조건 왜곡}

\kfPassageNote{문항 10 지문 / 조건 왜곡}{%
일반적으로 혼인은 남자와 여자가 부부가 되는 일을 말한다. 과거에는 특정한 법적 절차 없이도 결혼식과 같은 사회적 관습에 의한 의식을 치르면 혼인이 성립한 것으로 보았다. 하지만 현재 대부분의 국가는 법에서 규정한 절차를 따라야만 혼인이 성립되는 법률혼 제도를 두고 있으며, 사회적 관습보다는 혼인 신고와 같은 법적 절차가 더 중요하다. 혼인은 사회를 구성하는 기본 단위인 가족을 형성하는 일이기 때문에 국가의 근간을 이루는 중요한 과정이고 이에 따라 법을 통해 관계를 보호하는 것이다. 따라서 현대 사회에서 혼인은 법적 절차에 의해 이루어지게 되고 여기에 참여한 남녀 간에는 법적 효력이 발생하게 된다.

혼인이 성립하기 위해서는 법에서 규정하고 있는 실질적 요건과 형식적 요건을 모두 만족하여야 한다. 실질적 요건에서 가장 중요한 것은 혼인 의사의 합치이다. 혼인 의사는 부부 관계를 성립하겠다는 의사로 반드시 혼인 당사자 간의 합의가 필요하다. 또 다른 실질적 요건은 연령과 혼인 자격에 …(중략)
}
\begin{kfQBlock}{10}{}
\kfQStem{윗글을 읽고 알 수 있는 내용으로 적절하지 않은 것은?}
\begin{kfChoices}
\kfChoice 혼인 당사자들이 서로의 4촌 이내의 혈족 및 그 배우자와 인척 관계가 형성되는 것은 혼인 신고서를 관련 기관에 제출한 때를 기준으로 한다.
\kfChoice 혼인 당사자들이 서로의 4촌 이내의 혈족 및 그 배우자와 인척 관계가 형성되는 것은 혼인 신고서를 관련 기관에 제출한 때를 기준으로 한다.
\kfChoice 현재 대부분의 국가에서는 사회적 관습에 따른 의식을 치르지 않더라도 법적 절차만으로 혼인이 성립될 수 있다.
\kfChoice 당사자 간의 혼인 의사 합치에 따라 관계가 형성되었더라도 혼인 당시 한쪽이 중혼인 경우에는 혼인 취소의 대상이 될 수 있다.
\end{kfChoices}
\end{kfQBlock}

\kfSecHeader{kfAreaWeekly}{지문}{문항 11 지문 / 긍부정 반전}

\kfPassageNote{문항 11 지문 / 긍부정 반전}{%
\#\#\# [31 \textasciitilde{} 34] 다음 글을 읽고 물음에 답하시오.

(가)   지팡이 짚고 바람 쐬며 좌우를 돌아보니   누대의 맑은 경치 아마도 깨끗하구나.   ㉠ 물도 하늘 같고 하늘도 물 같으니   푸른 물과 긴 하늘이 한빛이 되었거든   물가에 갈매기는 오는 듯 가는 듯 그칠 줄을 모르네.   ㉡ 바위 위 산꽃은 수놓은 병풍 되었고   시냇가 버들은 초록 장막 되었는데,   좋은 날 좋은 경치 나 혼자 거느리고   ㉢ 꽃피는 시절 허송하지 말리라 하고   아이 불러 하는 말, 이 깊은 산속에서 해산물을 볼쏘냐.   ㉣ 살진 고사리, 향기로운 당귀를 돼지고기, 사슴고기 섞어서   크나큰 바구니에 흡족히 담아두고   붕어회에다 눌어, 꿩 섞어 먹음직하게 구워지거든   술동이의 맑은 술을 술잔에 가득 부어   한잔, 또 한잔 취토록 먹은 후에,   ㉤ 복숭아꽃 붉은 비 되어 취한 낯에 뿌리는데   낚시터 넓은 돌을 높이 베고 누우니   무회씨 때 사람인가, 갈천씨 때 백성*인가.*   *…(중략)
}
\begin{kfQBlock}{11}{}
\kfQStem{㉠ \textasciitilde{} ㉤에 대한 이해로 적절하지 않은 것은?}
\begin{kfChoices}
\kfChoice ㉢ : 의지적인 어조를 활용하여 학문 수양을 게을리하지 않으려는 자세를 드러내고 있다.
\kfChoice ㉢ : 의지적인 어조를 활용하여 학문 수양을 게을리하지 않으려는 자세를 드러내고 있다.
\kfChoice ㉠ : 유사한 문장 구조를 반복하여 자연물 간의 경계가 사라진 풍광을 묘사하고 있다.
\kfChoice ㉡ : 일상의 사물에 빗대어 화자를 둘러싼 자연의 모습을 표현하고 있다.
\end{kfChoices}
\end{kfQBlock}

\kfSecHeader{kfAreaWeekly}{지문}{문항 12 지문 / 비교대상 뒤바꿈}

\kfPassageNote{문항 12 지문 / 비교대상 뒤바꿈}{%
\# 
}
\begin{kfQBlock}{12}{}
\kfQStem{(가)와 (나)의 공통점으로 가장 적절한 것은?}
\begin{kfChoices}
\kfChoice 설의적 표현을 활용하여 전달하고자 하는 내용을 강조하고 있다.
\kfChoice 설의적 표현을 활용하여 전달하고자 하는 내용을 강조하고 있다.
\kfChoice 의인법을 통해 대상과의 친밀감을 표현하고 있다.
\kfChoice 언어유희를 통해 현실에 대한 태도를 간접적으로 드러내고 있다.
\end{kfChoices}
\end{kfQBlock}

\kfSecHeader{kfAreaWeekly}{지문}{문항 13 지문 / 주체/객체 혼동}

\kfPassageNote{문항 13 지문 / 주체/객체 혼동}{%
\#\#\# [18 \textasciitilde{} 23] 다음 글을 읽고 물음에 답하시오.

(가)   녹양방초 언덕에 소 먹이는 저 아이야   인간 영욕을 아는가 모르는가   인생 백년이 풀 끝에 이슬이라   삼만 육천 일이 그 아니 보잘것없는가   하물며 장수 단명이 운명이어니 사생(死生)을 정할쏘냐   여관 같은 세상에 하루살이같이 나왔다가   ㉠ 공명도 못 이루고 초목같이 썩어지면   ⓐ 공산 백골이 그 아니 느꺼우냐   하늘의 뜻을 이어 법칙을 세움은 옛 성인의 사업이요   아름다운 이름을 후세에 전함은 대장부의 할 일이라   생애는 유한하고 사일(死日)은 무궁하니   유한한 생애로 썩지 않을 이름을   영구히 전하여 ⓑ 천지와 함께 무궁하려고   (중략)   어와 그 뉘신고 어떠한 사람인고   형용이 초췌하니 초나라 대부 굴원이신가   잔혼이 영락하니 학사 유자후이신가*   눈썹을 찡그리시니 근심이 많으신가   발끝으로 서시니 어디를 바라는고   아름다운 기약을 바라는가 이별의 슬픔이 중하신가   ⓒ 해…(중략)
}
\begin{kfQBlock}{13}{}
\kfQStem{ⓐ \textasciitilde{} ⓔ에 대한 설명으로 가장 적절한 것은?}
\begin{kfChoices}
\kfChoice ⓑ : 삶의 유한성에 대한 인식에서 비롯한 욕망을 드러낸다.
\kfChoice ⓑ : 삶의 유한성에 대한 인식에서 비롯한 욕망을 드러낸다.
\kfChoice ⓐ : 자신의 삶을 성찰하며 느끼는 후회의 감정을 드러낸다.
\kfChoice ⓒ : 임과 이별한 자신의 외로운 처지를 나타낸다.
\end{kfChoices}
\end{kfQBlock}

\kfSecHeader{kfAreaWeekly}{지문}{문항 14 지문 / 내용과 방법 혼동}

\kfPassageNote{문항 14 지문 / 내용과 방법 혼동}{%
\#\#\# [32～34] 다음 글을 읽고 물음에 답하시오.

(가)   장풍에 돛을 달고 육선이 함께 떠나   삼현과 군악 소리 해산을 진동하니   물속의 어룡들이 응당히 놀라리라   해구를 얼른 나서 오륙도를 뒤 지우고   고국을 돌아보니 야색이 아득하여   아무것도 아니 뵈고 연해 각진포에   불빛 두어 점이 구름 밖에 뵐 만하다   배 방에 누워 있어 내 신세를 생각하니   가뜩이 심란한데 대풍이 일어나서   태산 같은 성난 물결 천지에 자욱하니   크나큰 만곡주가 나뭇잎 불리이듯   하늘에 올랐다가 지함에 내려지니   열두 발 쌍돛대는 차아처럼 굽어 있고   쉰두 폭 초석 돛은 반달처럼 배불렀네   (중략)   날이 마침 극열하고 석양이 비치어서   끓는 땅에 엎디어서 말씀을 여쭈오니   속에서 불이 나고 관대에 땀이 배어   물 흐르듯 하는지라 나라께서 보시고서   너희 더위 어려우니 먼저 나가 쉬라시니   곡배하고 사퇴하니 천은이 망극하다   더위를 장히 먹어 막힐 듯하는지라…(중략)
}
\begin{kfQBlock}{14}{}
\kfQStem{(가), (나)의 표현상 특징에 대한 설명으로 가장 적절한 것은?}
\begin{kfChoices}
\kfChoice (가)는 사물의 형태가 변화한 모습을 묘사하여 외부 환경의 영향력을 부각하고 있다.
\kfChoice (가)는 사물의 형태가 변화한 모습을 묘사하여 외부 환경의 영향력을 부각하고 있다.
\kfChoice (가)는 과거를 회상하는 표현을 통해 현재 상황에 대한 아쉬움을 드러내고 있다.
\kfChoice (나)는 계절을 나타내는 어휘를 활용해 애달픈 정서를 부각하고 있다.
\end{kfChoices}
\end{kfQBlock}

\kfSecHeader{kfAreaWeekly}{지문}{문항 15 지문 / 내용 왜곡}

\kfPassageNote{문항 15 지문 / 내용 왜곡}{%
[26\textasciitilde{}31] 다음 글을 읽고 물음에 답하시오.

(가)   십 년 종사 후에 고향으로 도라오니   산천 의구하되 인사(人事)는 달라졌구나   아마도 세간의 존멸을 못내 슬허 하노라 <제1수>

산화(山花)는 믈의 픠고 물새는 산의 운다   일신이 한가하야 산수간의 누어시니   세상의 어즈러은 긔별을 나는 몰라 하노라 <제4수>

거믄고 빗기 들고 산수를 희롱하니   청풍은 건듯 블고 명월도 도라온다   하물며 유신(有信)한 갈매기는 오명 가명 하나니 <제5수>

거믄고 흥진(興盡)커던 조대(釣臺)로 내려가니   도화 뜬 말근 믈 뛰노나니 고기로다   아이야 밋기 다지 마라 취적(取適)*이나 하오리라 <제7수>

* 신교, ｢귀산음(歸山吟)｣ -   * 취적 : 낚시질의 참뜻이 세상 생각을 잊고자 하는 데 있음.

(나)   백수(白首)에 산수 구경 늦은 줄 알지마는   평생 품은 뜻을 이루고야 말리라 여겨   병자년 봄에 봄옷을 새로 입고   죽장망혜(竹杖芒鞋)로 노계 깊은 골에 …(중략)
}
\begin{kfQBlock}{15}{}
\kfQStem{ⓐ와 ⓑ에 대한 이해로 가장 적절한 것은?}
\begin{kfChoices}
\kfChoice ⓐ는 하늘의 모습을 물에서 보게 된 것에 대한, ⓑ는 산의 모습이 평소와 달리 보이는 것에 대한 반응이다.
\kfChoice ⓐ는 하늘의 모습을 물에서 보게 된 것에 대한, ⓑ는 산의 모습이 평소와 달리 보이는 것에 대한 반응이다.
\kfChoice ⓐ는 하늘과 물의 변함없는 모습을 본 것에 대한, ⓑ는 선명하게 드러난 산의 모습을 본 것에 대한 반응이다.
\kfChoice ⓐ는 하늘이 물의 모습을 닮아 변해 가는 것에 대한, ⓑ는 산이 주변의 모습을 닮아 변해 가는 것에 대한 반응이다.
\end{kfChoices}
\end{kfQBlock}

\kfSecHeader{kfAreaWeekly}{지문}{문항 16 지문 / 과잉 해석}

\kfPassageNote{문항 16 지문 / 과잉 해석}{%
[26\textasciitilde{}28] 다음 글을 읽고 물음에 답하시오.

어머니는 동물적 본능에 가까울 정도로 생에 대한 집착이 강했다. 조금만 아프거나 배고픈 것도 참지 못했다. 노인정에서 점심 먹은 것이 조금 부실한 날은 해가 떨어지기도 전에 허기진 모습으로 집에 돌아와서 숟가락을 들고 밥통부터 찾곤 했다. 이 때문에 우리 집 전기밥통에는 언제나 밥이 준비되어 있게 마련이다. 밥이 없으면 아무렇지 않은 일에도 까탈을 부리며 심하게 며느리를 닦달했다. 어머니한테 밥은 곧 생명이며 에너지원이다. 어머니는 또 몸의 컨디션이 조금만 나빠도 아이들처럼 엄살을 떨며 당장 병원에 찾아가 주사 맞는 것을 좋아했다. ㉠노인네들이 항생제 주사를 많이 맞는 것이 좋지 않다는 말을 해도 듣지 않았다. 우리 가족들 중에서 해마다 가장 먼저 독감 예방주사를 맞는 것도 어머니다.   어머니가 젊었을 적에는 그렇지가 않았다. 배고픈 것도 잘 참았고 아무리 아파도 자리보전하거나 약을 먹지도 않았다. ㉡몸살이 나서 끙끙 앓으면서도 휘…(중략)
}
\begin{kfQBlock}{16}{}
\kfQStem{<보기>를 바탕으로 윗글을 감상한 것으로 적절하지 않은 것은? [3점]}
\begin{kfChoices}
\kfChoice ‘어머니의 삶’을 ‘궁핍과 땀과 희생과 인종’으로 보는 것에는 어머니의 인생을 부정적으로 여기는 ‘나’의 모습이 드러나 있군.
\kfChoice ‘어머니의 삶’을 ‘궁핍과 땀과 희생과 인종’으로 보는 것에는 어머니의 인생을 부정적으로 여기는 ‘나’의 모습이 드러나 있군.
\kfChoice ‘집’에 ‘밥이 없으면’ ‘까탈을 부리’는 어머니의 모습을 ‘나’가 ‘생에 대한 집착’이라 생각하는 것에는 부모 세대에 대한 그릇된 관념이 드러나 있군.
\kfChoice ‘특유한 어머니의 냄새’는 현실의 문제를 드러내는 소재로 사용되고 있군.
\end{kfChoices}
\end{kfQBlock}

\kfSecHeader{kfAreaWeekly}{지문}{문항 17 지문 / 범위 오류}

\kfPassageNote{문항 17 지문 / 범위 오류}{%
\#\#\# [42 \textasciitilde{} 45] 다음 글을 읽고 물음에 답하시오.

[앞부분 줄거리] 정 소저는 계모 박 씨의 모함을 의심 없이 받아들인 아버지 정공 때문에 위기에 처하고, 집에서 나와 숨어 다니던 중 도적을 만나 강물에 몸을 던진다. 이때, 정혼자 조무(용홍)와 동생 조성이 정 소저를 우연히 발견하여 구출한다.

소저가 매우 놀라며 말하였다.   “내가 외가로 가지 않고 구차하게 길가에서 분주하게 다닌 것은 조숙모에게 부끄럽고, 아버지의 허물을 드러내고 싶지 않아서였다. 뜻밖에 저 공자들을 만나니 내가 차마 사실을 말하여 부끄러움을 더하겠는가? 은인의 덕이 산과 바다 같으나 차마 근본을 아뢰게 되어 저 집에서 우리 집의 허물을 알게 되면 매우 부끄럽게 될 것이다. 모름지기 너는 다만 대답하기를 내가 타향에서 떠돌아다니다가 서울의 친척을 찾으러 왔다가 도적을 만나 물에 빠져 죽을 뻔했다고 말하여라. 조 공자가 이미 우리가 여자인 줄을 알았으니 남녀는 구별이 있는 것이다. 생명을 구해준 은혜…(중략)
}
\begin{kfQBlock}{17}{}
\kfQStem{[A]와 [B]에 대한 이해로 가장 적절하지 않은 것은?}
\begin{kfChoices}
\kfChoice [A]와 [B]는 모두 과거에 일어난 일을 상대에게 언급하고 있다.
\kfChoice [A]와 [B]는 모두 과거에 일어난 일을 상대에게 언급하고 있다.
\kfChoice [A]는 [B]와 달리 상대의 행동에 변화를 촉구하고 있다.
\kfChoice [B]는 [A]와 달리 상대에게 다른 인물의 말을 전하고 있다.
\end{kfChoices}
\end{kfQBlock}

\kfSecHeader{kfAreaWeekly}{지문}{문항 18 지문 / 시점/화자 혼동}

\kfPassageNote{문항 18 지문 / 시점/화자 혼동}{%
[26\textasciitilde{}31] 다음 글을 읽고 물음에 답하시오.

(가)   십 년 종사 후에 고향으로 도라오니   산천 의구하되 인사(人事)는 달라졌구나   아마도 세간의 존멸을 못내 슬허 하노라 <제1수>

산화(山花)는 믈의 픠고 물새는 산의 운다   일신이 한가하야 산수간의 누어시니   세상의 어즈러은 긔별을 나는 몰라 하노라 <제4수>

거믄고 빗기 들고 산수를 희롱하니   청풍은 건듯 블고 명월도 도라온다   하물며 유신(有信)한 갈매기는 오명 가명 하나니 <제5수>

거믄고 흥진(興盡)커던 조대(釣臺)로 내려가니   도화 뜬 말근 믈 뛰노나니 고기로다   아이야 밋기 다지 마라 취적(取適)*이나 하오리라 <제7수>

* 신교, ｢귀산음(歸山吟)｣ -   * 취적 : 낚시질의 참뜻이 세상 생각을 잊고자 하는 데 있음.

(나)   백수(白首)에 산수 구경 늦은 줄 알지마는   평생 품은 뜻을 이루고야 말리라 여겨   병자년 봄에 봄옷을 새로 입고   죽장망혜(竹杖芒鞋)로 노계 깊은 골에 …(중략)
}
\begin{kfQBlock}{18}{}
\kfQStem{<보기>를 참고하여 (다)를 감상한 내용으로 적절하지 않은 것은?}
\begin{kfChoices}
\kfChoice ‘팔베개를 하고 누워’ 하늘을 ‘무심히’ 바라보는 것은 지식을 멈추고 새로운 것을 경험하려는 행동으로 볼 수 있겠군.
\kfChoice ‘팔베개를 하고 누워’ 하늘을 ‘무심히’ 바라보는 것은 지식을 멈추고 새로운 것을 경험하려는 행동으로 볼 수 있겠군.
\kfChoice ‘사람’과 ‘사물’을 ‘일상적’으로 대하는 것은 미지의 것을 경험하는 데에 장애가 될 수 있겠군.
\kfChoice 어떤 대상에 대해 ‘아무개 하’는 것은 그 대상을 기존의 지식으로 해석하게 한다고 볼 수 있겠군.
\end{kfChoices}
\end{kfQBlock}



\clearpage

\kfChapterCover{3}{비트겐슈타인2 ch03}{Chapter 3}{이번 챕터: 문법 → 문학 5작품 → 비문학 5지문 → 패턴 워크북.}

\kfStartGrammar{이번 챕터 문법 영역 — 음운 / 비음화·유음화 + 복합 음운 변동}


\kfSecHeader{kfAreaGrammar}{문제}{}

\begin{multicols}{2}\raggedcolumns\setlength{\columnsep}{12pt}
\kfPassageFull{문항 1 지문 (concept)}{%
비음화란 비음이 아닌 자음이 비음의 영향을 받아 비음으로 바뀌는 음운 변동이다. 파열음 'ㄱ, ㄷ, ㅂ'이 비음 'ㄴ, ㅁ' 앞에서 각각 'ㅇ, ㄴ, ㅁ'으로 교체된다. 예를 들어 '국물'[궁물], '받는'[반는], '합니다'[함니다]가 이에 해당한다. 비음화는 인접한 비음에 동화되는 현상이므로 동화의 방향은 역행 동화(뒤 음소가 앞 음소에 영향)이다. 비음화의 결과 바뀌는 음소와 원래 음소는 조음 위치가 같고 조음 방법만 달라진다. 즉 'ㄱ→ㅇ'(연구개), 'ㄷ→ㄴ'(치조), 'ㅂ→ㅁ'(양순)은 각각 조음 위치가 보존된다.
}
\begin{kfQBlock}{1}{}
\kfQStem{윗글을 바탕으로 <보기>의 단어들에서 비음화가 적용되는 과정을 분석한 것으로 적절하지 않은 것은?

<보기> ⓐ 먹는[멍는]  ⓑ 십만[심만]  ⓒ 있는[인는]  ⓓ 꽃말[꼰말]  ⓔ 법률[범뉼]}
\begin{kfChoices}
\kfChoice ⓐ에서 'ㄱ'이 비음 'ㄴ' 앞에서 같은 조음 위치인 'ㅇ'으로 바뀌어 [멍는]이 된다.
\kfChoice ⓑ에서 'ㅂ'이 비음 'ㅁ' 앞에서 같은 조음 위치인 'ㅁ'으로 바뀌어 [심만]이 된다.
\kfChoice ⓒ에서 'ㅆ'이 비음 'ㄴ' 앞에서 바로 'ㄴ'으로 바뀌어 [인는]이 된다.
\kfChoice ⓓ에서 'ㅅ'이 끝소리 규칙으로 [ㄷ]이 된 후 비음 'ㅁ' 앞에서 [ㄴ]으로 바뀌어 [꼰말]이 된다.
\kfChoice ⓔ에서 'ㅂ'이 'ㄹ' 앞에서 비음화와 유음화가 적용되어 [범뉼]이 된다.
\end{kfChoices}
\end{kfQBlock}

\kfPassageFull{문항 2 지문 (concept)}{%
유음화란 'ㄴ'이 유음 'ㄹ' 앞이나 뒤에서 'ㄹ'로 바뀌는 현상이다. 'ㄴ'이 'ㄹ' 뒤에 올 때(순행 유음화)와 'ㄹ' 앞에 올 때(역행 유음화) 모두 유음화가 적용된다. '신라'[실라], '칼날'[칼랄]이 대표적인 예이다. 순행 유음화의 경우 '설날'[설랄]처럼 앞의 'ㄹ'이 뒤의 'ㄴ'을 동화시키고, 역행 유음화의 경우 '신라'[실라]처럼 뒤의 'ㄹ'이 앞의 'ㄴ'을 동화시킨다. 그러나 합성어에서 'ㄴ'이 'ㄹ' 앞에 올 때 유음화되지 않고 오히려 'ㄹ'이 [ㄴ]으로 바뀌는 경우도 있다. '담력'[담녁], '생산량'[생산냥]이 이에 해당한다.
}
\begin{kfQBlock}{2}{}
\kfQStem{윗글을 바탕으로 <보기>의 발음에 대한 설명으로 적절한 것은?

<보기> ⓐ 난로[날로]    ⓑ 진리[질리]    ⓒ 물난리[물랄리] ⓓ 의견란[의견난]    ⓔ 칼날[칼랄]}
\begin{kfChoices}
\kfChoice ⓐ와 ⓑ는 모두 순행 유음화가 적용된 예이다.
\kfChoice ⓐ와 ⓑ는 모두 역행 유음화가 적용된 예이다.
\kfChoice ⓒ에서 'ㄴ'이 모두 유음화되어 [물랄리]가 된다.
\kfChoice ⓓ에서 'ㄴ'이 유음화되지 않고 'ㄹ'이 [ㄴ]으로 바뀐 것은 역행 비음화에 해당한다.
\kfChoice ⓔ는 역행 유음화가 적용된 예이다.
\end{kfChoices}
\end{kfQBlock}

\kfPassageFull{문항 3 지문 (concept)}{%
하나의 단어에서 비음화와 유음화가 함께 적용되는 경우가 있다. '독립'[동닙]은 'ㄱ'이 'ㄹ' 앞에서 바로 비음화되지 않는다. 먼저 'ㄹ'이 'ㄱ' 뒤에서 비음 [ㄴ]으로 바뀌고(ㄹ의 비음화), 그 결과 'ㄱ'이 비음 'ㄴ' 앞에서 [ㅇ]으로 비음화되어 [동닙]이 된다. 반면 '설날'[설랄]은 'ㄹ' 뒤의 'ㄴ'이 유음화되어 [ㄹ]로 바뀐다. 이처럼 동일한 'ㄹ'과 다른 자음의 만남이라도 환경에 따라 유음화가 적용될 수도, 비음화가 적용될 수도 있다. 일반적으로 'ㄹ'이 비음이 아닌 자음 뒤에 올 때는 'ㄹ'의 비음화(ㄹ→ㄴ)가, 'ㄹ' 뒤에 'ㄴ'이 올 때는 유음화(ㄴ→ㄹ)가 적용된다.
}
\begin{kfQBlock}{3}{}
\kfQStem{윗글을 바탕으로 <보기>의 단어에 적용된 음운 변동을 분석한 것으로 적절하지 않은 것은?

<보기> ⓐ 협력[혐녁]    ⓑ 찰나[찰라]    ⓒ 백리[뱅니] ⓓ 실내[실래]    ⓔ 법률[범뉼]}
\begin{kfChoices}
\kfChoice ⓐ에서 'ㄹ'→[ㄴ](비음화) 후 'ㅂ'→[ㅁ](비음화)이 적용되어 [혐녁]이 된다.
\kfChoice ⓑ에서 'ㄹ' 뒤의 'ㄴ'이 유음화되어 [찰라]가 된다.
\kfChoice ⓒ에서 'ㄹ'→[ㄴ](비음화) 후 'ㄱ'→[ㅇ](비음화)이 적용되어 [뱅니]가 된다.
\kfChoice ⓓ에서 'ㄴ'이 앞의 'ㄹ'에 의해 유음화되어 [실래]가 된다.
\kfChoice ⓔ에서 'ㅂ'→[ㅁ](비음화) 후 'ㄹ'이 유음화되어 [범률]이 된다.
\end{kfChoices}
\end{kfQBlock}

\kfPassageFull{문항 4 지문 (concept)}{%
복합 음운 변동에서는 여러 규칙이 순차적으로 적용되므로, 각 단계를 정확히 파악하는 것이 중요하다. 음운 변동의 유형은 크게 교체(음운이 다른 음운으로 바뀜), 탈락(음운이 사라짐), 첨가(새 음운이 생김), 축약(두 음운이 하나로 합쳐짐)의 네 가지이다. 비음화와 유음화는 교체에 해당하며, 끝소리 규칙도 교체이다. 한 단어에 복합 음운 변동이 적용될 때는 변동의 유형과 순서를 함께 분석해야 한다. 예: '없는'에서 겹받침 'ㅄ'→[ㅂ](끝소리 규칙, 교체) → [ㅁ](비음화, 교체) → [엄는]. 이 과정에서 교체가 두 번 연속 적용된다.
}
\begin{kfQBlock}{4}{}
\kfQStem{윗글을 바탕으로 <보기>의 단어에 적용되는 음운 변동의 유형과 횟수를 분석한 것으로 적절한 것은?

<보기> 「닫는」, 「읊는」의 발음 과정을 각각 분석하시오.}
\begin{kfChoices}
\kfChoice '닫는'에는 교체가 1회 적용되고, '읊는'에는 교체가 2회 적용된다.
\kfChoice '닫는'에는 교체가 2회 적용되고, '읊는'에는 교체가 1회 적용된다.
\kfChoice '닫는'과 '읊는' 모두 교체가 2회씩 적용된다.
\kfChoice '닫는'에는 비음화만 적용되고, '읊는'에는 끝소리 규칙과 비음화가 적용된다.
\kfChoice '닫는'과 '읊는' 모두 끝소리 규칙과 비음화가 적용된다.
\end{kfChoices}
\end{kfQBlock}

\kfPassageFull{문항 5 지문 (concept)}{%
동화란 인접한 두 음운 중 하나가 다른 하나의 성질을 닮아가는 현상이다. 동화의 방향에 따라 순행 동화(앞 음소가 뒤 음소에 영향)와 역행 동화(뒤 음소가 앞 음소에 영향)로 나뉜다. 비음화는 뒤의 비음이 앞의 파열음에 영향을 미치므로 역행 동화이다. 순행 유음화는 앞의 'ㄹ'이 뒤의 'ㄴ'에 영향을 미치므로 순행 동화이고, 역행 유음화는 뒤의 'ㄹ'이 앞의 'ㄴ'에 영향을 미치므로 역행 동화이다. 한편 동화의 정도에 따라 완전 동화(두 음소가 완전히 같아짐)와 부분 동화(일부 자질만 같아짐)로 나뉜다. 비음화에서 'ㄱ→ㅇ'은 조음 방법만 바뀌고 조음 위치는 유지되므로 부분 동화이다. 유음화에서 'ㄴ→ㄹ'은 조음 위치(치조)가 같고 조음 방법이 바뀌므로 역시 부분 동화이다.
}
\begin{kfQBlock}{5}{}
\kfQStem{윗글을 바탕으로 <보기>의 각 단어에 적용된 동화의 방향과 정도를 분석한 것으로 적절하지 않은 것은?

<보기> ⓐ 국민[궁민]    ⓑ 신라[실라]    ⓒ 설날[설랄]    ⓓ 밥물[밤물]}
\begin{kfChoices}
\kfChoice ⓐ는 역행 동화이며 부분 동화이다.
\kfChoice ⓑ는 역행 동화이며 완전 동화이다.
\kfChoice ⓒ는 순행 동화이며 완전 동화이다.
\kfChoice ⓓ는 순행 동화이며 부분 동화이다.
\kfChoice ⓐ와 ⓓ는 동화의 방향이 같다.
\end{kfChoices}
\end{kfQBlock}

\end{multicols}


\kfStartLit{이번 챕터 문학 작품 5편}


\kfSecHeader{kfAreaLiterature}{지문}{성산별곡 — 정철 (가사) / 정철, 「성산별곡」, 『송강가사(松江歌辭)』}

\kfPassageNote{성산별곡 — 정철 (가사) / 정철, 「성산별곡」, 『송강가사(松江歌辭)』}{%
식영정(息影亭) 나린 물이 서석산(瑞石山) 빗겨 가니 \\[4pt]
큰 들이 거기로다. 아으 그 넘어 너른 들은 \\[4pt]
물아래 뵈는 것은 무어스 무어스 다 무엇이냐 \\[4pt]
어촌(漁村)에 해 비치니 물결이 금(金)비단이 \\[4pt]
사공은 어디 가고 빈 배만 매여 있다. \\[14pt]
장안(長安)을 도라보니 북극(北極)이 하 머니 \\[4pt]
세상 만사(萬事)를 생각하면 구름이로다. \\[4pt]
사람이 풍경(風景)을 몰라 하니 이 내 생애(生涯) 어떠한고? \\[4pt]
강산(江山)이 좋다 한들 내 분(分)인가 하노라. \\[14pt]
소동파(蘇東坡)의 적벽부(赤壁賦)를 여기 와 외와도다. \\[4pt]
물아일체(物我一體)를 어느 것이 다르리. \\[4pt]
무궁(無窮)한 이 강산에 무궁한 이내 흥(興)을 \\[4pt]
누구서 다 누리며 어떻게 다할소냐.
}
\kfSecHeader{kfAreaLiterature}{활동}{작품 정리표 — 성산별곡(정철)}

\noindent\textit{\small 빈칸에 알맞은 말을 채워 넣으시오. (초성 힌트 참고)}\par\medskip
\begin{kfActTable}{|>{\centering\arraybackslash}m{2.6cm}|m{6cm}|m{4cm}|}
\rowcolor{kfPrimaryLight!50}\textbf{\color{kfPrimary}항목} & \textbf{\color{kfPrimary}내용 (빈칸 채우기)} & \textbf{\color{kfPrimary}답안 작성}\\\hline
갈래 & ㄱㅅ ① & \kfTBlank \\\hline
작가 & ㅈㅊ(송강) ② & \kfTBlank \\\hline
배경 공간 & 담양 ㅅㅇㅈ ③ & \kfTBlank \\\hline
식영정 주인 & ㄱㅅㅇ ④ & \kfTBlank \\\hline
전개 방식 & ㅅㄱㅈ의 변화에 따른 전개 ⑤ & \kfTBlank \\\hline
인용 고사 & ㅅㄷㅍ의 ㅈㅂㅂ ⑥ & \kfTBlank \\\hline
핵심 개념 & 자연과 내가 하나 = ㅁㅇㅇㅊ ⑦ & \kfTBlank \\\hline
'장안'의 의미 & 서울(ㅈㅈ) ⑧ & \kfTBlank \\\hline
'북극'의 의미 & ㅇㄱ이 계신 곳 ⑨ & \kfTBlank \\\hline
작품 속 태도 & 세상 만사는 ㄱㄹ처럼 허무 ⑩ & \kfTBlank \\\hline
주제 1 & ㅈㅇ ㅇㄱ의 아름다움과 즐거움 ⑪ & \kfTBlank \\\hline
주제 2 & 유교적 ㄷㅎ 실천과 ㅎㅁ ㅅㅇ ⑫ & \kfTBlank \\\hline
\end{kfActTable}
\kfSecHeader{kfAreaLiterature}{문제}{}

\begin{multicols}{2}\raggedcolumns\setlength{\columnsep}{12pt}
\begin{kfQBlock}{1}{}
\kfQStem{윗글에 대한 설명으로 적절하지 않은 것은?}
\begin{kfChoices}
\kfChoice 자연의 경물을 열거하면서 풍경의 아름다움을 구체적으로 묘사하고 있다.
\kfChoice 중국의 고사(蘇東坡·赤壁賦)를 인용하여 자연 감상의 깊이를 더하고 있다.
\kfChoice '장안을 도라보니 북극이 하 머니'에서 정치적 현실에 대한 인식이 드러난다.
\kfChoice 화자는 자연 속의 삶을 부정하고 다시 관직에 나아가려는 의지를 표출하고 있다.
\kfChoice 설의법을 사용하여 자연 속 흥취의 무궁함을 강조하고 있다.
\end{kfChoices}
\end{kfQBlock}

\begin{kfQBlock}{2}{}
\kfQStem{윗글의 '장안(長安)을 도라보니 북극(北極)이 하 머니'에서 '장안'과 '북극'의 의미로 가장 적절한 것은?}
\begin{kfChoices}
\kfChoice '장안'은 중국의 옛 수도, '북극'은 밤하늘의 별자리를 가리킨다.
\kfChoice '장안'은 서울(조정), '북극'은 임금이 계신 곳을 비유한다.
\kfChoice '장안'은 화자의 고향, '북극'은 먼 이국땅을 의미한다.
\kfChoice '장안'은 번화한 시장, '북극'은 험준한 산을 비유한다.
\kfChoice '장안'은 학문의 전당, '북극'은 진리의 세계를 상징한다.
\end{kfChoices}
\end{kfQBlock}

\begin{kfQBlock}{3}{}
\kfQStem{'물아일체(物我一體)를 어느 것이 다르리'에서 화자가 말하고자 하는 바로 가장 적절한 것은?}
\begin{kfChoices}
\kfChoice 자연의 모든 만물은 인간과 무관하게 독립적으로 존재한다.
\kfChoice 이곳 성산에서도 소동파가 느꼈던 자연과의 합일을 체험할 수 있다.
\kfChoice 물아일체는 이론으로만 가능하고 현실에서는 실현이 불가능하다.
\kfChoice 소동파가 지은 적벽부는 이곳 성산의 풍경과 전혀 맞지 않는다.
\kfChoice 자연 속에서 학문에 몰두하여 입신양명을 이루어 내야만 한다.
\end{kfChoices}
\end{kfQBlock}

\begin{kfQBlock}{4}{}
\kfQStem{윗글에서 '사공은 어디 가고 빈 배만 매여 있다'라는 구절의 효과로 가장 적절한 것은?}
\begin{kfChoices}
\kfChoice 사공의 부재를 한탄하며 화자의 고립감을 드러낸다.
\kfChoice 어촌의 적막한 풍경을 통해 한적하고 평화로운 분위기를 형성한다.
\kfChoice 사공이 떠난 이유를 추리하게 하여 서사적 긴장감을 조성한다.
\kfChoice 빈 배를 통해 화자의 경제적 궁핍을 상징적으로 나타낸다.
\kfChoice 사공과 화자 사이의 갈등이 해소되었음을 암시한다.
\end{kfChoices}
\end{kfQBlock}

\begin{kfQBlock}{5}{}
\kfQStem{윗글을 바탕으로 <보기>를 감상한 내용으로 적절하지 않은 것은?

<보기> 「성산별곡」은 정철이 김성원의 식영정 생활을 노래한 것으로, 자연 속 은거와 유교적 도학의 실천이 결합된 작품이다. 조선 전기 사대부 문인들에게 자연은 단순한 풍류의 공간이 아니라, 도(道)를 체득하는 수양의 장이었다.}
\begin{kfChoices}
\kfChoice '이 내 생애 어떠한고'에서 화자가 자연 속 삶에 대해 자부심을 느끼는 것은 은거 생활의 가치를 인식한 것으로 볼 수 있겠군.
\kfChoice '강산이 좋다 한들 내 분인가 하노라'에서 화자가 겸양하는 것은 자연을 수양의 장으로 여기는 사대부의 태도와 연결되겠군.
\kfChoice 소동파의 적벽부를 인용한 것은 중국 고전을 통해 자연 감상의 전통을 계승하려는 의식으로 볼 수 있겠군.
\kfChoice '무궁한 이 강산에 무궁한 이내 흥을'은 자연의 영원함과 자신의 흥취를 대응시킨 것으로, 자연과의 합일 의식을 보여주겠군.
\kfChoice '세상 만사를 생각하면 구름이로다'는 세속의 일에 적극적으로 참여하겠다는 화자의 출사 의지를 드러낸 것으로 보아야겠군.
\end{kfChoices}
\end{kfQBlock}

\end{multicols}

\kfSecHeader{kfAreaLiterature}{지문}{사미인곡 — 정철 (가사) / 정철, 「사미인곡」, 『송강가사(松江歌辭)』}

\kfPassageNote{사미인곡 — 정철 (가사) / 정철, 「사미인곡」, 『송강가사(松江歌辭)』}{%
이 몸 삼기실 제 님을 조차 삼기시니 \\[4pt]
한생(一生) 연분(緣分)이며 하늘 모를 일이런가. \\[4pt]
나 하나 졈어 이셔 님을 하나 따랏더니 \\[4pt]
이제야 중년(中年)이 되도록 못 잊었어 하노라. \\[14pt]
아마도 이 님이 나인 줄도 모르시니 \\[4pt]
님을 향한 뜻은 나도 몰라 하노라. \\[4pt]
잠을 자니 꿈에나 보려 하되 \\[4pt]
눈을 감으면 생각이 새로워라. \\[14pt]
차라리 삭여 내여 곳비 되여 나으리라. \\[4pt]
향(香) 머금은 곳비 되여 님의 옷세 올으리라. \\[4pt]
님이야 날인 줄 모르셔도 \\[4pt]
내 님 뫼셔 옷깃 좇아 가리라.
}
\kfSecHeader{kfAreaLiterature}{활동}{작품 정리표 — 사미인곡(정철)}

\noindent\textit{\small 빈칸에 알맞은 말을 채워 넣으시오. (초성 힌트 참고)}\par\medskip
\begin{kfActTable}{|>{\centering\arraybackslash}m{2.6cm}|m{6cm}|m{4cm}|}
\rowcolor{kfPrimaryLight!50}\textbf{\color{kfPrimary}항목} & \textbf{\color{kfPrimary}내용 (빈칸 채우기)} & \textbf{\color{kfPrimary}답안 작성}\\\hline
갈래 & ㄱㅅ ① & \kfTBlank \\\hline
작가 & ㅈㅊ(송강) ② & \kfTBlank \\\hline
제목 의미 & 아름다운 사람(ㅇㄱ)을 ㄱㄹㅇ ③ & \kfTBlank \\\hline
핵심 정서 & 임금에 대한 그리움 = ㅇㄱㅈㅈ ④ & \kfTBlank \\\hline
화자 설정 & ㅇㅅ 화자를 내세움 ⑤ & \kfTBlank \\\hline
문학적 전통 & 굴원(屈原)의 ㅇㅅ 전통 계승 ⑥ & \kfTBlank \\\hline
변신 모티프 & ㄱㅂ가 되어 님의 옷에 오르겠다 ⑦ & \kfTBlank \\\hline
인연의 성격 & '삼기실 제 님을 조차 삼기시니' = ㅅㅁ ⑧ & \kfTBlank \\\hline
화자의 태도 & 님이 몰라도 따르는 ㅇㅂㅈ ㅎㅅ ⑨ & \kfTBlank \\\hline
그리움의 심화 & 눈을 감으면 ㅅㄱ이 ㅅㄹㅇ ⑩ & \kfTBlank \\\hline
짝 작품 & 「ㅅㅁㅇㄱ」과 쌍벽 ⑪ & \kfTBlank \\\hline
창작 배경 & ㅅㅈ에서 물러나 있던 시절 ⑫ & \kfTBlank \\\hline
\end{kfActTable}
\kfSecHeader{kfAreaLiterature}{문제}{}

\begin{multicols}{2}\raggedcolumns\setlength{\columnsep}{12pt}
\begin{kfQBlock}{1}{}
\kfQStem{윗글에 대한 설명으로 적절하지 않은 것은?}
\begin{kfChoices}
\kfChoice 여성 화자를 설정하여 님에 대한 그리움을 사랑의 감정으로 표현하고 있다.
\kfChoice '이 몸 삼기실 제 님을 조차 삼기시니'에서 님과의 인연이 태어날 때부터 정해진 것임을 강조하고 있다.
\kfChoice 화자는 님이 자신을 알아보지 못하는 것에 대해 원망하며 님을 떠나려 한다.
\kfChoice 꽃비(곳비)가 되어 님의 옷에 오르겠다는 것은 변신 모티프를 활용한 표현이다.
\kfChoice '님이야 날인 줄 모르셔도'에서 일방적 헌신의 태도가 드러난다.
\end{kfChoices}
\end{kfQBlock}

\begin{kfQBlock}{2}{}
\kfQStem{윗글에서 '곳비(꽃비)'가 지닌 문학적 기능으로 가장 적절한 것은?}
\begin{kfChoices}
\kfChoice 화자의 죽음과 영원한 소멸에 대한 깊은 두려움을 상징한다.
\kfChoice 형체를 버려서라도 님 곁에 있고 싶은 간절한 그리움이다.
\kfChoice 계절의 변화를 알리는 자연 현상에 대한 객관적 묘사이다.
\kfChoice 님이 화자를 다시 알아볼 수 있게 만드는 소통의 수단이다.
\kfChoice 화자의 경제적 궁핍과 일상의 고단한 삶을 비유한 표현이다.
\end{kfChoices}
\end{kfQBlock}

\begin{kfQBlock}{3}{}
\kfQStem{윗글의 화자가 '님'에 대해 취하는 태도를 가장 적절하게 설명한 것은?}
\begin{kfChoices}
\kfChoice 님의 무관심에 대해 분노하며 적극적으로 항의한다.
\kfChoice 님과의 인연을 숙명으로 받아들여 변함없이 따르려 한다.
\kfChoice 님을 잊기 위해 의도적으로 거리를 두려 한다.
\kfChoice 님의 잘못을 지적하고 반성을 촉구한다.
\kfChoice 님에 대한 그리움을 떨치고 새로운 삶의 길을 모색한다.
\end{kfChoices}
\end{kfQBlock}

\begin{kfQBlock}{4}{}
\kfQStem{'눈을 감으면 생각이 새로워라'라는 구절이 화자의 심리 상태를 드러내는 방식으로 가장 적절한 것은?}
\begin{kfChoices}
\kfChoice 꿈속에서 님을 직접 만나 기쁨을 느끼는 화자의 만족감을 나타낸다.
\kfChoice 잠을 자려 해도 님 생각이 더욱 또렷해져 잠들 수 없는 괴로움을 보인다.
\kfChoice 눈을 감는 행위를 통해 현실에서 벗어나 마음의 평화를 얻는 모습이다.
\kfChoice 꿈에서 깨어난 뒤 님에 대한 기억이 점점 희미해지는 아쉬움을 보인다.
\kfChoice 새로운 생각이 떠올라 님을 잊고 다른 관심사로 향하는 전환을 보여 준다.
\end{kfChoices}
\end{kfQBlock}

\begin{kfQBlock}{5}{}
\kfQStem{윗글을 바탕으로 <보기>를 감상한 내용으로 적절하지 않은 것은?

<보기> 「사미인곡」에서 정철은 굴원의 이소(離騷) 전통을 계승하여, 임금과 신하의 관계를 남녀의 관계에 빗대어 표현했다. 이 작품에서 '님'은 임금을, 화자인 '여성'은 신하인 정철 자신을 의미한다. 이러한 표현 방식은 직접적인 충성 표현보다 문학적 깊이를 더한다.}
\begin{kfChoices}
\kfChoice '이 몸 삼기실 제 님을 조차 삼기시니'는 신하가 태어날 때부터 임금과의 인연이 정해져 있다는 군신 관계의 숙명성을 표현한 것이겠군.
\kfChoice '잠을 자니 꿈에나 보려 하되'는 낙향한 신하가 임금을 만나기 어려운 현실적 처지를 반영한 것이겠군.
\kfChoice '곳비 되여 님의 옷세 올으리라'는 어떤 형태로든 임금 곁에 있고 싶은 충정을 비유적으로 표현한 것이겠군.
\kfChoice '님이야 날인 줄 모르셔도'는 임금에게 잊힌 신하의 처지를 드러낸 것이겠군.
\kfChoice '못 잊었어 하노라'는 임금에 대한 원한이 쌓여 더는 충성하지 않겠다는 결심을 나타낸 것이겠군.
\end{kfChoices}
\end{kfQBlock}

\end{multicols}

\kfSecHeader{kfAreaLiterature}{지문}{강호사시가 — 맹사성 (시조(연시조)) / 맹사성, 「강호사시가」, 수능·학평 기출 다수(현대어 정역)}

\kfPassageNote{강호사시가 — 맹사성 (시조(연시조)) / 맹사성, 「강호사시가」, 수능·학평 기출 다수(현대어 정역)}{%
강호에 봄이 드니 미친 흥이 절로 난다 \\[4pt]
시냇가 막걸리에 쏘가리 안주로다 \\[4pt]
이 몸이 한가한 것도 역시 임금의 은혜로다 \\[14pt]
강호에 여름이 드니 초당에 일이 없다 \\[4pt]
미더운 강 물결이 보내는 것은 바람이로다 \\[4pt]
이 몸이 서늘한 것도 역시 임금의 은혜로다 \\[14pt]
강호에 가을이 드니 고기마다 살져 있다 \\[4pt]
조그마한 배에 그물 실어 흐르게 던져두고 \\[4pt]
이 몸이 소일하는 것도 역시 임금의 은혜로다 \\[14pt]
강호에 겨울이 드니 눈 깊이 자가 넘다 \\[4pt]
삿갓 비껴쓰고 도롱이로 옷을 삼아 \\[4pt]
이 몸이 춥지 않은 것도 역시 임금의 은혜로다
}
\kfSecHeader{kfAreaLiterature}{활동}{작품 정리표 — 강호사시가(맹사성)}

\noindent\textit{\small 빈칸에 알맞은 말을 채워 넣으시오. (초성 힌트 참고)}\par\medskip
\begin{kfActTable}{|>{\centering\arraybackslash}m{2.6cm}|m{6cm}|m{4cm}|}
\rowcolor{kfPrimaryLight!50}\textbf{\color{kfPrimary}항목} & \textbf{\color{kfPrimary}내용 (빈칸 채우기)} & \textbf{\color{kfPrimary}답안 작성}\\\hline
갈래 & ㅇㅅㅈ(4수) ① & \kfTBlank \\\hline
작가 & ㅁㅅㅅ ② & \kfTBlank \\\hline
전개 구조 & ㅅㄱㅈ의 순서에 따른 전개 ③ & \kfTBlank \\\hline
후렴구 & ㅇㄱㅇ이샤 ㅇㅎ야 ㄷㄹ랴 ④ & \kfTBlank \\\hline
봄의 경물 & 탁료계변의 ㄱㄹㅇ ⑤ & \kfTBlank \\\hline
여름의 경물 & ㅊㄷ에 일이 없고, ㄱㅍ가 보내는 ㅂㄹ ⑥ & \kfTBlank \\\hline
가을의 행위 & ㅅㅈ에 그물 싣고 ㄷㅂ 아래 놓음 ⑦ & \kfTBlank \\\hline
겨울의 차림 & ㅅㄱ 쓰고 ㄴㅂㅇ 입고 ⑧ & \kfTBlank \\\hline
화자의 태도 & 자연 속 ㅍㄹ + ㅇㄱ에 대한 ㄱㅅ ⑨ & \kfTBlank \\\hline
문학사적 의의 & ㄱㅎㄱㄷ의 효시 ⑩ & \kfTBlank \\\hline
주제 1 & ㄱㅎ ㅎㄱ의 아름다움과 즐거움 ⑪ & \kfTBlank \\\hline
주제 2 & ㅇㄱ의 ㅇㅎ에 대한 감사 ⑫ & \kfTBlank \\\hline
\end{kfActTable}
\kfSecHeader{kfAreaLiterature}{문제}{}

\begin{multicols}{2}\raggedcolumns\setlength{\columnsep}{12pt}
\begin{kfQBlock}{1}{}
\kfQStem{윗시에 대한 설명으로 적절하지 않은 것은?}
\begin{kfChoices}
\kfChoice 사계절의 순서에 따라 시상이 전개되는 연시조이다.
\kfChoice 각 수의 종장에 '역군은이샤 은혜야 다르랴'가 반복되고 있다.
\kfChoice 봄·여름·가을·겨울의 대표적 경물을 통해 강호의 아름다움을 구체화하고 있다.
\kfChoice 화자는 자연 속에서 한가로운 삶을 즐기며 만족감을 드러내고 있다.
\kfChoice 화자는 임금의 은혜에 감사하기보다 정치적 소외에 대한 원망을 드러내고 있다.
\end{kfChoices}
\end{kfQBlock}

\begin{kfQBlock}{2}{}
\kfQStem{윗시에서 '역군은(亦君恩)이샤'가 지닌 의미로 가장 적절한 것은?}
\begin{kfChoices}
\kfChoice 임금의 은혜가 너무 적어 불만스럽다는 뜻이다.
\kfChoice 자연 속의 즐거움도 결국 임금의 은혜 덕분이라는 뜻이다.
\kfChoice 임금이 은혜를 베풀어 관직에 복귀시켜 달라는 요청이다.
\kfChoice 임금에게 자연의 아름다움을 알려 주고 싶다는 뜻이다.
\kfChoice 임금의 은혜와 자연의 아름다움은 서로 무관하다는 뜻이다.
\end{kfChoices}
\end{kfQBlock}

\begin{kfQBlock}{3}{}
\kfQStem{윗시의 '가을' 수에서 '소정에 그물 싣고 달빛 아래 놓으리라'가 드러내는 화자의 삶의 모습으로 가장 적절한 것은?}
\begin{kfChoices}
\kfChoice 생계를 위해 고된 어업 노동에 종사하는 모습
\kfChoice 자연 속에서 여유롭게 풍류를 즐기는 한가한 삶
\kfChoice 달빛 아래서 학문에 몰두하는 근면한 삶
\kfChoice 어업 기술을 연마하여 성공하려는 야심찬 모습
\kfChoice 자연 속에서 고독과 외로움을 견디는 쓸쓸한 삶
\end{kfChoices}
\end{kfQBlock}

\begin{kfQBlock}{4}{}
\kfQStem{윗시의 '겨울' 수에서 '삿갓 쓰고 누비옷 입고 눈 위에 홀로 가니'라는 구절의 효과로 가장 적절한 것은?}
\begin{kfChoices}
\kfChoice 겨울의 혹독한 추위 속에서 고통받는 화자의 비참한 처지를 강조한다.
\kfChoice 소박한 차림으로 눈 덮인 자연을 거니는 풍류인의 모습을 형상화한다.
\kfChoice 전쟁의 참혹함 속에서 홀로 살아남은 화자의 비장함을 드러낸다.
\kfChoice 관직에서의 쫓겨남을 상징적으로 나타낸다.
\kfChoice 눈 위를 홀로 걷는 모습을 통해 화자의 반사회적 태도를 보여준다.
\end{kfChoices}
\end{kfQBlock}

\begin{kfQBlock}{5}{}
\kfQStem{윗시를 바탕으로 <보기>를 감상한 내용으로 적절하지 않은 것은?

<보기> 조선 전기 사대부 문인들은 '강호가도(江湖歌道)'라 하여 자연 속 은거 생활을 시조로 노래하는 전통을 형성했다. 이들에게 자연은 정치에서 물러난 공간이면서도 군주에 대한 충성을 잊지 않는 공간이었다. 이러한 이중적 태도가 강호가도의 핵심이다.}
\begin{kfChoices}
\kfChoice 매 수마다 '역군은이샤'를 반복한 것은 자연 속에서도 군주에 대한 충성을 잊지 않는 강호가도의 태도를 잘 보여주겠군.
\kfChoice 각 계절의 경물을 구체적으로 제시한 것은 강호 생활의 실감을 높이기 위한 것이겠군.
\kfChoice '미친 듯이 즐거워라'에서 화자가 자연에 몰입하는 모습은 은거 생활의 만족을 보여주겠군.
\kfChoice 화자가 자연 속에서 소일하면서도 군은에 감사하는 것은 정치 참여와 은거 사이의 이중적 태도를 반영하겠군.
\kfChoice '역군은이샤'는 군주에 대한 풍자적 비판을 후렴구에 은밀히 담은 것이겠군.
\end{kfChoices}
\end{kfQBlock}

\end{multicols}

\kfSecHeader{kfAreaLiterature}{지문}{만흥 — 윤선도 (시조(연시조)) / 윤선도, 「만흥」, 수능·학평 기출 다수}

\kfPassageNote{만흥 — 윤선도 (시조(연시조)) / 윤선도, 「만흥」, 수능·학평 기출 다수}{%
산수간(山水間) 바위 아래 뫼집을 짓노라 하니 \\[4pt]
그 모른 남이 웃는다 말리라. \\[4pt]
어리고 하물며 일도 다 못하거든 \\[4pt]
글 닐거 무엇 하리. \\[14pt]
보리밥 풋나물을 알마초 먹을 것이 \\[4pt]
있거든 그만이지 부귀를 부러랴. \\[4pt]
내 생애 이러하되 설운 일이 하나 없다. \\[14pt]
내 벗이 멧이나 하니 수석(水石)이 사(四)이로다. \\[4pt]
구름 빗치 좋다 하나 검기를 자로 한다. \\[4pt]
바람이 맑다 하나 그칠 적이 하노매라. \\[4pt]
좋을씨고, 믿을씨고, 수석(水石)과 송죽(松竹)을. \\[14pt]
더우면 꽃 피고 추우면 잎 지거늘 \\[4pt]
솔아 너는 어이 눈서리를 모르는다. \\[4pt]
구태여 곧은 줄을 뉘라서 알랴마는 \\[4pt]
꼿 곧고 늘 푸르니 그를 좋아하노라. \\[14pt]
나무도 아닌 것이 풀도 아닌 것이 \\[4pt]
곧기는 뉘 시기며 속은 어이 비었는다. \\[4pt]
저러고 사시(四時)에 푸르니 그를 좋아하노라. \\[14pt]
작은 것이 높이 떠 만물을 다 비추니 \\[4pt]
밤중의 광명이 너만 한 이 또 있느냐. \\[4pt]
보고 듣고 먹는 것이 흥(興)이 아닌 것 없다.
}
\kfSecHeader{kfAreaLiterature}{활동}{작품 정리표 — 만흥(윤선도)}

\noindent\textit{\small 빈칸에 알맞은 말을 채워 넣으시오. (초성 힌트 참고)}\par\medskip
\begin{kfActTable}{|>{\centering\arraybackslash}m{2.6cm}|m{6cm}|m{4cm}|}
\rowcolor{kfPrimaryLight!50}\textbf{\color{kfPrimary}항목} & \textbf{\color{kfPrimary}내용 (빈칸 채우기)} & \textbf{\color{kfPrimary}답안 작성}\\\hline
갈래 & ㅇㅅㅈ(6수) ① & \kfTBlank \\\hline
작가 & ㅇㅅㄷ(고산) ② & \kfTBlank \\\hline
제목 의미 & 느긋이 이는 ㅎㅊ ③ & \kfTBlank \\\hline
벗의 수 & 수석이 ㅅ이로다 ④ & \kfTBlank \\\hline
네 벗 & ㅁ·ㅂㅇ·ㅅㄴㅁ·ㄷㄴㅁ ⑤ & \kfTBlank \\\hline
배제한 자연물 & ㄱㄹ(변덕), ㅂㄹ(그침) ⑥ & \kfTBlank \\\hline
솔의 미덕 & 꼿 ㄱㄱ ㄴ ㅍㄹ ⑦ & \kfTBlank \\\hline
대나무의 미덕 & ㄱㄱ + 속이 비고 ㅅㅅ에 푸름 ⑧ & \kfTBlank \\\hline
달의 미덕 & ㅂㅈ의 ㄱㅁ ⑨ & \kfTBlank \\\hline
화자의 식사 & ㅂㄹㅂ ㅍㄴㅁ ⑩ & \kfTBlank \\\hline
주제 1 & ㅈㅇ ㅅ ㅎㅊ ⑪ & \kfTBlank \\\hline
주제 2 & ㅇㅂㄴㄷ ⑫ & \kfTBlank \\\hline
\end{kfActTable}
\kfSecHeader{kfAreaLiterature}{문제}{}

\begin{multicols}{2}\raggedcolumns\setlength{\columnsep}{12pt}
\begin{kfQBlock}{1}{}
\kfQStem{윗시에 대한 설명으로 적절한 것은?}
\begin{kfChoices}
\kfChoice 자연물을 의인화하여 벗으로 삼고 교감하는 태도가 나타난다.
\kfChoice 화자는 부귀영화를 추구하며 입신양명의 의지를 드러내고 있다.
\kfChoice 계절의 변화에 따른 시상 전개가 후렴구로 통일되고 있다.
\kfChoice 화자는 구름과 바람을 가장 믿을 수 있는 벗으로 삼고 있다.
\kfChoice 관직 생활의 어려움을 한탄하며 자연 속 삶을 소망하고 있다.
\end{kfChoices}
\end{kfQBlock}

\begin{kfQBlock}{2}{}
\kfQStem{윗시에서 화자가 구름과 바람을 '벗'으로 삼지 않는 까닭으로 가장 적절한 것은?}
\begin{kfChoices}
\kfChoice 구름과 바람은 본디 자연이 아니라 인위적으로 만들어진 산물이기 때문이다.
\kfChoice 구름은 자주 검어지고 바람은 그치므로 변함없는 벗이 못 되기 때문이다.
\kfChoice 구름과 바람은 인간의 일상에 자주 해를 끼쳐 두려운 존재이기 때문이다.
\kfChoice 구름과 바람은 수석이나 송죽보다 그 아름다움이 부족하다 여겨지기 때문이다.
\kfChoice 구름과 바람은 겨울철에는 사라져 더 이상 존재하지 않기 때문이다.
\end{kfChoices}
\end{kfQBlock}

\begin{kfQBlock}{3}{}
\kfQStem{윗시에서 '솔'과 '대나무'를 좋아하는 공통된 이유로 가장 적절한 것은?}
\begin{kfChoices}
\kfChoice 아름다운 꽃을 피워 화자의 눈을 즐겁게 하기 때문이다.
\kfChoice 사시사철 푸르고 곧아 변함없는 절개를 상징하기 때문이다.
\kfChoice 화자에게 먹을 것을 제공하는 실용적 가치가 있기 때문이다.
\kfChoice 겨울에만 아름다워 다른 계절과 차별되기 때문이다.
\kfChoice 화자가 살고 있는 집의 건축 재료이기 때문이다.
\end{kfChoices}
\end{kfQBlock}

\begin{kfQBlock}{4}{}
\kfQStem{윗시의 마지막 수에서 '작은 것이 높이 떠 만물을 다 비추니'의 '작은 것'이 가리키는 대상과 그 효과로 적절한 것은?}
\begin{kfChoices}
\kfChoice 별 — 화자의 외로움을 심화하는 효과를 낸다.
\kfChoice 달 — 밤의 어둠 속에서 만물을 비추는 존재로, 자연의 경이로움을 부각한다.
\kfChoice 반딧불 — 작은 빛이 큰 어둠을 이기는 역설을 나타낸다.
\kfChoice 해 — 낮의 태양을 겸손하게 표현하여 화자의 겸양을 드러낸다.
\kfChoice 구름 — 하늘에 떠서 만물을 가리는 존재로, 세상의 혼탁함과 흐림을 상징한다.
\end{kfChoices}
\end{kfQBlock}

\begin{kfQBlock}{5}{}
\kfQStem{윗시를 바탕으로 <보기>를 감상한 내용으로 적절하지 않은 것은?

<보기> 조선 시대 사대부들은 유교적 덕목인 '안빈낙도(安貧樂道)'를 실천하고자 했다. 이는 가난한 삶 속에서도 도(道)를 즐기며 흔들리지 않는 마음가짐을 뜻한다. 자연 속의 삶은 이러한 정신을 실현하는 공간으로 여겨졌다.}
\begin{kfChoices}
\kfChoice '보리밥 풋나물을 알마초 먹을 것이'에서 소박한 식사는 안빈낙도의 물질적 측면을 보여 주겠군.
\kfChoice '부귀를 부러랴'에서 세속적 가치를 초월하는 태도는 안빈낙도의 정신적 측면을 보여 주겠군.
\kfChoice '설운 일이 하나 없다'에서 가난 속에서도 만족하는 화자의 태도는 안빈낙도와 깊게 연결되겠군.
\kfChoice 수석·송죽을 벗으로 삼는 모습은 자연 속에서 변함없는 도(道)를 추구하는 자세와 연결되겠군.
\kfChoice '보고 듣고 먹는 것이 흥이 아닌 것 없다'는 자연을 떠나 세속에서 즐거움을 찾으려는 의지겠군.
\end{kfChoices}
\end{kfQBlock}

\end{multicols}

\kfSecHeader{kfAreaLiterature}{지문}{고공가 — 허전 (가사) / 허전, 「고공가」}

\kfPassageNote{고공가 — 허전 (가사) / 허전, 「고공가」}{%
이 집 지은 지 몇 해나 되었으며 \\[4pt]
집 지은 연분(緣分)이 없다 할 수 있으랴. \\[4pt]
대들보 석가래며 기둥이 바로 서야 \\[4pt]
풍우(風雨)를 능히 막아 상하(上下) 평안하리라. \\[14pt]
집안의 고공(雇工)들아 내 말씀 들어 보오. \\[4pt]
혼자 힘 부족하여 너희를 의지하니 \\[4pt]
너희 맡은 소임(所任)을 성심껏 다하여라. \\[14pt]
밭갈 때 바로 갈고 씨뿌릴 때 골로 뿌려 \\[4pt]
김매기 정히 매고 거둘 때 다 거두소. \\[4pt]
농사(農事) 성패(成敗) 달렸으니 마음을 놓지 마라. \\[14pt]
네 직분(職分) 다하거든 품삯(品?)을 넉넉히 줄 것이요 \\[4pt]
직분을 게을리하면 종시(終始) 용서 못하리라. \\[4pt]
상벌(賞罰)이 분명하니 너희가 모를소냐.
}
\kfSecHeader{kfAreaLiterature}{활동}{작품 정리표 — 고공가(허전)}

\noindent\textit{\small 빈칸에 알맞은 말을 채워 넣으시오. (초성 힌트 참고)}\par\medskip
\begin{kfActTable}{|>{\centering\arraybackslash}m{2.6cm}|m{6cm}|m{4cm}|}
\rowcolor{kfPrimaryLight!50}\textbf{\color{kfPrimary}항목} & \textbf{\color{kfPrimary}내용 (빈칸 채우기)} & \textbf{\color{kfPrimary}답안 작성}\\\hline
갈래 & ㄱㅅ ① & \kfTBlank \\\hline
작가 & ㅎㅈ ② & \kfTBlank \\\hline
'고공'의 뜻 & ㅁㅅ(품팔이꾼) ③ & \kfTBlank \\\hline
비유 대응: 주인 & = ㅇㄱ ④ & \kfTBlank \\\hline
비유 대응: 머슴 & = ㅅㅎ ⑤ & \kfTBlank \\\hline
비유 대응: 농사 & = ㄴㄹ ㄷㅅㄹ ⑥ & \kfTBlank \\\hline
비유 대응: 집 & = ㄴㄹ ⑦ & \kfTBlank \\\hline
비유 대응: 풍우 & = ㄱㄴ(외부 위기) ⑧ & \kfTBlank \\\hline
보상 원칙 & 직분 다하면 ㅍㅅ 넉넉, 게으르면 ㅇㅅ 불가 ⑨ & \kfTBlank \\\hline
핵심 원칙 & ㅅㅂ이 분명 ⑩ & \kfTBlank \\\hline
주제 & ㅇㅁ ㅈㅊ의 정신 ⑪ & \kfTBlank \\\hline
작품 목적 & ㄱㅎ 목적의 가사 ⑫ & \kfTBlank \\\hline
\end{kfActTable}
\kfSecHeader{kfAreaLiterature}{문제}{}

\begin{multicols}{2}\raggedcolumns\setlength{\columnsep}{12pt}
\begin{kfQBlock}{1}{}
\kfQStem{윗글에 대한 설명으로 적절하지 않은 것은?}
\begin{kfChoices}
\kfChoice 주인이 머슴에게 일을 당부하는 형식을 취하고 있다.
\kfChoice 농사일의 비유를 통해 신하의 도리를 전달하고 있다.
\kfChoice 상벌의 분명함을 강조하여 경고와 약속을 병행하고 있다.
\kfChoice 화자는 머슴의 처지에서 주인의 부당한 대우를 고발하고 있다.
\kfChoice 집(나라)의 안정을 위해 각자의 소임을 다할 것을 촉구하고 있다.
\end{kfChoices}
\end{kfQBlock}

\begin{kfQBlock}{2}{}
\kfQStem{윗글에서 '대들보 석가래며 기둥이 바로 서야'가 비유하는 바로 가장 적절한 것은?}
\begin{kfChoices}
\kfChoice 건축 기술의 발전이 국가가 번영하는 기초가 된다.
\kfChoice 나라의 근간 신하들이 바르게 서야 국가가 안정된다.
\kfChoice 집을 지을 때는 재료의 품질이 무엇보다 중요하다.
\kfChoice 자연재해에 대비하여 견고한 방어 시설을 갖추어야 한다.
\kfChoice 주인이 직접 건축에 참여해야 집이 더욱 튼튼해진다.
\end{kfChoices}
\end{kfQBlock}

\begin{kfQBlock}{3}{}
\kfQStem{윗글의 화자가 머슴(고공)에게 전달하고자 하는 핵심 메시지로 가장 적절한 것은?}
\begin{kfChoices}
\kfChoice 머슴은 주인에게 충성하되 부당한 명령에는 불복해야 한다.
\kfChoice 맡은 소임을 성실히 다하면 보상을, 게을리하면 처벌이 따른다.
\kfChoice 머슴은 주인 없이도 스스로 농사일을 판단하고 결정해야 한다.
\kfChoice 농사보다 학문에 힘써야 진정한 보상을 받을 수 있다.
\kfChoice 머슴끼리 서로 감시하여 게으른 자를 내쫓아야 한다.
\end{kfChoices}
\end{kfQBlock}

\begin{kfQBlock}{4}{}
\kfQStem{'밭갈 때 바로 갈고 씨뿌릴 때 골로 뿌려 / 김매기 정히 매고 거둘 때 다 거두소'에서 사용된 표현 기법과 그 효과로 가장 적절한 것은?}
\begin{kfChoices}
\kfChoice 점층법 — 점차 강도가 높아지는 노동을 통해 긴장감을 조성한다.
\kfChoice 열거법 — 농사의 전 과정을 나열하여 소임의 구체적 내용을 전달한다.
\kfChoice 대구법 — 상반되는 내용을 대조하여 갈등을 부각한다.
\kfChoice 역설법 — 모순된 표현을 통해 진리를 드러낸다.
\kfChoice 반어법 — 겉뜻과 속뜻을 반대로 하여 풍자한다.
\end{kfChoices}
\end{kfQBlock}

\begin{kfQBlock}{5}{}
\kfQStem{윗글을 바탕으로 <보기>를 감상한 내용으로 적절하지 않은 것은?

<보기> 조선 시대 일부 가사 작품은 직접적인 정치적 발언 대신, 일상적 소재를 비유로 활용하여 정치적 메시지를 전달하는 방식을 취했다. 「고공가」는 주인과 머슴이라는 일상적 관계를 빌려 군신(君臣) 관계의 도리를 노래한 대표적 사례이다.}
\begin{kfChoices}
\kfChoice 농사의 성패가 고공에게 달렸다는 것은 나라 다스림의 성패가 신하의 성실함에 달렸다는 비유로 볼 수 있겠군.
\kfChoice 품삯을 넉넉히 주겠다는 약속은 성실한 신하에 대한 포상의 비유로 이해할 수 있겠군.
\kfChoice '풍우를 능히 막아'는 외부의 위기(국난)를 막는다는 비유로 읽을 수 있겠군.
\kfChoice 주인이 고공에게 당부하는 형식은 정치적 메시지를 부드럽게 전달하는 효과가 있겠군.
\kfChoice 이 작품에서 주인은 백성을, 고공은 임금을 비유한 것이므로, 백성이 임금에게 바른 통치를 촉구하는 것이겠군.
\end{kfChoices}
\end{kfQBlock}

\end{multicols}

\kfStartNonlit{이번 챕터 비문학 지문 5편}


\kfSecHeader{kfAreaNonlit}{지문}{노자와 한비자 — 도가와 법가의 분기 (인문)}

\kfPassageNote{노자와 한비자 — 도가와 법가의 분기 (인문)}{%
노자(老子)는 '도(道)'를 만물의 근원으로 보았다. 도는 언어로 규정할 수 없는 것으로, 이름 붙일 수 있는 도는 이미 참된 도가 아니다. 노자에 따르면 도의 본성은 '무위자연(無爲自然)', 곧 인위적 작위 없이 스스로 그러한 상태에 있는 것이다. 도는 만물을 낳되 소유하지 않고, 만물을 기르되 지배하지 않는다. 노자는 이러한 도의 원리를 인간 사회에 적용하여, 통치자 역시 백성의 삶에 인위적으로 간섭하지 않아야 한다고 주장했다. 이것이 '무위이치(無爲而治)', 곧 함이 없이 다스림이다.

노자의 정치철학에서 이상적 통치자는 백성이 그 존재를 거의 의식하지 못하는 자이다. 통치자가 법령을 남발하고 형벌을 강화할수록 도둑이 늘어나며, 지혜와 기교를 숭상할수록 거짓이 번성한다고 보았다. 따라서 노자는 소국과민(小國寡民)의 사회를 이상향으로 제시했다. 작은 나라에 적은 백성이 살면서 먼 곳까지 가지 않고, 문자를 쓰지 않으며, 이웃 나라의 닭 울음소리가 들리되 서로 왕래하지 않는 소박한 공동체가 그것이다.

한비자(韓非子)는 노자 사상의 영향을 받았으나 전혀 다른 방향으로 발전시켰다. 한비자가 노자에게서 계승한 것은 도의 형이상학적 측면이 아니라, 군주가 사적인 감정이나 편견에 흔들리지 않아야 한다는 '무위'의 원리였다. 그러나 한비자에게 무위란 아무것도 하지 않는 것이 아니라, 객관적 제도에 의해 국가가 저절로 운영되게 만드는 것을 의미했다. 이를 위해 한비자는 법(法)·술(術)·세(勢)의 세 가지 통치 수단을 제시했다.

법(法)은 명문화된 규범으로, 모든 사람에게 공평하게 적용되어야 한다. 법 앞에서 귀천의 차별이 없어야 하며, 공을 세운 자에게는 반드시 상을 주고 죄를 범한 자에게는 반드시 벌을 내려야 한다. 이것이 '신상필벌(信賞必罰)'이다. 술(術)은 군주가 신하를 부리는 기술이다. 군주는 자신의 속마음을 드러내지 않으며, 신하가 스스로 말한 것과 실제 업적을 대조하여 평가하는 '형명참동(形名參同)'의 방법을 쓴다. 세(勢)는 군주의 지위에서 나오는 권세이다. 아무리 현명한 군주라도 권세 없이는 통치할 수 없으며, 세는 법과 술이 작동하기 위한 토대이다.

노자와 한비자는 모두 인의(仁義)와 예(禮)를 강조하는 유가를 비판했다는 공통점이 있다. 그러나 그 비판의 방향은 달랐다. 노자는 인의와 예가 도를 잃은 뒤에 만들어진 인위적 산물이라 보아 도로 돌아갈 것을 주장했고, 한비자는 인의와 예만으로는 혼란한 시대를 다스릴 수 없으니 법과 형벌이라는 현실적 수단이 필요하다고 주장했다. 같은 뿌리에서 출발하되 하나는 자연으로의 회귀를, 다른 하나는 제도의 정비를 지향한 셈이다.
}
\kfSecHeader{kfAreaNonlit}{문제}{}

\begin{multicols}{2}\raggedcolumns\setlength{\columnsep}{12pt}
\begin{kfQBlock}{1}{}
\kfQStem{윗글의 내용과 일치하는 것은?}
\begin{kfChoices}
\kfChoice 노자는 통치자가 적극적으로 법령을 제정하여 백성을 이끌어야 한다고 보았다.
\kfChoice 한비자는 노자의 형이상학적 도의 개념을 그대로 계승하여 법가 사상을 확립했다.
\kfChoice 한비자의 '술(術)'은 신하의 말과 실적을 대조하여 평가하는 방법을 포함한다.
\kfChoice 노자와 한비자는 모두 유가의 인의(仁義)를 긍정적으로 수용하여 발전시켰다.
\kfChoice 노자의 소국과민은 강대한 군사력을 갖춘 이상 국가를 의미한다.
\end{kfChoices}
\end{kfQBlock}

\begin{kfQBlock}{2}{}
\kfQStem{윗글을 바탕으로 한비자가 '법(法)·술(術)·세(勢)'를 제시한 까닭으로 가장 적절한 것은?}
\begin{kfChoices}
\kfChoice 노자의 무위자연 사상을 도덕적으로 더욱 완성하려 했기 때문이다.
\kfChoice 유가의 인의와 예를 객관적 법적 제도로 보완하려 했기 때문이다.
\kfChoice 군주의 사적 감정에 기대지 않는 객관적 통치 체계를 세우려 했기 때문이다.
\kfChoice 노자가 말한 소국과민의 이상 사회를 현실에서 구현하려 했기 때문이다.
\kfChoice 도(道)의 형이상학적 원리를 정치 영역에서 증명하려 했기 때문이다.
\end{kfChoices}
\end{kfQBlock}

\begin{kfQBlock}{3}{}
\kfQStem{윗글을 바탕으로 <보기>를 이해한 내용으로 적절하지 않은 것은?

<보기> A국의 왕은 유능한 재상 갑(甲)에게 국정을 일임한 뒤 자신은 궁궐 깊숙이 물러나 아무런 지시도 내리지 않았다. 갑은 스스로의 판단으로 정책을 시행하여 나라를 잘 다스렸으나, 시간이 지나자 갑의 권위가 왕을 능가하게 되었다. 백성은 갑의 이름만 알 뿐 왕의 존재를 잊었다.}
\begin{kfChoices}
\kfChoice 노자의 관점에서 A국의 왕이 백성에게 잊힌 것은 이상적 통치에 가까운 모습이라 할 수 있다.
\kfChoice 한비자의 관점에서 왕이 갑에게 국정을 일임한 것은 '세(勢)'를 상실할 위험이 있는 행위이다.
\kfChoice 한비자의 관점에서 왕은 갑의 말과 실적을 '형명참동'으로 점검했어야 한다.
\kfChoice 한비자의 관점에서 왕의 행위는 '무위'를 올바르게 실현한 사례에 해당한다.
\kfChoice 노자와 한비자 모두 왕이 사적 감정에 흔들리지 않았다는 점에서는 긍정적으로 평가할 수 있다.
\end{kfChoices}
\end{kfQBlock}

\begin{kfQBlock}{4}{}
\kfQStem{윗글의 ㉠ '신상필벌'과 ㉡ '형명참동'에 대한 설명으로 적절한 것은?}
\begin{kfChoices}
\kfChoice ㉠은 법을 공평하게 적용하는 원칙을, ㉡은 신하에 대한 감성적 신뢰를 강조한다.
\kfChoice ㉠은 상벌의 확실성을 통한 법의 실효성을, ㉡은 말과 실적 대조를 통한 평가를 뜻한다.
\kfChoice ㉠과 ㉡은 모두 노자의 무위자연 원리를 직접적으로 계승하여 만든 통치 방법이다.
\kfChoice ㉠은 군주의 권세인 '세(勢)'에, ㉡은 객관적 제도인 '법(法)'에 해당하는 개념이다.
\kfChoice ㉠과 ㉡은 모두 공자와 맹자가 강조한 유가의 인의(仁義) 사상에 기반한 개념이다.
\end{kfChoices}
\end{kfQBlock}

\end{multicols}

\kfSecHeader{kfAreaNonlit}{지문}{장재·정이·주희의 기론 — 이기론의 전개 (인문)}

\kfPassageNote{장재·정이·주희의 기론 — 이기론의 전개 (인문)}{%
송대(宋代) 성리학의 핵심 문제 중 하나는 만물의 존재 근거를 어떻게 설명할 것인가 하는 것이었다. 이 물음에 대해 장재(張載), 정이(程頤), 주희(朱熹)는 각기 다른 방식으로 답했는데, 이들의 논의는 후대 성리학 전개의 근간이 되었다.

장재는 '기(氣)'를 만물의 유일한 실체로 보았다. 이를 '기일원론(氣一元論)'이라 한다. 장재에 따르면 기가 모이면 만물이 생겨나고, 기가 흩어지면 만물이 소멸한다. 기의 모이고 흩어짐은 끊임없이 반복되며, 이것이 우주 만물의 생성과 변화를 설명한다. 장재는 기가 흩어져도 완전히 소멸하는 것은 아니라고 보았다. 기가 흩어진 상태를 '태허(太虛)'라 하는데, 태허는 텅 비어 있는 것처럼 보이지만 실은 기의 본래 상태이다. 따라서 무(無)에서 유(有)가 생겨나는 것이 아니라, 기의 응취(凝聚)와 산일(散逸)이 있을 뿐이다.

정이는 장재와 달리 '이(理)'와 '기(氣)'를 구분하는 '이기이원론(理氣二元論)'을 제시했다. 정이에 따르면 이는 만물의 존재 법칙, 곧 사물이 그렇게 되는 까닭(所以然)이며, 기는 그 법칙이 실현되는 질료(質料)이다. 이와 기는 서로 다른 차원에 속하되, 현실에서는 분리될 수 없다. 이 없으면 기가 존재 근거를 갖지 못하고, 기 없으면 이가 현실에서 드러나지 못한다. 정이는 이를 형이상(形而上)의 것으로, 기를 형이하(形而下)의 것으로 규정하여 양자의 존재론적 위상을 구분했다.

주희는 장재의 기론과 정이의 이기이원론을 종합하여 성리학의 체계를 완성했다. 주희에 따르면, 이와 기는 개념적으로 분리할 수 있으나 현실적으로는 결코 떨어져 있지 않다. 이를 '불상리(不相離, 서로 떨어지지 않음)'와 '불상잡(不相雜, 서로 섞이지 않음)'이라 표현했다. 또한 주희는 존재론적으로 이가 기보다 앞선다는 '이선기후(理先氣後)'를 주장했다. 이것은 시간적 선후가 아니라 논리적·존재론적 우선성을 의미한다. 만물이 존재하려면 먼저 그 존재의 법칙(이)이 있어야 하기 때문이다.

주희는 이러한 이기론을 인간 심성론에도 적용했다. 인간의 본성(性)은 이에 근거하므로 본래 선(善)하다. 그러나 기질(氣質)의 차이에 따라 본성이 온전히 드러나지 못할 수 있다. 주희는 이것을 '이발기수(理發氣隨)', 곧 이(理)가 발(發)하면 기(氣)가 따르는 구조로 설명했다. 사단(四端: 측은·수오·사양·시비)은 이가 발한 것이며, 칠정(七情: 희·노·애·구·애·오·욕)은 기가 발한 것이다. 이처럼 주희는 인간의 마음 작용까지 이기의 구도로 체계화함으로써, 우주론과 심성론을 하나의 원리 아래 통합했다.
}
\kfSecHeader{kfAreaNonlit}{문제}{}

\begin{multicols}{2}\raggedcolumns\setlength{\columnsep}{12pt}
\begin{kfQBlock}{1}{}
\kfQStem{윗글의 내용과 일치하지 않는 것은?}
\begin{kfChoices}
\kfChoice 장재는 기가 흩어진 상태인 태허가 기의 본래 상태라고 보았다.
\kfChoice 정이는 이를 형이상의 것으로, 기를 형이하의 것으로 규정했다.
\kfChoice 주희는 이와 기가 시간적 선후 관계로 존재한다고 주장했다.
\kfChoice 장재의 기일원론에서는 무(無)에서 유(有)가 생겨나는 것이 아니다.
\kfChoice 주희는 인간의 본성(性)이 이에 근거하므로 본래 선하다고 보았다.
\end{kfChoices}
\end{kfQBlock}

\begin{kfQBlock}{2}{}
\kfQStem{윗글의 장재, 정이, 주희에 대한 이해로 가장 적절한 것은?}
\begin{kfChoices}
\kfChoice 장재는 이와 기를 명확하게 구분하여 송대 이원론 사상의 기초를 마련하였다.
\kfChoice 정이는 기(氣)만이 유일한 실체라는 기일원론의 사상적 입장을 주장하였다.
\kfChoice 주희는 장재의 기론을 전적으로 부정하고 오직 정이의 이론만을 그대로 수용했다.
\kfChoice 주희는 이와 기가 개념적으로 구분되나 현실적으로는 분리되지 않는다고 보았다.
\kfChoice 장재와 주희는 모두 기가 이보다 존재론적으로 앞선다고 보고 우선시하였다.
\end{kfChoices}
\end{kfQBlock}

\begin{kfQBlock}{3}{}
\kfQStem{윗글을 바탕으로 <보기>를 이해한 내용으로 적절하지 않은 것은?

<보기> 물은 높은 곳에서 낮은 곳으로 흐르는 법칙을 따른다. 그러나 실제로 물이 흐르려면 물이라는 물질이 있어야 한다. 법칙만 있고 물이 없으면 흐름은 일어나지 않으며, 물이 있어도 법칙이 없으면 어디로 흐를지 정해지지 않는다.}
\begin{kfChoices}
\kfChoice 정이의 관점에서 '물이 흐르는 법칙'은 이(理)에 해당한다고 볼 수 있다.
\kfChoice 정이의 관점에서 '물'이라는 물질은 기(氣)에 해당한다고 볼 수 있다.
\kfChoice 주희의 관점에서 법칙과 물이 서로 떨어질 수 없는 것은 '불상리'에 해당한다.
\kfChoice 장재의 관점에서 물이 증발하여 사라지면 기가 완전히 소멸한 것이다.
\kfChoice 주희의 관점에서 법칙이 물보다 논리적으로 앞서는 것은 '이선기후'에 해당한다.
\end{kfChoices}
\end{kfQBlock}

\begin{kfQBlock}{4}{}
\kfQStem{윗글에서 주희가 '불상리'와 '불상잡'이라는 개념을 동시에 사용한 이유로 가장 적절한 것은?}
\begin{kfChoices}
\kfChoice 이와 기가 완전히 동일한 것임을 강조하기 위해서이다.
\kfChoice 이와 기가 현실에서 함께하되 본질적으로 다른 것임을 보이려는 데 있다.
\kfChoice 이가 기에 종속된다는 점을 논증하기 위해서이다.
\kfChoice 장재의 기일원론을 비판하면서 정이의 이원론을 부정하기 위해서이다.
\kfChoice 이와 기의 시간적 선후를 설정하기 위해서이다.
\end{kfChoices}
\end{kfQBlock}

\end{multicols}

\kfSecHeader{kfAreaNonlit}{지문}{왕필과 곽상의 무(無) 해석 — 현학의 두 갈래 (인문)}

\kfPassageNote{왕필과 곽상의 무(無) 해석 — 현학의 두 갈래 (인문)}{%
위진(魏晉) 시대에 유행한 현학(玄學)은 유가와 도가의 경전을 새로운 시각에서 해석하려는 사상적 움직임이었다. 현학자들은 노자와 장자의 사상을 바탕으로 만물의 근본이 무엇인지를 탐구했는데, 이에 대한 대표적인 두 가지 답이 왕필(王弼)의 '귀무론(貴無論)'과 곽상(郭象)의 '독화론(獨化論)'이다.

왕필은 '무(無)'가 만물의 근본이라고 보았다. 여기서 무란 단순히 아무것도 없다는 뜻이 아니라, 어떤 특정한 형태나 성질로 규정될 수 없는 근원적 존재를 가리킨다. 왕필에 따르면, 특정한 형태를 가진 것(有)은 다른 것과 구별되기 때문에 제한적이다. 나무는 나무이지 돌이 아니며, 물은 물이지 불이 아니다. 이처럼 규정된 것은 다른 것이 될 수 없으므로 만물 전체의 근거가 될 수 없다. 반면 무는 어떤 형태로도 규정되지 않기 때문에 오히려 모든 것의 근거가 될 수 있다. 왕필은 이 논리로 노자의 '도(道)'를 해석했다. 도가 만물을 낳을 수 있는 것은 도 자체가 무, 곧 어떤 규정성도 갖지 않은 상태이기 때문이라는 것이다.

왕필은 이러한 존재론을 정치론에도 적용했다. 이상적인 군주는 자신의 개인적 지혜나 욕망을 앞세우지 않고 무위(無爲)의 자세를 취해야 한다. 군주가 특정한 선호를 드러내면 신하들이 거기에 영합하여 참된 모습을 숨기게 되기 때문이다. 따라서 군주는 자신을 비워 '무'의 자리에 서야 하며, 그렇게 할 때 신하와 백성이 각자의 본분에 따라 스스로 움직이게 된다.

곽상은 왕필과 근본적으로 다른 입장을 취했다. 곽상에 따르면, 무가 만물의 근본이라는 주장은 성립할 수 없다. 무는 말 그대로 없는 것이므로, 없는 것이 있는 것을 낳을 수는 없다. 곽상은 만물이 어떤 외부의 근원에 의해 만들어진 것이 아니라, 각각 스스로 그렇게 되어 존재한다고 보았다. 이를 '독화(獨化)', 곧 홀로 변화한다는 개념으로 표현했다. 만물은 타자에 의존하지 않고 저마다 자기 원인에 의해 존재하며, 이것이 곧 '자생(自生)'이자 '자화(自化)'이다.

곽상의 독화론은 만물 사이의 관계를 설명하는 방식에서도 왕필과 달랐다. 왕필에게 만물은 무라는 하나의 근본에 의해 통일되지만, 곽상에게 만물은 각각 독립적으로 존재한다. 그렇다면 만물 사이에 관계는 없는가? 곽상은 만물이 각자 독립적으로 존재하면서도 서로 자연스럽게 맞물려 조화를 이룬다고 보았다. 이를 '현명(玄冥)', 곧 어둡고 알 수 없는 조화라 했다. 누가 설계하거나 주재하는 것이 아니라 각자가 스스로 그러할 뿐인데, 결과적으로 전체가 조화를 이루는 것이다. 이처럼 왕필이 '하나의 근본(무)'에서 만물의 통일성을 찾았다면, 곽상은 '근본 없음' 속에서 만물의 자발적 조화를 발견했다.
}
\kfSecHeader{kfAreaNonlit}{문제}{}

\begin{multicols}{2}\raggedcolumns\setlength{\columnsep}{12pt}
\begin{kfQBlock}{1}{}
\kfQStem{윗글의 내용과 일치하는 것은?}
\begin{kfChoices}
\kfChoice 왕필의 '무'는 아무것도 존재하지 않는 공허한 상태를 의미한다.
\kfChoice 곽상은 만물이 무(無)라는 근원에서 비롯되었다고 주장했다.
\kfChoice 왕필과 곽상 모두 현학의 흐름 속에서 만물의 근본을 탐구했다.
\kfChoice 곽상의 독화론에서 만물은 타자의 설계에 의해 조화를 이룬다.
\kfChoice 왕필은 군주가 자신의 지혜를 적극적으로 드러내야 한다고 보았다.
\end{kfChoices}
\end{kfQBlock}

\begin{kfQBlock}{2}{}
\kfQStem{왕필이 '유(有)는 만물 전체의 근거가 될 수 없다'고 본 근거로 가장 적절한 것은?}
\begin{kfChoices}
\kfChoice 유는 감각으로 인식할 수 없어 존재를 확인할 방법이 없기 때문이다.
\kfChoice 유는 특정 형태로 규정되어 있어 다른 것이 될 수 없으므로 제한적이기 때문이다.
\kfChoice 유는 시간이 지나면 반드시 소멸하므로 영원한 근거가 될 수 없기 때문이다.
\kfChoice 유는 무에 의해 만들어진 것이므로 2차적 존재에 불과하기 때문이다.
\kfChoice 유는 도(道)의 가르침에 어긋나는 것이므로 근본이 될 자격이 없기 때문이다.
\end{kfChoices}
\end{kfQBlock}

\begin{kfQBlock}{3}{}
\kfQStem{윗글을 바탕으로 왕필의 '귀무론'과 곽상의 '독화론'을 비교한 것으로 적절한 것은?}
\begin{kfChoices}
\kfChoice 왕필은 만물이 하나의 근본에서 통일된다고 보았고, 곽상은 각각 독립 존재한다고 보았다.
\kfChoice 왕필은 만물의 근본이 본래 존재하지 않는다고 보았고, 곽상은 무가 그 근본이라고 보았다.
\kfChoice 왕필과 곽상은 모두 만물 사이에 어떤 조화도 결코 이루어질 수 없다고 보았다.
\kfChoice 왕필은 만물의 자발적 생성을, 곽상은 외부 근원에 의한 생성을 주장한 사상가이다.
\kfChoice 왕필과 곽상은 모두 유가의 관점에 입각해 도가의 경전을 비판적으로 해석하였다.
\end{kfChoices}
\end{kfQBlock}

\begin{kfQBlock}{4}{}
\kfQStem{윗글을 바탕으로 <보기>를 이해한 내용으로 적절하지 않은 것은?

<보기> 어떤 교향악단에서 지휘자 없이 연주를 했다. 각 단원은 악보와 자신의 음악적 감각에 따라 연주했을 뿐인데, 결과적으로 놀라울 정도로 조화로운 음악이 만들어졌다.}
\begin{kfChoices}
\kfChoice 곽상의 관점에서 각 단원이 자신의 감각에 따라 연주한 것은 '독화'에 비유할 수 있다.
\kfChoice 곽상의 관점에서 지휘자 없이 이루어진 조화는 '현명'의 개념으로 설명할 수 있다.
\kfChoice 왕필의 관점에서 조화로운 연주가 가능하려면 보이지 않는 근원적 통일 원리가 있어야 한다.
\kfChoice 왕필의 관점에서 지휘자가 없는 것은 군주가 무위의 자세를 취한 것에 해당한다.
\kfChoice 곽상의 관점에서 각 단원의 연주는 지휘자라는 외부 근원에 의해 비로소 가능해진 것이다.
\end{kfChoices}
\end{kfQBlock}

\end{multicols}

\kfSecHeader{kfAreaNonlit}{지문}{불교 유식학 — 팔식과 유식무경 (인문)}

\kfPassageNote{불교 유식학 — 팔식과 유식무경 (인문)}{%
유식학(唯識學)은 인도 대승불교의 한 학파로, 우리가 경험하는 세계가 마음(識) 바깥에 독립적으로 존재하는 것이 아니라 오직 마음의 작용에 의해 나타난 것이라고 주장한다. 이를 '유식무경(唯識無境)', 곧 오직 식(識)만 있을 뿐 외부 대상(境)은 없다는 말로 요약한다. 여기서 '무경'이란 외부 세계가 물리적으로 존재하지 않는다는 뜻이 아니라, 우리가 인식하는 대상이 마음의 구성 작용을 거친 것이므로 마음과 독립된 객관적 대상은 성립하지 않는다는 의미이다.

유식학에서는 인간의 마음을 여덟 가지 식(識)으로 분석하는데, 이를 '팔식(八識)'이라 한다. 먼저 전오식(前五識)은 눈·귀·코·혀·몸에 해당하는 안식(眼識)·이식(耳識)·비식(鼻識)·설식(舌識)·신식(身識)이다. 이들은 각각의 감각 기관을 통해 외부 자극을 받아들이는 역할을 한다. 그러나 전오식은 단순히 자극을 수용할 뿐, 그것이 무엇인지 판단하거나 의미를 부여하지는 못한다.

여섯 번째 식인 의식(意識, 제6식)은 전오식이 수용한 감각 자료를 종합하고 판단하는 역할을 한다. 우리가 일상에서 '생각한다', '판단한다'고 할 때 작동하는 것이 바로 이 의식이다. 의식은 대상을 분별하고 개념화하며, 과거를 회상하고 미래를 상상하는 기능도 수행한다. 전오식과 의식을 합하여 '전육식(前六識)'이라 하는데, 이 여섯 가지 식은 일상적 인식의 영역에 속한다.

일곱 번째 식인 말나식(末那識, 제7식)은 '자아'에 대한 집착의 근거이다. 말나식은 끊임없이 제8식인 아뢰야식을 '자기 자신'이라고 집착한다. 우리가 '나'라는 존재가 고정적으로 있다고 느끼는 것, 자기중심적으로 세계를 경험하는 것은 모두 말나식의 작용 때문이다. 말나식은 깨어 있을 때는 물론 잠들어 있을 때에도 쉬지 않고 작동하며, 수행을 통해 이 집착을 끊는 것이 유식 수행의 핵심 과제 중 하나이다.

여덟 번째 식인 아뢰야식(阿賴耶識, 제8식)은 유식학에서 가장 중요한 개념이다. 아뢰야식은 모든 경험의 흔적, 곧 종자(種子)를 저장하는 근본식이다. 과거의 행위와 경험에서 비롯된 종자들이 아뢰야식에 쌓이고, 이 종자들이 조건이 갖추어지면 현실의 인식으로 발현된다. 이 과정을 '식전변(識轉變)', 곧 식이 변하여 대상 세계를 드러낸다고 표현한다. 아뢰야식은 개인의 모든 경험을 담고 있으면서도 그 자체는 선악의 판단 없이 중립적으로 작동한다. 유식학에서는 이러한 팔식의 구조를 이해하고, 마음 바깥에 독립된 대상이 없음을 체득하는 것이 곧 깨달음의 길이라고 본다.
}
\kfSecHeader{kfAreaNonlit}{문제}{}

\begin{multicols}{2}\raggedcolumns\setlength{\columnsep}{12pt}
\begin{kfQBlock}{1}{}
\kfQStem{윗글의 내용과 일치하지 않는 것은?}
\begin{kfChoices}
\kfChoice 전오식은 감각 자극을 수용하지만 대상의 의미를 판단하지는 못한다.
\kfChoice 의식(제6식)은 감각 자료를 종합·판단하고 과거를 회상하는 기능을 수행한다.
\kfChoice 말나식은 깨어 있을 때만 작동하며 잠들면 활동을 중단한다.
\kfChoice 아뢰야식은 종자를 저장하되 그 자체는 선악 판단 없이 중립적으로 작동한다.
\kfChoice 유식무경에서 '무경'은 외부 대상이 마음과 독립하여 성립하지 않음을 뜻한다.
\end{kfChoices}
\end{kfQBlock}

\begin{kfQBlock}{2}{}
\kfQStem{윗글에서 팔식(八識)의 기능을 정리한 것으로 적절하지 않은 것은?}
\begin{kfChoices}
\kfChoice 전오식: 감각 기관을 통해 외부 자극을 수용한다.
\kfChoice 의식(제6식): 전오식의 감각 자료를 종합하고 대상을 분별·개념화한다.
\kfChoice 말나식(제7식): 아뢰야식을 자기 자신으로 집착하여 자아 의식의 근거가 된다.
\kfChoice 아뢰야식(제8식): 종자를 저장하고, 조건이 갖추어지면 현실의 인식으로 발현시킨다.
\kfChoice 의식(제6식): 모든 경험의 종자를 저장하는 근본식으로 선악 판단 없이 작동한다.
\end{kfChoices}
\end{kfQBlock}

\begin{kfQBlock}{3}{}
\kfQStem{윗글을 바탕으로 <보기>를 이해한 내용으로 적절하지 않은 것은?

<보기> 철수는 시장에서 과일을 보았다. 눈으로 빨간색을 인식하고, 코로 달콤한 향을 맡았다. 그는 '아, 이것은 사과구나'라고 판단한 뒤, 예전에 먹었던 사과의 맛을 떠올리며 구매를 결정했다.}
\begin{kfChoices}
\kfChoice 빨간색을 인식한 것은 전오식 중 안식(眼識)의 작용에 해당한다.
\kfChoice 달콤한 향을 맡은 것은 전오식 중 비식(鼻識)의 작용에 해당한다.
\kfChoice '이것은 사과구나'라고 판단한 것은 의식(제6식)의 분별 작용에 해당한다.
\kfChoice 예전 사과의 맛을 떠올린 것은 아뢰야식에 저장된 종자가 발현된 것으로 볼 수 있다.
\kfChoice 구매를 결정한 것은 말나식이 아뢰야식을 자기 자신으로 집착한 결과이다.
\end{kfChoices}
\end{kfQBlock}

\begin{kfQBlock}{4}{}
\kfQStem{'식전변(識轉變)'의 의미를 가장 정확하게 설명한 것은?}
\begin{kfChoices}
\kfChoice 외부 대상이 마음속으로 들어와 식(識)을 변화시키는 과정
\kfChoice 식(識)이 변하여 대상 세계를 드러내는 과정
\kfChoice 전오식이 의식(제6식)으로 전환되는 과정
\kfChoice 말나식이 아뢰야식의 저장 기능을 대신하는 과정
\kfChoice 아뢰야식에 저장된 종자가 영구히 소멸하는 과정
\end{kfChoices}
\end{kfQBlock}

\end{multicols}

\kfSecHeader{kfAreaNonlit}{지문}{퇴계 이황과 율곡 이이 — 사단칠정 논쟁 (인문)}

\kfPassageNote{퇴계 이황과 율곡 이이 — 사단칠정 논쟁 (인문)}{%
조선 성리학의 가장 중요한 철학적 논쟁은 사단칠정(四端七情) 논쟁이다. 이 논쟁은 퇴계 이황(李滉)과 고봉 기대승(奇大升) 사이에서 시작되었고, 이후 율곡 이이(李珥)가 자신의 체계를 제시하면서 조선 성리학의 두 큰 흐름이 형성되었다. 사단(四端)은 맹자가 말한 인간 본성의 네 가지 단서, 곧 측은지심(惻隱之心)·수오지심(羞惡之心)·사양지심(辭讓之心)·시비지심(是非之心)이며, 칠정(七情)은 『예기(禮記)』에 나오는 일곱 가지 감정, 곧 희(喜)·노(怒)·애(哀)·구(懼)·애(愛)·오(惡)·욕(欲)이다.

이황은 사단과 칠정의 발생 근거가 다르다고 보았다. 그에 따르면, 사단은 이(理)가 발하고 기(氣)가 따르는 것(理發氣隨)이며, 칠정은 기(氣)가 발하고 이(理)가 타는 것(氣發理乘)이다. 이를 '이기호발설(理氣互發說)'이라 한다. 이황이 이렇게 주장한 근거는 다음과 같다. 사단은 순선(純善)한 도덕 감정이므로 선한 원리인 이(理)에서 비롯되어야 하며, 칠정은 선할 수도 악할 수도 있는 감정이므로 기(氣)에서 비롯된다고 보았다. 이황에게 이(理)는 단순히 수동적인 법칙이 아니라 능동적으로 발할 수 있는 것이었다.

이이는 이황의 이기호발설에 반대하여 '기발이승일도설(氣發理乘一途說)'을 제시했다. 이이에 따르면, 발(發)하는 것은 오직 기(氣)뿐이며, 이(理)는 기 위에 타고 있을 뿐 스스로 발하지 않는다. 사단이든 칠정이든 모두 기가 발하고 이가 거기에 타는 구조라는 것이다. 이이는 이를 '일도(一途)', 곧 하나의 길로 표현했다. 이(理)가 스스로 발한다는 것은 이에 운동성을 부여하는 것인데, 이는 형이상의 원리로서 활동하는 것이 아니라 기를 통해서만 드러나므로, 이가 독자적으로 발한다고 말할 수 없다는 것이 이이의 논거였다.

그렇다면 이이의 체계에서 사단과 칠정의 차이는 어떻게 설명되는가? 이이는 사단을 칠정의 한 부분으로 보았다. 칠정 가운데 이(理)의 본래 모습이 온전히 드러난 것이 사단이며, 기질에 의해 가려진 것이 악한 감정이다. 비유하자면, 칠정이라는 큰 범주 안에 사단이 포함되는 구조이다. 이이는 이를 '사단은 칠정의 선한 한쪽(善一邊)'이라고 표현했다.

이황과 이이의 논쟁은 단순한 용어 싸움이 아니라, 이(理)의 성격을 어떻게 규정할 것인가 하는 근본 문제와 관련된다. 이황이 이에 능동성을 부여한 것은 도덕 실천의 근거를 확고히 하려는 의도에서였다. 이가 스스로 발할 수 있다면, 인간은 누구나 도덕적 주체로서 선을 실현할 수 있는 내적 근거를 갖게 되기 때문이다. 반면 이이는 이의 능동성을 인정하면 이와 기의 존재론적 구분이 무너진다고 보았다. 이가 기처럼 운동한다면 형이상과 형이하의 경계가 흐려지기 때문이다. 이처럼 두 학자의 논쟁은 형이상학적 원리와 도덕 실천의 관계라는 심오한 문제를 둘러싸고 전개된 것이었다.
}
\kfSecHeader{kfAreaNonlit}{문제}{}

\begin{multicols}{2}\raggedcolumns\setlength{\columnsep}{12pt}
\begin{kfQBlock}{1}{}
\kfQStem{윗글의 내용과 일치하는 것은?}
\begin{kfChoices}
\kfChoice 이이는 사단과 칠정이 모두 이(理)가 발하는 것이라고 주장했다.
\kfChoice 이황은 사단을 칠정의 선한 한쪽으로 보아 칠정 속에 포함시켰다.
\kfChoice 이이는 이(理)가 스스로 발하지 않고 기(氣)를 통해서만 드러난다고 보았다.
\kfChoice 사단칠정 논쟁은 이황과 율곡 이이 사이에서 직접 시작되었다.
\kfChoice 이황은 이(理)가 수동적 법칙에 불과하여 능동적으로 발할 수 없다고 보았다.
\end{kfChoices}
\end{kfQBlock}

\begin{kfQBlock}{2}{}
\kfQStem{이황이 이(理)에 능동성을 부여한 의도로 가장 적절한 것은?}
\begin{kfChoices}
\kfChoice 이와 기의 존재론적 구분을 해소하기 위해서이다.
\kfChoice 칠정이 사단보다 도덕적으로 우월함을 증명하기 위해서이다.
\kfChoice 형이상과 형이하의 경계를 무너뜨리기 위해서이다.
\kfChoice 인간의 도덕 실천 근거를 확고히 하기 위해서이다.
\kfChoice 기(氣)의 운동성을 부정하기 위해서이다.
\end{kfChoices}
\end{kfQBlock}

\begin{kfQBlock}{3}{}
\kfQStem{윗글의 이황과 이이의 견해를 비교한 것으로 적절하지 않은 것은?}
\begin{kfChoices}
\kfChoice 이황은 사단과 칠정의 발생 근거가 다르다고 보았고, 이이는 같다고 보았다.
\kfChoice 이황은 이가 스스로 발할 수 있다고 보았고, 이이는 이가 스스로 발하지 않는다고 보았다.
\kfChoice 이황은 사단을 이가 발한 결과로 보았고, 이이는 사단을 칠정의 선한 한쪽으로 보았다.
\kfChoice 이황과 이이 모두 사단과 칠정이 별개의 독립적 범주라고 보았다.
\kfChoice 이황은 이의 능동성을 통해 도덕 실천의 근거를 마련하고자 했고, 이이는 이의 능동성이 이기 구분을 무너뜨린다고 우려했다.
\end{kfChoices}
\end{kfQBlock}

\begin{kfQBlock}{4}{}
\kfQStem{윗글을 바탕으로 <보기>를 이해한 내용으로 적절하지 않은 것은?

<보기> 갑(甲)은 길에서 다친 아이를 보고 측은한 마음이 일어나 달려가 도와주었다. 을(乙)은 불의를 보고 분노하여 항의했으나, 분노가 지나쳐 폭력을 행사했다.}
\begin{kfChoices}
\kfChoice 이황의 관점에서 갑의 측은한 마음은 이(理)가 발한 사단에 해당한다.
\kfChoice 이이의 관점에서 갑의 측은한 마음은 기(氣)가 발하고 이(理)가 탄 결과이다.
\kfChoice 이황의 관점에서 을의 분노는 기(氣)가 발한 칠정에 해당한다.
\kfChoice 이이의 관점에서 을의 폭력은 기질에 의해 이의 본래 모습이 가려진 경우이다.
\kfChoice 이이의 관점에서 갑의 측은한 마음과 을의 분노는 발생 근거가 근본적으로 다르다.
\end{kfChoices}
\end{kfQBlock}

\end{multicols}

\kfStartWeekly{패턴 워크북 — 총 ?문항}


\kfSecHeader{kfAreaWeekly}{지문}{문항 1 지문 / 섞어치기}

\kfPassageNote{문항 1 지문 / 섞어치기}{%
컴퓨터는 사람의 언어를 직접 이해할 수 없다. 따라서 컴퓨터가 사람의 언어를 이해할 수 있도록 표현해 주어야 한다. 컴퓨터가 사람의 언어를 이해할 수 있도록 표현하는 방법에는 여러 가지가 있을 수 있는데, 그중 ㉠원-핫 인코딩은 표현하고 싶은 단어의 색인 값에 1을 부여하고 나머지에는 0을 부여하는 방법이다. 원-핫 인코딩을 수행하기 위해서는 먼저 대상이 되는 텍스트의 단어 집합을 만든다. 예를 들어 ‘오늘 날씨 정말 좋다.’라는 한 문장으로 이루어진 텍스트를 대상으로 한다고 가정하면, 단어 집합은 \{오늘, 날씨, 정말, 좋다\}가 되고 집합의 원소가 되는 단어의 개수가 4개이므로 단어 집합의 크기는 4가 된다. 단어 집합의 크기에 따라 각각의 단어에 [1, 0, 0, 0], [0, 1, 0, 0], [0, 0, 1, 0], [0, 0, 0, 1]과 같이 4차원의 고유한 숫자를 부여하는 것이 원-핫 인코딩이다. 원-핫 인코딩은 직관적이고 단순하지만 단어들이 개별적, 독립적으로 표현되기

(중략)
}
\begin{kfQBlock}{1}{}
\kfQStem{㉮의 목적으로 가장 적절한 것은?}
\begin{kfChoices}
\kfChoice 단어들 간의 유사성이나 관계성을 표현하기 위해서
\kfChoice 텍스트를 대상으로 단어 집합을 만들기 위해서
\kfChoice 단어들을 직관적이고 단순하게 표현하기 위해서
\end{kfChoices}
\end{kfQBlock}

\kfSecHeader{kfAreaWeekly}{지문}{문항 2 지문 / 부정상대어}

\kfPassageNote{문항 2 지문 / 부정상대어}{%
최근까지 일반적인 컴퓨터 환경에서는 운영 체제를 제어하거나 사용자 인터페이스를 처리하고, 문서를 작성하거나 웹을 탐색하는 등 다양한 범용 작업이 주요 연산 대상이었다. 이러한 작업은 명령어의 종류가 많고 그 순서가 수시로 바뀌며, 연산 간 종속성이 크다는 특징이 있다. 컴퓨터의 연산은 주로 여러 개의 코어로 구성된 ㉠중앙 처리 장치(CPU)에서 수행된다. 코어는 컴퓨터의 두뇌 역할을 하는 것으로 명령어를 해석하고 각종 연산을 실행한다. CPU는 여러 개의 고성능 코어로 구성되어 복잡한 제어 흐름을 정밀하게 조율할 수 있도록 설계되었다. 그러나 인공 지능 학습, 고해상도 영상 처리, 과학 기술 관련 계산 등 대규모 데이터를 대상으로 동일한 연산을 반복 수행해야 하는 작업이 증가하면서 기존의 연산 환경에 변화가 생겼다. CPU는 코어 수에 한계가 있고 개별 코어가 복잡한 제어 흐름에 맞추어 다양한 연산을 처리하도록 설계되어 있기 때문에 동일한 연산을 대량으로 반복 수행해야 하는 현대의 

(중략)
}
\begin{kfQBlock}{2}{}
\kfQStem{<보기>의 시스템 설계를 윗글과 연결해 이해한 내용으로 적절하지 않은 것은?}
\begin{kfEvidence}<보기> 자율주행 시스템은 센서 융합 전처리와 제어 판단, 대규모 행렬 기반 객체 인식, 예외 상황 판단을 동시에 처리한다. 시스템은 CPU·GPU·NPU에 작업을 분담하고, 상황에 따라 장치 간 연산 흐름을 조정한다.\end{kfEvidence}
\begin{kfChoices}
\kfChoice 분기 상황이 많아질수록 GPU의 모든 스레드는 서로 다른 경로를 동시에 병목 없이 처리하므로 효율이 항상 오른다.
\kfChoice 대규모 행렬 기반 반복 계산 구간은 GPU 또는 NPU 활용 우선순위를 검토할 수 있다.
\kfChoice 예외 상황 판단처럼 분기 빈도가 높은 구간은 CPU 분담이 유리할 수 있다.
\end{kfChoices}
\end{kfQBlock}

\kfSecHeader{kfAreaWeekly}{지문}{문항 3 지문 / 주체객체}

\kfPassageNote{문항 3 지문 / 주체객체}{%
(가) 인과 관계에 관한 ㉮조건설에 의하면 어떤 범죄적 행위는 그것이 특정한 범죄 결과를 필연적으로 야기하는 전체 조건에 포함된 하나의 조건일 때, 그 결과에 대한 원인이 된다. 조건설은 이와 같이 조건의 전체적 복합을 이루는 개개의 조건들을 모두 동등하게 결과 발생의 원인으로 취급하므로 등가설이라고도 불리는데, 이러한 조건들 중에 어느 하나라도 결여된다면 문제의 결과는 발생하지 않게 된다.

조건설에서는 이른바 조건 관계의 발견을 위해 선행 사실이 없었다고 가정할 경우 후행 사실이 발생했을 것인가를 묻는 조건 공식을 사용하고 있다. 당해 행위를 제거하고 생각해 봤을 때 문제의 결과가 발생하지 않을 것으로 판단된다면, 행위와 결과 사이에는 조건 관계가 성립한다고 보는 것이다. 이러한 ⓐ반사실적 사고 실험의 절차는 일종의 가설적 소거법을 따르는 것이라 할 수 있다.

(중략)
}
\begin{kfQBlock}{3}{}
\kfQStem{(가)와 (나)를 바탕으로 <보기>를 이해한 내용으로 가장 적절한 것은?}
\begin{kfEvidence}ㄱ. X의 물건 위에 그것이 더 이상 버티지 못하고 비가역적인 손상을 입기 시작하게 만드는 무게의 95\%에 해당하는 하중을 A가 가하는 순간, 그러한 무게의 나머지 5\%에 해당하는 하중을 B가 가했다. 이후 X의 물건은 곧 손상되었다. ㄴ. A는 X가 평소 무게 50kg의 추를 여러 개 보관하던 창고에 그것과 모양이 비슷하지만 무게가 90kg에 달하는 추 하나를 놓아두었는데, 마침 추를 사용하려던 X는 60kg까지 버틸 수 있는 기계로 A가 놓아둔 추를 들어 올렸다. 당시 X가 고를 수 있었을 나머지 모든 추들은 B가 A와 마찬가지로 모양이 비슷한 무게 90kg의 추들로 바꿔 둔 상태였다. X가 추를 들어 올리자 곧 기계가 추의 무게를 견디지 못하고 고장을 일으켰다. ㄷ. A는 X가 평소 무게 50kg의 추를 보관하던 창고에 그것과 모양이 비슷하지만 무게가 90kg에 달하는 추를 놓아두었다. 마침 추를 사용하려던 X가 60kg까지 버틸 수 있는 기계로 A가 놓아둔 추를 집으려는 순간, 전날 B가 창고 지붕을 수리하며 허술하게 쌓아 두었던 자재가 기계 위로 쏟아지면서 기계가 고장을 일으켰다. (모든 경우에 A와 B는 서로 알지 못하며, 협력하거나 공모한 적도 없다.)\end{kfEvidence}
\begin{kfChoices}
\kfChoice ㄱ: 조건설에 따르면 A의 행위를 X의 물건이 손상되게 한 원인으로 보지 않을 것이다.
\kfChoice ㄴ: 조건설에 따르면 A의 행위를 X의 기계가 고장 나게 한 원인으로 보지 않을 것이다.
\end{kfChoices}
\end{kfQBlock}

\kfSecHeader{kfAreaWeekly}{지문}{문항 4 지문 / 순서방향}

\kfPassageNote{문항 4 지문 / 순서방향}{%
소비란 가계가 재화와 서비스를 구입해 사용하는 행위이며, 얼마나 소비할 것인지를 결정하는 요인 중에서 가장 중요한 것은 소득의 크기이다. 특히 개인이 마음대로 쓸 수 있는 가처분 소득의 크기와 밀접한 관계를 맺고 있다. 소비와 가처분 소득 간에는 양(+)의 상관관계가 있어 가처분 소득이 증가하면 소비도 증가한다. 그리고 가처분 소득이 늘어나면서 커지는 소비의 증가분은 가처분 소득의 증가분보다 작다. 소비자들은 증가된 소득을 전부 소비하는 것이 아니라 그중 일부를 저축하기 때문이다.

(중략)

어떤 국민 경제의 가처분 소득이 Yd, 그리고 소비가 C라고 할 때, ㉠둘 사이의 관계는 ‘C=a+bYd (단, a>0, 0와 같이 나타낼 수 있다. 이때 ａ는 생존을 위해 필요한 최소한의 소비이고, ｂ는 한계 소비 성향이다. 이 식처럼 소비와 가처분 소득 사이의 관계를 단순한 1차 함수로 표현한 것을 케인스의 소비 함수라고 한다. 이 함수에 나타난 소비와 가처분 소득 사이의 관계에서 평균 소비 성향과 한계

(중략)
}
\begin{kfQBlock}{4}{}
\kfQStem{[01] 윗글에 대한 이해로 적절하지 않은 것은?}
\begin{kfChoices}
\kfChoice 평균 소비 성향과 한계 소비 성향이라는 개념을 활용하여 소비 함수를 도출해 볼 수 있다.
\kfChoice 소비자들은 가처분 소득으로 소비도 하고 저축도 한다.
\kfChoice 소비와 가처분 소득 간의 관계에 따를 때 가처분 소득이 감소하면 소비도 감소하게 된다.
\end{kfChoices}
\end{kfQBlock}

\kfSecHeader{kfAreaWeekly}{지문}{문항 5 지문 / 인과도치}

\kfPassageNote{문항 5 지문 / 인과도치}{%
현대 경제는 과거와는 다른 양상으로 움직인다. 기술은 시시각각 변하고 세계 금융 시장은 얽히고설켜 예측하기 어려운 변동성을 보인다. 이러한 복잡한 경제 상황을 이해하기 위해 미국의 산타페 연구소를 중심으로 등장한 복잡계 경제학은 경제 주체가 언제나 가장 합리적 선택을 한다는 명제는 진리가 아니라고 주장한다. 그러면서 경제에 대한 근본적인 생각의 전환을 요구한다.

(중략)

복잡계 경제학에서는 현대 경제를 수학적 계산이나 도식적 모델의 결과물이 아니라, 동적으로 변화하는 복잡한 시스템이라고 본다. 또한 현대 경제의 특징은 다양한 주체들이 제한된 정보 속에서 상호 작용하며 예측 불가능한 질서를 만들어 내는 비정형성이라고 설명한다. 어떤 기술이 표준이 될지, 어떤 기업이 시장을 주도할지, 어떤 정책이 장기적으로 성공할지는 누구도 확신할 수 없다는 것이다. 예를 들어 기술의 선택에서는 효율성보다는 경로 의존성이 더 크게 작용하기 때문에 새로운 기술 개발의 효과를 예측하는 것은 매우 어렵다. 여기서 경로

(중략)
}
\begin{kfQBlock}{5}{}
\kfQStem{[02] 윗글을 바탕으로 <보기>를 이해한 내용으로 적절하지 않은 것은?}
\begin{kfEvidence}<보기> 쿼티(QWERTY) 자판 배열은 타자기의 기계적 충돌을 방지하기 위해 고안된 것이다. 오늘날 컴퓨터 키보드는 타자기에서와 같은 기계적 충돌을 염려할 필요가 없기 때문에 입력 효율 면에서 쿼티보다 더 우수하다고 평가되는 여러 새로운 자판 배열이 개발되었다. 그러나 이미 쿼티 자판이 전 세계적으로 표준화되어 널리 보급된 상태이기 때문에, 새로운 배열 자판의 사용률은 매우 저조하다.\end{kfEvidence}
\begin{kfChoices}
\kfChoice 쿼티 자판 배열보다 입력 효율이 높은 자판의 사용률이 저조한 것은 기술 선택에 경로 의존성이 작용한 것이라고 볼 수 있겠군.
\kfChoice 쿼티보다 입력 효율이 우수한 자판 배열이 개발되었다는 것은 경제 주체의 합리적 선택이 경제의 비정형성을 해소한다는 것을 말해 준다고 볼 수 있겠군.
\end{kfChoices}
\end{kfQBlock}

\kfSecHeader{kfAreaWeekly}{지문}{문항 6 지문 / 의도/목적}

\kfPassageNote{문항 6 지문 / 의도/목적}{%
(가) 타인에 대한 힘과 영향력의 근원으로서 사회 활동에서 다양한 이윤을 확보할 수 있는 자원이 되는 ‘자본’은 여러 형태로 존재한다. \uline{㉠부르디외}는 세계화를 기득권 세력이 지배력을 유지 및 확장하기 위해 수행하는 전략이라고 보고, 세계화로 인해 자본을 바탕으로 한 세력 간의 위계가 강화됨으로써 불평등 구조가 심화된다고 보았다. 그는 자본을 돈 또는 돈으로 즉각 전환 가능한 경제 자본, 학연·지연·혈연 등과 같은 사회적 관계망인 사회 자본, 개인이 보유한 문화적 지식, 기술, 교육 및 문화적 자산인 문화 자본, 그리고 상징 자본으로 나누었다. 상징 자본은 소유한 자본이 세간으로부터 인정되었을 때 발생하는 자본으로서, 명예, 권위 등과 같이 보이지 않는 상징적 가치가 자본으로서의 힘을 발휘하는 것이다. 부르디외는 상징 자본이 국가 간의 관계에도 작용하며, 그 결과 불평등한 국제 질서가 형성될 수 있다고 보았다.

부르디외는 상징 폭력과 아비투스를 바탕으로 자본과 세계화에 관…(중략)
}
\begin{kfQBlock}{6}{}
\kfQStem{㉠과 ㉡에 대한 설명으로 적절하지 않은 것은?}
\begin{kfChoices}
\kfChoice ㉡은 국가가 물질적 풍요를 바탕으로 국민의 후생을 증대하는 것을 우선적인 목표로 삼아야 한다고 보았다.
\kfChoice ㉡은 국가가 물질적 풍요를 바탕으로 국민의 후생을 증대하는 것을 우선적인 목표로 삼아야 한다고 보았다.
\kfChoice ㉠은 시민 교육을 통해 사람들이 세계화의 숨은 본질을 알게 해야 한다고 보았다.
\kfChoice ㉡은 ㉠과 달리 세계화를 통해 부유하지 않은 국가들이 부유해질 수 있는 기회를 얻을 수 있다고 보았다.
\end{kfChoices}
\end{kfQBlock}

\kfSecHeader{kfAreaWeekly}{지문}{문항 7 지문 / 상관관계 반전}

\kfPassageNote{문항 7 지문 / 상관관계 반전}{%
(가) \uline{건축 설계}는 공간을 계획하거나 건물의 외형을 아름답게 만드는 것만을 의미하지 않는다. \uline{건축 설계}는 \uline{구조적 안정성}, \uline{기능적 효율성}, \uline{미적 가치}를 위한 과학이자 예술이다. \uline{건축 설계}가 과학이자 예술로 기능하는 데 있어 가장 중요한 두 가지는 \uline{건축 구조}와 \uline{건축 재료}이다.

\uline{건축 구조}는 건물에 가해지는 \uline{하중}을 \uline{㉠지탱}하며 형체를 유지하도록 설계된다. \uline{하중}이란 구조물에 작용하는 모든 힘으로, \uline{정하중}과 \uline{활하중}으로 구분된다. \uline{정하중}은 건축물 자체의 무게, 즉 기둥, 보, 슬래브, 벽체 등 영구적으로 작용하는 무게를 뜻한다. \uline{활하중}은 사람이 건축물 안에서 이동하거나 가구와 같은 \uline{가변적}인 물체가 추가되면서 발생하는 하중을 의미한다. 건축 구조는 이 하중을 지탱할 수 있도록 설계되어야 …(중략)
}
\begin{kfQBlock}{7}{}
\kfQStem{(가)를 바탕으로 추론한 내용으로 가장 적절하지 않은 것은?}
\begin{kfChoices}
\kfChoice 철근 콘크리트는 콘크리트의 인장 약점을 보완하므로, 동일 조건에서 콘크리트 단독보다 균열 위험을 낮출 수 있다.
\kfChoice 철근 콘크리트는 콘크리트의 인장 약점을 보완하므로, 동일 조건에서 콘크리트 단독보다 균열 위험을 낮출 수 있다.
\kfChoice 응력은 단위 면적당 힘이므로, 같은 힘에서 단면적이 커지면 응력도 커진다.
\kfChoice 좌굴은 주로 인장력에 의해 부재가 늘어나며 생기는 현상이다.
\end{kfChoices}
\end{kfQBlock}

\kfSecHeader{kfAreaWeekly}{지문}{문항 8 지문 / 필요조건}

\kfPassageNote{문항 8 지문 / 필요조건}{%
(가) 조세 법률주의는 세금의 부과와 징수가 반드시 법률에 근거해 이루어져야 한다는 원칙이다. 법률에 명시되지 않은 세금을 부과하거나, 법률의 규정을 넘어 세율과 과세 기준을 자의적으로 바꾸는 것은 원칙적으로 허용되지 않는다. 이 원칙은 정부의 과세 권한을 제한해 납세자의 재산권을 보호하고 과세의 공정성과 예측 가능성을 높인다. 조세 법률주의는 실무에서 조세 열거주의와 조세 포괄주의라는 두 해석 방식으로 구체화된다. 조세 열거주의는 법조문에 적힌 사항만 엄격히 적용해 과세하는 입장으로, 법에 없는 새로운 거래 유형에는 과세하지 못한다. 이는 납세자 권리 보호와 조세권 남용 방지에 강점이 있다. 반면 조세 포괄주의는 법 문언뿐 아니라 입법 취지와 과세 목적까지 종합해, 법에 직접 적시되지 않은 유사 거래에도 과세 가능성을 검토한다. 이 방식은 과세 회피를 줄이고 경제적 실질이 유사한 행위 사이의 부담 형평을 맞추는 데 유리하다. 현대 조세 실무는 두 입장의 균형을 요구한다. 문언만 고…(중략)
}
\begin{kfQBlock}{8}{}
\kfQStem{(가)의 논지를 바탕으로 추론한 내용으로 적절하지 않은 것은?}
\begin{kfChoices}
\kfChoice 포괄주의는 법률 근거 없이도 과세 목적만 타당하면 과세를 정당화한다.
\kfChoice 포괄주의는 법률 근거 없이도 과세 목적만 타당하면 과세를 정당화한다.
\kfChoice 문언만 고수하면 신종 거래에서 과세 공백이 생길 수 있다.
\kfChoice 포괄주의를 무제한 허용하면 납세자 권리 침해 우려가 커질 수 있다.
\end{kfChoices}
\end{kfQBlock}

\kfSecHeader{kfAreaWeekly}{지문}{문항 9 지문 / 결합/분리}

\kfPassageNote{문항 9 지문 / 결합/분리}{%
우리는 일반적으로 마음속 상태를 겉으로 분명하게 드러내거나 나타낼 때 ‘\uline{표현}’이라는 말을 사용한다. 그러나 표현이라는 말이 단지 자신의 심리적 상태를 드러내는 데에만 쓰이는 것은 아니다. 때로는 자연을 대하는 우리의 태도나 반응에도 표현이라는 말이 적용된다. 예를 들어 축 늘어진 나뭇가지를 보고 “애처롭다.”라고 말하는 경우가 있다. 이는 나뭇가지 자체가 감정을 지닌 것은 아니지만, 우리가 그것을 바라보며 내면의 정서를 \uline{투영}하고 외부로 드러내는 행위를 통해 표현이라는 말을 사용하게 되는 것이다. 하지만 표현이라는 말이 가장 널리 사용되는 분야는 다름 아닌 예술이며, 예술과 표현의 관계에 대해서는 오래전부터 다양한 견해가 논의되어 왔다.

첫째로, \uline{㉠표현}을 예술가가 창작 과정 속에서 행하는 내적 활동으로 보는 견해가 있다. 이 입장에서 표현은 예술가의 마음속에서 일어나는 정신적 과정, 즉 내적인 작용으로 이해된다. 이러한 입장을 대표하는 사상가가 …(중략)
}
\begin{kfQBlock}{9}{}
\kfQStem{㉮의 이유로 가장 적절한 것은?}
\begin{kfChoices}
\kfChoice 일상적 경험과 예술 모두 내적 이미지가 자발적으로 생겨나는 과정을 통해 이루어지므로
\kfChoice 일상적 경험과 예술 모두 내적 이미지가 자발적으로 생겨나는 과정을 통해 이루어지므로
\kfChoice 예술은 일상의 정신적 삶과 유리된 특수한 기능의 산물이므로
\kfChoice 기존의 미학이나 예술학이 일상과 예술 간의 연속성을 중시했으므로
\end{kfChoices}
\end{kfQBlock}

\kfSecHeader{kfAreaWeekly}{지문}{문항 10 지문 / 조건 왜곡}

\kfPassageNote{문항 10 지문 / 조건 왜곡}{%
일반적으로 혼인은 남자와 여자가 부부가 되는 일을 말한다. 과거에는 특정한 법적 절차 없이도 결혼식과 같은 사회적 관습에 의한 의식을 치르면 혼인이 성립한 것으로 보았다. 하지만 현재 대부분의 국가는 법에서 규정한 절차를 따라야만 혼인이 성립되는 법률혼 제도를 두고 있으며, 사회적 관습보다는 혼인 신고와 같은 법적 절차가 더 중요하다. 혼인은 사회를 구성하는 기본 단위인 가족을 형성하는 일이기 때문에 국가의 근간을 이루는 중요한 과정이고 이에 따라 법을 통해 관계를 보호하는 것이다. 따라서 현대 사회에서 혼인은 법적 절차에 의해 이루어지게 되고 여기에 참여한 남녀 간에는 법적 효력이 발생하게 된다.

혼인이 성립하기 위해서는 법에서 규정하고 있는 실질적 요건과 형식적 요건을 모두 만족하여야 한다. 실질적 요건에서 가장 중요한 것은 혼인 의사의 합치이다. 혼인 의사는 부부 관계를 성립하겠다는 의사로 반드시 혼인 당사자 간의 합의가 필요하다. 또 다른 실질적 요건은 연령과 혼인 자격에 …(중략)
}
\begin{kfQBlock}{10}{}
\kfQStem{㉮의 이유로 가장 적절한 것은?}
\begin{kfChoices}
\kfChoice 협의상 이혼은 이혼 당사자 간의 자율적인 합의에 의한 것이기 때문에
\kfChoice 협의상 이혼은 이혼 당사자 간의 자율적인 합의에 의한 것이기 때문에
\kfChoice 이혼 사유의 타당성 여부는 법원이 판단하는 것이 아니기 때문에
\kfChoice 이혼 당사자 간의 합의만으로도 법적 효력이 발생하게 되기 때문에
\end{kfChoices}
\end{kfQBlock}

\kfSecHeader{kfAreaWeekly}{지문}{문항 11 지문 / 긍부정 반전}

\kfPassageNote{문항 11 지문 / 긍부정 반전}{%
\#\#\# [31 \textasciitilde{} 34] 다음 글을 읽고 물음에 답하시오.

(가)   지팡이 짚고 바람 쐬며 좌우를 돌아보니   누대의 맑은 경치 아마도 깨끗하구나.   ㉠ 물도 하늘 같고 하늘도 물 같으니   푸른 물과 긴 하늘이 한빛이 되었거든   물가에 갈매기는 오는 듯 가는 듯 그칠 줄을 모르네.   ㉡ 바위 위 산꽃은 수놓은 병풍 되었고   시냇가 버들은 초록 장막 되었는데,   좋은 날 좋은 경치 나 혼자 거느리고   ㉢ 꽃피는 시절 허송하지 말리라 하고   아이 불러 하는 말, 이 깊은 산속에서 해산물을 볼쏘냐.   ㉣ 살진 고사리, 향기로운 당귀를 돼지고기, 사슴고기 섞어서   크나큰 바구니에 흡족히 담아두고   붕어회에다 눌어, 꿩 섞어 먹음직하게 구워지거든   술동이의 맑은 술을 술잔에 가득 부어   한잔, 또 한잔 취토록 먹은 후에,   ㉤ 복숭아꽃 붉은 비 되어 취한 낯에 뿌리는데   낚시터 넓은 돌을 높이 베고 누우니   무회씨 때 사람인가, 갈천씨 때 백성*인가.*   *…(중략)
}
\begin{kfQBlock}{11}{}
\kfQStem{(나)의 ‘빛나는 생명의 예술가’가 갖추어야 할 태도로 가장 적절한 것은?}
\begin{kfChoices}
\kfChoice 직관을 통해 자연에 대한 솔직한 감각을 드러낼 수 있어야 한다.
\kfChoice 직관을 통해 자연에 대한 솔직한 감각을 드러낼 수 있어야 한다.
\kfChoice 자연의 모든 것을 알아낼 수 있다는 확신으로 탐구에 임해야 한다.
\kfChoice 여러 기록을 참고하며 자연의 새로운 경지를 소개할 수 있어야 한다.
\end{kfChoices}
\end{kfQBlock}

\kfSecHeader{kfAreaWeekly}{지문}{문항 12 지문 / 비교대상 뒤바꿈}

\kfPassageNote{문항 12 지문 / 비교대상 뒤바꿈}{%
\#\#\# [24 \textasciitilde{} 26] 다음 글을 읽고 물음에 답하시오.

(가)   팽이가 돈다   어린아해이고 어른이고 살아가는 것이 신기로워   물끄러미 보고 있기를 좋아하는 나의 너무 큰 눈 앞에서   아이가 팽이를 돌린다 / 살림을 사는 아해들도 아름다웁듯이   노는 아해도 아름다워 보인다고 생각하면서   손님으로 온 나는 ㉠ 이 집 주인과의 이야기도 잊어버리고   또 한번 팽이를 돌려주었으면 하고 원하는 것이다   도회 안에서 쫓겨다니는 듯이 사는   나의 일이며 / 어느 소설보다도 신기로운 나의 생활이며   모두 다 내던지고   점잖이 앉은 나의 나이와 나이가 준 나의 무게를 생각하면서   ㉡ 정말 속임 없는 눈으로 / 지금 팽이가 도는 것을 본다   그러면 팽이가 까맣게 변하여 서서 있는 것이다   누구 집을 가보아도 나 사는 곳보다는 여유가 있고   바쁘지도 않으니 / 마치 별세계(別世界)같이 보인다   팽이가 돈다 / 팽이가 돈다   팽이 밑바닥에 끈을 돌려 매이니 이상하고   …(중략)
}
\begin{kfQBlock}{12}{}
\kfQStem{(가)와 (나)에 대한 설명으로 가장 적절한 것은?}
\begin{kfChoices}
\kfChoice (가)는 동일한 시구를 반복하여 주제 의식을, (나)는 공감각적 이미지를 사용하여 대상의 특성을 드러내고 있다.
\kfChoice (가)는 동일한 시구를 반복하여 주제 의식을, (나)는 공감각적 이미지를 사용하여 대상의 특성을 드러내고 있다.
\kfChoice (가)는 현재형 표현으로 현장감을, (나)는 음성 상징어를 사용하여 대상의 생동감을 나타내고 있다.
\kfChoice (가)는 공간의 이동에 따라, (나)는 계절의 변화에 따라 시상을 전개하여 대상이 변화하는 모습을 나타내고 있다.
\end{kfChoices}
\end{kfQBlock}

\kfSecHeader{kfAreaWeekly}{지문}{문항 13 지문 / 주체/객체 혼동}

\kfPassageNote{문항 13 지문 / 주체/객체 혼동}{%
\#\#\# [43 \textasciitilde{} 45] 다음 글을 읽고 물음에 답하시오.

(가)   사월이 돌아와 다사로운 봄볕에   목련이 꽃망울지기 시작하면   내 슬픔은 비롯하나보다.   경운동집 앞마당에   목련이 ㉠ 가지마다 꽃등을 달면   병석의 어머님은 방문을 열고   사월 팔일이 온 것 같다고 웃고 계셨다.

옛날을 꽃피우던   늙은 나무는 죽은 지 오래이고   남은 가지가 자라난 지 스물 두 해   오늘은 ㉡ 아침부터 바람이 불고   연약한 가지에 매어달린 목련은   떠나가는 몸짓을 한다.

목련이 지면 어머님은 떠나가시고   삼백 예순 날이 또 지나가겠지   아 새봄이 와서   가지마다 새싹이 움틀 때까지   나는 서서 나무가 되고 싶다.

* 김광균, 「목련나무 옆에서」 -

(나)   그리운 곳에는 우리를 부르는 소리가 있네   헐벗은 영혼들도 귀의할 안식이 있듯   상처뿐인 삶들도 돌아가 잠들 그리운 집은 있네   천상의 사랑은 ㉢ 이미 빗장을 풀고 달아나버려   보리밭 위로 부는 바람에…(중략)
}
\begin{kfQBlock}{13}{}
\kfQStem{㉠ \textasciitilde{} ㉤에 대한 설명으로 적절하지 않은 것은?}
\begin{kfChoices}
\kfChoice ㉤: 세상의 마지막 저녁을 거듭해서 겪고 있음을 부각하고 있다.
\kfChoice ㉤: 세상의 마지막 저녁을 거듭해서 겪고 있음을 부각하고 있다.
\kfChoice ㉠: 목련이 가지에 낱낱이 피어 있는 형상을 부각하고 있다.
\kfChoice ㉡: 화자가 인식한, 바람이 불기 시작한 시점을 드러내고 있다.
\end{kfChoices}
\end{kfQBlock}

\kfSecHeader{kfAreaWeekly}{지문}{문항 14 지문 / 내용과 방법 혼동}

\kfPassageNote{문항 14 지문 / 내용과 방법 혼동}{%
\#\#\# [22 \textasciitilde{} 25] 다음 글을 읽고 물음에 답하시오.

“김기사장의 인격이나 기술을 우리 사에서는 믿고 맡기고 있었소. 한 회사라는 것은 그 회사의 사업을 위주로 해서 사람을 쓴다는 것은 두말도 할 필요 없겠지요. 김기사는 우리 회사가 환도 후 재건에 있어서 가장 큰 공로를 세운 사람이라고 생각하고 있소. 그러나 금년 삼월에 들어 발전기와 제23호 인쇄기의 고장은 우리 회사의 치명적인 타격이었소. 이렇게 되면 회사에서는 그 기계를 다루는 기사장의 책임으로 돌리지 않을 수 없소. 기사라는 것은 기계의 고장을 사전에 발견하는 것이라고, 아니 고장이 나지 않게 하는 것이 제일 큰 직책이라고 회사에서는 생각하고 있소. 한두 번이 아닌 고장의 수리가 이삼일이 멀다 해서 또 고장 또 고장이라면 결국 기사는 고장의 원인을 모르는 것이라고 해석하지 않을 수 없겠지요. 그 책임을 기사장이 져야 한다면 현명한 기사장은 자기 처신을 어떻게 해야 한다는 것은 여기서 민망한 말을 하지 않아도 잘 이해…(중략)
}
\begin{kfQBlock}{14}{}
\kfQStem{윗글의 서술상의 특징으로 가장 적절하지 않은 것은?}
\begin{kfChoices}
\kfChoice 역순행적 구성을 활용하여 사건이 일어난 이유를 밝히고 있다.
\kfChoice 역순행적 구성을 활용하여 사건이 일어난 이유를 밝히고 있다.
\kfChoice 시대적 배경을 설명하여 갈등 해결의 실마리를 제공하고 있다.
\kfChoice 등장인물이 관찰자 입장에서 인물들의 말과 행동을 전달하고 있다.
\end{kfChoices}
\end{kfQBlock}

\kfSecHeader{kfAreaWeekly}{지문}{문항 15 지문 / 내용 왜곡}

\kfPassageNote{문항 15 지문 / 내용 왜곡}{%
[44 \textasciitilde{} 45] 다음 글을 읽고 물음에 답하시오.   [앞부분 줄거리] 동물원의 코끼리들이 도심으로 탈출했다. 근처 선거 유세장에서는 정치인이 부상을 당하였고, 일대는 쑥대밭이 되었다. 조련사는 유세를 방해하기 위해 일부러 코끼리를 풀어 준 혐의로 경찰서에 붙잡혀 와 조사를 받는다. 참고인 자격의 의사와 아들의 면회를 온 어머니도 함께 있다.

조련사 : 정말인데. 코끼리들은 공연하면서 많이 우는데. 답답하다고 우는데. 슬퍼서 우는데. 난 다 알고 있었는데. 코끼리들이 며칠 전서부터 도망갈 조짐을 보인 것도 알았는데. 도망가려고 의논하는 소릴 들었는데. 그리고 그날은 공원에 갈 때 다른 날과 다르게 빨리 걸었는데. 난 눈치를 챘는데. 오늘이구나. 다른 조련사들이 나한테 다 맡기고 매점에 갔을 때, 코끼리들이 주위를 살피기 시작했는데. 거위들이 꽥꽥댈 때 서로 눈을 마주쳤는데. 나도 코끼리랑 눈이 마주쳤지만 휘파람을 불었는데. 못 본 척 휘파람만 불었는데. 도망가라고. 가서 가족들 …(중략)
}
\begin{kfQBlock}{15}{}
\kfQStem{윗글을 이해한 내용으로 적절하지 않은 것은?}
\begin{kfChoices}
\kfChoice 형사, 의사, 어머니는 서로 의견을 교환하며 조련사를 설득할 방법을 모색했다.
\kfChoice 형사, 의사, 어머니는 서로 의견을 교환하며 조련사를 설득할 방법을 모색했다.
\kfChoice 조련사는 코끼리들이 동물원에서 탈출하려는 모습을 보고도 방관했다고 말했다.
\kfChoice 형사는 조련사에게 배후 세력의 지시를 받았다는 것을 인정하라고 다그쳤다.
\end{kfChoices}
\end{kfQBlock}

\kfSecHeader{kfAreaWeekly}{지문}{문항 16 지문 / 과잉 해석}

\kfPassageNote{문항 16 지문 / 과잉 해석}{%
\#\#\# [24 \textasciitilde{} 26] 다음 글을 읽고 물음에 답하시오.

(가)   팽이가 돈다   어린아해이고 어른이고 살아가는 것이 신기로워   물끄러미 보고 있기를 좋아하는 나의 너무 큰 눈 앞에서   아이가 팽이를 돌린다 / 살림을 사는 아해들도 아름다웁듯이   노는 아해도 아름다워 보인다고 생각하면서   손님으로 온 나는 ㉠ 이 집 주인과의 이야기도 잊어버리고   또 한번 팽이를 돌려주었으면 하고 원하는 것이다   도회 안에서 쫓겨다니는 듯이 사는   나의 일이며 / 어느 소설보다도 신기로운 나의 생활이며   모두 다 내던지고   점잖이 앉은 나의 나이와 나이가 준 나의 무게를 생각하면서   ㉡ 정말 속임 없는 눈으로 / 지금 팽이가 도는 것을 본다   그러면 팽이가 까맣게 변하여 서서 있는 것이다   누구 집을 가보아도 나 사는 곳보다는 여유가 있고   바쁘지도 않으니 / 마치 별세계(別世界)같이 보인다   팽이가 돈다 / 팽이가 돈다   팽이 밑바닥에 끈을 돌려 매이니 이상하고   …(중략)
}
\begin{kfQBlock}{16}{}
\kfQStem{<보기>를 바탕으로 (가), (나)를 감상한 내용으로 적절하지 않은 것은? [3점]}
\begin{kfChoices}
\kfChoice (나)의 ‘댓가지가 바람에 흔들리면서 나는 소리’를 ‘물소리처럼 들리는 거’라고 하는 말에서 ‘나’가 자연에 내재한 풍부한 생명력을 느끼고 있음을 알 수 있군.
\kfChoice (나)의 ‘댓가지가 바람에 흔들리면서 나는 소리’를 ‘물소리처럼 들리는 거’라고 하는 말에서 ‘나’가 자연에 내재한 풍부한 생명력을 느끼고 있음을 알 수 있군.
\kfChoice (가)의 ‘도회 안에서 쫓겨다니는 듯이 사는’과 ‘누구 집을 가보아도 나 사는 곳보다는 여유가 있고’에서 ‘나’의 생활이 고단함을 알 수 있군.
\kfChoice (가)의 ‘팽이가 돌면서 나를 울린다’와 ‘팽이는 나를 비웃는 듯이 돌고 있다’에서 ‘나’가 팽이가 도는 모습을 보며 주체성이 결여된 자신의 삶을 성찰하고 있음을 알 수 있군.
\end{kfChoices}
\end{kfQBlock}

\kfSecHeader{kfAreaWeekly}{지문}{문항 17 지문 / 범위 오류}

\kfPassageNote{문항 17 지문 / 범위 오류}{%
\#\#\# 2024년 11월 고3 수능

[22-27] 다음 글을 읽고 물음에 답하시오.

(가)   배를 민다   배를 밀어보는 것은 아주 드문 경험   희번덕이는 잔잔한 가을 바닷물 위에   배를 밀어넣고는   온몸이 아주 추락하지 않을 순간의 한 허공에서   밀던 힘을 한껏 더해 밀어주고는   아슬아슬히 배에서 떨어진 손, 순간 환해진 손을   허공으로부터 거둔다

사랑은 참 부드럽게도 떠나지   뵈지도 않는 길을 부드럽게도

배를 한껏 세게 밀어내듯이 슬픔도   그렇게 밀어내는 것이지

배가 나가고 남은 빈 물 위의 흉터   잠시 머물다 가라앉고

그런데 오, 내 안으로 들어오는 배여   아무 소리 없이 밀려들어오는 배여   - 장석남, <배를 밀며>

(나)   당신……, 당신이라는 말 참 좋지요, 그래서 불러봅니다 킥킥거리며 한때 적요로움의 울음이 있었던 때, 한 슬픔이 문을 닫으면 또 한 슬픔이 문을 여는 것을 이만큼 살아옴의 상처에 기대, 나 킥킥……, 당신을 부릅니다 단…(중략)
}
\begin{kfQBlock}{17}{}
\kfQStem{문항별}
\begin{kfChoices}
\kfChoice 번은 '통제할 수 없는 익명의 욕구'를 상대를 향한 사랑이 운명적이어서 멈출 수 없다는 뜻으로 해석했습니다.
\kfChoice 번은 '통제할 수 없는 익명의 욕구'를 상대를 향한 사랑이 운명적이어서 멈출 수 없다는 뜻으로 해석했습니다.
\kfChoice 번은 배가 들어오는 것을 '대상과의 재회가 예상대로 이루어짐'이라고 해석했습니다.
\kfChoice 번은 '당신'을 화자의 내면에 사는 '병자'로서 연민의 대상이라고 설명했습니다.
\end{kfChoices}
\end{kfQBlock}

\kfSecHeader{kfAreaWeekly}{지문}{문항 18 지문 / 시점/화자 혼동}

\kfPassageNote{문항 18 지문 / 시점/화자 혼동}{%
[38\textasciitilde{}42] 다음 글을 읽고 물음에 답하시오.

(가)   장마 가뭄에 피해 입은 백성이 관찰사 가을 순행 기다림은   가을걷이 부족함을 채워줄까 해서인데 지나는 곳마다 죄를 묻는 폐단 있네   무논 재해도 감췄는데 목화밭이야 거론할까   백 묘(畝)나 되는 벌건 땅에 백지징세 하는구나   인자한 우리 임금 곡식 한 묶음도 모래 덮일까 염려하는데   불쌍한 백성 논밭에다 좁은 길 넓히란다   각읍 관리 독촉하니 채찍 몽둥이 낭자하다   허다한 관인들이 대호(大戶) 소호(小戶)에 분담시켜   사방(四方) 부근 십 리 안에 닭과 개가 멸종하네   부자는 괜찮지만 가련한 이 가난한 자로다   해는 기울고 이정*은 저녁밥 재촉할 때   텅 빈 부엌에서 우는 아낙 발 구르며 하는 말이   방아품에 얻은 양식 한두 되 있건마는   채소도 있건마는 그릇은 누구에게 빌릴꼬   앞뒷집 돌아보니 섣달그믐에 시루 빌리는 격이로다   한 마을 닭과 개 다 먹어 치우고 집집마다 또 거둔단 말인가   대호…(중략)
}
\begin{kfQBlock}{18}{}
\kfQStem{<보기>를 바탕으로 (가)\textasciitilde{}(다)를 감상한 내용으로 적절하지 않은 것은?}
\begin{kfChoices}
\kfChoice (나)의 산천은 인사로 인해 변해 버린다는 점에서 변함없는 자연에 대한 화자의 소망을 투영한 공간이라고 할 수 있다.
\kfChoice (나)의 산천은 인사로 인해 변해 버린다는 점에서 변함없는 자연에 대한 화자의 소망을 투영한 공간이라고 할 수 있다.
\kfChoice (가)의 논밭은 지배층을 위해 길로 넓혀진다는 점에서 백성들이 빼앗긴 삶의 터전을 의미하는 공간이라고 할 수 있다.
\kfChoice (가)의 텅 빈 부엌은 방아품으로 얻은 양식을 담을 그릇조차 없는 곳이라는 점에서 아낙이 자신의 처지에 슬픔을 느끼는 공간이라고 할 수 있다.
\end{kfChoices}
\end{kfQBlock}



\clearpage

\kfChapterCover{4}{비트겐슈타인2 ch04}{Chapter 4}{이번 챕터: 문법 → 문학 5작품 → 비문학 5지문 → 패턴 워크북.}

\kfStartGrammar{이번 챕터 문법 영역 — 음운 / 구개음화·된소리되기 + 음운 변동 유형 분류 종합}


\kfSecHeader{kfAreaGrammar}{문제}{}

\begin{multicols}{2}\raggedcolumns\setlength{\columnsep}{12pt}
\kfPassageFull{문항 1 지문 (concept)}{%
구개음화란 끝소리가 'ㄷ, ㅌ'인 형태소가 모음 'ㅣ'나 반모음 'ㅣ'로 시작하는 형식 형태소와 만날 때 'ㅈ, ㅊ'으로 바뀌는 현상이다. '굳이'[구지], '같이'[가치], '해돋이'[해도지]가 대표적인 예이다. 구개음화의 핵심 조건은 두 가지이다. 첫째, 앞 형태소의 받침이 'ㄷ' 또는 'ㅌ'이어야 한다. 둘째, 뒤에 오는 형식 형태소가 'ㅣ' 모음이나 반모음 'ㅣ'로 시작해야 한다. 따라서 실질 형태소끼리의 결합에서는 구개음화가 일어나지 않는다. '디딤돌'의 '디'에서 'ㄷ' 뒤에 'ㅣ'가 오지만, 이는 하나의 실질 형태소 내부이므로 구개음화가 적용되지 않는다.
}
\begin{kfQBlock}{1}{}
\kfQStem{윗글을 바탕으로 <보기>에서 구개음화가 적용된 것만을 모두 고른 것은?

<보기> ㄱ. 굳이[구지]    ㄴ. 디디다[디디다]    ㄷ. 붙이다[부치다] ㄹ. 닫히다[다치다]    ㅁ. 티끌[티끌]}
\begin{kfChoices}
\kfChoice ㄱ, ㄷ
\kfChoice ㄱ, ㄷ, ㄹ
\kfChoice ㄱ, ㄴ, ㄷ
\kfChoice ㄱ, ㄷ, ㅁ
\kfChoice ㄱ, ㄴ, ㄷ, ㄹ
\end{kfChoices}
\end{kfQBlock}

\kfPassageFull{문항 2 지문 (concept)}{%
된소리되기(경음화)란 예사소리가 된소리로 바뀌는 현상이다. 된소리되기가 적용되는 대표적 환경은 다음과 같다. ① 받침 'ㄱ, ㄷ, ㅂ'(파열음) 뒤에서 예사소리 'ㄱ, ㄷ, ㅂ, ㅅ, ㅈ'이 된소리로 바뀐다. '국밥'[국빱], '앞산'[압싼]. ② 어간 받침 'ㄴ, ㅁ' 뒤에서 어미 첫소리 'ㄱ, ㄷ, ㅅ, ㅈ'이 된소리로 바뀐다. '신고'[신꼬], '감다'[감따]. ③ 관형사형 어미 '-(으)ㄹ' 뒤에서 된소리로 바뀐다. '할 것'[할 껏], '갈 데'[갈 떼]. ④ 한자어에서 'ㄹ' 받침 뒤 'ㄷ, ㅅ, ㅈ'이 된소리로 바뀐다. '발달'[발딸], '물질'[물찔]. 각 환경을 정확히 구분하는 것이 중요하다.
}
\begin{kfQBlock}{2}{}
\kfQStem{윗글을 바탕으로 <보기>의 단어들에서 된소리되기가 적용되는 환경을 분류한 것으로 적절한 것은?

<보기> ⓐ 학교[학꾜]    ⓑ 안다[안따]    ⓒ 할 수[할 쑤]    ⓓ 갈등[갈뜽]    ⓔ 문법[문뻡]}
\begin{kfChoices}
\kfChoice ⓐ와 ⓔ는 같은 환경(①)에서 된소리되기가 적용된다.
\kfChoice ⓑ와 ⓔ는 같은 환경(②)에서 된소리되기가 적용된다.
\kfChoice ⓐ는 환경①, ⓑ는 환경②, ⓒ는 환경③에 해당한다.
\kfChoice ⓒ와 ⓓ는 같은 환경(③)에서 된소리되기가 적용된다.
\kfChoice ⓓ와 ⓔ는 같은 환경(④)에서 된소리되기가 적용된다.
\end{kfChoices}
\end{kfQBlock}

\kfPassageFull{문항 3 지문 (concept)}{%
음운 변동은 교체, 탈락, 첨가, 축약의 네 유형으로 분류된다. 교체는 하나의 음운이 다른 음운으로 바뀌는 것으로, 비음화·유음화·구개음화·된소리되기·끝소리 규칙이 이에 해당한다. 탈락은 음운이 없어지는 것으로, 자음군 단순화(겹받침에서 하나 탈락), 'ㅎ' 탈락, 'ㅡ' 탈락, 동일 모음 탈락 등이 있다. 첨가는 없던 음운이 새로 생기는 것으로, 사이시옷에 의한 'ㄴ' 첨가가 대표적이다. 축약은 두 음운이 하나로 합쳐지는 것으로, 'ㅎ' 축약('놓다'+'고'→[노코], 'ㄷ'+'ㅎ'→[ㅌ])과 모음 축약('ㅗ'+'ㅏ'→'ㅘ' 등)이 있다.
}
\begin{kfQBlock}{3}{}
\kfQStem{윗글을 바탕으로 <보기>의 각 예에 적용된 음운 변동의 유형을 바르게 짝지은 것은?

<보기> ⓐ 놓고[노코]    ⓑ 꽃잎[꼰닙]    ⓒ 쓰어→써    ⓓ 맏이[마지]}
\begin{kfChoices}
\kfChoice ⓐ 축약, ⓑ 교체+첨가, ⓒ 축약, ⓓ 교체
\kfChoice ⓐ 교체, ⓑ 교체, ⓒ 탈락, ⓓ 교체
\kfChoice ⓐ 축약, ⓑ 첨가, ⓒ 탈락, ⓓ 축약
\kfChoice ⓐ 교체, ⓑ 교체+첨가, ⓒ 축약, ⓓ 첨가
\kfChoice ⓐ 축약, ⓑ 교체+첨가, ⓒ 탈락, ⓓ 교체
\end{kfChoices}
\end{kfQBlock}

\kfPassageFull{문항 4 지문 (concept)}{%
음운 변동이 적용되면 음운의 개수가 변하거나 유지된다. 교체는 하나의 음운이 다른 하나로 바뀌므로 음운의 개수가 변하지 않는다. 탈락은 음운이 사라지므로 개수가 줄어든다. 첨가는 새 음운이 생기므로 개수가 늘어난다. 축약은 두 음운이 하나로 합쳐지므로 개수가 줄어든다. 이를 활용하면 특정 단어에 적용된 음운 변동의 유형을 역으로 추론할 수 있다. 예를 들어 '놓다'[노타]에서 음운의 수가 'ㄴ, ㅗ, ㅎ, ㄷ, ㅏ'(5개) → 'ㄴ, ㅗ, ㅌ, ㅏ'(4개)로 줄었으므로 축약이 적용되었음을 알 수 있다.
}
\begin{kfQBlock}{4}{}
\kfQStem{윗글을 바탕으로 <보기>의 단어에서 음운 변동 전후의 음운 개수 변화를 분석한 것으로 적절한 것은?

<보기> ⓐ 꽃을[꼬츨] → 연음    ⓑ 닫히다[다치다] → 축약+구개음화 ⓒ 솜이불[솜니불] → 첨가    ⓓ 국물[궁물] → 비음화}
\begin{kfChoices}
\kfChoice ⓐ에서 음운의 개수가 줄어든다.
\kfChoice ⓑ에서 축약과 구개음화가 일어나 음운 개수가 1개 줄어든다.
\kfChoice ⓒ에서 'ㄴ'이 새로이 첨가되어 음운 개수가 2개 늘어난다.
\kfChoice ⓓ에서 비음화가 적용되어 음운 개수가 1개 줄어든다.
\kfChoice ⓐ\textasciitilde{}ⓓ 중 음운 개수가 변하지 않는 것은 ⓓ뿐이다.
\end{kfChoices}
\end{kfQBlock}

\kfPassageFull{문항 5 지문 (concept)}{%
하나의 단어에 여러 음운 변동이 복합적으로 적용되는 경우, 각 변동의 유형·순서·결과를 체계적으로 분석해야 한다. '꽃히다'(존재하지 않는 단어이나 예시로)가 아닌 실제 단어 '맞히다'를 분석해 보자. '맞-'+'히다'에서 먼저 'ㅈ'과 'ㅎ'이 축약되어 'ㅊ'이 된다. 그 결과 [마치다]가 된다. 또 다른 예로 '넓히다'에서는 겹받침 'ㄼ' 중 'ㄹ'이 탈락(자음군 단순화)하고 'ㅂ'+'ㅎ'이 축약되어 'ㅍ'이 되므로 [널피다]가 된다. 이렇게 탈락과 축약이 연쇄적으로 적용되는 경우도 있다.
}
\begin{kfQBlock}{5}{}
\kfQStem{윗글을 바탕으로 <보기>의 단어에 적용된 음운 변동의 유형과 순서를 바르게 분석한 것은?

<보기> 「읽히다」의 발음[일키다]를 도출하는 과정을 분석하시오.}
\begin{kfChoices}
\kfChoice 'ㄺ'에서 'ㄹ' 탈락(탈락) → 'ㄱ'+'ㅎ'→'ㅋ'(축약) → [이키다]
\kfChoice 'ㄺ'에서 'ㄱ' 탈락(탈락) → 'ㄹ'+'ㅎ'→'ㄹ'(축약?) → 적용 불가
\kfChoice 'ㄺ'에서 'ㄱ'+'ㅎ'→'ㅋ'(축약) → 'ㄹ'이 받침에 남음 → [일키다]
\kfChoice 'ㄺ' → [ㄱ](끝소리 규칙) → 'ㄱ'+'ㅎ'→'ㅋ'(축약) → [익키다]
\kfChoice 'ㄺ'에서 'ㄹ' 탈락(탈락) → 구개음화 → [이치다]
\end{kfChoices}
\end{kfQBlock}

\end{multicols}


\kfStartLit{이번 챕터 문학 작품 5편}


\kfSecHeader{kfAreaLiterature}{지문}{춘향전 — 작자 미상 (고전소설(판소리계 소설)) / 「춘향전」(완판 84장본 계열), 수능·학평 기출 다수}

\kfPassageNote{춘향전 — 작자 미상 (고전소설(판소리계 소설)) / 「춘향전」(완판 84장본 계열), 수능·학평 기출 다수}{%
변학도 분부하되, \\[4pt]
"네 수청 거역하니 관장(官杖) 맞아 죽기가 서러우냐, 내 말 좇아 영화가 서러우냐. 바로 대답하여라." \\[4pt]
춘향이 기가 차서, \\[4pt]
"일부종사(一夫從事) 굳은 마음 물에 빠져 죽는 한이 있어도 변할 마음 추호도 없사오니, 천만 번 죽여 주십시오." \\[14pt]
(중략) \\[14pt]
이때 어사또 출도(出道)하니, \\[4pt]
"암행어사 출도야!" \\[4pt]
좌우의 나졸이 벌여 서며 금부도사(禁府都事) 앞세우니, \\[4pt]
변학도 혼비백산(魂飛魄散)하여 기절하여 넘어진다. \\[4pt]
어사또 변학도 죄상을 낱낱이 읽어 내니, \\[4pt]
첫째 백성의 재물을 빼앗은 죄, \\[4pt]
둘째 무죄한 백성을 가두어 때린 죄, \\[4pt]
셋째 관기(官妓) 아닌 양인(良人)의 딸에게 수청을 강요한 죄, \\[4pt]
이 세 죄를 엮어 문초(問招)하니, \\[4pt]
변학도 아무 말도 못하더라.
}
\kfSecHeader{kfAreaLiterature}{활동}{작품 정리표 — 춘향전}

\noindent\textit{\small 빈칸에 알맞은 말을 채워 넣으시오. (초성 힌트 참고)}\par\medskip
\begin{kfActTable}{|>{\centering\arraybackslash}m{2.6cm}|m{6cm}|m{4cm}|}
\rowcolor{kfPrimaryLight!50}\textbf{\color{kfPrimary}항목} & \textbf{\color{kfPrimary}내용 (빈칸 채우기)} & \textbf{\color{kfPrimary}답안 작성}\\\hline
갈래 & ㅍㅅㄹㄱ ㅅㅅ ① & \kfTBlank \\\hline
주인공 & ㅅㅊㅎ과 ㅇㅁㄹ ② & \kfTBlank \\\hline
악인 & ㅂㅎㄷ ③ & \kfTBlank \\\hline
춘향의 태도 & ㅇㅂㅈㅅ로 수청 거부 = ㅈㄱ ④ & \kfTBlank \\\hline
반전 장치 & ㅇㅎㅇㅅ 출도 ⑤ & \kfTBlank \\\hline
변학도 죄상 1 & 백성의 ㅈㅁ을 빼앗음 ⑥ & \kfTBlank \\\hline
변학도 죄상 2 & 무죄한 백성을 ㄱㄷ ⑦ & \kfTBlank \\\hline
변학도 죄상 3 & ㅇㅇ의 딸에게 ㅅㅊ 강요 ⑧ & \kfTBlank \\\hline
대비 구성 & 호화 ㅈㅊ vs 옥중 ㅊㅎ ⑨ & \kfTBlank \\\hline
주제 1 & ㅅㄹ의 성취 ⑩ & \kfTBlank \\\hline
주제 2 & 부당한 권력에 대한 ㅈㅎ ⑪ & \kfTBlank \\\hline
주제 3 & ㅅㅁ ㅈㄷ의 모순 비판 ⑫ & \kfTBlank \\\hline
\end{kfActTable}
\kfSecHeader{kfAreaLiterature}{문제}{}

\begin{multicols}{2}\raggedcolumns\setlength{\columnsep}{12pt}
\begin{kfQBlock}{1}{}
\kfQStem{윗글에 대한 설명으로 적절하지 않은 것은?}
\begin{kfChoices}
\kfChoice 변학도와 춘향의 대화에서 권력자와 피지배자 사이의 갈등이 드러난다.
\kfChoice 암행어사 출도 장면에서 극적 반전이 이루어지고 있다.
\kfChoice 변학도의 죄상을 열거하여 부패한 관리의 모습을 구체적으로 폭로하고 있다.
\kfChoice 춘향은 변학도의 제안을 수락하여 신분 상승을 꾀하고 있다.
\kfChoice '혼비백산'이라는 한자성어를 사용하여 변학도의 공포를 생생하게 표현하고 있다.
\end{kfChoices}
\end{kfQBlock}

\begin{kfQBlock}{2}{}
\kfQStem{윗글에서 변학도의 죄상 중 '관기 아닌 양인의 딸에게 수청을 강요한 죄'가 지닌 의미로 가장 적절한 것은?}
\begin{kfChoices}
\kfChoice 관기에 대한 수청 강요는 본래 합법이었으므로 춘향의 경우에만 문제가 된다.
\kfChoice 양인의 딸을 관기 취급한 것으로 신분 제도의 모순과 권력 남용을 함께 보여 준다.
\kfChoice 양인의 딸이라는 점에서 춘향의 집안이 가진 경제적 부유함을 부각해 보여 준다.
\kfChoice 변학도가 관기와 양인을 구별하지 못한 단순한 행정상 실수를 지적해 드러낸다.
\kfChoice 수청이라는 관행 자체가 법적으로 전혀 허용되지 않았음을 명확히 보여 준다.
\end{kfChoices}
\end{kfQBlock}

\begin{kfQBlock}{3}{}
\kfQStem{윗글에서 춘향이 '일부종사(一夫從事) 굳은 마음'이라고 말한 까닭으로 가장 적절한 것은?}
\begin{kfChoices}
\kfChoice 이몽룡에 대한 사랑과 절개를 지키겠다는 굳은 의지를 표현하기 위해서이다.
\kfChoice 변학도라는 인물에 대한 깊은 존경심을 우회적으로 드러내기 위해서이다.
\kfChoice 자신이 갖춘 학식과 교양을 변학도 앞에서 과시하고 싶었기 때문이다.
\kfChoice 관기로서의 본분을 끝까지 다하겠다는 굳은 다짐을 보여 주기 위해서이다.
\kfChoice 변학도의 수청 제안을 일부 수락하되 까다로운 조건을 달기 위해서이다.
\end{kfChoices}
\end{kfQBlock}

\begin{kfQBlock}{4}{}
\kfQStem{윗글에서 암행어사 출도 장면이 독자(청중)에게 주는 효과로 가장 적절한 것은?}
\begin{kfChoices}
\kfChoice 변학도에 대한 깊은 동정심을 유발하여 비극적 결말을 예고하고 있다.
\kfChoice 정의 실현의 통쾌함과 권선징악의 쾌감을 청중에게 제공한다.
\kfChoice 조선 시대 암행어사 제도의 여러 문제점을 비판하는 계기를 마련한다.
\kfChoice 춘향이 지킨 절개가 결국 무의미했음을 깨닫게 하여 허탈감을 안긴다.
\kfChoice 이몽룡의 출세에 대한 강한 욕망을 부정적으로 평가하게 만든다.
\end{kfChoices}
\end{kfQBlock}

\begin{kfQBlock}{5}{}
\kfQStem{윗글을 바탕으로 <보기>를 감상한 내용으로 적절하지 않은 것은?

<보기> 「춘향전」은 조선 후기 사회에서 신분 제도의 모순과 부패한 관리의 횡포를 비판하는 작품이다. 동시에 사랑의 성취와 정의의 실현이라는 보편적 가치를 추구한다. 특히 암행어사라는 제도적 장치를 통해 '법과 정의'의 문제를 서사적으로 풀어낸다.}
\begin{kfChoices}
\kfChoice 변학도의 죄상 열거는 부패한 관리의 횡포를 구체적으로 폭로하는 장면으로 볼 수 있겠군.
\kfChoice 춘향의 수청 거부는 부당한 권력에 굴복하지 않는 저항의 의미를 지니겠군.
\kfChoice 암행어사 출도는 법과 정의가 실현되는 제도적 장치로 기능하겠군.
\kfChoice '양인의 딸에게 수청을 강요한 죄'는 신분 제도의 모순을 드러내는 대목이겠군.
\kfChoice 변학도의 처벌은 춘향의 사랑이 신분 제도에 의해 좌절되었음을 보여주겠군.
\end{kfChoices}
\end{kfQBlock}

\end{multicols}

\kfSecHeader{kfAreaLiterature}{지문}{이춘풍전 — 작자 미상 (고전소설) / 작자 미상, 「이춘풍전」 — 학평 기출 수록}

\kfPassageNote{이춘풍전 — 작자 미상 (고전소설) / 작자 미상, 「이춘풍전」 — 학평 기출 수록}{%
춘풍 아내 곁에 앉아 하는 말이 “마오 마오 그리 마오. 청루미색 좋아 마오. 자고로 이런 사람이 어찌 망하지 않을까? 내 말을 자세히 들어보소. 미나리골 박화진이라는 이는 청루미색 즐기다가 나중에는 굶어 죽고, 남산 밑에 이 패두는 소년 시절 부자였으나 주색에 빠져 다니다가 늙어서는 상거지 되고, 모시전골 김 부자는 술 잘 먹기 유명하여 누룩 장수가 도망을 다니기로 장안에 유명터니 수만금을 다 없애고 끝내 똥 장수가 되었다니, 이것으로 두고 볼지라도 청루잡기 잡된 마음 부디부디 좋아 마소.”

춘풍이 대답하되, “자네 내 말 들어보게. 그 말이 다 옳다 하되, 이 앞집 매갈쇠는 한잔 술도 못 먹어도 돈 한 푼 못 모으고, 비우고개 이도명은 오십이 다 되도록 주색을 몰랐으되 남의 집만 평생 살고, 탁골 사는 먹돌이는 투전 잡기 몰랐으되 수천 금 다 없애고 나중에는 굶어 죽었으니, 이런 일을 두고 볼지라도 주색잡기 안 한다고 잘 사는 바 없느니라. 내 말 자네 들어보게. 술 잘 먹던 이태백은 호사스런 술잔으로 매일 장취 놀았으되 한림학사 다 지내고, 투전에 으뜸인 원두표는 잡기를 방탕히 하여 소년부터 유명했으나 나중에 잘되어서 정승 벼슬 하였으니, 이로 두고 볼진대 주색잡기 좋아하기는 장부의 할 바라. 나도 이리 노닐다가 나중에 일품 정승 되어 후세에 전하리라.”

아내의 말을 아니 듣고 수틀리면 때리기와 전곡 남용 일삼으니 이런 변이 또 있을까? 이리저리 놀고 나니 집안 형용 볼 것 없다.

[중략 부분 줄거리] 춘풍 아내가 열심히 품을 팔아 집안을 일으키자 춘풍은 다시 교만해지고, 아내의 만류에도 호조에서 이천 냥을 빌려 평양으로 장사를 떠나게 된다. 춘풍이 평양에서 기생 추월의 유혹에 넘어가 장사는 하지 않고 재물을 모두 탕진한 채 추월의 하인이 되었다는 소식을 듣고 춘풍의 아내가 통곡한다.

이리 한참 울다가 도로 풀고 생각하되, ‘우리 가장 경성으로 데려다가 호조 돈 이천 냥을 한 푼 없이 다 갚은 후에 의식 염려 아니하고 부부 둘이 화락하여 백 년 동락하여 보자. 평생의 한이로다.’

마침 그때 김 승지 댁이 있으되 승지는 이미 죽고, 맏자제가 문장을 잘해 소년 급제하여 한림 옥당 다 지내고 도승지를 지낸 고로, 작년에 평양 감사 두 번째 물망에 있다가 올해 평양 감사 하려고 도모한단 말을 사환 편에 들었것다. 승지 댁이 가난하여 아침저녁으로 국록을 타서 많은 식구들이 사는 중에 그 댁에 노부인 있다는 말을 듣고, 바느질품을 얻으려고 그 댁에 들어가니, 후원 별당 깊은 곳에 도승지의 모부인이 누웠는데 형편이 가난키로 식사도 부족하고 의복도 초췌하다. 춘풍 아내 생각하되, ‘이 댁에 붙어서 우리 가장 살려내고 추월에게 복수도 할까.’

천만 의외로 김 승지가 평양 감사가 되었구나. 춘풍 아내, 부인 전에 문안하고 여쭈되, “승지 대감, 평양 감사 하였사오니 이런 경사 어디 있사오리까?” 부인이 이른 말이, “나도 평양으로 내려 갈 제, 너도 함께 따라가서 춘풍이나 찾아보아라.” 하니 춘풍 아내 여쭈되, “소녀는 고사하옵고 오라비가 있사오니 비장으로 데려가 주시길 바라나이다.” 대부인이 이른 말이, “네 청이야 아니 듣겠느냐? 그리하라.” 허락하고 감사에게 그 말을 하니 감사도 허락하고, “회계 비장 하라.” 하니 좋을시고, 좋을시고. 춘풍의 아내 없던 오라비를 보낼쏜가? 제가 손수 가려고 여자 의복 벗어놓고 남자 의복 치장한다.
}
\kfSecHeader{kfAreaLiterature}{활동}{작품 정리표 — 이춘풍전(작자 미상)}

\noindent\textit{\small 빈칸에 알맞은 말을 채워 넣으시오. (초성 힌트 참고)}\par\medskip
\begin{kfActTable}{|>{\centering\arraybackslash}m{2.6cm}|m{6cm}|m{4cm}|}
\rowcolor{kfPrimaryLight!50}\textbf{\color{kfPrimary}항목} & \textbf{\color{kfPrimary}내용 (빈칸 채우기)} & \textbf{\color{kfPrimary}답안 작성}\\\hline
갈래 & ㅍㅅㄹㄱ ㅅㅅ, ㅍㅈ ㅅㅅ, ㅅㅌ ㅅㅅ ① & \kfTBlank \\\hline
남편 & ㅇㅊㅍ ② & \kfTBlank \\\hline
아내 & 춘풍의 ㅊ(이름은 본문에 직접 등장하지 않음) ③ & \kfTBlank \\\hline
기생 & 평양 기생 ㅊㅇ ④ & \kfTBlank \\\hline
춘풍의 합리화 & ㅇㅌㅂ·ㅇㄷㅍ의 사례를 ㅂㄹ로 듦 ⑤ & \kfTBlank \\\hline
탕진 자금 & 부모 유산 + ㅎㅈ 돈 ㅇㅊ ㄴ ⑥ & \kfTBlank \\\hline
아내의 결심 & ‘우리 가장 살려내고 ㅊㅇ에게 ㅂㅅ’ ⑦ & \kfTBlank \\\hline
접근 대상 & ㄱ ㅅㅈ 댁 ㄴㅂㅇ ⑧ & \kfTBlank \\\hline
수단 & ㅂㄴㅈㅍ + ㅊㄷㅅ ⑨ & \kfTBlank \\\hline
결정적 청 & ㅇㄹㅂ를 ㅂㅈ으로 데려가 달라 ⑩ & \kfTBlank \\\hline
반전 & ‘없던 오라비’ — 본인이 ㄴㅈ 후 ㅎㄱ ㅂㅈ ⑪ & \kfTBlank \\\hline
주제 & ㅁㄴ한 가장에 대한 풍자와 ㅈㅊㅈ 여성상의 제시 ⑫ & \kfTBlank \\\hline
\end{kfActTable}
\kfSecHeader{kfAreaLiterature}{문제}{}

\begin{multicols}{2}\raggedcolumns\setlength{\columnsep}{12pt}
\begin{kfQBlock}{1}{}
\kfQStem{윗글을 이해한 내용으로 적절하지 않은 것은?}
\begin{kfChoices}
\kfChoice 춘풍은 호조 돈 이천 냥을 빌려 평양으로 떠났다.
\kfChoice 춘풍 아내는 바느질품을 팔며 생계를 이었다.
\kfChoice 춘풍 아내는 춘풍의 잘못에도 가정의 화목을 바라고 있다.
\kfChoice 도승지는 평양 감사직을 연이어 두 번 맡게 되었다.
\kfChoice 대부인은 춘풍 아내의 정성을 흡족히 받아들이고 그의 청을 들어준다.
\end{kfChoices}
\end{kfQBlock}

\begin{kfQBlock}{2}{}
\kfQStem{윗글에서 춘풍의 아내가 ‘김 승지 댁’ 노부인에게 정성껏 차담상을 차리는 행위의 의도로 가장 적절한 것은?}
\begin{kfChoices}
\kfChoice 노부인의 환심을 얻어 평양 감사 댁의 도움으로 남편을 구하려는 계획
\kfChoice 어려운 이웃을 보살펴 마을의 도덕적 모범이 되려는 자선의 표현
\kfChoice 도승지에게 자신의 음식 솜씨를 자랑하기 위한 시연
\kfChoice 노부인을 사적인 양어머니로 삼아 신분 상승을 꾀하려는 시도
\kfChoice 기생 추월에게 같은 음식을 보내어 화해하려는 사전 준비
\end{kfChoices}
\end{kfQBlock}

\begin{kfQBlock}{3}{}
\kfQStem{윗글에서 춘풍이 자신의 주색잡기를 합리화한 방식으로 가장 적절한 것은?}
\begin{kfChoices}
\kfChoice 성리학적 권위를 빌려 자신의 행동에 도덕적 정당성을 부여한다.
\kfChoice 춘풍 아내의 사례를 그대로 인정하면서 그 결론만 다르게 받아들인다.
\kfChoice 이태백·원두표의 사례를 들어 춘풍 아내가 든 사례의 반례로 제시한다.
\kfChoice 기생 추월의 입장에서 자신의 행위를 바라보며 변호한다.
\kfChoice 관료의 권위를 내세워 일방적으로 명령한다.
\end{kfChoices}
\end{kfQBlock}

\begin{kfQBlock}{4}{}
\kfQStem{윗글에서 춘풍의 아내가 ‘오라비를 비장으로 데려가 주시길’ 청한 까닭으로 가장 적절한 것은?}
\begin{kfChoices}
\kfChoice 오라비가 평양 풍경을 구경하기를 평소 간절히 원했기 때문
\kfChoice 오라비를 통해 도승지의 정치적 지지 세력을 키우기 위해
\kfChoice 자신이 직접 남장을 하고 회계 비장이 되어 평양으로 가기 위해
\kfChoice 오라비가 회계와 무역에 능통한 전문가이기 때문
\kfChoice 춘풍이 평양에서 외로워하지 않도록 동행자를 보내려고
\end{kfChoices}
\end{kfQBlock}

\begin{kfQBlock}{5}{}
\kfQStem{윗글을 바탕으로 <보기>를 감상한 내용으로 적절하지 않은 것은?

<보기> 이 작품은 남편이 저지른 일을 아내가 수습하는 서사가 중심이 된다. 춘풍은 가장이지만 경제관념 없이 현실적 쾌락만을 추구하며 자신이 초래한 문제를 해결하려 하지 않는다. 반면 춘풍 아내는 적극적으로 현실의 문제를 해결하려는 의지를 갖고 주도면밀하게 목적을 달성한다. 이러한 두 인물의 대비를 통해 무능한 가장의 모습과 주체적인 아내의 역할 및 능력이 부각된다.}
\begin{kfChoices}
\kfChoice 춘풍이 부모 유산을 다 탕진한 뒤 아내에게 집안일에 대한 모든 권리를 넘기는 것에서 무책임한 가장의 모습을 엿볼 수 있군.
\kfChoice 춘풍이 전곡을 남용하고 주색잡기에 빠져 있는 것에서 경제관념 없이 현실적 쾌락만 추구하는 모습을 엿볼 수 있군.
\kfChoice 춘풍 아내가 사환에게 정보를 얻고 김 승지 댁 노부인에게 의도적으로 접근한 것에서 주도면밀한 모습을 엿볼 수 있군.
\kfChoice 춘풍 아내가 비장 지위를 획득하고 남장을 하는 것에서 적극적인 문제 해결 의지를 엿볼 수 있군.
\kfChoice 춘풍이 각서를 쓰고 춘풍 아내가 차담상을 차리는 행위에서 두 인물 모두 신분 상승을 통해 목적을 이루려 했음을 엿볼 수 있군.
\end{kfChoices}
\end{kfQBlock}

\end{multicols}

\kfSecHeader{kfAreaLiterature}{지문}{유충렬전 — 작자 미상 (고전소설(군담소설·영웅소설)) / 「유충렬전」, 수능·학평 기출}

\kfPassageNote{유충렬전 — 작자 미상 (고전소설(군담소설·영웅소설)) / 「유충렬전」, 수능·학평 기출}{%
정한담이 천자(天子)께 아뢰되, "유심(유충렬의 아버지)은 반역을 꾀하옵고, 그 아들 유충렬은 도적의 무리와 내통하여 사직을 위태롭게 하옵니다. 속히 잡아 죄를 다스리소서." 천자 크게 노하여 유심을 하옥(下獄)하고, 유충렬을 잡아들이라 하시니,

이때 유충렬이 천리 밖에서 이 소식을 듣고 하늘을 우러러 통곡하되, "아버지 충성이 하늘 같거늘 간신의 모함에 옥에 갇히시니, 이 원수를 갚지 않으면 사람이 아니로다."

(중략)

유충렬이 대군을 이끌고 도성(都城)으로 돌아오니, 천자 친히 성문 밖에 나와 맞으시며, "경의 충성이 아니었으면 사직이 위태로웠을 것이라." 하시고 정한담의 죄를 낱낱이 밝히시니, 정한담이 땅에 엎드려 두 손을 묶인 채 용서를 빌되, "신의 죄 만 번 죽어 마땅하오나 목숨만 살려 주소서." 천자 노하여, "충신을 모함하고 사직을 어지럽힌 죄, 용서할 수 없노라." 하시고 참형(斬刑)에 처하시니, 만백성이 환호하더라.
}
\kfSecHeader{kfAreaLiterature}{활동}{작품 정리표 — 유충렬전}

\noindent\textit{\small 빈칸에 알맞은 말을 채워 넣으시오. (초성 힌트 참고)}\par\medskip
\begin{kfActTable}{|>{\centering\arraybackslash}m{2.6cm}|m{6cm}|m{4cm}|}
\rowcolor{kfPrimaryLight!50}\textbf{\color{kfPrimary}항목} & \textbf{\color{kfPrimary}내용 (빈칸 채우기)} & \textbf{\color{kfPrimary}답안 작성}\\\hline
갈래 & ㄱㄷ ㅅㅅ / ㅇㅇ ㅅㅅ ① & \kfTBlank \\\hline
충신 & ㅇㅊㄹ ② & \kfTBlank \\\hline
간신 & ㅈㅎㄷ ③ & \kfTBlank \\\hline
핵심 갈등 & ㅊ과 ㅇ의 대립 ④ & \kfTBlank \\\hline
정한담의 악행 & 거짓 ㅅㅅ로 충신 ㅁㅎ ⑤ & \kfTBlank \\\hline
유충렬의 결심 & '이 ㅇㅅ를 갚지 않으면 사람이 아니로다' ⑥ & \kfTBlank \\\hline
천자의 예우 & ㅅㅁ 밖에서 직접 맞이 ⑦ & \kfTBlank \\\hline
정한담의 최후 & ㅊㅎ에 처해짐 ⑧ & \kfTBlank \\\hline
서사 결말 & ㅁㅂㅅ이 환호 ⑨ & \kfTBlank \\\hline
영웅 서사 구조 & 고귀한 출생 → ㄱㄴ → 조력 → ㅌㄷ → 성취 ⑩ & \kfTBlank \\\hline
법의 문제 & 법이 ㄱㅅ에 의해 ㅇㅇ될 수 있음 ⑪ & \kfTBlank \\\hline
주제 & ㄱㅅㅈㅇ — 충신 승리, 간신 몰락 ⑫ & \kfTBlank \\\hline
\end{kfActTable}
\kfSecHeader{kfAreaLiterature}{문제}{}

\begin{multicols}{2}\raggedcolumns\setlength{\columnsep}{12pt}
\begin{kfQBlock}{1}{}
\kfQStem{윗글에 대한 설명으로 적절하지 않은 것은?}
\begin{kfChoices}
\kfChoice 간신 정한담이 충신 유심을 모함하여 옥에 가두게 한다.
\kfChoice 유충렬은 아버지의 충성을 믿으며 원수를 갚겠다고 결심한다.
\kfChoice 천자가 정한담의 죄를 밝히고 처벌함으로써 정의가 실현된다.
\kfChoice 정한담은 자신의 죄를 인정하지 않고 끝까지 결백을 주장한다.
\kfChoice '만백성이 환호하더라'에서 권선징악의 서사적 결말이 확인된다.
\end{kfChoices}
\end{kfQBlock}

\begin{kfQBlock}{2}{}
\kfQStem{윗글에서 정한담의 행위가 '법의 악용'에 해당하는 근거로 가장 적절한 것은?}
\begin{kfChoices}
\kfChoice 정한담이 자신이 직접 군대를 이끌고 반란을 일으켰기 때문이다.
\kfChoice 정한담이 거짓 상소를 올려 법적 처벌 절차를 악용했기 때문이다.
\kfChoice 정한담이 백성들의 재물을 강제로 빼앗는 경제 범죄를 저질렀기 때문이다.
\kfChoice 정한담이 유충렬과 벌인 개인적 싸움에서 일방적으로 패배했기 때문이다.
\kfChoice 정한담이 외국 세력과 내통하여 자신의 나라를 배신했기 때문이다.
\end{kfChoices}
\end{kfQBlock}

\begin{kfQBlock}{3}{}
\kfQStem{윗글에서 천자가 유충렬을 성문 밖에서 직접 맞이한 행위가 지닌 서사적 의미로 가장 적절한 것은?}
\begin{kfChoices}
\kfChoice 천자의 권위가 유충렬에 의해 점차 약화되었음을 보여 준다.
\kfChoice 충신에 대한 군주의 예우와 정의 회복의 확인을 의미한다.
\kfChoice 천자가 자신의 잘못을 인정하고 정한담에게 용서를 구하는 장면이다.
\kfChoice 유충렬이 마침내 천자의 자리를 빼앗으려는 야심을 드러낸 것이다.
\kfChoice 천자가 전쟁의 참혹함을 직접 깨닫고 깊이 반성하는 장면이다.
\end{kfChoices}
\end{kfQBlock}

\begin{kfQBlock}{4}{}
\kfQStem{윗글을 바탕으로 <보기>를 감상한 내용으로 적절하지 않은 것은?

<보기> 「유충렬전」은 충(忠)과 역(逆)의 대립을 서사의 중심에 놓고, 충신의 승리와 간신의 몰락을 통해 권선징악의 주제를 구현한다. 이 과정에서 법과 제도가 간신에 의해 악용되었다가, 군주의 올바른 판단에 의해 회복되는 구조를 보여준다.}
\begin{kfChoices}
\kfChoice 정한담의 모함은 법과 제도가 간신에 의해 악용될 수 있음을 보여주겠군.
\kfChoice 천자가 처음에 정한담의 말을 믿은 것은 법적 판단이 잘못될 수 있음을 드러내겠군.
\kfChoice 유충렬의 귀환과 정한담의 처벌은 충과 역의 대립이 충의 승리로 귀결된 것이겠군.
\kfChoice 만백성의 환호는 권선징악의 서사가 독자(청중)에게 주는 카타르시스를 대변하겠군.
\kfChoice 정한담의 처벌은 법 자체가 불필요하며 군사력만이 정의를 실현할 수 있음을 보여주겠군.
\end{kfChoices}
\end{kfQBlock}

\end{multicols}

\kfSecHeader{kfAreaLiterature}{지문}{골목 안 — 박태원 (현대소설) / 박태원, 「골목 안」(1937)}

\kfPassageNote{골목 안 — 박태원 (현대소설) / 박태원, 「골목 안」(1937)}{%
골목 안에는, 어디서나, 사람들이 살고 있었다.

셋집 주인 허 씨는, 석 달째 밀린 셋돈을 받으러 다시 한 번 돌아왔다. 그러나 안에서는 기척도 없었다. 아마 일을 나간 모양이었다. 아니, 셋돈 받으러 오는 줄 알고 숨은 것인지도 몰랐다.

한편 골목 건너편, 갓난아이를 업은 열세 살쯤 되어 보이는 계집아이가 코를 훌쩍거리며 서 있었다. 아이의 울음을 그치게 하려는 것인지, 아이를 업은 채 제자리에서 뒤뚱뒤뚱 몸을 흔들었다. 어머니는 삯바느질을 하러 나간 것이다.

골목 끝에서는, 실업자 김 군이 벽에 기대어 서서 멍하니 하늘을 올려다보고 있었다. 그는 오늘도 일자리를 구하지 못했다. 누렇게 뜬 얼굴에 어떤 표정도 없었다.
}
\kfSecHeader{kfAreaLiterature}{활동}{작품 정리표 — 골목 안(박태원)}

\noindent\textit{\small 빈칸에 알맞은 말을 채워 넣으시오. (초성 힌트 참고)}\par\medskip
\begin{kfActTable}{|>{\centering\arraybackslash}m{2.6cm}|m{6cm}|m{4cm}|}
\rowcolor{kfPrimaryLight!50}\textbf{\color{kfPrimary}항목} & \textbf{\color{kfPrimary}내용 (빈칸 채우기)} & \textbf{\color{kfPrimary}답안 작성}\\\hline
갈래 & ㅎㄷ ㅅㅅ ① & \kfTBlank \\\hline
작가 & ㅂㅌㅇ ② & \kfTBlank \\\hline
시대 배경 & 1930년대 ㅅㅁㅈ ㄱㅅ ③ & \kfTBlank \\\hline
서술 기법 & ㅋㅁㄹ ㅇㅇ ④ & \kfTBlank \\\hline
구성 방식 & 여러 인물의 일상을 나란히 보여주는 ㅅㅎㅈ 구성 ⑤ & \kfTBlank \\\hline
인물 1 & 셋돈 ㅊㄴ을 받으러 오는 허 씨 ⑥ & \kfTBlank \\\hline
인물 2 & 갓난아이를 업은 ㄱㅈ아이 ⑦ & \kfTBlank \\\hline
인물 3 & 일자리를 구하지 못한 ㅅㅇㅈ 김 군 ⑧ & \kfTBlank \\\hline
인물들의 공통점 & 경제적 ㄱㅍ과 사회적 ㅅㅇ ⑨ & \kfTBlank \\\hline
소외의 원인 & ㅅㅁㅈ ㅈㄷ의 구조적 모순 ⑩ & \kfTBlank \\\hline
문학사적 의의 & ㅁㄷㄴㅈ 소설 실험 ⑪ & \kfTBlank \\\hline
주제 & 식민지 도시 빈민의 ㅈㄷㅈ ㅅㅇ ⑫ & \kfTBlank \\\hline
\end{kfActTable}
\kfSecHeader{kfAreaLiterature}{문제}{}

\begin{multicols}{2}\raggedcolumns\setlength{\columnsep}{12pt}
\begin{kfQBlock}{1}{}
\kfQStem{윗글에 대한 설명으로 적절하지 않은 것은?}
\begin{kfChoices}
\kfChoice 특정 사건의 발단·전개·절정이 아니라 여러 인물의 일상을 나란히 보여주는 구성이다.
\kfChoice 인물의 내면 심리를 깊이 분석하기보다 외부의 행동과 상황을 객관적으로 묘사하고 있다.
\kfChoice 골목 안 인물들의 경제적 궁핍과 사회적 소외가 드러나고 있다.
\kfChoice 작가는 인물들의 고통에 감정적으로 개입하여 직접적인 비판을 가하고 있다.
\kfChoice '어디서나, 사람들이 살고 있었다'는 보편적 존재를 환기하면서도 그들의 고단한 삶을 암시한다.
\end{kfChoices}
\end{kfQBlock}

\begin{kfQBlock}{2}{}
\kfQStem{윗글에서 '누렇게 뜬 얼굴에 어떤 표정도 없었다'가 지닌 문학적 기능으로 가장 적절한 것은?}
\begin{kfChoices}
\kfChoice 어려움 속에서도 잃지 않은 김 군의 낙천적 성격을 보여 준다.
\kfChoice 반복된 실패로 감정이 메마른 인물의 무기력함을 전달한다.
\kfChoice 김 군이 새로운 일을 도모할 계획을 마음에 품고 있음을 암시한다.
\kfChoice 고된 노동에도 불구하고 김 군의 건강 상태가 양호함을 보여 준다.
\kfChoice 골목의 아름다운 풍경과 자연스럽게 조화를 이루는 평화로운 장면이다.
\end{kfChoices}
\end{kfQBlock}

\begin{kfQBlock}{3}{}
\kfQStem{윗글에 등장하는 인물들의 공통점으로 가장 적절한 것은?}
\begin{kfChoices}
\kfChoice 모두 자발적으로 골목에 들어와 공동체를 이루고 있다.
\kfChoice 모두 경제적 궁핍과 사회적 소외라는 공통된 처지에 놓여 있다.
\kfChoice 모두 식민 정부에 적극적으로 저항하는 인물들이다.
\kfChoice 모두 가족 관계가 원만하여 심리적 안정을 누리고 있다.
\kfChoice 모두 교육을 통해 신분 상승을 꾀하는 의지적 인물들이다.
\end{kfChoices}
\end{kfQBlock}

\begin{kfQBlock}{4}{}
\kfQStem{윗글을 바탕으로 <보기>를 감상한 내용으로 적절하지 않은 것은?

<보기> 박태원은 1930년대 경성(서울)을 배경으로 모더니즘 소설을 실험한 작가이다. 그의 작품에서는 작가의 직접적 개입 없이 인물의 외부 행동을 관찰자 시점으로 묘사하는 '카메라 아이(camera eye)' 기법이 자주 사용된다. 이를 통해 식민지 근대 도시의 모습을 담담히 기록한다.}
\begin{kfChoices}
\kfChoice 셋돈을 내지 못하는 세입자의 모습은 식민지 도시 빈민의 경제적 궁핍을 보여주겠군.
\kfChoice 갓난아이를 업은 계집아이의 모습은 어른의 역할을 떠맡아야 하는 아동의 소외를 드러내겠군.
\kfChoice 실업자 김 군의 무표정한 얼굴 묘사는 카메라 아이 기법으로 무기력함을 전달하겠군.
\kfChoice 작가가 인물의 고통을 직접적으로 비판하지 않는 것은 모더니즘 소설의 객관적 서술 태도와 관련되겠군.
\kfChoice 이 작품은 식민지 도시의 활기차고 희망찬 면모를 카메라 아이 기법으로 포착한 것이겠군.
\end{kfChoices}
\end{kfQBlock}

\end{multicols}

\kfSecHeader{kfAreaLiterature}{지문}{저당 잡힌 사내 — 이동하 (현대소설) / 이동하, 「저당 잡힌 사내」 — (가공 의심 본문. 실제 작품 대조 + 본문 교체 필요)}

\kfPassageNote{저당 잡힌 사내 — 이동하 (현대소설) / 이동하, 「저당 잡힌 사내」 — (가공 의심 본문. 실제 작품 대조 + 본문 교체 필요)}{%
사내는 아침마다 출근을 한다. 그것은 그가 살아 있기 때문이 아니라, 갚아야 할 것이 있기 때문이다.

은행 대출, 아파트 할부금, 자동차 할부금, 보험료, 아이들 학원비……. 그는 달마다 들어오는 월급을 받자마자 나누어 보낸다. 월급은 그의 손을 거치지 않고 각처로 흩어진다. 남는 것은 아무것도 없다.

어느 날 그는 문득 생각했다. '나는 무엇인가.' 그리고 답을 찾았다. '나는 저당 잡힌 물건이다.'

그렇다. 그는 저당 잡힌 물건이었다. 은행이 그를 잡고 있었고, 보험 회사가 그를 잡고 있었고, 학원이 그를 잡고 있었다. 그는 갚기 위해 일하고, 일하기 위해 살고, 살기 위해 갚았다. 그 순환에는 시작도 끝도 없었다.

그는 가끔 생각했다. 이 모든 것을 놓아 버리면 어떨까. 그러나 그것은 불가능했다. 저당 잡힌 물건은 스스로 빠져나올 수 없는 것이다.
}
\kfSecHeader{kfAreaLiterature}{활동}{작품 정리표 — 저당 잡힌 사내(이동하)}

\noindent\textit{\small 빈칸에 알맞은 말을 채워 넣으시오. (초성 힌트 참고)}\par\medskip
\begin{kfActTable}{|>{\centering\arraybackslash}m{2.6cm}|m{6cm}|m{4cm}|}
\rowcolor{kfPrimaryLight!50}\textbf{\color{kfPrimary}항목} & \textbf{\color{kfPrimary}내용 (빈칸 채우기)} & \textbf{\color{kfPrimary}답안 작성}\\\hline
갈래 & ㅎㄷ ㅅㅅ(단편) ① & \kfTBlank \\\hline
작가 & ㅇㄷㅎ ② & \kfTBlank \\\hline
제목의 비유 & 경제적 ㄷㅂㅁ처럼 취급당하는 인간 ③ & \kfTBlank \\\hline
출근 동기 & 살기 위해가 아니라 ㄱㅇ야 할 것이 있어서 ④ & \kfTBlank \\\hline
경제적 의무 & ㄷㅊ·ㅎㅂㄱ·ㅂㅎㄹ·ㅎㅇㅂ ⑤ & \kfTBlank \\\hline
월급의 행방 & 손을 거치지 않고 각처로 ㅎㅇ ⑥ & \kfTBlank \\\hline
자기 인식 & '나는 ㅈㄷ 잡힌 ㅁㄱ이다' ⑦ & \kfTBlank \\\hline
악순환 구조 & 갚기→ㅇ→살기→갚기, ㅅㅈ도 ㄲ도 없음 ⑧ & \kfTBlank \\\hline
탈출 불가 & 저당 잡힌 물건은 ㅅㅅ 빠져나올 수 없음 ⑨ & \kfTBlank \\\hline
인간 존재의 물음 & '나는 ㅁㅇ인가' ⑩ & \kfTBlank \\\hline
주제 1 & 경제적 ㅅㅂ에 의한 ㅇㄱ ㅈㅇ 훼손 ⑪ & \kfTBlank \\\hline
주제 2 & 자본주의 체제의 ㄱㅈㅈ ㅂㅈㄹ ⑫ & \kfTBlank \\\hline
\end{kfActTable}
\kfSecHeader{kfAreaLiterature}{문제}{}

\begin{multicols}{2}\raggedcolumns\setlength{\columnsep}{12pt}
\begin{kfQBlock}{1}{}
\kfQStem{윗글에 대한 설명으로 적절하지 않은 것은?}
\begin{kfChoices}
\kfChoice 사내의 출근 동기가 삶의 보람이 아니라 경제적 의무임이 드러난다.
\kfChoice '나는 저당 잡힌 물건이다'에서 인간의 물화(物化) 현상이 표현된다.
\kfChoice 월급이 각처로 흩어지는 구조에서 경제적 속박의 구체적 양상이 나타난다.
\kfChoice 사내는 경제적 속박에서 벗어나기 위한 구체적 계획을 세우고 실행한다.
\kfChoice '시작도 끝도 없는 순환'에서 악순환의 구조가 드러난다.
\end{kfChoices}
\end{kfQBlock}

\begin{kfQBlock}{2}{}
\kfQStem{윗글의 제목 '저당 잡힌 사내'에서 '저당 잡히다'가 지닌 상징적 의미로 가장 적절한 것은?}
\begin{kfChoices}
\kfChoice 경제적 의무에 묶여 자유와 존엄을 잃은 인간의 상태
\kfChoice 법적 분쟁에 휘말려 오랫동안 재판을 받고 있는 상황
\kfChoice 스스로 경제 활동에 자발적으로 참여하는 적극적 태도
\kfChoice 부동산 투자를 통해 큰 부를 축적해 가려는 전략
\kfChoice 도덕적인 타락으로 사회 공동체에서 버림받은 상태
\end{kfChoices}
\end{kfQBlock}

\begin{kfQBlock}{3}{}
\kfQStem{윗글에서 '갚기 위해 일하고, 일하기 위해 살고, 살기 위해 갚았다'라는 구절이 드러내는 바로 가장 적절한 것은?}
\begin{kfChoices}
\kfChoice 사내가 삶의 목표를 향해 단계적으로 나아가는 성장 과정
\kfChoice 시작과 끝이 없는 경제적 악순환 속에 갇힌 인간의 존재 양상
\kfChoice 근면과 성실이 보상받는 건전한 경제 활동의 모습
\kfChoice 가족을 위한 희생이 보람 있는 삶의 원동력이 되는 모습
\kfChoice 사회 구성원으로서 의무를 다하는 시민 의식의 발현
\end{kfChoices}
\end{kfQBlock}

\begin{kfQBlock}{4}{}
\kfQStem{윗글을 바탕으로 <보기>를 감상한 내용으로 적절하지 않은 것은?

<보기> 현대 자본주의 사회에서 인간은 경제적 의무(대출·이자·세금 등)에 묶여 자유로운 삶을 영위하기 어려운 경우가 많다. 이동하의 「저당 잡힌 사내」는 이러한 경제적 속박이 인간의 존엄과 자유를 어떻게 훼손하는지를 '저당'이라는 법적·경제적 비유를 통해 형상화한다.}
\begin{kfChoices}
\kfChoice '살아 있기 때문이 아니라 갚아야 할 것이 있기 때문'이라는 대목은 삶의 목적이 존재 자체가 아니라 경제적 의무로 전도되었음을 보여주겠군.
\kfChoice 월급이 '손을 거치지 않고 각처로 흩어진다'는 것은 노동의 대가가 자신에게 돌아오지 않는 소외를 드러내겠군.
\kfChoice '저당 잡힌 물건은 스스로 빠져나올 수 없다'는 것은 개인의 노력만으로는 극복할 수 없는 구조적 문제를 암시하겠군.
\kfChoice '나는 무엇인가'라는 질문은 경제적 속박 속에서 인간 존재의 의미를 되묻는 실존적 물음이겠군.
\kfChoice 사내의 경제적 속박은 사내 개인의 방탕한 생활 때문이므로, 자본주의 체제와는 무관한 개인적 문제이겠군.
\end{kfChoices}
\end{kfQBlock}

\end{multicols}

\kfStartNonlit{이번 챕터 비문학 지문 5편}


\kfSecHeader{kfAreaNonlit}{지문}{법률행위의 무효와 취소 (사회(법))}

\kfPassageNote{법률행위의 무효와 취소 (사회(법))}{%
법률행위란 당사자의 의사 표시를 핵심 요소로 하여 법적 효과를 발생시키는 행위이다. 그런데 의사 표시에 하자(瑕疵)가 있을 경우, 그 법률행위의 효력이 문제된다. 민법은 의사 표시의 하자 유형으로 착오, 사기, 강박을 규정하고 있으며, 하자의 정도에 따라 법률행위를 '무효'로 하거나 '취소'할 수 있도록 한다.

무효란 법률행위가 처음부터 법적 효력을 갖지 못하는 것을 말한다. 무효인 법률행위는 당사자가 따로 주장하지 않아도 처음부터 효력이 없으며, 누구든지 무효를 주장할 수 있다. 또한 시간이 지나도 유효한 행위로 전환되지 않는다. 무효 사유의 대표적 예로는 반사회적 법률행위(선량한 풍속에 반하는 행위), 불공정한 법률행위(궁박·경솔·무경험을 이용한 폭리 행위) 등이 있다.

취소란 일단 유효하게 성립한 법률행위의 효력을 나중에 소급하여 소멸시키는 것이다. 취소가 이루어지면 법률행위는 처음부터 무효였던 것으로 된다. 이를 '소급 무효(遡及無效)'라 한다. 그러나 취소 전까지는 법률행위가 유효하므로, 취소권자가 취소하지 않으면 법률행위는 그대로 효력을 유지한다. 취소 사유에 해당하는 것으로는 미성년자의 법률행위, 착오에 의한 의사 표시, 사기·강박에 의한 의사 표시 등이 있다.

착오(錯誤)란 의사와 표시가 일치하지 않은 것을 표의자(의사 표시를 한 사람) 스스로 알지 못하는 경우를 말한다. 착오가 '법률행위의 내용의 중요 부분'에 해당하고 표의자에게 중대한 과실이 없을 때에 한하여 취소할 수 있다. 사기(詐欺)란 타인의 기망 행위에 의해 착오에 빠져 의사 표시를 한 경우이며, 강박(强迫)이란 타인의 협박에 의해 공포심을 느끼고 의사 표시를 한 경우이다.

무효와 취소는 모두 법률행위의 효력을 부정하는 제도이지만, 중요한 차이가 있다. 첫째, 주장 권한이 다르다. 무효는 누구나 주장할 수 있지만, 취소는 취소권자(표의자, 법정 대리인 등)만 주장할 수 있다. 둘째, 효력 발생 시점이 다르다. 무효는 처음부터 효력이 없지만, 취소는 취소할 때까지는 유효하다가 취소하면 소급하여 무효가 된다. 셋째, 선의의 제3자 보호 규정이 있다. 사기에 의한 취소는 선의의 제3자에게 대항하지 못한다. 즉, 사기를 당한 사람이 계약을 취소하더라도, 사기 사실을 모르고 거래한 제3자의 권리는 보호된다. 이는 거래의 안전과 신뢰를 보호하기 위한 것이다.
}
\kfSecHeader{kfAreaNonlit}{문제}{}

\begin{multicols}{2}\raggedcolumns\setlength{\columnsep}{12pt}
\begin{kfQBlock}{1}{}
\kfQStem{윗글의 내용과 일치하는 것은?}
\begin{kfChoices}
\kfChoice 무효인 법률행위는 시간이 지나면 유효한 행위로 전환될 수 있다.
\kfChoice 취소는 취소권자가 아닌 누구든지 주장할 수 있다.
\kfChoice 취소된 법률행위는 취소 시점부터 효력을 잃으며 그 이전 효력은 유지된다.
\kfChoice 착오에 의한 의사 표시는 표의자에게 중대한 과실이 있어도 취소할 수 있다.
\kfChoice 사기에 의한 취소는 사기 사실을 모르고 거래한 제3자에게 대항할 수 없다.
\end{kfChoices}
\end{kfQBlock}

\begin{kfQBlock}{2}{}
\kfQStem{윗글의 '무효'와 '취소'를 비교한 것으로 적절하지 않은 것은?}
\begin{kfChoices}
\kfChoice 무효는 처음부터 효력이 없고, 취소는 취소 전까지는 유효하다.
\kfChoice 무효는 누구나 주장 가능하고, 취소는 취소권자만 주장 가능하다.
\kfChoice 무효와 취소 모두 법률행위의 효력을 부정하는 제도이다.
\kfChoice 취소는 소급하여 무효가 되므로 결과적으로 무효와 동일한 법적 상태가 된다.
\kfChoice 무효와 취소 모두 선의의 제3자에 대해 대항할 수 없다는 규정을 갖고 있다.
\end{kfChoices}
\end{kfQBlock}

\begin{kfQBlock}{3}{}
\kfQStem{윗글을 바탕으로 <보기>를 이해한 내용으로 적절한 것은?

<보기> 갑(甲)은 을(乙)의 사기에 속아 자신의 토지를 을에게 매도하였다. 을은 이 토지를 사기 사실을 전혀 모르는 병(丙)에게 다시 매도하였다. 이후 갑은 을의 사기를 알게 되어 매매 계약을 취소하였다.}
\begin{kfChoices}
\kfChoice 갑의 매매 계약은 처음부터 무효이므로 갑은 별도의 취소 절차가 필요 없다.
\kfChoice 갑이 계약을 취소하면 병은 사기 사실을 몰랐더라도 토지를 반환해야 한다.
\kfChoice 갑이 취소하기 전까지 갑-을 간 매매 계약은 유효한 상태였다.
\kfChoice 을이 병에게 토지를 매도한 행위도 사기에 해당하므로 무효이다.
\kfChoice 갑은 취소권자가 아니므로 매매 계약을 취소할 수 없다.
\end{kfChoices}
\end{kfQBlock}

\begin{kfQBlock}{4}{}
\kfQStem{윗글에서 '선의의 제3자 보호' 규정이 존재하는 이유로 가장 적절한 것은?}
\begin{kfChoices}
\kfChoice 사기를 당한 사람의 권리를 최우선으로 보장하기 위해서이다.
\kfChoice 사기를 행한 사람에게 이익을 돌려주기 위해서이다.
\kfChoice 모든 법률행위를 무효로 처리하여 법적 안정성을 확보하기 위해서이다.
\kfChoice 거래의 안전과 신뢰를 보호하기 위해서이다.
\kfChoice 사기 사실을 알고도 거래한 제3자의 권리를 보호하기 위해서이다.
\end{kfChoices}
\end{kfQBlock}

\end{multicols}

\kfSecHeader{kfAreaNonlit}{지문}{법 해석의 본질 (사회(법))}

\kfPassageNote{법 해석의 본질 (사회(법))}{%
법은 추상적인 규범이므로, 구체적 사건에 적용하기 위해서는 해석이 필요하다. 법 해석이란 법 조문의 의미와 적용 범위를 밝히는 작업이다. 법 해석의 방법은 크게 문리 해석, 목적론적 해석, 체계적 해석, 유추 해석, 반대 해석으로 나뉜다.

문리 해석은 법 조문의 문언(文言), 곧 글자 그대로의 의미에 따라 해석하는 방법이다. 이 방법은 법 해석의 출발점이 되며, 법적 안정성과 예측 가능성을 높이는 데 기여한다. 그러나 법 조문의 문언만으로는 입법자의 의도나 사회적 변화를 충분히 반영하기 어려운 경우가 있다.

목적론적 해석은 법 조문이 달성하고자 하는 목적이나 취지를 고려하여 해석하는 방법이다. 법 조문의 글자에 얽매이지 않고, 그 조문이 왜 만들어졌는지, 어떤 가치를 실현하고자 하는지를 탐구한다. 이 방법은 사회적 변화에 유연하게 대응할 수 있다는 장점이 있지만, 해석자의 주관이 개입될 여지가 있다.

체계적 해석은 특정 법 조문을 그 법 조문이 속한 법률 전체 또는 법 체계 안에서의 위치와 관련 조문과의 관계를 고려하여 해석하는 방법이다. 하나의 조문을 고립적으로 읽지 않고, 다른 관련 조문과의 정합성(整合性)을 유지하는 방향으로 해석한다.

유추 해석은 법에 직접적인 규정이 없는 사항에 대해, 유사한 사항을 규정한 법 조문을 적용하는 방법이다. 예를 들어, A에 대해서는 법 규정이 있지만 B에 대해서는 규정이 없을 때, A와 B가 본질적으로 유사하다면 A에 관한 규정을 B에도 적용하는 것이다. 그러나 유추 해석에는 중요한 제한이 있다. 형법에서는 '죄형법정주의(罪刑法定主義)' 원칙에 따라 유추 해석이 금지된다. 죄형법정주의란 범죄와 형벌은 반드시 법률에 명시되어 있어야 한다는 원칙으로, 법에 규정되지 않은 행위를 유추로 처벌하는 것은 국민의 자유와 권리를 침해할 수 있기 때문이다.

반대 해석은 법 조문에 규정된 내용의 반대 상황에는 반대의 결론이 적용된다고 보는 해석 방법이다. 예를 들어, '미성년자는 법정 대리인의 동의가 필요하다'라는 규정에서, 반대 해석을 하면 '성년자는 법정 대리인의 동의가 필요하지 않다'라는 결론을 도출할 수 있다. 이처럼 다양한 법 해석 방법은 구체적 사건의 성격과 관련 법률의 목적에 따라 적절히 선택되어야 하며, 하나의 방법만으로는 올바른 해석에 이르기 어려운 경우가 많다.
}
\kfSecHeader{kfAreaNonlit}{문제}{}

\begin{multicols}{2}\raggedcolumns\setlength{\columnsep}{12pt}
\begin{kfQBlock}{1}{}
\kfQStem{윗글의 내용과 일치하지 않는 것은?}
\begin{kfChoices}
\kfChoice 문리 해석은 법 조문의 글자 그대로의 의미에 따라 해석하는 방법이다.
\kfChoice 목적론적 해석은 법 조문의 입법 목적이나 취지를 고려하는 해석이다.
\kfChoice 체계적 해석은 특정 조문을 관련 조문과의 관계 속에서 해석하는 방법이다.
\kfChoice 유추 해석은 모든 법 분야에서 제한 없이 허용된다.
\kfChoice 반대 해석은 규정된 내용의 반대 상황에 반대의 결론을 적용하는 것이다.
\end{kfChoices}
\end{kfQBlock}

\begin{kfQBlock}{2}{}
\kfQStem{윗글에서 형법에서 유추 해석이 금지되는 이유로 가장 적절한 것은?}
\begin{kfChoices}
\kfChoice 형법 조문이 다른 일반 법률에 비해 더 명확하게 작성되기 때문이다.
\kfChoice 규정에 없는 행위를 유추로 처벌하면 국민의 자유·권리가 침해되기 때문이다.
\kfChoice 형법은 다른 어떤 법률과도 관련 없는 완전히 독립적 영역이기 때문이다.
\kfChoice 유추 해석이 형법에 적용되면 사회적 범죄가 줄어드는 효과가 있기 때문이다.
\kfChoice 형법은 오로지 문리 해석으로만 해석할 수 있도록 규정되어 있기 때문이다.
\end{kfChoices}
\end{kfQBlock}

\begin{kfQBlock}{3}{}
\kfQStem{윗글을 바탕으로 <보기>의 상황에 적용되는 법 해석 방법으로 적절한 것은?

<보기> '공원에서 자동차를 운전하는 것을 금지한다'라는 규정이 있다. 이 규정에 전동 킥보드에 대한 언급은 없다. 그런데 전동 킥보드도 자동차와 마찬가지로 동력으로 움직이며 보행자에게 위험을 줄 수 있으므로, 이 규정을 전동 킥보드에도 적용할 수 있는지가 문제된다.}
\begin{kfChoices}
\kfChoice 문리 해석 — 조문에 '전동 킥보드'가 명시되어 있으므로 적용 가능하다.
\kfChoice 목적론적 해석 — 보행자 안전이라는 입법 목적에 비추어 적용 가능하다고 볼 수 있다.
\kfChoice 반대 해석 — 자동차가 아닌 것은 허용되므로 전동 킥보드는 금지 대상이 아니다.
\kfChoice 체계적 해석 — 전동 킥보드 관련 다른 법 조문과의 정합성에 따라 판단한다.
\kfChoice 유추 해석 — 자동차와 본질적으로 유사하므로 자동차 규정을 유추 적용한다.
\end{kfChoices}
\end{kfQBlock}

\begin{kfQBlock}{4}{}
\kfQStem{윗글에서 '죄형법정주의(罪刑法定主義)'의 의미로 가장 적절한 것은?}
\begin{kfChoices}
\kfChoice 범죄의 경중에 따라 형벌의 종류를 법관이 자유롭게 결정하는 원칙
\kfChoice 범죄와 형벌은 반드시 법률에 명시되어 있어야 한다는 원칙
\kfChoice 모든 범죄는 동일한 형벌로 처벌해야 한다는 원칙
\kfChoice 범죄자의 사회 복귀를 최우선으로 고려해야 한다는 원칙
\kfChoice 범죄 피해자의 의견에 따라 형벌을 결정하는 원칙
\end{kfChoices}
\end{kfQBlock}

\end{multicols}

\kfSecHeader{kfAreaNonlit}{지문}{행정 대집행의 요건과 절차 (사회(법))}

\kfPassageNote{행정 대집행의 요건과 절차 (사회(법))}{%
행정 대집행이란 행정상 의무를 이행하지 않는 자가 있을 때, 행정청이 스스로 또는 제3자로 하여금 그 의무를 대신 이행하게 하고, 그 비용을 의무자에게 징수하는 제도이다. 대집행은 행정의 실효성을 확보하기 위한 강제 수단 중 하나로, 「행정대집행법」에 그 요건과 절차가 규정되어 있다.

대집행이 허용되려면 다음의 요건을 충족해야 한다. 첫째, 의무의 성격이 '대체적 작위의무'이어야 한다. 대체적 작위의무란 의무자 본인이 아니더라도 다른 사람이 대신 이행할 수 있는 의무를 말한다. 예를 들어, 위법 건축물의 철거 의무는 전문 업체가 대신 이행할 수 있으므로 대체적 작위의무에 해당한다. 반면, 출석하여 진술하라는 의무는 본인만 이행할 수 있으므로 비대체적 작위의무이며, 이 경우에는 대집행이 허용되지 않는다. 둘째, 의무자가 의무를 이행하지 않아야 한다. 셋째, 다른 수단으로는 의무 이행을 확보하기 어렵고 그 불이행을 방치하면 공익에 현저히 반하는 경우여야 한다. 이를 '보충성의 원칙'이라 한다.

대집행의 절차는 네 단계로 진행된다. 제1단계는 '계고(戒告)'이다. 행정청은 의무자에게 일정한 기한을 정하여 의무를 이행할 것을 문서로 명하고, 이행하지 않으면 대집행을 하겠다는 뜻을 미리 통지한다. 제2단계는 '대집행 영장의 통지'이다. 계고에도 불구하고 의무자가 이행하지 않으면, 행정청은 대집행의 시기, 책임자, 비용 견적 등을 기재한 대집행 영장을 의무자에게 발부한다. 제3단계는 '대집행의 실행'이다. 통지된 시기에 행정청 또는 제3자가 의무를 대신 이행한다. 제4단계는 '비용 징수'이다. 대집행에 든 비용은 의무자로부터 국세 체납 처분의 예에 따라 징수한다.

대집행에 대해 의무자가 불복하고자 할 경우, 행정심판이나 행정소송을 통해 권리를 구제받을 수 있다. 계고 처분, 대집행 영장 통지, 대집행 실행 등 각 단계의 행위가 각각 독립된 처분으로서 행정 쟁송의 대상이 될 수 있다. 또한, 대집행이 위법하게 실행되어 의무자에게 손해가 발생한 경우에는 국가배상을 청구할 수도 있다.

이처럼 대집행은 행정 목적을 달성하면서도 국민의 재산권과 자유를 침해할 수 있는 강제 수단이므로, 엄격한 요건과 절차를 거쳐야 하고 사후적 권리 구제도 보장되어야 한다. 대집행 제도는 행정의 실효성 확보와 국민의 권리 보호라는 두 가치를 조화시키는 데 그 의의가 있다.
}
\kfSecHeader{kfAreaNonlit}{문제}{}

\begin{multicols}{2}\raggedcolumns\setlength{\columnsep}{12pt}
\begin{kfQBlock}{1}{}
\kfQStem{윗글의 내용과 일치하는 것은?}
\begin{kfChoices}
\kfChoice 비대체적 작위의무에 대해서도 행정 대집행이 허용된다.
\kfChoice 대집행에 드는 모든 비용은 행정청이 전액 부담하게 된다.
\kfChoice 대집행은 계고 절차 없이도 곧바로 실행할 수 있는 제도이다.
\kfChoice 의무자는 대집행에 대해 행정심판이나 행정소송으로 불복할 수 있다.
\kfChoice 대집행 영장에는 비용 견적이나 실행 계획을 기재하지 않아도 된다.
\end{kfChoices}
\end{kfQBlock}

\begin{kfQBlock}{2}{}
\kfQStem{윗글에서 '보충성의 원칙'의 의미로 가장 적절한 것은?}
\begin{kfChoices}
\kfChoice 대집행은 다른 수단으로 의무 이행 확보가 어려울 때만 허용된다는 원칙
\kfChoice 대집행에 든 모든 비용을 의무자에게 분할 납부시키도록 하는 원칙
\kfChoice 대집행의 절차를 보충 서류가 없이도 곧바로 진행할 수 있다는 원칙
\kfChoice 의무자가 대집행에 든 비용 가운데 일부를 보충해 납부해야 한다는 원칙
\kfChoice 행정청이 대집행을 시작하기 전에 의무자를 보충 교육해야 한다는 원칙
\end{kfChoices}
\end{kfQBlock}

\begin{kfQBlock}{3}{}
\kfQStem{윗글을 바탕으로 <보기>를 이해한 내용으로 적절하지 않은 것은?

<보기> 갑(甲)은 무허가 건축물을 지어 행정청으로부터 철거 명령을 받았으나 이행하지 않았다. 행정청은 대집행 절차를 밟기로 결정했다.}
\begin{kfChoices}
\kfChoice 갑의 철거 의무는 다른 사람이 대신 이행할 수 있으므로 대체적 작위의무에 해당한다.
\kfChoice 행정청은 먼저 갑에게 일정 기한을 정하여 철거를 명하는 계고를 해야 한다.
\kfChoice 계고 후에도 갑이 이행하지 않으면 행정청은 대집행 영장을 발부해야 한다.
\kfChoice 대집행에 든 비용은 행정청이 부담하고 갑에게 청구할 수 없다.
\kfChoice 갑은 대집행에 불복하여 행정심판이나 행정소송을 제기할 수 있다.
\end{kfChoices}
\end{kfQBlock}

\begin{kfQBlock}{4}{}
\kfQStem{윗글에서 대집행 절차에 엄격한 요건과 절차가 요구되는 이유로 가장 적절한 것은?}
\begin{kfChoices}
\kfChoice 행정청의 업무량을 줄이기 위해서이다.
\kfChoice 대집행이 국민의 재산권과 자유를 침해할 수 있는 강제 수단이기 때문이다.
\kfChoice 의무자가 대집행 비용을 미리 납부할 시간을 주기 위해서이다.
\kfChoice 행정청이 대집행 실행 인력을 확보하는 데 시간이 필요하기 때문이다.
\kfChoice 의무자에게 재심 기회를 부여하기 위해서이다.
\end{kfChoices}
\end{kfQBlock}

\end{multicols}

\kfSecHeader{kfAreaNonlit}{지문}{특허권의 성립과 효력 (사회(법))}

\kfPassageNote{특허권의 성립과 효력 (사회(법))}{%
특허권이란 새로운 발명을 한 자에게 국가가 일정 기간 독점적으로 그 발명을 실시할 수 있는 권리를 부여하는 제도이다. 특허 제도의 목적은 발명자에게 독점적 이익을 보장함으로써 기술 혁신을 장려하고, 동시에 발명의 내용을 공개하게 하여 사회 전체의 기술 수준을 높이는 데 있다.

특허를 받으려면 발명이 세 가지 요건을 갖추어야 한다. 첫째, 신규성(新規性)이다. 출원 전에 이미 알려진 기술과 동일한 발명은 특허를 받을 수 없다. 이미 공개된 기술을 독점시키면 오히려 기술 발전을 저해할 수 있기 때문이다. 둘째, 진보성(進步性)이다. 기존 기술로부터 쉽게 생각해 낼 수 있는 수준의 발명은 특허를 받을 수 없다. 해당 기술 분야에서 통상의 지식을 가진 사람이 쉽게 발명할 수 있는 것이라면, 그 발명에 독점권을 줄 필요가 없기 때문이다. 셋째, 산업이용가능성이다. 산업에 실제로 이용할 수 있는 발명이어야 하며, 단순한 학문적 발견이나 자연법칙 자체는 특허의 대상이 되지 않는다.

특허 출원은 발명의 내용을 상세히 기술한 명세서와 청구범위를 특허청에 제출하는 것으로 시작된다. 특허청은 출원된 발명을 심사하여 요건을 갖추었다고 판단하면 특허 등록을 결정한다. 특허권은 등록에 의해 발생하며, 출원일로부터 20년간 존속한다. 존속 기간이 만료되면 해당 발명은 누구나 자유롭게 실시할 수 있는 공유 영역으로 들어간다.

특허권자는 타인이 정당한 권원 없이 자신의 특허 발명을 실시하는 것을 금지할 수 있는 독점적 권리를 가진다. 타인이 특허권을 침해한 경우, 특허권자는 침해 행위의 금지를 청구하는 '금지 청구권'과 침해로 인한 손해의 배상을 청구하는 '손해배상 청구권'을 행사할 수 있다. 또한, 침해 행위에 사용된 물건의 폐기나 침해 행위에 제공된 설비의 제거를 청구할 수도 있다.

그러나 특허권에도 제한이 있다. 연구 또는 시험을 위한 특허 발명의 실시, 특허 출원 전부터 선의(善意)로 동일한 발명을 실시하고 있던 자의 계속 실시(선사용권) 등은 특허권 침해에 해당하지 않는다. 이러한 제한은 기술 발전을 위한 연구 활동의 자유와, 선의의 기존 실시자의 기대 이익을 보호하기 위한 것이다. 특허 제도는 이처럼 발명자의 독점적 이익과 사회 전체의 기술 발전이라는 두 가치의 균형 위에 성립한다.
}
\kfSecHeader{kfAreaNonlit}{문제}{}

\begin{multicols}{2}\raggedcolumns\setlength{\columnsep}{12pt}
\begin{kfQBlock}{1}{}
\kfQStem{윗글의 내용과 일치하지 않는 것은?}
\begin{kfChoices}
\kfChoice 특허권은 출원일로부터 20년간 존속하며, 기간 만료 후에는 공유 영역이 된다.
\kfChoice 자연법칙 자체는 특허의 대상이 되지 않는다.
\kfChoice 연구 또는 시험 목적의 특허 발명 실시는 침해에 해당하지 않는다.
\kfChoice 특허 출원 전부터 동일 발명을 실시하던 자는 항상 침해로 처벌된다.
\kfChoice 특허를 받으려면 신규성, 진보성, 산업이용가능성을 갖추어야 한다.
\end{kfChoices}
\end{kfQBlock}

\begin{kfQBlock}{2}{}
\kfQStem{윗글에서 진보성(進步性) 요건이 필요한 이유로 가장 적절한 것은?}
\begin{kfChoices}
\kfChoice 이미 일반에 공개된 기술의 독점을 방지하기 위해서이다.
\kfChoice 기존 기술에서 쉽게 도출되는 발명에 독점권을 줄 필요가 없기 때문이다.
\kfChoice 특허 출원 시 작성해야 할 서류의 부담을 줄이기 위해서이다.
\kfChoice 현재의 특허 존속 기간을 더 단축시키기 위한 정책적 이유에서이다.
\kfChoice 발명이 실제 산업에 이용 가능한지 여부를 판단하기 위한 기준이다.
\end{kfChoices}
\end{kfQBlock}

\begin{kfQBlock}{3}{}
\kfQStem{윗글을 바탕으로 <보기>를 이해한 내용으로 적절한 것은?

<보기> 갑(甲)은 새로운 배터리 기술을 발명하여 특허 등록을 마쳤다. 을(乙)은 갑의 허락 없이 동일한 기술을 사용하여 배터리를 제조·판매하고 있다. 한편 병(丙)은 갑의 특허 기술을 연구 목적으로 실험실에서 시험하고 있다.}
\begin{kfChoices}
\kfChoice 을의 행위는 선사용권에 해당하므로 침해가 아니다.
\kfChoice 병의 행위는 정당한 권원 없이 실시한 것이므로 침해에 해당한다.
\kfChoice 갑은 을에게 금지 청구와 손해배상 청구를 할 수 있다.
\kfChoice 갑의 특허권은 등록 시점부터 25년간 존속한다.
\kfChoice 을이 갑보다 먼저 해당 기술을 개발했다면 갑의 특허는 무효가 된다.
\end{kfChoices}
\end{kfQBlock}

\begin{kfQBlock}{4}{}
\kfQStem{윗글에서 특허 제도가 발명의 내용을 '공개'하게 하는 이유로 가장 적절한 것은?}
\begin{kfChoices}
\kfChoice 다른 발명자가 동일한 발명을 하지 못하도록 경고하기 위해서이다.
\kfChoice 사회 전체의 기술 수준을 높이기 위해서이다.
\kfChoice 특허 심사 비용을 줄이기 위해서이다.
\kfChoice 발명자의 명예를 사회적으로 인정하기 위해서이다.
\kfChoice 특허 존속 기간을 연장하기 위한 조건이기 때문이다.
\end{kfChoices}
\end{kfQBlock}

\end{multicols}

\kfSecHeader{kfAreaNonlit}{지문}{법적 의제(擬制)의 개념 (사회(법))}

\kfPassageNote{법적 의제(擬制)의 개념 (사회(법))}{%
법에서 '의제(擬制)'란 실제로는 성질이 다른 것을 법률상 같은 것으로 취급하는 법기술적 장치이다. 의제는 '\textasciitilde{}으로 본다'라는 표현으로 나타나며, 사실과 다르더라도 법률이 그렇게 규정한 이상 반대의 증거를 제시하여 번복할 수 없다. 이와 대비되는 개념이 '추정(推定)'이다. 추정은 '\textasciitilde{}으로 추정한다'라는 표현을 쓰며, 반대 증거(반증)가 제시되면 그 효과가 뒤집힐 수 있다.

의제와 추정의 가장 핵심적인 차이는 '반증 가능성'에 있다. 의제는 반증을 허용하지 않는다. 법이 A를 B로 본다고 규정하면, 실제로 A가 B가 아닐지라도 법적으로는 B로 확정된다. 반면, 추정은 반증을 허용한다. 법이 A를 B로 추정한다고 규정하더라도, A가 실제로 B가 아님을 증명하면 그 추정은 깨진다. 추정의 경우 반증의 부담은 추정을 번복하려는 측에 있다.

의제의 대표적 사례로는 '사망 의제'가 있다. 민법은 전쟁이나 재해 등으로 생사가 불분명한 사람에 대해 법원이 실종 선고를 하면, 실종 기간이 만료된 때에 사망한 것으로 본다고 규정한다. 이때 '본다'는 의제의 표현이므로, 실종자가 실제로는 살아 있다 하더라도 법적으로는 사망한 것으로 취급된다. 물론 실종자가 나중에 나타나면 실종 선고 취소를 청구할 수 있지만, 이는 의제 자체를 반증으로 뒤집는 것이 아니라 별도의 법적 절차(실종 선고 취소 심판)를 통해 이루어지는 것이다.

추정의 대표적 사례로는 '친생 추정'이 있다. 민법은 혼인 중에 출생한 자녀는 부(夫)의 자녀로 추정한다고 규정한다. 이는 추정이므로, 부(夫)가 유전자 검사 등 반증을 제출하여 혈연관계가 없음을 증명하면 이 추정을 번복할 수 있다.

법이 의제를 사용하는 이유는 법적 안정성의 확보에 있다. 만약 의제 사항에 대해서도 언제든 반증을 통해 번복이 가능하다면, 이미 형성된 법률관계가 불안정해질 수 있다. 예를 들어, 실종 선고로 사망 의제가 이루어지면 상속이 개시되고 혼인이 해소되는 등 다양한 법률관계가 형성되는데, 이러한 법률관계가 반증만으로 무너진다면 관련 당사자들의 법적 지위가 극도로 불안해진다. 따라서 의제는 법적 확실성을 높이기 위해 반증 가능성을 차단하는 제도적 장치인 것이다. 이처럼 의제와 추정은 모두 사실 관계에 대한 법률적 판단이라는 점에서 공통되지만, 반증의 허용 여부라는 결정적 차이로 인해 법적 효과가 크게 달라진다.
}
\kfSecHeader{kfAreaNonlit}{문제}{}

\begin{multicols}{2}\raggedcolumns\setlength{\columnsep}{12pt}
\begin{kfQBlock}{1}{}
\kfQStem{윗글의 내용과 일치하는 것은?}
\begin{kfChoices}
\kfChoice 의제는 반대 증거를 제시하면 번복할 수 있다.
\kfChoice 추정은 반대 증거가 제시되어도 효과가 유지된다.
\kfChoice 의제와 추정 모두 반증 가능성이 동일하다.
\kfChoice 사망 의제에서 실종자가 나타나면 반증으로 의제가 자동 소멸한다.
\kfChoice 친생 추정은 유전자 검사 등 반증을 제출하면 번복할 수 있다.
\end{kfChoices}
\end{kfQBlock}

\begin{kfQBlock}{2}{}
\kfQStem{윗글의 '의제'와 '추정'을 비교한 것으로 적절한 것은?}
\begin{kfChoices}
\kfChoice 의제는 '\textasciitilde{}으로 추정한다'는 표현을, 추정은 '\textasciitilde{}으로 본다'는 표현을 사용하는 차이를 보인다.
\kfChoice 의제는 반증이 가능하고, 추정은 반증이 불가능하다.
\kfChoice 의제는 반증을 허용하지 않아 법적 확실성이 높고, 추정은 반증이 가능해 유연성이 크다.
\kfChoice 의제와 추정 모두 사실 관계와 무관한 순수한 법 창조 행위이다.
\kfChoice 의제는 민법에만, 추정은 형법에만 적용되는 개념이다.
\end{kfChoices}
\end{kfQBlock}

\begin{kfQBlock}{3}{}
\kfQStem{윗글을 바탕으로 <보기>를 이해한 내용으로 적절하지 않은 것은?

<보기> 갑(甲)은 10년 전 해난 사고로 실종되어 법원의 실종 선고를 받았다. 갑의 재산은 상속인 을(乙)에게 상속되었다. 그런데 최근 갑이 살아 있는 것이 확인되었다.}
\begin{kfChoices}
\kfChoice 갑에 대한 사망 의제는 '\textasciitilde{}으로 본다'는 규정에 따른 것이다.
\kfChoice 갑이 살아 있다는 사실만으로 사망 의제가 자동으로 소멸하지는 않는다.
\kfChoice 갑이 나타나면 별도의 법적 절차(실종 선고 취소 심판)를 통해 사망 의제를 해소해야 한다.
\kfChoice 을이 상속받은 것은 사망 의제에 따라 적법하게 형성된 법률관계이다.
\kfChoice 갑이 살아 있다는 반증을 제출하면 의제가 즉시 번복되어 상속이 자동 취소된다.
\end{kfChoices}
\end{kfQBlock}

\begin{kfQBlock}{4}{}
\kfQStem{윗글에서 법이 의제를 사용하여 반증 가능성을 차단하는 이유로 가장 적절한 것은?}
\begin{kfChoices}
\kfChoice 법원의 복잡한 심사 업무를 줄여 행정적 효율성을 높이기 위해서이다.
\kfChoice 이미 형성된 법률관계의 불안정을 막아 법적 안정성을 확보하기 위해서이다.
\kfChoice 객관적 사실과 다른 결론을 내리는 것 자체가 법의 본래 목적이기 때문이다.
\kfChoice 반증 자료를 직접 제출할 능력이 부족한 사회적 약자를 보호하기 위해서이다.
\kfChoice 기존의 추정 제도를 새로운 의제 제도로 단계적으로 대체하기 위해서이다.
\end{kfChoices}
\end{kfQBlock}

\end{multicols}

\kfStartWeekly{패턴 워크북 — 총 ?문항}


\kfSecHeader{kfAreaWeekly}{지문}{문항 1 지문 / 섞어치기}

\kfPassageNote{문항 1 지문 / 섞어치기}{%
컴퓨터는 사람의 언어를 직접 이해할 수 없다. 따라서 컴퓨터가 사람의 언어를 이해할 수 있도록 표현해 주어야 한다. 컴퓨터가 사람의 언어를 이해할 수 있도록 표현하는 방법에는 여러 가지가 있을 수 있는데, 그중 ㉠원-핫 인코딩은 표현하고 싶은 단어의 색인 값에 1을 부여하고 나머지에는 0을 부여하는 방법이다. 원-핫 인코딩을 수행하기 위해서는 먼저 대상이 되는 텍스트의 단어 집합을 만든다. 예를 들어 ‘오늘 날씨 정말 좋다.’라는 한 문장으로 이루어진 텍스트를 대상으로 한다고 가정하면, 단어 집합은 \{오늘, 날씨, 정말, 좋다\}가 되고 집합의 원소가 되는 단어의 개수가 4개이므로 단어 집합의 크기는 4가 된다. 단어 집합의 크기에 따라 각각의 단어에 [1, 0, 0, 0], [0, 1, 0, 0], [0, 0, 1, 0], [0, 0, 0, 1]과 같이 4차원의 고유한 숫자를 부여하는 것이 원-핫 인코딩이다. 원-핫 인코딩은 직관적이고 단순하지만 단어들이 개별적, 독립적으로 표현되기

(중략)
}
\begin{kfQBlock}{1}{}
\kfQStem{문맥상 ⓐ와 의미가 가장 가까운 것은?}
\begin{kfChoices}
\kfChoice 흐린 날보다 맑은 날에 사진이 잘 나온다.
\kfChoice 나는 웃음이 나오면 참을 수가 없다.
\kfChoice 이 상품은 시장에 나온 후 큰 인기를 끌었다.
\end{kfChoices}
\end{kfQBlock}

\kfSecHeader{kfAreaWeekly}{지문}{문항 2 지문 / 부정상대어}

\kfPassageNote{문항 2 지문 / 부정상대어}{%
에너지 대사는 생명 유지를 위해 필요한 에너지를 얻고 사용하는 과정으로 인체에 필요한 에너지는 대부분 음식물 섭취를 통해 이루어진다. 섭취를 통해 체내에 ⓐ유입된 음식물은 그대로 에너지로 소비되지 않고 소화 과정을 거쳐 체내 각 조직에 공급된다. 음식물에 함유된 탄수화물, 단백질, 지방은 각각 일련의 과정을 거쳐 세포 활동에 직접 사용되는 에너지원인 ATP로 전환된 후 필요한 조직에 사용된다. 에너지로 바로 사용되지 않은 탄수화물 일부는 포도당의 형태로 간이나 근육에 저장되었다가 필요할 때 다시 ATP로 전환된다. 섭취된 에너지와 인체가 소비하는 에너지는 균형을 유지해야 하지만, 섭취한 에너지의 양이 소비되는 양보다 많아지는 상태가 지속되면 남는 에너지는 지방산으로 전환되어 지방 세포에 저장된다. 이러한 상태가 ⓑ심화되면 지방 세포의 크기가 커지게 된다. 지방산이 과도하게 축적되면 체내의 에너지 조절 체계가 정상적으로 작동하지 않게 되고 비만을 비롯한 다양한 대사 질환의 발생 가능성

(중략)
}
\begin{kfQBlock}{2}{}
\kfQStem{윗글을 읽고 알 수 있는 내용으로 적절하지 않은 것은?}
\begin{kfChoices}
\kfChoice 에너지 과잉이 지속될수록 지방산 저장은 줄고 대사 질환 위험은 낮아진다.
\kfChoice 기초 대사량이 낮으면 같은 섭취에서도 에너지 과잉 위험이 커질 수 있다.
\kfChoice 갈색 지방 활성 저하는 소비 감소로 이어질 수 있다.
\end{kfChoices}
\end{kfQBlock}

\kfSecHeader{kfAreaWeekly}{지문}{문항 3 지문 / 주체객체}

\kfPassageNote{문항 3 지문 / 주체객체}{%
(가) 유추는 통상적인 법의 해석 및 적용 과정에서 이른바 법의 흠결을 직면하게 된 해석자가 취할 수 있는 보조 수단의 하나로 여겨져 왔다. 본래 법의 해석은 일반적이고 추상적인 규범으로부터 개별적이고 구체적인 사례에 대한 처리 방안을 도출해 내는 연역적 절차로 이해되었으며, 유추는 특정 사례에 직접적으로 적용될 규범이 존재하지 않을 때 그 사례와 유사한 다른 사례에 적용되는 규범을 써서 문제를 해결하는 일종의 흠결 보충의 논리였던 것이다.

(중략)

그러나 이 같은 전통적인 이해 방식은 법 해석의 본질을 제대로 파악하지 못하고 있는 것으로 보인다. 우선 개별적이고 구체적인 사례에 대한 처리 방안으로서 법이 순수한 연역적 절차를 거쳐 발견되는 것은 아니다. 흔히 흠결 보충이 필요 없어 보이는 사례에 대해서도 그 사실 관계를 규범에 맞춰 보고, 다시 규범을 사실 관계에 맞춰 봄으로써 규범이 정제되는 과정을 필요로 하기 때문이다. 애초에 일반적이고 추상적인 형태로 주어진 규범은 다양한 사례들과 그 

(중략)
}
\begin{kfQBlock}{3}{}
\kfQStem{(가)의 유추 구조를 이해한 내용으로 가장 적절한 것은?}
\begin{kfChoices}
\kfChoice 유추는 비교점 없는 직관에 의존할수록 정확해진다.
\kfChoice 유추는 개별 사례에서 출발해 제3의 비교점을 담은 규범을 거쳐 결론에 이른다.
\end{kfChoices}
\end{kfQBlock}

\kfSecHeader{kfAreaWeekly}{지문}{문항 4 지문 / 순서방향}

\kfPassageNote{문항 4 지문 / 순서방향}{%
법은 사회 구성원 간의 갈등을 조정하고 질서를 유지하기 위한 제도적 장치로서, 주로 국가에 의해 제정되고 시행되는 ⓐ실정법을 중심으로 구성된다. 실정법은 입법 기관이 사회의 구체적 상황에 맞추어 만든 규범으로, 강제력과 명확한 절차를 통해 사회 질서를 실현한다. 그러나 실정법만으로는 인간 사회의 근본적 정의와 도덕적 정당성을 충분히 설명하기 어렵다. 이러한 한계를 보완하는 개념이 바로 인간의 본성과 이성에 근거한 보편적 규범인 ‘자연법’이다. ⓑ자연법은 시대나 사회에 따라 변하지 않는 도덕적 원리로, 인간이라면 누구나 인식할 수 있는 옳고 그름의 기준으로 여겨진다. 따라서 자연법은 법 제도의 정당성을 판단하는 궁극적 기준으로 기능하며, 실정법이 정당성을 가지는 것 역시 자연법의 원리에 부합할 때 비로소 가능한 것이다.

(중략)

자연법의 철학적 기초는 고대 그리스 스토아학파에 의해 체계화되었으나, 그 이전에도 헤라클레이토스나 소피스트, 아리스토텔레스 등은 인간과 자연의 보편적 질서 속에서 법의 

(중략)
}
\begin{kfQBlock}{4}{}
\kfQStem{윗글의 주제로 가장 적절한 것은?}
\begin{kfChoices}
\kfChoice 자연법과 실정법의 관계 및 자연법 사상의 전개와 현대적 재구성
\kfChoice 실정법의 제정 절차와 처벌 규정의 중요성
\kfChoice 스토아학파 윤리와 로마법의 조문 분석
\end{kfChoices}
\end{kfQBlock}

\kfSecHeader{kfAreaWeekly}{지문}{문항 5 지문 / 인과도치}

\kfPassageNote{문항 5 지문 / 인과도치}{%
리더는 목표 달성을 위해 구성원들을 이끌고 그들에게 영향력을 미치는 조직의 핵심 인물이다. 오래전부터 리더로서의 능력을 나타내는 리더십은 무엇이며, 어떠한 리더십이 효과적인지에 관한 연구가 행해졌다. 피들러는 리더를 업무의 목표 달성을 중시하는 업무 지향적 스타일을 가진 리더와 구성원들과의 관계, 소통 등을 중시하는 관계 지향적 스타일을 가진 리더로 구분하였다. 그리고 리더가 자신의 집단에 영향을 미칠 수 있는 정도를 나타내는 상황의 호의성이 ‘리더와 구성원들의 관계, 업무 구조, 리더의 권한’이라는 세 가지 변수로 결정된다고 보았다. 즉 리더와 구성원들이 서로 믿고 존중하는 관계일수록, 업무가 명확하게 구조화되어 있을수록, 채용 및 승진 등에 대한 리더의 공식적 권한이 클수록 상황의 호의성이 높고, 반대의 경우는 호의성이 낮다고 본 것이다. 그는 상황의 호의성에 따라 적합한 리더십의 종류가 다르다고 보았는데, 상황의 호의성이 매우 낮거나 매우 높을 때는 업무 지향적 리더가, 중간 

(중략)
}
\begin{kfQBlock}{5}{}
\kfQStem{윗글을 바탕으로 <보기>를 이해한 반응으로 적절하지 않은 것은?}
\begin{kfEvidence}<보기>\textbackslash{}n갑이 경영하고 있는 기업 K는 얼마 전 회사에 오랫동안 근무하던 구성원들이 대거 퇴직하여 대부분이 입사한 지 오래되지 않은 이들로 구성되어 있다. 이들은 업무에 임하는 태도는 매우 적극적이지만, 아직 업무에 숙달되지 않은 상태이다. 한편 을이 경영하고 있는 기업 P는 업무에 숙달된 구성원들이 매우 많지만 전반적으로 구성원들이 자신감이 부족하고 근로 의욕이 저하된 상태이다.\textbackslash{}n갑은 각 구성원이 수행해야 할 업무와 성과에 따른 보상과 처벌을 명확히 정하여 제시하였다. 을은 구성원에게 기업이 추구하는 목표를 제시하고 구성원들이 회사의 의사 결정에 적극적으로 참여할 수 있게 하였다. 그리고 목표를 달성한 구성원에게 추가 수당이라는 보상을 제공하였다. 기업 K의 구성원들은 대부분 목표를 달성하지 못했고, 기업 P의 구성원들은 대부분 목표를 쉽게 달성하였다.\end{kfEvidence}
\begin{kfChoices}
\kfChoice 로크는 을이 구성원들의 동기를 자극하기 위해서는 현재보다 더 어려운 목표를 제시해야 한다고 생각하겠군.
\kfChoice 데시와 라이언은 기업 P의 구성원들이 목표를 달성한 이유를 을이 구성원에게 제공한 보상 때문이라고 여기겠군.
\end{kfChoices}
\end{kfQBlock}

\kfSecHeader{kfAreaWeekly}{지문}{문항 6 지문 / 의도/목적}

\kfPassageNote{문항 6 지문 / 의도/목적}{%
(가) 타인에 대한 힘과 영향력의 근원으로서 사회 활동에서 다양한 이윤을 확보할 수 있는 자원이 되는 ‘자본’은 여러 형태로 존재한다. \uline{㉠부르디외}는 세계화를 기득권 세력이 지배력을 유지 및 확장하기 위해 수행하는 전략이라고 보고, 세계화로 인해 자본을 바탕으로 한 세력 간의 위계가 강화됨으로써 불평등 구조가 심화된다고 보았다. 그는 자본을 돈 또는 돈으로 즉각 전환 가능한 경제 자본, 학연·지연·혈연 등과 같은 사회적 관계망인 사회 자본, 개인이 보유한 문화적 지식, 기술, 교육 및 문화적 자산인 문화 자본, 그리고 상징 자본으로 나누었다. 상징 자본은 소유한 자본이 세간으로부터 인정되었을 때 발생하는 자본으로서, 명예, 권위 등과 같이 보이지 않는 상징적 가치가 자본으로서의 힘을 발휘하는 것이다. 부르디외는 상징 자본이 국가 간의 관계에도 작용하며, 그 결과 불평등한 국제 질서가 형성될 수 있다고 보았다.

부르디외는 상징 폭력과 아비투스를 바탕으로 자본과 세계화에 관…(중략)
}
\begin{kfQBlock}{6}{}
\kfQStem{㉠과 ㉡의 대응 전략을 비교한 설명으로 적절하지 않은 것은?}
\begin{kfChoices}
\kfChoice ㉡은 삶의 정치보다 해방의 정치만을 목표로 삼는다.
\kfChoice ㉡은 삶의 정치보다 해방의 정치만을 목표로 삼는다.
\kfChoice ㉠은 시민 교육과 국제주의 연대를 통해 대응 전략을 제시한다.
\kfChoice ㉡은 제3의 길과 사회 투자 국가를 통해 위험 대응을 모색한다.
\end{kfChoices}
\end{kfQBlock}

\kfSecHeader{kfAreaWeekly}{지문}{문항 7 지문 / 상관관계 반전}

\kfPassageNote{문항 7 지문 / 상관관계 반전}{%
전류란 \uline{전하를 띤 입자들이 일정한 방향으로 이동하는 현상}이다. 전류의 세기는 일정한 시간 동안 도선이나 회로의 단면을 통과한 전하의 양으로 측정되며, 단위는 \uline{암페어(A)}이다. 그런데 도선이나 회로 안에서 전류가 흐를 때 \uline{전자와 원자 간의 충돌}로 열이 발생한다. 전자는 도선이나 회로 안에서 자유롭게 이동하지만, 완전히 자유롭게는 흐르지 않는다. 도선이나 회로 안의 원자들과 충돌을 반복하면서 전자의 운동 에너지는 \uline{열에너지}로 변환되고 전류의 세기가 커지면 열에너지가 더욱 많이 발생한다. 이때 같은 전류라도 도선의 단면적이 좁으면 단위 면적을 통과하는 전류의 양인 \uline{전류 밀도}는 높아지고, 넓으면 전류 밀도는 낮아진다. 따라서 전류 밀도가 높을수록 단위 부피당 열 발생량이 증가한다. 이러한 열 발생은 불필요한 전력 손실을 야기할 뿐만 아니라, 과도한 열로 인해 기기의 성능이 저하되거나 회로가 손상될 수 있으므로 설계 시 이에 대한 …(중략)
}
\begin{kfQBlock}{7}{}
\kfQStem{표피 효과가 도선의 발열을 증가시키는 이유로 가장 적절하지 않은 것은?}
\begin{kfChoices}
\kfChoice 표피 효과는 전류를 표면층으로 몰아 유효 단면적을 줄이고 저항을 키울 수 있기 때문에
\kfChoice 표피 효과는 전류를 표면층으로 몰아 유효 단면적을 줄이고 저항을 키울 수 있기 때문에
\kfChoice 표피 효과는 전류가 흐르는 유효 단면적을 늘려 저항을 낮추기 때문에
\kfChoice 표피 효과는 전류의 방향을 직류로 바꾸어 열 발생을 멈추게 하기 때문에
\end{kfChoices}
\end{kfQBlock}

\kfSecHeader{kfAreaWeekly}{지문}{문항 8 지문 / 필요조건}

\kfPassageNote{문항 8 지문 / 필요조건}{%
(가) 역사적으로 예술 형식을 구현하는 수단인 매체의 자리를 두고 경쟁한 문자와 이미지의 관계는 매체 연구의 주된 주제였다. 레싱은 시와 회화의 관계를 논하면서 당시 지배적이던 호라티우스의 경구 '시는 그림과 같이'를 대체할 새로운 관점을 제시했다. 호라티우스의 경구는 시와 회화의 공통 미메시스를 설명하는 말로 읽히기도 하고, 시가 회화를 모범으로 삼아야 한다는 말로 해석되기도 했다. 레싱은 회화의 원칙을 시에 확장 적용해 두 예술 형식의 경계를 모호하게 만드는 고전적 태도를 비판했다. 그는 각 예술 형식에서 매체와 매체 효과를 구분해 보아야 한다고 주장했다. 레싱에 따르면 시와 회화는 정보를 처리하는 방식이 달라 근본적 차이가 발생한다. 그는 미메시스라 통칭되던 공통 기능을 회화적 묘사에서는 표현으로, 서술적 묘사에서는 재현으로 나누어 매체 차별성을 강조했다. 레싱은 「라오콘 군상」을 사례로 들어, 조각은 고통을 표정으로 직접 제시하지만 시는 독자가 장면을 상상해 구성하게 만든다고…(중략)
}
\begin{kfQBlock}{8}{}
\kfQStem{(나)의 페히 관점에 대한 설명으로 가장 적절한 것은?}
\begin{kfChoices}
\kfChoice 페히는 매체를 형식과 분리하되 위치 전환 가능성을 인정한다.
\kfChoice 페히는 매체를 형식과 분리하되 위치 전환 가능성을 인정한다.
\kfChoice 페히는 매체와 예술 형식을 완전히 동일한 개념으로 본다.
\kfChoice 페히는 통합 관점만으로 상호 작용 분석이 충분하다고 본다.
\end{kfChoices}
\end{kfQBlock}

\kfSecHeader{kfAreaWeekly}{지문}{문항 9 지문 / 결합/분리}

\kfPassageNote{문항 9 지문 / 결합/분리}{%
우리는 일반적으로 마음속 상태를 겉으로 분명하게 드러내거나 나타낼 때 ‘\uline{표현}’이라는 말을 사용한다. 그러나 표현이라는 말이 단지 자신의 심리적 상태를 드러내는 데에만 쓰이는 것은 아니다. 때로는 자연을 대하는 우리의 태도나 반응에도 표현이라는 말이 적용된다. 예를 들어 축 늘어진 나뭇가지를 보고 “애처롭다.”라고 말하는 경우가 있다. 이는 나뭇가지 자체가 감정을 지닌 것은 아니지만, 우리가 그것을 바라보며 내면의 정서를 \uline{투영}하고 외부로 드러내는 행위를 통해 표현이라는 말을 사용하게 되는 것이다. 하지만 표현이라는 말이 가장 널리 사용되는 분야는 다름 아닌 예술이며, 예술과 표현의 관계에 대해서는 오래전부터 다양한 견해가 논의되어 왔다.

첫째로, \uline{㉠표현}을 예술가가 창작 과정 속에서 행하는 내적 활동으로 보는 견해가 있다. 이 입장에서 표현은 예술가의 마음속에서 일어나는 정신적 과정, 즉 내적인 작용으로 이해된다. 이러한 입장을 대표하는 사상가가 …(중략)
}
\begin{kfQBlock}{9}{}
\kfQStem{㉮의 주장에 대한 근거로 가장 적절한 것은?}
\begin{kfChoices}
\kfChoice 예술가는 특별한 직관 능력을 지니지만, 내적 이미지가 자발적으로 형성되는 방식은 일상에서도 본질적으로 동일하므로
\kfChoice 예술가는 특별한 직관 능력을 지니지만, 내적 이미지가 자발적으로 형성되는 방식은 일상에서도 본질적으로 동일하므로
\kfChoice 예술은 일상의 정신적 삶과 분리된 특수한 기능의 산물이므로
\kfChoice 기존의 미학이나 예술학이 일상과 예술의 연속성을 일관되게 중시해 왔으므로
\end{kfChoices}
\end{kfQBlock}

\kfSecHeader{kfAreaWeekly}{지문}{문항 10 지문 / 조건 왜곡}

\kfPassageNote{문항 10 지문 / 조건 왜곡}{%
일반적으로 혼인은 남자와 여자가 부부가 되는 일을 말한다. 과거에는 특정한 법적 절차 없이도 결혼식과 같은 사회적 관습에 의한 의식을 치르면 혼인이 성립한 것으로 보았다. 하지만 현재 대부분의 국가는 법에서 규정한 절차를 따라야만 혼인이 성립되는 법률혼 제도를 두고 있으며, 사회적 관습보다는 혼인 신고와 같은 법적 절차가 더 중요하다. 혼인은 사회를 구성하는 기본 단위인 가족을 형성하는 일이기 때문에 국가의 근간을 이루는 중요한 과정이고 이에 따라 법을 통해 관계를 보호하는 것이다. 따라서 현대 사회에서 혼인은 법적 절차에 의해 이루어지게 되고 여기에 참여한 남녀 간에는 법적 효력이 발생하게 된다.

혼인이 성립하기 위해서는 법에서 규정하고 있는 실질적 요건과 형식적 요건을 모두 만족하여야 한다. 실질적 요건에서 가장 중요한 것은 혼인 의사의 합치이다. 혼인 의사는 부부 관계를 성립하겠다는 의사로 반드시 혼인 당사자 간의 합의가 필요하다. 또 다른 실질적 요건은 연령과 혼인 자격에 …(중략)
}
\begin{kfQBlock}{10}{}
\kfQStem{윗글과 <보기>를 바탕으로 판단한 내용으로 적절하지 않은 것은?}
\begin{kfEvidence}<보기> A와 B는 결혼식을 올렸지만 혼인 신고를 하지 않은 채 4년간 동거했다. 이후 A는 일방적으로 관계를 종료했고, B는 사실혼 파기로 인한 손해배상을 검토하고 있다. 한편 C와 D는 혼인 신고를 마친 법률혼 부부로, 이혼에 합의한 뒤 가정법원 확인을 받고 숙려 기간을 거쳤다. 다만 두 사람은 확인 후 4개월이 지나서야 이혼 신고를 하려 한다.\end{kfEvidence}
\begin{kfChoices}
\kfChoice C와 D는 법원 확인만 있었으므로 신고 시점과 무관하게 이미 법적 이혼이 확정되었다고 볼 수 있다.
\kfChoice C와 D는 법원 확인만 있었으므로 신고 시점과 무관하게 이미 법적 이혼이 확정되었다고 볼 수 있다.
\kfChoice A와 B 관계는 신고가 없으므로 법률혼이 아니라 사실혼 성격으로 다뤄질 수 있다.
\kfChoice 사실혼이라도 정당한 사유 없는 일방 파기에 대해 손해배상 문제가 발생할 수 있다.
\end{kfChoices}
\end{kfQBlock}

\kfSecHeader{kfAreaWeekly}{지문}{문항 11 지문 / 긍부정 반전}

\kfPassageNote{문항 11 지문 / 긍부정 반전}{%
\#\#\# [31 \textasciitilde{} 33] 다음 글을 읽고 물음에 답하시오.

세상이 린 몸이 밭이랑의 늘거 가니   바깥일 내 모고  일 무 일인고   이 중의 우국성심은 년풍을 원노라   <제1수 - 원풍>

새벽 밝아 오자 지빠귀가 소 다   일어나거라 아들아 ㉠ 밭 보러 가자꾸나   밤 이 이슬 긔운에 얼마나 기런고 노라   <제6수 - 신>

보리밥 지어 담고 풀로 끓인 국을 여   골 ㉡ 농부들을 제때에 먹이자꾸나   아야  그릇 올녀라 친히 맛보아 보내리라   <제7수 - 오>

서산에  지고 풀 끝에 이슬 난다   호미를 둘러 메고  등에 지고 가자꾸나   이 중의 즐거운 을 닐러 무엇 리오   <제8수 - 석>

* 이휘일, ｢ 전가팔곡 ｣ -


}
\begin{kfQBlock}{11}{}
\kfQStem{윗글에 대한 설명으로 가장 적절한 것은?}
\begin{kfChoices}
\kfChoice 권유의 의미를 나타내는 표현을 반복하여 화자의 태도를 강조하고 있다.
\kfChoice 권유의 의미를 나타내는 표현을 반복하여 화자의 태도를 강조하고 있다.
\kfChoice 반어적인 표현을 사용하여 시적 상황을 부각하고 있다.
\kfChoice 색채어를 활용하여 시적 대상을 감각적으로 표현하고 있다.
\end{kfChoices}
\end{kfQBlock}

\kfSecHeader{kfAreaWeekly}{지문}{문항 12 지문 / 비교대상 뒤바꿈}

\kfPassageNote{문항 12 지문 / 비교대상 뒤바꿈}{%
[26\textasciitilde{}28] 다음 글을 읽고 물음에 답하시오.

(가)   거미란 놈이 흉한 심보로 병원 뒤뜰 난간과 꽃밭 사이 사람 발이 잘 닿지 않는 곳에 그물을 쳐 놓았다. 옥외 요양을 받는 젊은 사나이가 누워서 치어다보기 바르게-

나비가 한 마리 꽃밭에 날아들다 그물에 걸리었다. 노오란 날개를 파득거려도 파득거려도 나비는 자꾸 감기우기만 한다. 거미가 쏜살같이 가더니 끝없는 끝없는 실을 뽑아 나비의 온몸을 감아 버린다. 사나이는 긴 한숨을 쉬었다.

나이보담 무수한 고생 끝에 때를 잃고 병을 얻은 이 사나이를 위로할 말이-거미줄을 헝클어 버리는 것밖에 위로의 말이 없었다.

* 윤동주, 「위로」 -

(나)   누가 와서 나를 부른다면   내 보여 주리라   저 얼은 들판 위에 내리는 달빛을.   얼은 들판을 걸어가는 한 그림자를   지금까지 내 생각해 온 것은 모두 무엇인가.   친구 몇몇 친구 몇몇 그들에게는   이제 내 것 가운데 그중 외로움이 아닌 길을   보여 주게 되리.   오…(중략)
}
\begin{kfQBlock}{12}{}
\kfQStem{(가)와 (나)의 공통점으로 가장 적절한 것은?}
\begin{kfChoices}
\kfChoice 동일한 시어를 반복하여 시적 의미를 강조하고 있다.
\kfChoice 동일한 시어를 반복하여 시적 의미를 강조하고 있다.
\kfChoice 명사로 시상을 마무리하여 시적 여운을 드러내고 있다.
\kfChoice 반어적 표현을 활용하여 화자의 태도를 부각하고 있다.
\end{kfChoices}
\end{kfQBlock}

\kfSecHeader{kfAreaWeekly}{지문}{문항 13 지문 / 주체/객체 혼동}

\kfPassageNote{문항 13 지문 / 주체/객체 혼동}{%
\#\#\# [22～27] 다음 글을 읽고 물음에 답하시오.

(가)   흰 벽에는 ――   어련히 해들 적마다 나뭇가지가 그림자 되어 떠오를 뿐이었다.   그러한 정밀*이 천년이나 머물렀다 한다.

단청은 연년(年年)이 빛을 잃어 두리기둥에는 틈이 생기고, 볕과 바람이 쓰라리게 스며들었다. 그러나 험상궂어 가는 것이 서럽지 않았다.

기왓장마다 푸른 이끼가 앉고 세월은 소리없이 쌓였으나 ㉠문은 상기 닫혀진 채 멀리 지나가는 바람 소리에 귀를 기울이는 밤이 있었다.

주춧돌 놓인 자리에 가을풀은 우거졌어도 봄이면 돋아나는 푸른 싹이 살고, 그리고 한 그루 진분홍 꽃이 피는 나무가 자랐다.

유달리도 푸른 높은 하늘을 눈물과 함께 아득히 흘러간 별들이 총총히 돌아오고 사납던 비바람이 걷힌 낡은 처마 끝에 찬란히 빛이 쏟아지는 새벽, 오래 닫혀진 문은 산천을 울리며 열리었다.

―― 그립던 깃발이 눈뿌리에 사무치는 푸른 하늘이었다.

* 김종길, 「문」 -   * 정밀 : 고요하고 편안함. …(중략)
}
\begin{kfQBlock}{13}{}
\kfQStem{ⓐ～ⓔ에 대한 설명으로 적절하지 않은 것은?}
\begin{kfChoices}
\kfChoice ⓒ : 잊음에 대해 ‘나’가 제시한 가정적 상황이 틀리지 않았음을 강조하기 위한 물음이다.
\kfChoice ⓒ : 잊음에 대해 ‘나’가 제시한 가정적 상황이 틀리지 않았음을 강조하기 위한 물음이다.
\kfChoice ⓐ : 잊는 것에 대한 ‘나’의 생각을 전개하기 위한 물음이다.
\kfChoice ⓑ : 잊음에 대한 ‘나’의 생각이 어디에서 비롯된 것인지에 대한 답을 제시하기 위해 던지는 물음이다.
\end{kfChoices}
\end{kfQBlock}

\kfSecHeader{kfAreaWeekly}{지문}{문항 14 지문 / 내용과 방법 혼동}

\kfPassageNote{문항 14 지문 / 내용과 방법 혼동}{%
\#\#\# [16\textasciitilde{}19] 다음 글을 읽고 물음에 답하시오.

‘부산 부두에 발을 올려 딛는 때부터 내 고향이다. 내 고향은 나에겐 편안히 쉴 자리를 줄 리가 없다. 그것을 바라고 그것을 꾀할 나도 아니다. 그곳에는 여러 동무들이 있을 것이다. 어서 신들메를 끄르지 말고 그대로 뛰어나오시오. 당신만은 온몸을 사리고 저편에 붙지 말고 용감하게 우리 속에 와 끼어 주시오. 이렇게 부르짖는 힘차고 씩씩한 친구들이 나를 맞아 줄 것이다. 오, 어서 달려가다오!’   윤건은 차 속이 좁고 갑갑한 듯이 땀에 절은 학생복 저고리는 벗어 걸어 놓고 셔츠 바람으로 몇 번이나 승강대에 나와서 날아가는 이국의 밤경치를 내다보곤 하였다.   그 이튿날 아침, 차가 고베 플랫폼에서 쉬게 되었음에 윤건은 도시락을 사러 나왔다가 어떤 낯익은 조선 청년을 만나게 되었다. 그 청년도 윤건을 얼른 알아보고 마주 와서 손을 잡았다.   “귀국하시는 길입니까?”   “네.”   “저도 이 찻간에 탔습니다.”   그 청년은 …(중략)
}
\begin{kfQBlock}{14}{}
\kfQStem{윗글의 서술상 특징으로 가장 적절하지 않은 것은?}
\begin{kfChoices}
\kfChoice 이야기 외부의 서술자가 특정 인물의 관점에서 사건과 인물의 심리를 서술하고 있다.
\kfChoice 이야기 외부의 서술자가 특정 인물의 관점에서 사건과 인물의 심리를 서술하고 있다.
\kfChoice 외부 이야기의 서술자가 자신이 겪은 내부 이야기의 의미를 밝히고 있다.
\kfChoice 서술자가 여러 인물의 내면을 서술하여 인물의 다양한 특성을 드러내고 있다.
\end{kfChoices}
\end{kfQBlock}

\kfSecHeader{kfAreaWeekly}{지문}{문항 15 지문 / 내용 왜곡}

\kfPassageNote{문항 15 지문 / 내용 왜곡}{%
\#\#\# [38\textasciitilde{}42] 다음 글을 읽고 물음에 답하시오.

(가)   어와 저 낭자야 내 말씀 들어보소   속세에 묻혔다고 오랜 인연 잊을쏘냐   낙포선녀* 보랴 하면 전생에 네 아닌가   남관* 포의(布衣) 백면서생도 신선(神仙)인 줄 뉘 알리오   요지의 잔치에서 복숭아를 훔친 사람 너이건만   주고받은 같은 죄라 너와 내가 귀양 왔네   넓고 아득한 천하에 동서로 나누이니   넓고 넓은 푸른 바다 은하수가 되어 있다   너도 나를 보랴 하면 여덟 고개 첩첩(疊疊)하고   나도 너를 보랴 하면 한라산이 아득하다   평생에 한이 되고 자나깨나 원하더니   옥황상제 감동한지 선관(仙官)이 두둔한지   태을선(太乙仙)의 연잎 배에 돛을 높이 달아   자라 수염에 배를 매고 제주 땅에 들어오니   아름다운 꽃과 나무 선계의 경치로다   풍경도 좋거니와 좋은 인연 더욱 좋다   연꽃 얼굴 버들눈썹 전생 모습 그대로요   검은 머리 흰 피부는 세속 모습 전혀 없다   (중략)   꽃다운 맹세…(중략)
}
\begin{kfQBlock}{15}{}
\kfQStem{(나)에 대한 이해로 적절하지 않은 것은?}
\begin{kfChoices}
\kfChoice ‘나’는 사람들이 ‘외손녀’에 대해 아는 것을 바라지 않았다.
\kfChoice ‘나’는 사람들이 ‘외손녀’에 대해 아는 것을 바라지 않았다.
\kfChoice ‘나’는 ‘외손녀’에게 자신의 마음을 전하고 싶어 글을 썼다.
\kfChoice ‘나’는 ‘외손녀’가 밝은 성품과 슬기로움을 지녔다고 생각했다.
\end{kfChoices}
\end{kfQBlock}

\kfSecHeader{kfAreaWeekly}{지문}{문항 16 지문 / 과잉 해석}

\kfPassageNote{문항 16 지문 / 과잉 해석}{%
[25 \textasciitilde{} 28] 다음 글을 읽고 물음에 답하시오.   [앞부분 줄거리] 설렁탕집 주인 ‘달평 씨’는 선행은 아무도 모르게 해야 한다는 신념을 가진 인물이다. 그러나 우연히 신문 기자들에 의해 선행이 과장되어 세상에 알려지면서 달평 씨는 대중들의 시선을 의식하게 되고, 본래 자신의 모습을 잃어버리는 첫 번째 죽음을 맞게 된다.

그러나 어쩐 일인지 세상 사람들의 관심은 달평 씨에게서 자꾸 멀어져가고 있었다. 그것을 눈치 못 챌 매스컴들이 아니었다. 달평 씨의 미담이 세상 사람들에게 알려지는 기회가 부쩍 줄어들었다.   그러나 달평 씨는 거기서 물러설 위인이 아니었다. 그가 입을 더 크게 벌렸다.   “나는 전과잡니다. 용서 못 받을 죄를 수없이 지고도 뻔뻔스럽게 살아온 흉악무도한 죄인입니다.”   달평 씨는 듣기에 끔찍한 지난날 자기의 악행을 요목요목 들추어 만천하에 공개하기 시작했다. 치한, 사기, 모리배, 폭력…… 등등, 그는 초빙되어 간 그 강단에 서서 꾸벅꾸벅 조는 사람들의 …(중략)
}
\begin{kfQBlock}{16}{}
\kfQStem{<보기>를 참고하여 윗글을 감상한 내용으로 적절하지 않은 것은? [3점]}
\begin{kfChoices}
\kfChoice ‘달평 씨에게 씌워’진 ‘친선 단체의 회장직 감투’를 거부하지 않은 것은 불우한 사람들까지도 철저하게 속이려는 달평 씨의 허위의식을 보여 주는군.
\kfChoice ‘달평 씨에게 씌워’진 ‘친선 단체의 회장직 감투’를 거부하지 않은 것은 불우한 사람들까지도 철저하게 속이려는 달평 씨의 허위의식을 보여 주는군.
\kfChoice ‘세상 사람들에게 알려지는 기회가 부쩍 줄어들’자 ‘입을 더 크게 벌’리는 달평 씨의 모습에서 대중의 관심을 얻고자 하는 인물의 욕심이 드러나는군.
\kfChoice ‘끔찍한 지난날 자기의 악행’을 공개하자 ‘다시 달평 씨를 입에 올리기 시작’하는 사람들을 통해 자극적인 정보에만 반응하는 대중들의 모습을 보여 주는군.
\end{kfChoices}
\end{kfQBlock}

\kfSecHeader{kfAreaWeekly}{지문}{문항 17 지문 / 범위 오류}

\kfPassageNote{문항 17 지문 / 범위 오류}{%
\#\#\# 2024년 11월 고3 수능

[22-27] 다음 글을 읽고 물음에 답하시오.

(가)   배를 민다   배를 밀어보는 것은 아주 드문 경험   희번덕이는 잔잔한 가을 바닷물 위에   배를 밀어넣고는   온몸이 아주 추락하지 않을 순간의 한 허공에서   밀던 힘을 한껏 더해 밀어주고는   아슬아슬히 배에서 떨어진 손, 순간 환해진 손을   허공으로부터 거둔다

사랑은 참 부드럽게도 떠나지   뵈지도 않는 길을 부드럽게도

배를 한껏 세게 밀어내듯이 슬픔도   그렇게 밀어내는 것이지

배가 나가고 남은 빈 물 위의 흉터   잠시 머물다 가라앉고

그런데 오, 내 안으로 들어오는 배여   아무 소리 없이 밀려들어오는 배여   - 장석남, <배를 밀며>

(나)   당신……, 당신이라는 말 참 좋지요, 그래서 불러봅니다 킥킥거리며 한때 적요로움의 울음이 있었던 때, 한 슬픔이 문을 닫으면 또 한 슬픔이 문을 여는 것을 이만큼 살아옴의 상처에 기대, 나 킥킥……, 당신을 부릅니다 단…(중략)
}
\begin{kfQBlock}{17}{}
\kfQStem{문항별}
\begin{kfChoices}
\kfChoice 번은 '통제할 수 없는 익명의 욕구'를 상대를 향한 사랑이 운명적이어서 멈출 수 없다는 뜻으로 해석했습니다.
\kfChoice 번은 '통제할 수 없는 익명의 욕구'를 상대를 향한 사랑이 운명적이어서 멈출 수 없다는 뜻으로 해석했습니다.
\kfChoice 번은 배가 들어오는 것을 '대상과의 재회가 예상대로 이루어짐'이라고 해석했습니다.
\kfChoice 번은 '당신'을 화자의 내면에 사는 '병자'로서 연민의 대상이라고 설명했습니다.
\end{kfChoices}
\end{kfQBlock}

\kfSecHeader{kfAreaWeekly}{지문}{문항 18 지문 / 시점/화자 혼동}

\kfPassageNote{문항 18 지문 / 시점/화자 혼동}{%
[38\textasciitilde{}42] 다음 글을 읽고 물음에 답하시오.

(가)   장마 가뭄에 피해 입은 백성이 관찰사 가을 순행 기다림은   가을걷이 부족함을 채워줄까 해서인데 지나는 곳마다 죄를 묻는 폐단 있네   무논 재해도 감췄는데 목화밭이야 거론할까   백 묘(畝)나 되는 벌건 땅에 백지징세 하는구나   인자한 우리 임금 곡식 한 묶음도 모래 덮일까 염려하는데   불쌍한 백성 논밭에다 좁은 길 넓히란다   각읍 관리 독촉하니 채찍 몽둥이 낭자하다   허다한 관인들이 대호(大戶) 소호(小戶)에 분담시켜   사방(四方) 부근 십 리 안에 닭과 개가 멸종하네   부자는 괜찮지만 가련한 이 가난한 자로다   해는 기울고 이정*은 저녁밥 재촉할 때   텅 빈 부엌에서 우는 아낙 발 구르며 하는 말이   방아품에 얻은 양식 한두 되 있건마는   채소도 있건마는 그릇은 누구에게 빌릴꼬   앞뒷집 돌아보니 섣달그믐에 시루 빌리는 격이로다   한 마을 닭과 개 다 먹어 치우고 집집마다 또 거둔단 말인가   대호…(중략)
}
\begin{kfQBlock}{18}{}
\kfQStem{㉠과 ㉡에 대한 이해로 가장 적절한 것은?}
\begin{kfChoices}
\kfChoice ㉠은 화자가 임금님의 은혜에 감사를 느끼는 시간이고, ㉡은 글쓴이가 새로운 것을 시도하는 시간이다.
\kfChoice ㉠은 화자가 임금님의 은혜에 감사를 느끼는 시간이고, ㉡은 글쓴이가 새로운 것을 시도하는 시간이다.
\kfChoice ㉠은 화자가 인생의 덧없음을 느끼는 시간이고, ㉡은 글쓴이가 이웃의 친절에 고마움을 느끼는 시간이다.
\kfChoice ㉠은 화자가 내적 갈등을 해결하는 시간이고, ㉡은 글쓴이가 자신의 삶의 가치를 새롭게 인식하게 되는 시간이다.
\end{kfChoices}
\end{kfQBlock}



\end{document}
